\documentclass[11pt]{book}

\usepackage{graphicx}
\usepackage{eso-pic}
\usepackage{tcolorbox}
\usepackage{polyglossia}
\setdefaultlanguage[indentfirst=false,forceheadingpunctuation=false]{russian}
\setotherlanguages{english}

\RequirePackage[a4paper, headsep = 0.5 \headsep, left=2.5cm, right=2.1cm, top=2cm, bottom=2.1cm]{geometry}

\setmainfont{Times New Roman}
\newfontfamily\cyrillicfont{Times New Roman}[Script=Cyrillic, Ligatures=TeX]

\usepackage{enumitem}
\setlist[itemize, 1]{label =\raisebox{-0.3\height}{\scalebox{1.8}{\textbullet}}}

\title{Философия эмоций}
\author{Кристина Тапполет}
\date{2023}
\begin{document}

\AddToShipoutPictureBG*{\includegraphics[width=\paperwidth]{title.png}}

\null
    \thispagestyle{empty}%
    \addtocounter{page}{-1}%
    \newpage

\tableofcontents

\part{Философия и наука об эмоциях}

\chapter{Философия эмоций}

\section{Введение}

Представьте, что пытаетесь набирать текст будучи запертым в иглу, и ''комок слякоти из перегретого купола уже в третий раз падает в пишущую машинку, заканчивая вашу работу на сегодня'' (Бриггс 1970, 259). Фрустрация и гнев в такой ситуации вполне понятны. Но антрополог Джин Л. Бриггс обнаружила, что понятны они не для всех --- в своей книге ''Без гнева'' (1970) она описывает эмоциональную жизнь уткухихалингмиутов (утку, для краткости) --- общины инуитов в арктической тундре к северо-западу от Гудзонова залива, где она провела полтора года. На собственном горьком опыте она узнала, что утку выражают неодобрение остракизмом --- им запрещается злится. Конечно, к вспышкам гнева у детей они терпимы, но взрослые должны контролировать свои эмоции. Потому злость, а вместе с ней грусть и любовь они выражают куда реже, чем это делается, скажем, в средней североамериканской семье. Казалось бы --- печальный факт, но в их суровой природной среде жесткий контроль сильных эмоций и сопровождающего их насилия вполне оправдан.

Подобные нормы есть и в других культурах --- идеал жизни, лишенной эмоциональных потрясений многим показался привлекательным. Например, согласно Стоицизму, страх, гнев и зависть --- иррациональные импульсы, связанные с ложными суждениями, и от эмоций, значит, нужно освободиться, чтобы достичь спокойствия ума. Буддизм также призывает отказаться от ненависти, и даже от желания и любви --- состояний, влекущих привязанности, удерживающие от освобождения. Да и в современных обществах к эмоциям относятся как минимум с осторожностью --- страх, гнев, зависть, отвращения и даже сострадание или любовь, считается, мешают ясно мыслить.

Но указывается также, что эмоции могут быть и полезны --- страх помогает избежать опасности, отвращение --- употребления вредных веществ. Иногда мы даже стремимся испытать негативные эмоции без, казалось бы, какой-то пользы: прыгаем с парашютом, слушаем грустную музыку или смотрим фильмы ужасов. А гнев помогает справиться с противником, и иногда кажется оправданным --- ''праведным гневом''. Уж тем более трудно отрицать пользу положительных эмоций --- сострадания, любви, радости, восхищения и надежды --- к ним и удовольствию от них мы обычно стремимся, и они, похоже, ключевой компонент счастливой жизни. Принято считать, что жизнь, лишенная эмоций --- и положительных, и отрицательных --- была бы бедна и едва ли человечна. Кто бы выбрал таковую жизнь андроида Дейты из ''Звездного пути'' --- владеющего мощнейшим позитронным мозгом, но вовсе не испытывающего эмоции?

Итак, наше отношение к эмоциям неоднозначно --- кажется, что следует быть не слишком им подверженным, и все же хочется вести полноценную, наполненную ими, жизнь. Эта амбивалентность проявляется не только в противоположных оценках эмоций, но и, собственно, эмоционально --- мы боимся гнева, отвращения, любви, и самого страха порой тоже боимся. И все же эти переживания нас привлекают --- и позитивные, такие как любовь и радость, и негативные, такие как гнев и страх. Но так думают большинство людей независимо от жизненного уклада и профессии, а что насчет философов? Конечно, философы тоже люди, и чувства их не могут принципиально отличаться. Так что и мысли их по поводу эмоций не отличаются фундаментально.

Чтобы дать представление о деятельности философов в этом направлении, я сделаю краткий исторический обзор философии эмоций. Обзор этот, конечно, избирателен и фокусируется лишь на нескольких основополагающих фигурах, тем самым игнорируя, быть может, не менее значительных --- Баруха Спинозу (1632-1677) и Адама Смита (1723-1780), например, а также не-европейских мыслителей --- таких как конфуцианский философ Мэн-Цзы (ок. 1040-221 гг. до н.э.) (Мэн-Цзы 2008, см. Вираг 2017; Ван Норден 2019; Вонг 2019) или буддийский философ Буддхагхоша (V-ый век) (Буддхагхоша 1984, см. Хейм 2013, Таске 2016), не говоря уже о женщинах-мыслителях, которым не предоставилась возможность оставить след в истории идей (см. Суперсон 2020).

\section{Эмоции: от Платона до Фомы Аквинского}

Современная западная философия берет свое начало в Древней Греции, и философия эмоций не исключение. И хотя ни Платон (ок. 429-347 гг. до н.э.), ни Аристотель (384-322 гг. до н.э.) не разработали законченную теорию эмоций (или страстей --- \textit{pathê} --- того, что мы претерпеваем, будучи по отношению к нему пассивны), они придерживались особых на них взглядов, оказавших огромное влияние --- и продолжающих влиять --- на мыслителей и по сей день (см. Прайс 2009). Оба они --- и Платон, и Аристотель --- считали человеческий разум (или ''душу'' --- \textit{psyche}) разделенным на части. Согласно ''Государству'' Платона, душа состоит из рациональной части --- разума, и низшей, нерациональной, делящейся на яростную (тумос) и страстную (эпитумия). Яростная часть ответственна за гнев, стыд, негодование, страх, гордость и восхищение, а страстная связана с телесными желаниями, такими как жажда и голод. Можно подумать, что эмоций касается только яростная часть --- но нет --- и Платон прямо утверждал, что эмоции сопровождают и некоторые телесные желания. Более того, он относил любовь к истине и мудрости к разумной части. Примечательно, что, согласно Платону, эмоции --- не исключительно психическое явление --- и они, и телесные позывы, он считал, соответствуют телесным нарушениям. Так, гнев, например, сопровождается кипением крови от мысли о несправедливости.

Центральное место в концепции души Платона занимает возможность конфликта между разными ее частями. Например, особый недостаток рациональности, называемый акрасией (слабость воли), Платон определял как совершаемое несмотря на веру в его неправильность, и объяснял в терминах конфликта между страстной и рациональной установкой. И чтобы таких конфликтов не было, рациональная часть, конечно, должна преобладать --- в ''Федре'' он сравнивает разум с возничим, ведущим двух тянущих в разные стороны крылатых коней. Платон, впрочем, считал яростную часть ближе и податливей разуму (в сравнении со страстной) и допускал возможность научиться, например, любить и ненавидеть так, чтобы конфликта с разумом не возникало.

У Аристотеля была схожая иерархия в его концепции, только он делил душу на две части: рациональную и не-рациональную --- ответственную за эмоции. В ''Никомаховой этике'' он включает в список эмоций телесные желания (эпитумию), а также гнев, страх, уверенность, зависть, радость, любовь, ненависть, тоску, соперничество и жалость. Как и Платон, Аристотель говорит об эмоциях как о затрагивающих не только душу, но и тело, описывая разгневанного при мысли об оскорблении как вскипающего. Он дает определения ряда эмоций, среди которых гнев, негодование, стыд, жалость, зависть и страх. Гнев --- это желание мести, сопровождающееся болью из-за действительного или воображаемого пренебрежения к себе или близкому. Примечательно, что для гнева характерна не только боль, но и удовольствие [предвкушения мести]. Эмоции, считал Аристотель, включают боль или удовольствие, оценку (например, веру в незначительность чего-то или его опасность), мотивацию, а также телесные изменения.

Как и Платон, Аристотель указывал на эмоции как на частый источник иррациональности --- конфликт между ними и разумом может привести к действию вопреки здравому смыслу --- акрасии. Но он также признавал за эмоциями возможность играть положительную роль --- считал их ключевыми для моральных/этических добродетелей и счастья. Моральные добродетели --- привычки, тесно связанные с ощущением нужных эмоций в нужное время по отношению к нужным предметам и людям (см. главу 10), так что и гнев может быть правильной реакцией. Добродетели, считал Аристотель, требуют воспитания, и эмоциональные реакции зависят от паттернов, усвоенных в общественных взаимодействиях. Как становятся хорошими игроками на лире, тренируясь на ней играть, так и становятся храбрыми, встречаясь лицом с опасностью и приучаясь испытывать при этом правильные эмоции. Идея воспитания эмоций оказалась, вообще очень влиятельной (см. главу 12).

Аристотель считал эмоции необходимыми для добродетели и счастья --- но в рамках Стоицизма, важной школе мысли, основанной в Афинах в III веке до нашей эра и просуществовавшей вплоть до римского периода, к эмоциям относились куда более критично. В работах стоиков можно найти подробное рассмотрение таких негативных эмоций как гнев, страх, горе, печаль, стыд, зависть, ревность и жалость. В отличие от Платона и Аристотеля, они не верили в разделенную душу, представляя ее как физическую субстанцию (пневму), смешанную с телом. Душа, однако, обладает управляющей способностью, через которую реализуется, среди прочего, способность к рассуждению. Эмоции же --- согласно стоикам --- иррациональны, их надо не просто контролировать, но --- в идеале --- полностью искоренить. Согласно Сенеке (ок. 1 до н.э. - 65 н.э.), ''Гнев следует не контролировать, но совершенно уничтожить, ибо какой контроль возможен на по сети своей злобным?'' (Сенека, ''О гневе'', iii. 42). То есть разделяя душу на части, стоики не согласовывали, но противопоставляли их. Наличие эмоций --- источник акрасии, влекущей скорые и разрушительные изменения в душе.

Что же не так с эмоциями, согласно стоикам? Эмоции подразумевают суждения о добре и зле, которые глубоко ошибочны. Горе, например, следствие убеждения в действительном наличии нечта плохого --- согласие с наличием такого и заставляет считать реакцию правильной. Однако такого рода убеждения основаны на ошибочном представлении о человеке и его месте в мире. Хороша же только добродетель, и только соответствующие ей действия обеспечивают истинное счастье. Считать же важным что-нибудь еще --- свою удовольствие или боль, например --- не просто плохо, но иррационально, ведь ни к чему кроме несчастья это не приведет. Стоики предлагали разные методы избавления от эмоций --- от когнитивной терапии, включающей философские рассуждения, до развития соответствующих привычек и даже музыкальной терапии.

Стоиков часто высмеивали за их идеал --- существо, начисто лишенное эмоций, кажется недосягаемым для нормальных людей, но, справедливости ради, стоики и не рекомендовали жизнь, лишенную даже намека на эмоции. Допустимо чувствовать своего рода протоэмоциональную реакцию --- не соглашаясь с ней, не позволяя ей стать полноценной эмоцией, движущей к определенным действиям. Некоторые психические состояния, сейчас считаемые эмоциями, стоиками даже поощрялись --- например, восторг (если он вполне обоснован). Впрочем, даже с учетом этих уступок, идеал их может показаться чрезмерно требовательным и в целом нежелательным.

Средневековых философов тоже интересовали эмоции (их называли \textit{passio} или \textit{affectus}). Фома Аквинский (1225-1274), на подход которого значительно повлиял Аристотель, предложил наиболее обширную трактовку эмоций в своей ''Сумме теологии'' (Часть II -- 1.22--48 и 2.17--36) (см. Кнууттила 2004, Перлер 2018). Подобно стоикам, современники Фомы и он сам были сосредоточены на нормативном вопросе --- хороши или плохи эмоции, являются ли любовь, ненависть, надежда, отчаяние, страх, радость или гнев, добродетелью или грехом. Фома Аквинский, считавший эмоции движениями души, подчеркивал их побудительную силу. Эмоции --- это мотивация к действию: гнев, например --- желание мести. Мотивация эта связана с телесными изменениями и возникает из оценки чего-либо как приятного или болезненного или, вообще, значимого. Фома Аквинский называет эти категории ''формальными объектами'', и понятие это (как будет видно в Главе 2) до сих пор играет в философии эмоций важную роль. Поскольку оценки, на которых основаны эмоции, могут быть ошибочными, их, считал Фома, следует проверять разумом. Но искоренять эмоции не стоит --- достаточно держать их под контролем. Итак, Фома Аквинский следовал Аристотелю, утверждая необходимость для добродетели соответствующих разуму эмоций, и --- раз эмоции могут командам разума сопротивляться, контроль разума должен быть подобен политической власти над свободными сопротивляться. Рефлексия, считал он, способна не только ослабить, например, гнев, но и полностью его прекратить, и для этого достаточно помыслить вызывающее его в ином свете. Как будет видно из Главы 11, стратегия такой переоценки занимает центральное место в современных подходах к регуляции эмоций.

\section{Современная философия и эмоции}

Рене Декарт (1596-1650) (см. Шапиро 2020) обычно считается отцом современной философии. Декарт открыто презирал средневековые теории эмоций в своей книге ''Страсти души'' (1649/1989). Его занимала проблема отношения разума и тела, которые он считал отдельными субстанциями. Эмоции, согласно Декарту --- это восприятия, ощущения или волнения души, вызванные движениями животных духов --- мельчайших частиц, сообщающих душе готовность тела к определенным действиям. То есть эмоции затрагивают не только ум, но и тело. Примечательно, что животные (кроме человека), считал Декарт, не имеют разума и потому неспособны испытывать эмоции и даже просто чувствовать боль или удовольствие --- это резко контрастирует со взглядами его предшественников, допускающими способность животных испытывать эмоции даже не имея разума.

Но как взаимодействуют разум и тело? Это сложный вопрос, отвечая на который, Декарт разработал теорию, согласно которой, тело влияет на разум передачей своих движений посредством животных духов через шишковидную железу в основании мозга. Но эта теория, конечно, не может быть корректной --- если разум состоит из нефизической субстанции, то такой ''факс в душу'' не объясняет воздействие физического на разум.

Хотя Декарт и уделяет много внимания такому механистическому объяснению эмоций, его теория весьма схожа как с аристотелевской, так и со стоической. Так, говоря про эмоции как про волнения души, вызванные движением животных духов, он все же указывал, что некоторые из них включают оценку вещей с точки зрения вреда и пользы. Любовь и радость, например --- являясь одной из основных эмоций наряду с удивлением, ненавистью, желанием и печалью --- вызваны представлением своего объекта как хорошего и полезного[, при этом любовь, в отличие от радости подразумевает неизменность этих свойств со временем]. Эмоции обладают мотивационной силой и отличаются от оценочных суждений тем, что суждения требуют также активного согласия с оценкой. То есть можно выносить суждение о вещи как о хорошей, и все же невольно реагировать на нее как на плохую, не соглашаясь, впрочем, с этим представлением. Как будет видно в Главе 6, такие ''непокорные эмоции'' (Д'Армс и Джейкобсон, 2003) занимают видное место в современных дебатах о природе эмоций.

Декарт отличал эмоции от суждений еще тем, что эмоции могут вводить в заблуждение. Они часто искажают наше понимание реальности, заставляя верить, что все лучше или хуже, чем на самом деле, так что эмоции следует контролировать. Напрямую управлять эмоцией нельзя, но, считал Декарт, с первоначальной эмоциональной реакцией можно не соглашаться --- согласие же даст ей импульс к дальнейшему развитию. Косвенный контроль над эмоциями возможен посредством другой эмоции --- чувства собственного достоинства, способности свободно осуществлять свою волю. Декарт также выделял несколько способов изменить свои [первичные] эмоциональные реакции. Можно, во-первых, пытаться представить объект своих эмоций в ином свете. Или представить последствия эмоциональных реакций --- если они негативные, это будет постепенно притуплять порождающие их первичные реакции. В конце концов, частое столкновение с тем, что вызывает эмоциональные реакции, также постепенно притупляет их.

Другой известный философ этого периода --- Дэвид Юм (1711-1776) --- взял на себя задачу пересмотреть традиционную рационалистическую концепцию эмоций (см. Кохон 2018). Вместо поиска способов их контроля он прямо заявил, что эмоции играют --- и должны играть --- ведущую роль. В своем ''Трактате о человеческой природе'' он указал, что ''разум должен быть --- и был всегда --- только служанкой чувств'' (Юм 1739-1740/2000, Книга 2, Часть 3, Раздел 3). Роль разума --- полезная, но вспомогательная: он предоставляет информацию о мире и средствах к достижению наших целей, но цели эти предоставляются эмоциями. Разум сам по себе не обладает мотивационной силой и, значит, не в состоянии действительно противостоять эмоциям --- любая попытка такого мнимого противостояния предполагала бы своей причиной другие эмоции. В своем ''Трактате'' Юм утверждал, что эмоции не состоят в представлении вещей, а значит и не могут, подобно суждениям, оцениваться разумом как рациональные или иррациональные. Впрочем, в более поздней работе ''О стандарте вкуса'' (1757/1985a) Юм от этого тезиса отказался. Но в ''Трактате'' он также утверждал, что ''нет никакого противоречия в том, чтобы предпочесть смерть всего мира, царапине'' (Юм 1739-1740/2000, Книга 2, Часть 3, Раздел 3). Однако, Юм, конечно, решительно осудил бы такое предпочтение --- размышление о смерти мира вызывает ужас --- но рациональность или иррациональность здесь не при чем.

Юм выделяет два вида эмоций. Первый --- ''прямые страсти'' --- включает желание, отвращение, надежду, горе, радость, страх, отчаяние и противоположность страха --- чувство безопасности. Они вызываются соответствующим опытом или мыслью о добре и зле, что для Юма то же, что переживание удовольствия или боли или мысль о них. К так называемым же ''косвенным страстям'' относятся гордость, смирение, стыд, честолюбие, тщеславие, любовь, ненависть, зависть, жалость, злоба и щедрость --- помимо переживания боли или удовольствия, они требуют также дополнительных мыслей. Гордость за красоту своего дома, например --- удовольствие, направленное на самого себя и вызванное мыслью о принадлежности вам этого дома. Вообще, гордость соответствует положительному представлению о себе, а стыд --- отрицательному. То есть, согласно Юму, и прямые, и косвенные эмоции подразумевают оценку своих объектов, так что различие его теории с таковыми предшественников не так уж велико --- Юма отличает не столько описание природы эмоций, сколько понимание их взаимосвязи с разумом.

Такая инверсия отношения разума к эмоциям является ключевой для морального сентиментализма --- для Юма, моральные оценки берут свое начало из эмоций, а не из разума, и именно на основе удовольствия или неудовольствия в от восприятия людей и их действий, мы выносим моральные суждения. То есть добродетельным считается человек, характер которого вызывает одобрение, а одобрение --- это определенный род удовольствия. Удовольствие это опосредованно ''симпатией'' --- спонтанном ощущении удовольствия/неудовольствия при виде его у других. Так, доброжелательность одобряется потому что приносит удовольствие другим. Юм, как не странно, указывал на способность эмоциональных реакций вводить в заблуждение. Например, при рассмотрении личности врага мы склонны быть предвзятыми, из-за чего не можем дать ей справедливую оценку. Нам, значит, стоит принять нейтральную точку зрения, чтобы не допускать подобных отклонений.

Другой знаменитый философ --- Иммануил Кант (1724-1804) --- выступает против концепции морали Юма (см. Уилсон и Дэнис 2018). Юм указывал, что основание морали --- эмоции по отношению к другим, Кант же полагал ее основой практический разум --- способность определять волю и приводить к действию. Мораль, согласно Канту, возможна благодаря автономной воле --- той, что устанавливает сама себе закон. Такая воля должна игнорировать всякие привлекательные черты своих объектов, поскольку руководство ими значило бы руководство ''склонностями'', которые нам в действительности чужды. Следовать склонностям значило бы рабство --- такая воля не автономна. Поэтому эмоции не следует учитывать в принятии решения --- только так мы можем быть автономными агентами и соблюдать свои моральные обязательства.

Кант критиковал эмоции как основу для морали, указывая, что они слишком шатки для столь фундаментальной роли. В своей работе ''Основы метафизики морали'' (1785/1996) он утверждал, что такие ''эмпирические принципы'' не могут поддержать моральные законы, справедливые для всех [разумных] существ, ведь они привязаны к конкретной --- человеческой --- природе. Кант также отмечал ''бесконечные различия в степенях и видах чувств'', из-за чего они ''не могут предоставить единый стандарт хорошего и плохого'' (Кант 1785/1996, 4:442). Можно, впрочем --- независимо от вопроса об универсальности моральных принципов --- усомниться в справедливости такого осуждения эмоций, ведь наши рациональные способности так же зависят от нашей природы, а наши выводы могут различаться в разных обстоятельствах и даже сбивать нас с верного пути.

Кант, однако, не столь плохо относится к эмоциям, как можно было подумать. Так, одной конкретной эмоции --- уважению морального закона --- он отводит центральную роль. Эмоция эта возникает при рассмотрении морального закона и приводит к подчинению ему. Для таких существ, как мы, согласно Канту, требуется чувство удовольствия или восторга от выполнения своего долга. В более поздних работах (см., например, Кант 1797/2018) он даже указывал на обязанность культивировать такие моральные эмоции, как сострадание , любовь к людям и уважение к себе. В работе об эстетических суждениях Кант также отводил эмоциям центральную роль понимании красоты и возвышенного. И все же, в отличие от Юма, Кант считал подлинно моральными лишь действия, совершаемые исключительно из чувства долга.

Итак, видно, что интерес философов к эмоциям привел к появлению и развитию все более тонких и сложных их теорий. Но здесь имеет смысл задать вопрос: в чем, собственно, должно состоять философское исследование эмоций? Иначе говоря, какие основные вопросы предполагает эта область?

\section{Основные вопросы}

Большинство современных философов, конечно, сочли бы центральным философии эмоций вопрос ''Что такое эмоция?'' (выступающий также заголовком знаменитой статьи Уильяма Джеймса (1884)), и большая часть работы в этой области на нем и сосредоточена. Для ответа нужно разработать теорию, раскрывающую сущность эмоций, отделяющую их от прочих состояний, таких как зуд или мигрень, например. Многие философы прошлого (да и современные тоже) рассматривали эмоции как состоящие из нескольких компонентов. Так, согласно Аристотелю, они включают боль/удовольствие, оценки, мотивацию, а также телесные изменения. Есть ли среди этих компонентов те, что присущи всякой эмоции? Как вообще компоненты эмоций связаны друг с другом? Некоторые эмоции очень различаются --- страх и гордость, например --- так что имеет ли смысл говорить обо всех эмоциях сразу, или лучше рассмотреть только одну или две? Одна из задач философии эмоций --- выявить и объяснить различия между видами эмоций.

Впрочем --- и это видно из исторического очерка --- философов волнует не только вопрос о природе эмоций. Эмоции могут быть хороши или плохи, некоторые из них, значит, следует воспитывать, некоторые --- контролировать, от некоторых желательно вовсе отказаться (согласно стоикам --- почти от всех). Но даже радикальные критики эмоций соглашались, что есть эмоции, необходимые для счастья --- обоснованный восторг у стоиков и уважение к моральному у Канта, например. С другой стороны, даже те, кто отводил эмоциям главенствующие роли, отмечали вред от некоторых из них --- даже Юм признал, что эмоции могут приводить к ошибкам.

Такая двойственность в отношении эмоций --- признак их сложности: эмоции принимают самые разные формы и оттенки. Это также признак важности нормативного вопроса ''Хороши или плохи эмоции для нас?''. Иначе говоря, какова их функция? Из-за разнообразия эмоций, речь обычно заходит о нормативном статусе лишь нескольких или даже одной из них. Более того, обычно эмоция хороша в чем-то одном и плоха в другом: страх, например, помогает избежать опасности, но должен быть направлен на действительно опасный объект и быть не слишком сильным чтобы не перерасти в панику. Так или иначе, философы прошлого склонны были рассматривать эмоции именно с нормативной точки зрения.

Один из важных аспектов, рассматриваемых теориями эмоций --- их регулирование, будь то путем их воспитания, контроля или даже искоренения. Можно ли изменить наши эмоциональные реакции, и если да, то как? Такой вопрос, конечно, тесно связан с нормативным, ведь регулировать имеет смысл лишь то, что может быть хорошим или плохим. Более того, вопрос о природе эмоций можно рассматривать как инструментальный по отношению к вопросу их регулирования. Философы прошлого предоставили несколько способов такой регуляции, представляющих интерес не только для философов, но и для психологов.

Итак, философов [касательно эмоций] занимали три взаимосвязанных вопроса:

\smallskip

a) Вопрос о сущности: какова природа эмоций?

б) Нормативный вопрос: хороши или плохи эмоции [для нас]?

в) Вопрос регулирования: можем ли мы управлять эмоциями и если да, то как?

\smallskip

Наибольшее практическое значение имеет, конечно, нормативный вопрос, ведь от ответа на него зависит интерес к природе и регулированию эмоций. Конечно, остальные два вопроса тоже важны, и вопрос о природе --- наиболее фундаментальный. Зная, что такое эмоции, мы можем понять, как они могут быть плохи или хороши [для нас] и как мы можем --- если можем --- ими управлять. Вопрос же регулирования важен потому, что ответы на остальные два еще не способны дать практических рекомендаций.

До сих пор речь шла об интересах философов касательно эмоций, но они, конечно, не единственные, кого они занимают. Многие дисциплины --- от социальных наук до биологических наук, вносят свой вклад в понимание эмоций, не говоря уже об искусстве вообще и литературе в частности. В свете этого возникает вопрос о связи этих различных направлений друг с другом в рассмотрении эмоций.

\section{Актуальность результатов научных исследований}

Эмоции так или иначе изучаются очень многими научными дисциплинами, среди которых биология, психология, психиатрия, нейронауки, антропология, социология и история, а потому важна связь между этими разными к ним подходами. Проще говоря: как философам реагировать на результаты касательно эмоций из научных дисциплин? Должны ли философы передать тему эмоций социальным и биологическим так же, как передали биологии тему функционирования организмов? Или, может, философия способна выступить здесь интегратором результатов рассмотрения эмоций в разных их аспектах, предложив так единое их объяснение и, возможно, даже обнаружив недоступные в рамках отдельных дисциплин, выводы? Многие философы не согласны ни с первым, ни со вторым вариантом. Но что тогда может философия предложить в качестве особого подхода к эмоциям?

Во-первых, философия использует особые методы исследования эмоций. Во-вторых, ее интересуют другие касательно них вопросы. То есть у философов касательно эмоций есть и собственный предмет, и собственные методы (эту линию защиты см. в Робертс 2003). А значит, кажется, философам не обязательно уделять много внимания тому, что делается в прочих дисциплинах --- достаточно вежливого любопытства. Как же защитить эти два тезиса?

Начнем с методологического. Метафилософский вопрос о методах философии --- спорный, но большинство философов согласны, что основными инструментами являются концептуальный анализ, мысленные эксперименты, интроспекция и наблюдение. Но достаточно ли они точны и надежны? Во всяком случае, они сильно отличаются от экспериментальных и статистических методов, используемых в большинстве социальных и биологических наук. Даже философы, занимающиеся экспериментальной философией и использующие анкетирование, не используют микроскопы или сканеры МРТ. И все же, разрыв между философскими и научными методами не следует преувеличивать.

Возможно, научные методы не столь уж отличаются от философских, по крайней мере не так, как это кажется на первый взгляд. Научное наблюдение, конечно, гораздо систематичней любого используемого философами метода, но ведь и философии в той или иной мере свойственна систематичность. Интроспекция же может считаться своего рода наблюдением, что делает разрыв между философской и научной методологией еще меньше. Наконец, концептуальный анализ и мысленные эксперименты в некоторой мере необходимы и для научных исследований. Проверяемые гипотезы, пожалуй, невозможно представить без соответствующих концептуальных разъяснений, и в ходе разработки этих пояснений кажется необходимым рассмотреть примеры и контрпримеры. То есть различия между научным и философским методом --- если они есть --- касаются акцентов, и неясно, достаточно ли их для рассмотрения философского взгляда на эмоции как совершенно самостоятельного.

Эмоции --- сложное явление, а потому разные дисциплины изучают разные их аспекты. Нейронауки сосредоточены на нервной системе вообще и мозге в частности, и ищут основания эмоций в них; эволюционная биология рассматривает механизмы реагирования, которые помогли нашим предкам выжить и распространить свои гены; антропологию занимают социальные нормы, регулирующие эмоции, и т.д.. Из этого легко сделать вывод, что и у философии должен быть особый, интересующий только ее, аспект эмоций. Но давайте вернемся к трем центральным вопросам философии эмоций: имеет ли какой-нибудь из них исключительно философский характер?

Вопрос о сущности, казалось бы, только философов всегда и волновал. Разве сам термин ''сущность'' не сугубо философский? Однако используемая терминология может вводить здесь в заблуждение. Вопрос о сущности эмоций централен для философии эмоций, но представляет интерес и для прочих дисциплин. Ведь как, вообще, можно проводить эксперименты, посвященные эмоциям, не предполагая о них хоть что-то? Разные дисциплины предлагают разное понимание эмоций, по-разному их объясняя, так что вопрос о сущности не чисто философский.

Что же касается нормативного вопроса, то ученые часто относятся к таким вопросам с некоторой опаской, рассматривая их как сомнительные и даже ненаучные и способные заинтересовать только философов. Отчасти так оно и есть: вопросы о том, хороши эмоции для нас или плохи, --- это вопросы, занимавшие философов на протяжении веков, и философы, похоже, лучше всех к таким вопросам подготовлены, ведь этика является традиционной философской дисциплиной. Но прогресс в этих вопросах, очевидно, вряд ли возможен вне вклада из нефилософских дисциплин. Вопросом влияния эмоций на суждения, например, занимается психология, но они совершенно точно имеют отношение к пониманию влияния эмоций на моральный облик. Так же как и помощь эмоций в решении проблем должна быть подвергнута эмпирической, научной, оценке. Итак, нормативный вопрос, хотя и специфичен для философии, решаться должен с учетом результатов анализа эмпирических данных в рамках нефилософских дисциплин.

Аналогично дело обстоит с вопросом о регулировании эмоций. Он так же интересен не только философии --- им занимается психология, психиатрия и нейробиология. И философам, похоже, придется сотрудничать с учеными из этих дисциплин для эффективного его решения.

Итак, даже если нормативный вопрос специфичен для философии --- даже тогда для полноценного ответа на центральные вопросы философам нужно учитывать результаты из нефилософских дисциплин, и ни предполагаемые методологические различия, ни различные интересующие аспекты не гарантируют обратного. Подход, значит, должен быть междисциплинарным.

Здесь можно возразить, что даже включение философии в такой междисциплинарный подход не обязательно --- есть опасения касательно надежности ее инструментов. Возможно, стоит обойтись без них --- без, значит, концептуального анализа, мысленных экспериментов, интроспекции и обыденного наблюдения? Но они, похоже, в некоторой мере присущи и прочим дисциплинам, так что здесь нет однозначных причин отвергнуть философский подход.

В защиту философских методов скажу несколько слов о концептуальном анализе. Его основная идея в разбиений сложных концепций на простые с целью определения фундаментальных и построения уже из них новых сложных. Это схоже с химическим анализом, в рамках которого определяются фундаментальные химические элементы, из которых составлены сложные. Большинство обыденных понятий --- понятие знания, например --- сопротивляются такому анализу, и тогда зачастую лучше попытаться показать концептуальные связи между ними. Концептуальный анализ --- это важно --- не подразумевает редуцирование, \textit{сведение} одних концепций к другим, более фундаментальным. Концептуальные разъяснения позволяют выявить концептуальные истины, то есть утверждения, истинные в силу только лишь включаемых концепций. Желтый, например, является цветом, а то, что тяжелее слона, не может быть легче него. Или: если некто знает, что Земля круглая, то Земля круглая. Интересная особенность концептуальных истин --- направленность не на просто таки нейтральную регистрацию употребления терминов, а на понятия и мысли. Для разрешения противоречий между предполагаемыми концептуально истинными утверждениями и ответить на контрпримеры, нужен акт концептуального балансирования, в ходе которого некоторые утверждения также могут быть и отклонены.

Метод концептуального балансирования схож с методом рефлексивного равновесия, предложенным Джоном Ролзом в своей знаменитой работе ''Теория справедливости'' (1971). Ролз утверждал, что для получения обоснованных этических убеждений следует начать с некоторых приемлемых этических суждений и последовательно уточнять и исправлять их для получения теории. Аналогично, для получения концептуальных истин нужно начать с некоторых банальностей, непосредственно следующих из самих понятий и, сталкиваясь с противоречиями или контрпримерами, эти банальности корректировать --- модифицировать или отвергать --- постепенно получая все более связный их набор (Джексон 1998). Этот метод, как надеюсь продемонстрировать в следующих главах, успешно применим к связям между понятиями эмоций и оценочными понятиями.

Вообще, философское мышление может быть продуктивным в применении ко всем трем основным вопросам об эмоциях. Как следствие, философами были разработаны различные сложные теории эмоций, сказаны проницательные вещи об их влиянии --- хорошем или плохом --- на наши мысли и действия и об их связи с ценностями и добродетелями. Не так много сказано философами по вопросу регулирования эмоций, но и это немногое оказалось очень влиятельным. Философы не потеряли интерес к эмоциям, и, что очевидно из текущих их дебатов, продолжают вносить важный вклад в их понимание.

\section{Подведение итогов}

Эмоции играют важную роль в нашей жизни, но, впрочем, неясно --- положительную или отрицательную. Философов прошлого --- от Платона и Аристотеля до Декарта и Канта --- беспокоило негативное влияние эмоций и потому они стремились понять, что такое эмоции и как их можно регулировать. Современные философы не менее заинтересованы в вопросе о сущности эмоций --- наиболее фундаментальном --- но также активно рассматривают нормативный вопрос --- хороши или плохи эмоции [для нас], негативно или позитивно они влияют на нашу способность рационально мыслить. Вопрос регулирования эмоций все менее занимает философов, но имеет большое практическое значение. Философия --- не единственная дисциплина, интересующаяся эмоциями, и потому прогресс в соответствующих вопросах --- даже в нормативном --- требует от философов учета результатов работы ученых-эмпириков. Следующая глава будет посвящена ключевым концепциям философии эмоций.

\section{Краткое содержание}

\begin{itemize}
  \item\ Философы, как и все люди, неоднозначно относятся к эмоциям: они могут всячески вредить нам, и все же необходимы для счастливой и этически правильной жизни.
  \item\ Философы прошлого были весьма заинтересованы в выяснении природы эмоций и их связи с разумом. Часто они заключали, что эмоции иррациональны и оттого нуждаются в регулировании.
  \item\ Философия эмоций разрабатывается вокруг трех вопросов: (а) вопрос о сущности, природе эмоций, (б) нормативный вопрос о пользе или вреде эмоций и (в) вопрос регулирования, связанный с воспитанием, контролем или искоренением эмоций.
  \item\ Результаты научных дисциплин касательно эмоций имеют значение, потому что науки занимаются выяснением сущности эмоций и способами их регулирования, а также исследованиями, релевантными для нормативного вопроса.
  \item\ Методы философии --- концептуальный анализ, мысленные эксперименты, интроспекция и наблюдение --- отличаются от научных лишь акцентами, но не качественно.

\end{itemize}

\begin{tcolorbox}

  \section{Учебные вопросы.}

  \begin{enumerate}
    \item\ Вспомните свой последний приступ гнева. Он был неприятным, приятным или, может, одновременно приятным и неприятным?
    \item\ Как вы, в целом, относитесь к эмоциям? Есть ли среди них особенно желательные или наоборот --- те, которых хочется избегать?
    \item\ В чем специфика философских методов применительно к эмоциям и достаточно ли они надежны?
  \end{enumerate}

\end{tcolorbox}

\chapter{Аффективная сфера}

\section{Введение}

Вспомните свою реакцию на пожар в Соборе Парижской Богоматери 15 апреля 2019 года. Я очень хорошо запомнила это утро, начавшееся как совершенно обычное. Была хорошая, ясная погода, и, будучи еще немного сонной и уже немного голодной, я приняла душ, позавтракала и стала читать новости, попутно думая, как начать новую главу. Наверное, подумала я, использую пример. И тут меня поразила новость: пожар частично разрушил Собор Парижской Богоматери. Я отреагировала с удивлением и ужасом и внезапно окончательно проснулась. Глядя на фотографии падающего шпиля и разрушенных крыш, я чувствовала печаль от потери таких красивых и исторически значимых сооружений, представляя, как брожу среди развалин от пожара.

Как ясно видно из приведенного примера, наша жизнь состоит из потока переплетающихся психических явлений: ощущение хорошего настроения, радость, сонливость, растерянность, голод, воспоминание, удовольствие, размышление, беспокойство, решение, удивление и ужас и чувство полного пробуждения, наблюдение, печаль, воображение. И с первого взгляда не так уж очевидно, как их можно четко классифицировать.

Впрочем, как должно быть ясно еще из исторического обзора в прошлой главе, такие потоковость и разнообразие психической жизни не помешали философам прошлого довольно строго отделять, например, рациональную часть от не-рациональной. Для современной философии, вообще, довольно характерно противопоставление разума эмоциям и прочим аффектам.

Относительно недавно это избитое противопоставление получило новую, более эмпирическую обертку. Психологи предложили разделение когнитивных процессов на относящиеся к Интуитивной Системе (Система 1) --- предполагающие быструю, автоматическую, но довольно негибкую реакцию --- и относящиеся к Рассудочной Системе (Система 2) --- аккуратные и обдуманные, но медленные. Некоторые считают Рассудочную систему в целом более ''умной'', по крайней мере когда ей не мешает атавистическая Интуитивная (см., например, Канеман 2011), но это весьма спорно.

Эти так называемые \textit{двупроцессные теории познания} не лишены критики, но --- в той или иной форме --- принимаются подавляющим большинством психологов и экономистов. Споры вызывает не столько само разделение, сколько конкретная его формулировка, точное различение между этими двумя видами. Согласно одной из концепций, следует разделять когнитивные процессы на \textit{рефлексивные} --- подразумевающие значительную степень контролируемого внимания и манипулирование информацией ''в уме'' --- и \textit{интуитивные} --- предъявляющие минимум требований к такому манипулированию и оттого почти что не требующие внимания. Такие характеристики, как быстрота, бессознательность и автоматизм, конечно, коррелируют с итуитивностью процесса, но не обязательны. Важно понимать, что процессы обоих видов способны приводить к правильным реакциям, и правильность эта зависит во многом от ситуации. То есть процессы каждого из этих видов по-своему ''умны''.

Эмоциональные процессы --- как всякие аффективные --- обычно относят к интуитивным. Но, возможно, все не так уж просто. Эмоции, хотя и считаются пассивными, зависящими от быстрых автоматических процессов, состояниями, все же довольно тесно связаны с сознательной деятельностью. Но поскольку это верно и для чувственного восприятия, очевидно состоящего из интуитивных процессов, эмоциональные процессы можно рассматривать и как интуитивные, хотя отнесение всех таких процессов к интуитивным не столь уж очевидно. Такие эмоции, как, например, негодование и вина, похоже, зависят от суждений, а суждения предполагают контролируемое внимание. Но как соотносятся эмоции и соответствующие им суждения? Если суждения не нужны всякий раз для переживания эмоции, то эти эмоции, пожалуй, можно и не связывать с суждениями (см. Главу 3 и Главу 6). Вопрос, вообще, в том, насколько процессы одного типа взаимодействуют с процессами другого типа --- например, в ходе учебы игры на фортепиано. Вопрос этот особенно важен для понимания возможности и способов регулирования эмоций (см. Главу 11 и Главу 12).

Для ясности в этой главе будут представлены стандартные разграничения внутри аффективной сферы. Так, один из важных вопросов --- чем эмоции отличаются от прочих аффектов. Но что такое, вообще, аффекты? Термин ''аффект'' происходит от латинского effectus --- ''подвергаться воздействию, влиянию'' и использовался средневековыми философами для обозначения эмоций (см. Главу 1 Раздел 1). В современных обсуждениях концепция аффекта --- центральная в так называемых ''аффективных науках'' --- относительно молодой области исследований, объединяющей психологов, нейробиологов, антропологов и философов --- является общей для эмоций и прочих подобных им явлений. То есть под аффектами понимают не только эмоции, но и настроения, чувства, желания, удовольствие и неудовольствие и боль. Аффекты считаются психическими явлениями, характеризующимися как чувствами, так и мотивацией. Чем же эмоции особенны среди других аффектов? Давайте начнем с рассмотрения основных их характеристик, чувства и мотивации, например.

\section{Чувство и мотивация}

Глядя вниз с высокой скалы, мы обычно испытываем страх. Страх этот, несомненно, можно считать эмоцией, так же как и, например, отвращение при виде разлагающейся туши краба, гнев на мусорящих туристов, радость при мысли о встрече с друзьями и грусть когда оказалось, что они не придут или удивление при виде, что они уже прибыли. Эмоции, как видно, бывают самые разные.

Для эмоций есть множество терминов, начиная с уже использованных в прошлых примерах --- ''страх'', ''отвращение'', ''гнев'', ''радость'', ''печаль'', ''удивление'' --- и которые, похоже, можем чувствовать не только мы, но другие животные, и продолжая характерными, похоже, только для человека --- ''возмущение'', ''стыд'', ''вина'', ''уважение'', ''восхищение'', ''трепет'', ''горе'', ''надежда'', ''развлечение'' и ''облегчение''. И это ведь еще далеко не полный список таких терминов даже на русском языке, не говоря уже о специфичных для других языков и культур --- такой, пожалуй, будет очень длинным. Анна Вежбицка (1999), например, рассматривает более 50 терминов английского лексикона, связанных с эмоциями, а также множество терминов из других языков и культур. Поскольку то, что верно для некоторых из этих эмоций, не обязательно верно для всех них, а также потому, что некоторые эмоции не столь прототипичны (веселье или трепет, например), стоит начать с наиболее простых и распространенных. В этой главе речь пойдет главным образом о страхе, отвращении, гневе, счастье/радости, печали и удивлении. (Как будет видно в Главе 3, это список базовых эмоций, который предложил Пол Экман (1972).) Что же верно для эмоциональных реакций вроде страха при взгляде вниз со скалы?

Если эмоции --- аффективные явления, характеризующиеся чувствами и мотивацией, то, наверное, стоит начать с чувств --- тезис об их связи с эмоциями принимается многими философами прошлого и современности (см., например, Деонна и Терони 2012). Чувства считались столь ключевыми для эмоций, что некоторые решались даже их отождествить (см. Главу 4).

Действительно: когда мы боимся, мы \textit{чувствуем} страх, когда злимся --- злость, когда радуемся --- радость и т.д.. Чувство страха, например, здесь отвечает на вопрос, какого это --- бояться, на что это похоже. То есть эмоции обычно обладают тем, что философы называют \textit{феноменальными свойствами}. Считается, что чувства от эмоций связаны с физиологическими изменениями --- учащенным сердцебиением, например, а также с психологическими изменениями --- смещением фокуса внимания, например. Ощущения зависят, конечно, от испытываемой эмоции, но каждый ли вид эмоций имеет особые феноменальные свойства? И что общего у негативных эмоций вроде страха, отвращения и печали, что отличает их от, скажем, радости? Во всяком случае очевидно, что эмоции валентны --- бывают либо отрицательными, либо положительными. Впрочем, удивление, кажется, может быть иногда отрицательным, а иногда положительным --- в зависимости от оценки его объекта. Возможно, удивление --- умеренно негативная эмоция, за которой уже могут следовать как положительные, так и отрицательные. Вопрос в том, как объяснить валентность и природу отличительных феноменальных свойств эмоций вообще. В любом случае, кажется верным, что

\smallskip

(Феноменология) Эмоциональные эпизоды обычно обладают феноменальными свойствами.

\smallskip

Заметьте, что неочевидна концептуальная истинность этого утверждения. У птиц обычно есть крылья, но верно ли, что эмоции обычно связаны с чувствами? Если вы найдете воробья без крыльев, это будет значить, что что-то не так с воробьем, но не с концепцией воробья или концепцией птицы. Бескрылый воробей --- скорее исключения, как и, например, пингвин. Из концепции же эмоций не следует, что эмоция без чувств --- такое исключение. Можно отметить, что овладение понятием \textit{эмоции} вообще и конкретных эмоций в частности, зависит от их опыта, их переживания. Чтобы объяснить ребенку, что такое страх, мы называем страхом испытываемое в соответствующих обстоятельствах другими и им самим --- например, когда он слышит гром или видит рычащую собаку. Так что утверждение, что эмоции обычно обладают феноменальными свойствами, похоже, все же концептуальная истина.

Но даже тогда все еще возможны эмоции без таких свойств, то есть, типичность не значит обязательность. Ведь, кажется, нет противоречия в том, чтобы испытывать страх со всеми присущими ему характеристиками, но без соответствующих чувств. В связи с этим возникает вопрос о возможности бессознательных эмоций, к которому мы еще вернемся (см. Главу 4). Бессознательные эмоции, впрочем --- даже если они возможны --- не типичны, а значит, не противоречат исходному тезису.

Второй правдоподобный тезис касательно эмоций как аффектов, заключается в их связи с мотивацией. Как еще будет обсуждаться в Главе 5, этот тезис особенно популярен в психологии (Фрида 1986), хотя не чужд и философам (Скарантино 2014). Глядя вниз со скалы и чувствуя при этом страх, мы склонны вернуться в более безопасное место, отвращение к разлагающемуся крабу заставляет отвернуться, гнев на туристов побуждает обругать их или сделать что-нибудь еще, что заставит их вести себя культурней. То есть, похоже, верно, что

\smallskip

(Мотивация) Эмоциональные эпизоды обычно включают мотивацию

\smallskip

Примечательно, что концептуальная истинность этого тезис, как и предыдущего, не столь уж очевидна --- не очевидно, что отрицающий обыкновенность связи эмоций с мотивацией не понимает концепцию эмоций. Возможны ли эмоции без такого мотивационного компонента? Кажется, что нет противоречия в страхе или гневе без соответствующей мотивации. Можно, конечно возразить, указав на связь этимологий слов ''эмоция'' и ''мотивация''. Термин ''эмоция'' восходит к старофранцузскому emouvoir, что означает ''стряхивать'', который, в свою очередь --- к латинскому emovere, значащему ''выдвигать, удалять, волновать'' от ''ex-'' (выходить) и movere (двигаться). Однако движение изнутри можно понимать и как выражение эмоций, да и аргументы на основе этимологии не столь уж убедительны, учитывая, что значение терминов меняется со временем. К тому же, некоторые эмоции --- восхищение, например --- похоже, не имеют тесной связи с мотивацией. Так что рассматриваемый тезис верен не для всех эмоций. Но он верен, по крайней мере, для страха, отвращения и гнева --- эти эмоции сопровождаются тенденциями к определенным образам действия.

\section{Объекты эмоций}

Третий тезис касается направленности эмоций на что-то. Так, мы боимся \textit{упасть со скалы}, испытываем отвращение \textit{к разлагающемуся крабу}, злимся \textit{на туристов}, испытываем радость \textit{от мысли о встрече с друзьями}, грустим, \textit{что они не пришли}, или удивляемся, \textit{что они уже приехали}. Для обозначения того, на что направлена эмоция, философы обычно используют термин ''интенциональный объект'' (не в том смысле, в котором используется ''интенция'' как ''намерение'', а в том же, что ''направленность''). С наличием интенциональных объектов у эмоций согласны многие философы (см., например, Брентано 1874/1995, Кенни 1963/2003; Лион 1980). Третий тезис, значит, таков:

\smallskip

(Интенциональность) Эпизоды эмоций обычно имеют интенциональные объекты

\smallskip

Согласно этому, эмоции имеют важную общую с убеждениями и желаниями, черту. Например, если мы верим, что идет дождь, то эта наша вера заключается в том, что идет дождь, а если мы хотим, чтобы светило солнце, то это наше желание --- про ясную, солнечную погоду.

Однако, в отличие от убеждений, эмоции могут иметь разные интенциональные объекты (Деонна и Терони, 2012 и Скарантино и де Соуза, 2018). Конечно, если мы злимся, например, что на пляже мусор, то интенциональный объект --- некоторое положение дел. Но эмоции могут быть направлены и на конкретные объекты или события, как в случае с отвращением к разлагающемуся крабу или удивлением прибытию друзей. Такие конкретные объекты не обязательно реальны --- можно бояться возможной, но --- во всяком случае, пока что --- не реальной, бури, например. Объекты эти могут быть абстрактными, например, когда вы удивлены новой гипотезой. Наконец, эмоции могут быть направлены на целые классы вещей: бояться можно, скажем, не конкретного волка, но всех волков вообще.

Интенциональный объект эмоции часто является и ее причиной --- страх перед встретившимся на тропе волком, например, этим самым волком и вызван. Но такое совпадение не обязательно (см. де Соуза 1987): гнев по поводу шутки коллеги, например, может быть вызван не столько самой шуткой, сколько фрустрацией из-за проблем на работе. Еще более очевидный пример --- эмоции по поводу несуществующих объектов, как в случае с возможной бурей.

Обычно, однако, считается, что эмоции причинно связаны со своими интенциональными объектами. В случае с волком страх явно причинно с ним связан, ведь не встреть вы волка --- не было бы и страха. Но зависимость возможно и не от чувственного восприятия: достаточно, например, верить, что по тропе идет волк, даже его не видя, и от этого бояться. Действительно, воспоминания и ожидания могут легко вызывать эмоции. Поскольку такие причины связаны с такими когнитивными состояниями, как восприятия и убеждения, их часто называют ''когнитивной основой'' эмоций (см. Муллиган 1998; Деона и Терони 2012).

Как видно из рассмотренных примеров, отрицать тезис об интенциональных объектах эмоций едва ли возможно. Так что, похоже, это концептуальная истина. Возможно, впрочем, это верно только для таких эмоций, как страх и отвращение, ведь концепция эмоций, казалось бы, сама по себе не исключает отсутствие интенционального объекта --- возможно, казалось бы, переживать эмоцию, ни к чему не относящуюся. Например, можно просто чувствовать себя счастливым, без явного повода или того, ''о чем'' это счастье. Тревога также может не касаться сколь-нибудь конкретного объекта. Есть два способа ответа на такие возражения.

Согласно первому, отсутствие интенционального объекта --- скорее исключение. Но тогда наличие интенционального объекта --- уже не концептуальная истина, ведь есть такие ''беспредметные'' эмоции (Голди 2000 и Прайс 2006), как, скажем, чувство страха при резком пробуждении посреди ночи: вы смотрите в темноту, на форму занавесок, прислушиваетесь к шуму, но чего-то конкретного, чего бы вы боялись, нет.

Можно, впрочем, возразить, что в этом примере вы боитесь именно что темноты и формы занавесок, или даже указать, что объект страха есть, но просто еще не идентифицирован. То есть эмоция то направлена на что-то --- но вы пока просто не знаете, на что. Но такой ответ годится не для всех случаев эмоций --- счастье, например, явно не о каком-то не пока идентифицированном объекте, оно не направлено на что-либо вообще.

Второй же, и стандартный, способ ответа --- выделять такие ''беспредметные эмоции'' в отдельную категорию: это не эмоции, а \textit{настроения}. Эмоции же, значит, всегда направлены на что-то, и это, согласно разделению, концептуально истинно. Такое разделение позволило бы выделить класс аффективных состояний, имеющих общую особенность --- интенциональность. Но, как будет видно в Разделе 2.7, вопрос этот сложен, и некоторые настаивают на понимании настроений как лишь определенного \textit{вида} эмоций, и даже вопрос о все же наличии предмета у таких, казалось бы, беспредметных эмоций спорен.

\section{Эмоции и оценки}

Четвертая характеристика эмоций заключается в их связи с оценками. Этот тезис, как было видно, принимался многими философами прошлого (см. Главу 1). Он также популярен в психологии (см. Лазарь 1991), но еще больше --- среди современных философов (см., например, Нуссбаум 2001), некоторые из которых уверены, что оценки являются ядром эмоций. Так, страх при взгляде вниз со скалы, кажется, подразумевает оценку такого положения в пространстве как опасного, а отвращение к гнилому крабу, похоже, связано с его оценкой как заразного и несъедобного. Вообще довольно очевидна связь эмоций и оценок. Но какую именно роль играют оценки по отношению к эмоциям? Являются ли они их причинами или составляют сущностные их компоненты? К этому вопросы мы еще вернемся в Главе 6, а пока, во всяком, случае, верно, что

\smallskip

(Оценка) Эпизоды эмоций обычно включают оценку

\smallskip

Как и в случае с предыдущими тезисами, сразу не ясно, является ли он концептуальной истиной. Кажется, можно представить страх без соответствующей оценки --- возможно, это скорее исключение, но все же. Тезис, впрочем, даже тогда верен. К тому же, как будет обсуждаться в Главе 9, есть концептуальные связи между некоторыми видами эмоций и соответствующими оценочными концепциями. Например, есть очевидная связь между понятиями страха и устрашающего, отвращения и отвратительного. Такие примеры могут показаться слишком тривиальными и даже круговыми, чтобы вызвать интерес, но, как я буду утверждать, есть основания, что они указывают на важные свойства эмоций и их отношений к ценностям. В любом случае, такие концептуальные связи подтверждают рассматриваемый тезис.

Интересно, что иногда мы используем такие термины, как ''страх'', ''гнев'' и т.д. производным способом --- акцентируя внимание на оценочном содержании и оставляя прочие, эмоциональные, аспекты в стороне. Мы можем, например, сказать, что боимся соседской собаки, имея в виду ее оценку как опасной, но не испытывая при этом соответствующей эмоции. Или, мы можем сказать, что злимся на какого-то политика, подразумевая убеждение в его вредоносности, но не испытывая злости как эмоции. Приписываемое состояние здесь --- просто \textit{оценочное убеждение}, а не полноценная эмоция.

Тезис об оценках обычно понимался философами как указывающий на некие критерии эмоций. Страх, например, по определению направлен на опасность, а отвращение --- на что-то грязное (или просто таки отвратительное), и т.д.. Эти оценочные свойства принято называть ''формальными объектами''. Понятие формального объекта восходит к Фоме Аквинскому (см. Главу 1, Раздел 2), а в современной философии было популяризировано Энтони Кенни (1963/2003; см. также Лион 1980; де Соуза 1987; Терони 2007).

Итак, формальные объекты эмоций --- это то, что определяет их возникновение. Так, страх возникает тогда и только тогда, когда нечто опасно, отвращение --- тогда и только тогда, когда нечто отвратительно. Есть, значит, тесная связь между оценками и формальными объектами, причем оценка предполагает приписывание статуса формального объекта тому, на что направлена эмоция. Здесь, конечно, возникает ряд вопросов касательно понимания этой связи в зависимости от роли оценок в эмоциональной жизни (см. Главу 6), но наличие этой связи редко отрицается. Таким образом,

\smallskip

(Формальные объекты) Эмоциональные эпизоды обычно имеют формальные объекты

\smallskip

Концептуальная истинность этого тезиса, как и всех прошлых, может быть неочевидна: понятия эмоции и формального объекта могут показаться не столь близкими. Но раз понятие формального объекта --- сугубо техническое, в него легко встроить обычность этих самых объектов для эмоциональных эпизодов.

Эмоции связаны с оценками еще и тем, что сами могут подвергаться им. Эмоция может быть подходящей (целесообразной) или неподходящей, в зависимости от природы своего интенционального объекта. Так, отвращение, похоже, подразумевает отвратительность этого объекта. Вопрос такого ''соответствия'' довольно спорен, и будет подробно рассмотрен в Главе 9. Эмоции можно оценивать очень по-разному. Так, стоики (см. Глава 1, Раздел 2) указывали на их иррациональность --- и это, несомненно, отрицательная оценка, но возможна и иногда встречается и положительная, когда эмоции хвалят за их рациональность. Так же можно оценить полезность эмоций для испытывающего их, то есть их соответствие его интересам. Наконец, эмоции могут быть рассмотрены с точки зрения морали, исходя из последствий или из особенностей самих эмоций (гнев, например, можно считать аморальным и уродливым независимо от его последствий). То есть можно говорить не только о целесообразности эмоции, но и о ее рациональности и о ее благоразумности (наличия за ней пруденциальной ценности), а также о ее моральности (см. Главу 7 и Главу 8).

Есть, конечно, и другие способы оценивать эмоции, но те четыре, что я привела --- наиболее существенны. И даже в свете приведенных возникает несколько вопросов. Как сопоставлять эти все измерения, если они не сводятся друг к другу? Зависят ли эти оценки от культуры или даже отдельных в одинаковой культуре, индивидов? Уместность определенной эмоции, например, определяется и культурными нормами, и особенностями испытывающего эту эмоцию индивида. Во многих культурах от людей ожидают то, что соответствует их социальным ролям. Так, в Северной Америке выражение печали и страха считается недостойным для мужчины поведением, а от женщин ожидают меньше гнева (Броди 1999). Один из сложных вопросов здесь --- наличие [объективных] фактов, определяющих уместность эмоций вне зависимости от культурных норм.

\section{Темпоральный аспект эмоций}

Эмоциональные эпизоды вроде вызываемого взглядом вниз с большой высоты, видом гнилого краба, нерадивыми туристами и т.д. --- довольно кратковременны и длятся несколько секунд или минут. Похоже, это касается всех эмоций, то есть

\smallskip

(Темпоральность) Эпизоды эмоций обычно скоротечны

\smallskip

Психологами (см., например, Экман 1994) и философами (см., например, Бен-Зеев 2001) этот тезис воспринимается как нечто само собой разумеющееся. Впрочем, и его концептуальная истинность не очевидна, по крайней мере сразу. Конечно, есть эмоции, для которых он однозначно верен: удивление, например, едва ли можно ощущать по несколько часов подряд. Но страх, гнев, или --- тем более --- печаль, могут продолжаться несколько дней, месяцев и даже лет, и такая продолжительность не выглядит как исключение. Так, если придется длительное время находится рядом с источником большой опасности, страх может не отпускать ни днем, ни ночью, мешая даже спать. А печаль от смерти близких, например, может продлиться недели и даже месяцы. Впрочем, эмоции все же имеют тенденцию быть скоротечными, особенно в сравнении со, скажем, чертами характера и темпераментом, которые неизменны на протяжении десятков лет или всей жизни.

Вопрос темпоральности эмоциональных эпизодов касается тонких, но фундаментальных вопросов онтологии эмоций. Есть как минимум три взгляда на их онтологический статус, разнящихся в зависимости от представлений об их темпоральности. Согласно первому взгляду, испытывать, например, страх при столкновении с волком --- значит находиться в определенном психическом состоянии (см.: Робертс 2003; Принц 2004). У этого состояния, значит, есть начало и конец, а также разные способы, которыми оно проявляется --- чувства и мотивации, например. Но состояния статичны, а эмоции могут усиливаться и ослабевать, и включать прочие изменения, поэтому некоторые теоретики считают эмоции не состояниями, а событиями (второй взгляд) или процессами (третий взгляд). Согласно второму взгляду, страх перед волком следует понимать как событие, так же как мы понимаем извержение вулкана, землетрясение, или поход в горы (см. Яворский 2018). События, так же как и состояния, имеют ту или иную продолжительность, но могут также быть простыми или сложными, характеризоваться разными фазами и иметь разные причины и последствия. Эмоции, значит, предлагается рассматривать как события, связанные с чувствами и физиологическими изменениями, вызванными оценками и приводящими к мотивации. Третий же взгляд на эмоции предполагает представление их как процессов, наряду с, например, коррозией, фотосинтезом или пищеварением. Процессы также происходят во времени, включают разные компоненты и имеют определенный результат. Эмоции тогда --- сложные процессы, могущие включать оценочный, мотивационный и физиологический компоненты, которые каузально взаимодействуют (см. Робинсон 2005; Голди 2011).

Для простоты, впрочем, буду рассматривать эмоции как состояния. Но допускает ли такое понимание эмоций их ослабление или усиление, а также взаимодействие различных компонентов эмоциональных эпизодов? Есть основания полагать положительный ответ. Психические состояния взаимодействуют с рядом событий и процессов, таких как физиологические изменения, чувства или мотивации, и это взаимодействие определяет динамику эмоциональных эпизодов (см. Сотириу 2018).

Впрочем, чем бы мы не считали, например, страх --- состоянием, событием или процессом --- он явно отличается от прочих эмоциональных и аффективных явлений. Прежде чем приступить к рассмотрению таких различий эмоциональных эпизодов, подведем итоги. Мы обсудили несколько утверждений касательно эпизодов базовых эмоций, таких как страх, отвращение, гнев, счастье, печаль и удивление, и выяснили, что такие эпизоды обыкновенно скоротечны, включают оценки, чувства, мотивации интенциональные и формальные объекты. Теперь, с учетом этих характеристик, перейдем к рассмотрению эмоциональных диспозиций.

\section{Эмоциональные диспозиции}

В философии эмоций принято разделение на актуальные эмоции (те, что называют также эмоциональными эпизодами) и эмоциональными диспозициями (см. Кенни 1963/2003; Лион 1980; Деонна и Терони 2012). Есть разница между эпизодом страха, происходящим в определенное время и при определенных обстоятельствах, и склонностью переживать такие эпизоды [в подобных обстоятельствах]. Отвращение при виде разлагающегося краба, например, явно отличается от склонности такое отвращение испытывать при виде трупов или испорченных продуктов: первое --- эмоция, испытываемая в определенном контексте, второе --- то, что есть у вас всегда, некоторое состояние, которое уже порождает эмоции в соответствующих контекстах, но само, конечно, одним из порождаемых эпизодов не являющееся.

Эмоциональные эпизоды, значит, всегда зависят от эмоциональных диспозиций: испытываемые эмоции --- это проявление глубинных диспозиций к ним, так же как, например, разбитие вазы --- проявление ее хрупкости. Такие диспозиции могут затруднять жизнь, если приводят к патологическим эмоциям --- как при клаустрофобии, например. Но, конечно, диспозиции не обязательно проблематичны, ведь возможна склонность бояться только того, чего и стоит бояться.

Одно из ключевых отличий эпизодов от диспозиций --- темпоральность. Эмоциональные эпизоды, как указывалось в прошлом разделе, кратковременны, тогда как диспозиции не имеют определенной длительности, отчего их можно считать стационарными состояниями. Диспозиции, впрочем, могут приобретаться и утрачиваться, характеризуя своего носителя в течение пусть и длительных, но не обязательно равных всей жизни, периодов. Так, клаустрофобия как диспозиция может возникнуть после ночи в сломанном лифте, и может отравлять жизнь своего носителя в течение многих лет, если не всей жизни. Однако, от нее, возможно, удастся избавиться --- диспозиции могут меняться. Эмоциональные диспозиции, значит, несмотря на устойчивость, достаточно пластичны.

Еще одно важное отличие состоит в наличии у эпизодов феноменальных свойств, тогда как диспозиции ими, похоже, не обладают. Ведь диспозиция к отвращению к гниющим трупам вызывает чувства только тогда, когда проявляет себя в эмоциональном эпизоде. Таким же образом эмоциональные диспозиции можно считать включающими прочие характеристики эмоций --- но только косвенно, конечно. Они могут включать мотивации, интенциональные объекты и оценки в той же мере, что и их проявления. Согласно же другому взгляду, эмоциональные диспозиции подразумевают представления о типах объектов и соответствующих реакциях. Диспозиция к отвращению, например, включает представление о корректности соответствующей реакции на трупы и испорченные продукты. Эмоциональные диспозиции в таком случае --- устойчивые состояния, включающие репрезентации, а, значит, и собственные интенциональные объекты.

Обычно, впрочем, рассматриваются эмоциональные эпизоды, ведь через них можно понять также и диспозиции. Диспозиции, однако, тоже представляют интерес, ведь именно через их изменение возможно изменение первичных эмоциональных реакций (см. Главу 12). Я далее буду использовать термин ''эмоция'' как относящийся к эмоциональным эпизодам (другое использование см. в Вльхейм 1999; Голди 2000).

Эмоциональные диспозиции разделяют на несколько видов (см. Броуд 1954, Фрида 1994a, Деонна и Терони 2012). Те, что затрагивают только один вид эмоций --- называются однонаправленными, те же, что затрагивают несколько --- многонаправленными или ''сентиментами''. Диспозиция к отвращению к испорченным продуктам и клаустрофобия --- примеры однонаправленных, хорошим же примером многонаправленной является чувство любви, подразумевающее склонность испытывать целый ряд разных эмоций в зависимости от ситуации: радость от нахождения рядом с любимым человеком, печаль от разлуки с ним, беспокойство, когда он в опасности, злость, когда его кто-то оскорбил, гордость за его достижения и т.д. (см. Наар 2013b). Часто полагают, что моральные сентименты --- моральное неодобрение, например --- состоят из таких многонаправленных диспозиций или из включающих их устойчивых состояний. Так, ваше моральное неодобрение лжи предполагает склонность чувствовать вину, когда лжете, испытывать гнев, если лжет кто-то другой, восхищаться тем, кто никогда не лжет, и т.д. (Принц 2007).

Такие чувства (''сентименты''), как любовь, направлены на людей или типы вещей, но есть также диспозиции, направленные на ценности. Честный человек, например, ценит истину, и потому склонен огорчаться или злиться из-за нечестности, презирать нечестно добившихся успеха, восхищаться тем, кто честен, даже несмотря на риск, радоваться ''победе'' честности и т.д. (см. Главу 10). Честность обычно рассматривается как морально достойная черта, как добродетель. Согласно взгляду, берущему начало у Аристотеля, добродетели (как и пороки) включают такие многонаправленные диспозиции в дополнение к прочим --- касающимся действий, например. Направленные на ценности эмоциональные диспозиции, впрочем, не обязательно считать добродетелями или пороками. Шутливый человек, например, будет склонен относиться с юмором даже к серьезным вещам, испытывая веселье при посещении серьезных мероприятий, например, и сожалея, что другими они воспринимаются серьезно. Такие диспозиции как честность и шутливость, считаются чертами характера --- то есть тем, что определяет личность человека.

Итак, есть основания различать эмоциональные эпизоды и диспозиции, а так же однонаправленные и многонаправленные и объектные и ценностные диспозиции. Еще одно выделяемое здесь понятие --- настроение, его рассмотрим в следующем разделе.

\section{Настроения}

Базовые примеры настроений --- тревога, раздражительность, восторг и печаль. Из перечисленного очевидна корреляция между такими настроениями и эмоциональными эпизодами, можно составить пары: страх и тревога, гнев и раздражительность, счастье и восторг, эмоция печали и печальное настроение. Но всем ли настроениям соответствуют эмоции? Похоже, что нет, ведь нет настроения, соответствующего отвращению или удивлению, например. И наоборот, есть настроения, не имеющие эмоционального аналога --- например, лень или нетерпение. Но настроения и эмоции все равно довольно тесно связаны и даже имеют общие свойства.

Во-первых, настроения и эмоции характеризуются сходными феноменальными свойствами. Тревога, например, напоминает страх, а раздражение --- гнев. Часто указывают на присущесть настроениям некоторой рассеянности, разбавленности в сравнении с эмоциями. Действительно: настроения обычно менее интенсивны, но длятся дольше. Конечно, раздражение предполагает не столь сильные чувства, как гнев, но следует не забывать, что и настроения могут варьироваться в интенсивности: тревога, например, может доходить до паники, эмоции же, напротив, могут быть довольно слабыми. Некоторые настроения также могут длиться не столь долго, как некоторые эмоции.

Одно из ключевых отличий настроений от эмоций заключается в том, что настроения, похоже, не способны к непосредственной мотивации. Так, переживание страха тесно связано с соответствующей мотивацией, но тревога, как правило, обходится без нее, хотя некоторые настроения могут и вызывать экспрессивное поведение --- стук пальцем по столу от раздражение, например (см. Голди 2000). Принято считать, что настроения в перспективе существенно влияют на нашу психическую жизнь (см. Прайс 2006). Они направляют наши мысли, провоцируя соответствующие суждения, а также сказываются на внимании, памяти и способностях к классификации.

Но самое разительное отличие эмоций от настроений заключается, конечно же, в интенциональности: настроения ни на что не направлены. Раздражительность не направлена на что-либо, она ни к чему конкретному не относится. Впрочем, возможно, настроения не совершенно лишены таких объектов, а просто касаются вещей глобально, будучи одновременно ни о чем конкретном и обо всем.

Некоторые теоретики, в свете описанных особенностей настроений, предлагают рассматривать их как диспозиционные состояния (см. Гриффитс 1997). Но проблема в том, что настроения включают чувства, поэтому принято все же считать, что они не диспозиционны, но актуальны --- так же как эмоциональные эпизоды, которые отличаются интенциональностью. Эмоции направлены на конкретные объекты, тогда как настроения --- на мир в целом или большинство вещей вокруг (см. Соломон 1976/1993; Принц 2004). Настроения, значит, как и эмоции, предполагают оценки и формальные объекты. Формальные объекты соответствующих настроений и эмоций коррелируют: и тревога, и страх, например, направлены на опасность, и присущие им оценки также касаются опасности. Одна из проблем с этим подходом в том, что большинство настроений тогда будут ошибочны (Деонна и Терони 2012): мир, например, редко опасен целиком и полностью. Другая проблема --- возможность направленности эмоций на довольно общие объекты, такие как стая волков или все человечество, например, из-за чего может быть неясно, к эмоции или настроению отнести страх, направленный на все вокруг.

Интенциональность настроений можно также показать через неопределенные объекты (Росси 2021), вероятности (Прайс 2006), или оценочные возможности (Тапполет 2018), но раз настроения подразумевают чувства, то, похоже, должны быть такими же актуальными состояниями, как эмоциональные эпизоды. А значит, возможны и соответствующие диспозициональные состояния --- склонности к определенным настроениям в определенного рода контекстах. Можно быть склонным, например, к тревоге, и, похоже, такие диспозиции и составляют \textit{темперамент}. Склонность к определенному настроению, конечно, порождает и склонность к соответствующим эмоциям --- через актуализацию этого настроения.

Учитывая все эти особенности, давайте вернемся к вопросу интенциональности эмоций. Ранее было указано, что такие, казалось бы, эмоции, как счастье и тревога, не направлены на что-либо конкретное. Эти контрпримеры можно исключить, назвав такие ненаправленные ''эмоции'' --- настроениями --- и эмоции тогда по определению будут интенциональны. Представляют ли указанные здесь особенности настроений угрозу для такого хода? Быть может, настроения это всего лишь диспозиции к эмоциям? Нет, уже была показана неудовлетворительность такой характеристики. Однако, можно считать, например, счастье --- настроениями, направленными, но не на что-то конкретное, а на весь мир.

Помимо тех, что обсуждались, есть и другие аффективные явления --- желания, удовольствия, неудовольствия и боль, например. При рассмотрении мотивационных теорий мы увидим, что желания и эмоции обычно различаются направленностью соответствия (Серль 1983): желания подразумевают соответствие мира уму, т.е. мир должен измениться, чтобы удовлетворить их, тогда как эмоции направлены от ума к миру, т.е. должны отражать существующее в нем (см. Главу 5). Что же касается удовольствий и неудовольствий, то некоторые из них очевидно подразумевают эмоции: можно, скажем, радоваться погоде или надеется что она улучшится. Удовольствия и неудовольствия, похоже, обладают признаками эмоций: они подразумевают интенциональные объекты и их оценку (а значит, и мотивацию), а также чувства. Более того, приятное и неприятное вполне можно считать формальными объектами удовольствия и неудовольствия. Это все верно, конечно, и для чувственных удовольствий, таких как удовольствие от вкуса эспрессо, например. Похоже, что такие удовольствия являются аффективными состояниями, предполагающими ощущения и их приятность (Росси 2018). Возможно, чувственные удовольствия и неудовольствия и не эмоции в полной мере, но, во всяком случае, имеют с ними много общего. Что же касается этих заключений относительно боли, многое зависит от ее связи с неудовольствием: поскольку некоторые боли (якобы) не болезненны, то есть основания не спешить относить их к чувственным неудовольствиям (см. Граек 2007, Бейн 2014).

Прежде чем завершить эту главу, хочу указать на то, что может создать впечатление пробела в описании --- страсти. Термин ''страсть'', как и большинство рассматривавшихся в этой главе --- ''аффект'', ''эмоция'' и ''сентимент'' --- имеет долгую историю, на протяжении которой менял значение. Современные теоретики, впрочем, склоняются к тому, что это лишь такое древнее название эмоций, подчеркивающее пассивность переживающего их к ним. В обыденном употреблении ''страсть'' значит сильную, едва контролируемую эмоцию или желание, но может также означать устойчивый интерес, играющий важную роль в жизни.

\section{Подведение итогов}

Рассматривая аффективные состояния вообще и эмоции в частности, мы обсудили основные их особенности и пришли к выводу, что, похоже, базовые эмоции вроде страха и гнева обычно включают феноменальные свойства и мотивации, а также направлены на [интенциональные] объекты. Объекты эти соответствующе оцениваются, например то, чего боимся --- опасное, и оценки эти определяют формальные объекты. Наконец, эмоции довольно скоротечны. Все эти утверждения не обязательно концептуально истинны, но достаточно правдоподобны, чтобы считаться хотя бы \textit{предполагаемыми} концептуальными истинами. Мы также увидели, что эмоциональные диспозиции или ''сентименты'' отличаются от эмоциональных эпизодов отсутствием феноменальных свойств и большей длительностью. Настроения же отличаются от эмоций отсутствием или, возможно, тотальностью/неопределенностью интенционального объекта --- они обо всем и не о чем одновременно.

Теперь, улучшив свое представление о том, что такое эмоции, мы можем приступить к решению вопроса отнесения их к культуре или природе.

\section{Краткое содержание}

\begin{itemize}
  \item\ Двупроцессные теории предлагают различать интуитивные и рефлективные процессы. Эмоции обычно затрагивают только интуитивную сторону, но куда относить те, что зависят от мыслей --- сложный вопрос.
  \item\ Эпизоды таких эмоций как страх, гнев и т.д. обычно включают чувства, мотивацию, интенциональные объекты, оценки и формальные объекты, а также относительно скоротечны.
  \item\ Эмоциональные эпизоды представляют собой проявления эмоциональных диспозиций (склонностей).
  \item\ Эмоциональная диспозиция может быть однонаправленной --- склонность испытывать определенную эмоцию в определенном классе случаев, или многонаправленной (''сентиментом'') --- включающей в себя несколько однонаправленных, также диспозиции могут касаться не только вещей и людей, но и ценностей.
  \item\ Настроения характеризуются чувствами, не направленными на что-либо конкретное. Некоторые философы уверены, что настроения можно считать направленными на такие глобальные объекты, как весь мир.
  \item\ Диспозиции к настроениям составляют темперамент.
\end{itemize}

\section{Учебные вопросы}

\begin{tcolorbox}
  \begin{enumerate}
    \item\ Можно ли отнести голод к аффективным состояниям? Если да, то следует ли считать его эмоцией? Если нет, то чем он отличается от прототипических эмоций вроде страха и гнева?
    \item\ Похоже, что страх можно ощущать в течение нескольких дней или даже недель. Но что в этом отношении следует сказать о времени, когда вы спите? Можно ли, вообще, испытывать эмоции во время сна?
    \item\ Каждой ли эмоции соответствует некоторое настроение? Какое, например, соответствует эмоции отвращения?
  \end{enumerate}
\end{tcolorbox}

\chapter{Эмоции как социальные конструкты}

\section{Введение}

Представьте небольшой итальянский город Верону, каким он был еще в эпоху Возрождения, когда в нем жили Ромео и Джульетта. Их история --- история любви с первого взгляда, напряженного и трагического романа. Когда Ромео Монтекки и Джульетта Капулетти, оба в подростковом возрасте, впервые видят друг друга, тут же влюбляются. Менее чем через 24 часа они тайно женятся, несмотря на вражду, разделяющую их семьи. Они успевают провести вместе одну ночь: вынуждаемая выйти замуж за другого, Джульетта принимает зелье и впадает в кому. Решив, что она мертва, Ромео убивает себя, отчего Джульетта, очнувшись, также кончает жизнь самоубийством. Конец.

Во-первых, странно, что описания столь драматичных, ужасных событий могут привлекать нас (к этому мы вернемся в Главе 12). Во-вторых, такие истории демонстрируют нам представления о любви и то, как они меняются со временем и как отличаются в разных культурах. Представления о любви, конечно, имеют свою историю, но как насчет самой любви? Она, казалось бы --- естественное явление, присущее всем людям, но так ли это однозначно? Возможно, любовь --- лишь социальное изобретение, конструкт. Согласно К.С.Льюису (1936, см. также Блох 1992), например, романтическая любовь как феномен зародилась в конце XI века на юге Франции, в специфическом социальном и культурном контексте, где придворные трубадуры сочиняли песни, прославляющие храбрость рыцарей, стремящихся угодить своим возлюбленным. Такая любовь предполагала ряд правил, определяющих, о чем следует думать, что следует говорить и делать влюбленному. Например, исключалась любовь к более чем одному человеку, а так же предполагалась ревность --- ''тот, кто не ревнует, не может любить''. И эти правила, согласно Льюису, выполнялись не только Ромео и Джульеттой, но и выполняются нами до сих пор --- мы так невольно подражаем ведем себя по сценарию тех трубадуров.

Вопрос, значит, в том, являются ли сами аффективные состояния --- помимо представлений о них --- продуктом культуры. Возможно, они --- как деньги или социальная роль --- совокупности базовых элементов, соединяемых в соответствии с во многом произвольными культурно установленными правилами. Вопрос о культурной сконструированности может быть поднят для каждого из аффективных состояний: как для страха при взгляде вниз с большой высоты, так и для эмоциональных диспозиций --- однонаправленных, как склонность к страху высоты, или многонаправленных, как любовь. В этой главе, впрочем, я сосредоточусь на эмоциональных эпизодах.

Социальный конструктивизм в отношении аффективных состояний противопоставляется биологическому детерминизму, согласно которому наши эмоции зависят от врожденных особенностей нашей психологии, являющихся результатом естественного отбора. Как мы вскоре увидим, обе эти позиции неубедительны, а истина, похоже, лежит где-то между ними, то есть наши эмоции --- продукт как природы, так и культуры, и вопрос лишь в том, как эти две силы взаимодействуют.

\section{Биологический детерминизм}

В противовес социальному конструктивизму, биологический детерминизм утверждает существование панкультурных видов эмоций --- присутствующих во всех человеческих популяциях. Эти эмоции вызваны психологическими диспозициями, заложенными в нашей генетической структуре. Дарвин в своей работе об универсальности эмоционального выражения (Дарвин 1872) утверждал, что панкультурные эмоциональные диспозиции --- это эволюционные адаптации, возникшие в результате естественного отбора. И даже будучи подверженными культурному влиянию, наши эмоциональные предрасположенности сохраняют свое ядро. Эмоциональные эпизоды, значит, представляют собой эффективные способы решения основных проблем, с которыми столкнулись наши предки --- ''борьба, влюбленность, побег от хищников, сексуальная неверность, неудача и потеря статуса, смерть члена семьи'' (Туби и Космидес 2000, 92; см. также Экман 1999) --- и преодолевая которые, могли значительно успешней распространять свои гены.

Есть две основные теории биологического детерминизма. Первая --- теория базовых эмоций, восходящая к 1970-м годам и связанная с психологами Полом Экманом (1972) и Кэроллом Изардом (1971), неврологом Яаком Панксеппом (1998), а также философами Полом Гриффтсом (1997) и Андреа Скарантино (2014). Согласно этой теории, биологический детерминизм верен только в отношении некоторых видов эмоций --- базовых --- которые мы разделяем с прочими животными. Вторая же теория --- развиваемая эволюционными психологами, такими как Роберт Франк (1988) или Джон Туби и Леда Космидес (1990) --- отдает предпочтение эволюционным объяснениям не только базовых, но и так называемых высших эмоций --- например, любви, ревности или вины --- считающихся уникальными для человека.

Рассмотрим сначала теорию базовых эмоций. Первоначальный список таких эмоций, составленный Экманом, ограничен шестью эмоциями: гнев, отвращение, страх, счастье, печаль и удивление (Экман, 1972). Примечательно, что позже Экман предложил расширенный список из 15 таких эмоций: веселье, гнев, презрение, покой, отвращение, смущение, восхищение, страх, вина, гордость, облегчение, печаль, удовлетворение, чувственное удовольствие и стыд (Экман, 1999). Панксепп предложил менее ортодоксальный такой список: забота, страх, похоть, паника, игра, ярость и поиск (Панксепп 2000), список же Изарда состоял из 10 пунктов: гнев, презрение, отвращение, стресс, страх, вина, интерес, радость, стыд и удивление (Изард 1971).

Крайне важно, как при этом различаются базовые и остальные эмоции, достаточно ли схожи, значит, состояния, чтобы относить их к одному виду. Пол Гриффс (1997), например, считает базовые эмоции слишком отличными от прочих, чтобы относить их к одному естественному виду. Большинство теоретиков, впрочем, сходятся во мнении, что базовые и не-базовые эмоции имеют достаточно общих черт, чтобы принадлежать к хотя бы одной психологической категории (вне зависимости от принадлежности к одному естественному виду) (см. Робертс 2003; Принц 2004; Деонна и Терони 2012). Принадлежность к такой категории зависит, конечно, от соотношения базовых и не-базовых эмоций, и к этому вопросу мы еще вернемся в Разделе 3.7.

Согласно теории базовых эмоций, соответствующие эмоциональные эпизоды представляют собой каскады сложных скоординированных изменений, запускаемых автоматической оценкой. Эти изменения касаются мимики, интонаций, тенденций к действию, а также затрагивают вегетативную нервную систему и могут провоцировать определенные воспоминания и ожидания. И основные аффективные программы, воспроизводимые в ходе эмоциональны эпизодов, строго определены отдельными нейронными связями (hardwired neural circuits). В рамках теории базовых эмоций приводятся тенденции для каждого из типов базовых эмоций. Так, страх подразумевает борьбу или замирание, тогда как гнев --- нападение.

Считается, что эти программы по большей части фиксированы и воспроизводятся автоматически, хотя определенная гибкость все же допустима. Это не удивительно, ведь она подразумевает большую адаптивность в сравнении с жестко фиксированными реакциями (Нессе 1990, Гриффитс 1997). То есть биологические детерминисты согласны, что окружающая среда влияет на формирование наших эмоций. Экман, например, считает, что выражения лиц, соответствующие эмоциям, культурно обусловлены. Существует исследование, показывающее, что выражения лиц японцев и американцев, смотрящих фильм о хирургических вмешательствах (провоцирующий отвращение) отличались в зависимости от публичности: японцы оказались менее склонны к выражению эмоций в присутствии других людей, и, значит, выражения лиц, вызываемые эмоциями, зависят от социальных факторов (Экман и Фризен, 1971).

Биологические детерминисты обычно признают связь триггеров базовых эмоций как с окружающей средой, так и с культурой. Ведь, например, то, что наши далекие предки боялись --- мамонты и дикие животные, очевидно, отличается от того, чего боимся мы. Нам зато приходится иметь дело с неосторожными водителями, ядерным оружием и изменением климата. И все же --- утверждают биологические детерминисты --- каждая базовая эмоция зависит от механизма оценки, настроенного на определенную небольшую группу стимулов: младенцы, например, боятся громких звуков (Дамасио 1994). Поскольку люди склонны обучаться реакциям на определенные условия, к врожденным стимулам легко добавляется ряд вызывающих условий (Охман и Минека, 2001). В результате мы легче научаемся (и труднее разучаемся) связывать со страхом пауков и высоту, чем, например, электрические розетки и оружие. Точно так же и макаки-резусы легко учатся бояться змей, наблюдая за соответствующей на них реакцией у сородичей, но то же самое не срабатывает для цветов или игрушечного кролика (Охман и Минека 2001).

Давайте рассмотрим свидетельства в пользу теории базовых эмоций.

\section{Свидетельства в пользу теории базовых эмоций}

В пользу существования панкультурных эмоций наиболее явно свидетельствуют исследования выражений лица. Те же исследования, что касаются неврологии и генетической наследственности, все еще спорны (см. Принц 2004, глава 5; Баррет 2006). В известном исследовании Экман и его сотрудник Уоллес Фризен (1971) изучали выражения лица, встречающиеся у народа Форе --- дописьменных людей, живущих в высокогорьях Папуа-Новой Гвинеи. Испытуемым --- среди них были 189 взрослых и 130 детей --- рассказывали разные провоцирующие эмоции истории --- о смерти ребенка или встрече с другом, например --- и предлагали выбрать одну из трех фотографий лиц актеров, имитирующих эмоции (ту которая наиболее соответствует истории). Оказалось, что выбор испытуемых из народа Форе незначительно отличается от выбора испытуемых американцев. Испытуемые из народа Форе, впрочем, иногда путали выражение удивления и выражение страха: немногим меньше половины из взрослых выбрали выражение страха, когда среди трех вариантов было и выражение удивления, но при этом 80\% из них сделали ''правильный'' выбор, когда среди вариантов не было выражения удивления. 68\% выбрали ''правильные'' выражения для историй, которым соответствовало удивление, 79\% --- которым соответствовала печаль, и 92\% --- счастье. На основании этого и ряда подобных исследований Экманом был сделан вывод, что выражения базовых эмоций на лице --- панкультурны и, значит, сами эти эмоции тоже. Таким образом, эти результаты, кажется, указывают на врожденный характер базовых эмоций, потому что эта врожденность лучше всего объясняет полученные данные.

Эти выводы, впрочем, активно критиковались, и большая часть критики была направлена на используемую в исследованиях методологию. Психолог Джеймс Рассел (1994; см. также Гендрон, Кривелли и Баррет 2018), например, указывал на основанность исследования на принудительном выборе: испытуемому приходилось выбирать между набором из всего лишь трех выражений лица, что, конечно, существенно повышало вероятность выбора ''правильного'' из них, и учитывая такую ограниченность, следовало бы ожидать гораздо более высокого процента согласия, чем тот, что был получен. И действительно, исследования, в которых испытуемым предлагалось выбрать любой [используемый для указания на эмоцию] термин для описания выражения лица, показали гораздо меньшую межкультурную корреляцию.

В ответ на возражения Экман (1994) указывал, что корреляция эта все еще куда выше, чем если бы эмоции и выражения лица сопоставлялись случайно --- и это, конечно, так, ряд эмоций обычно сопровождается характерной мимикой. Но вывод, который из этого можно сделать, гораздо слабее нужного для обоснования теории базовых эмоций. Джесси Принц поясняет (2004; см. также Райзенцейн, Штудтманн и Хорстманн 2013), что из корреляции между эмоциями и выражениями лица следует лишь частичная ограниченность развития эмоций специфичными для них нервными структурами, и что это утверждение куда слабее требуемого теорией базовых эмоций.

Другая критика направлена на понимании функции мимики. Приверженцы теории базовых эмоций предполагают, что функция улыбки, например --- показать эмоцию счастья. Но почему естественный отбор давал предпочтение такому автоматическому выражению эмоций? Так, нам --- и, значит, нашим генам --- только вредит демонстрация страха своим врагам. Выражения лица различаются сильнее, чем должны бы согласно теории базовых эмоций: олимпийские спортсмены, например, не склонны улыбаться сразу после завоевания золотой медали, и делают это только когда уже стоят на подиуме и смотрят на публику (Фернандес-Долс и Руис-Бельда, 1995). По мнению Алана Фридланда (1994), это указывает, что улыбка не обязательно выражает счастье, и сходства в выражениях лиц можно объяснить социальным обучением и подражанием.

Экман в ответ утверждал (1997), что социальное одобрение конечно может быть источником счастья, и потому спортсмен его на подиуме и выражает. Главное, что выражение лица может служить своего рода сообщением о намерениях, ведь тесно связано с переживаниями и, значит, мотивацией. И если бы улыбка не понималась уже заранее как выражающая счастье, то было бы загадкой, что чувствуют спортсмены на подиуме. Это, впрочем, разумный, но недостаточный для защиты теории базовых эмоций, ответ. Дело в том, что такое объяснение совместимо с в том числе культурной обусловленностью выражений лиц --- такая двойственная обусловленность не исключает связи между переживаемым и выражаемым и коммуникативную функцию мимики.

Защитники теории базовых эмоций предполагают, что возникновение эмоции каждого из типов сопровождается определенным нейронным паттерном и телесными изменениями на уровне вегетативной нервной системы. Экман и его коллеги (1983) пытались доказать наличие обратной связи --- способность выражения лица вызывать соответствующие телесные изменения. Они дали испытуемым инструкции по формированию нужных выражений лиц и просили их каждое такое выражение удерживать в течение десяти секунд. И действительно, такое подражание приводило к телесным изменениям: выражения гнева, страха и печали коррелировали с более быстрым сердцебиением, чем выражения счастья, удивления и отвращения, а выражения страха и печали сопровождались понижением температуры пальцев. Подобные результаты были подтверждены и в других исследованиях, и метаанализ показал, что влияние такой обратной связи стабильно, хотя и невелико (Коулз, Ларсен и Ленч, 2019).

Какой же вывод следует сделать из этих исследований? Как отмечает Лиза Фельдман Баррет (2017), исследование Экмана не показывает наличие уникальных телесных паттернов для каждого вида эмоций. Баррет утверждает, что сравнение исследований, вызывающих эмоции с помощью фильмов ужасов, например, показывает отсутствие корреляции видов эмоций с телесными паттернами: разные такие исследования смогли выявить разные закономерности.

Вопрос, конечно, в том, насколько большого разнообразия здесь можно ожидать. Страх, например, бывает разный: паника от неминуемого падения и тревога по поводу надвигающейся опасности сопровождаются совершенно разными, зависящими от контекста, реакциями, включая борьбу, замирание и бегство (Нессе 1990, Принц 2004). Столь разные реакции подразумевают в качестве способствующих совершенно разные телесные изменения, так что не удивительно, что страх может сопровождаться разными такими изменениями. И классификация с помощью машинного обучения показывает наличие отдельных паттернов для таких эмоций как удивление, страх, гнев и печаль (Крагель и ЛаБар 2013; но см. также Сигел и др. 2018).

То же верно и касательно исследований, пытающихся показать определенные закономерности работы мозга при переживании эмоции. Миндалевидное тело долгое время считалось структурой, отвечающей за страх (Леодукс 1996), но, как показал метаанализ ученицы Баррета Кристен Линдквист (2012), она активна и при переживании гнева, отвращения, печали и даже счастья. Имеющиеся данные, значит, показывают отсутствие внятных корреляций базовых эмоций с определенными паттернами мозговой активности.

Возможно, для каждой эмоции следует искать не единый паттерн, а лишь закономерности паттернов. Но даже при таком подходе результаты исследований неоднозначны. Методы классификации паттернов использовались, как и в случае с телесными паттернами, для демонстрации существования таковых для основных эмоций, таких как удивление, страх, гнев и печаль (Крагель и ЛаБар 2016), Но, как указывает Диас 2019, вывод об обоснованности биологического детерминизма делать все еще преждевременно. Дело в том, что наличие таких паттернов согласуется с возможным не только поверхностным культурным влиянием, но и таким влиянием на саму сущность тех или иных эмоций.

Еще один источник свидетельств в пользу теории базовых эмоций --- исследование эмоций [нечеловеческих] животных и человеческих младенцев: человеческие младенцы --- даже слепые --- могут выражать мимикой такие эмоции, как страх или счастье (Мацумото и Уиллингем, 2009). Более того, мимика животных и людей при переживании эмоций очень схожа. Еще Дарвин (1872) отмечал, что шерсть и животных и волоски у людей встают дыбом при переживании ими страха. Переживание страха людьми и животными характеризуется также тремя возможными реакциями --- борьба, бегство или замирание --- причем последняя, особенно у животных, включает существенное замедление сердечной деятельности и дыхания, а также снижение температуры тела и близкий к смерти паралич (сознание, однако, сохраняется). Такая тоническая неподвижность, встречающаяся у людей при столкновении с серьезной опасностью, наблюдалась также у ряда млекопитающих, птиц, рыб и насекомых (Арчер, 1979).

Но даже преемственность между выражениями лица и поведением животных и людей не доказывает врожденный характер базовых эмоций. Потому что, во-первых, лица младенцев, возможно, функционируют не так же, как у взрослых --- выражение удивления наблюдается у младенцев, рассматривающих знакомые предметы (Камрас 1994). Да и выражаемое мимика и поведение животными может плохо соотноситься с выражаемым схожей мимикой и поведением у людей.

Итак, как отмечает Принц (2004), биология оказывает значительное влияние на наши эмоции, но ничто не исключает столь же значительное влияние культуры, и вполне возможно, что даже казалось бы фундаментальные свойства эмоций подвержены культурному влиянию и формируются в существенной степени под действием него.

Давайте теперь перейдем к аргументам со стороны эволюционных психологов.

\section{Свидетельства из эволюционной психологии}

Программа эволюционных психологов, как уже было сказано, более амбициозна: они указывают на определенность эволюцией не только базовых эмоций, но и всех остальных --- в том числе тех, что уникальны для людей. Так, Джон Теби и Леда Космидес пишут о следствиях адаптивности мышления: ''каждое функционально отличающееся эмоциональное состояние --- страх перед хищниками, вина, сексуальная ревность, ярость, горе и т.д. --- будет соответствовать интегрированному режиму работы, предназначенному для получения преимущества при повторяющейся ситуации/условии, которому эта эмоция соответствует'' (2000, 101). Даже если Туби и Космидес симпатизируют теории базовых эмоций, их подход все же направлен на эмоции, интегрированные в когнитивные процессы и запланированные действия.

Этот подход хорошо проиллюстрирован в описании романтической любви экономистом Робертом Фрэнком (1988). Согласно Фрэнку, эмоции --- это разрушительные мотивации, которые, однако, помогают нам преодолевать проблемы сотрудничества и этим служат нашим интересам. Дело в том, что эмоции побуждают нас совершать казалось бы иррациональные --- с точки зрения личной выгоды --- действия, которые в долгосрочной перспективе оказываются более выгодными. Фрэнк видит романтическую любовь природным решением так называемой ''проблемы обязательств'': мы никогда не можем быть уверены, что потенциальный партнер --- лучший из тех, кого мы сможем найти, и, значит, предпочли бы не вступать в отношения вовсе --- и не передали бы свои гены --- если бы не иррациональные порывы. То есть склонные к романтической любви и соответствующей ей иллюзии идеальности потенциального партнера, распространяли свои гены куда успешней, и потому многие к ней склонны. Таким образом Фрэнк приходит к выводу о врожденном характере сентименту романтической любви. Аналогичную точку зрения поддерживает и Хелен Фишер (2004), указывая на фундаментальность романтической любви как стремления, развившегося благодаря той экономии времени и сил при выборе партнера, которую оно дает.

Проблема эволюционного подхода Фрэнка в том, что вывод об эволюционной природе любви делается на основе ее полезности. Пол Гриффис (1997) указывает на необоснованность такого вывода --- объяснить, как некоторая полезная особенность могла возникнуть в результате естественного отбора так же легко, как трудно объяснить, почему она на самом деле так возникла. То есть теории на основе эволюционных адаптаций поддерживаются лишь своей простотой: лучшее объяснение наличия у нас полезной особенности --- ее возникновение в результате естественного отбора. Но альтернативные объяснения часто оказываются столь же удачными. Эволюционные объяснения, значит (так называемые ''потому что так'') могут объяснить некоторую особенность только если уже известно, что она возникла благодаря естественному отбору. Но пока существуют альтернативные объяснения, это не работает.

У подхода Фрэнка есть и другие проблемы. Может показаться, что он принижает любовь, сводя ее к лишь полезному средству, но гораздо важнее, что, как отмечает Джесси Принц (2004), Фрэнк предполагает западную концепцию любви. Что любовь подразумевает долгосрочные обязательства --- не общепризнанно: в китайской культуре любовь ассоциируется с кратковременными связями, тогда как в браке поощряются доверие и дружба. Возможно, западная концепция [романтической] любви --- лишь одно, и, пожалуй, специфичное, решение проблемы обязательств, но нет никаких оснований полагать, что оно единственно и является продуктом естественного отбора.

Эволюционные психологи испытывают трудности с доказательством, что их объяснения --- наилучшие, ведь во многих случаях возможны по крайней мере столь же подходящие и простые объяснения. Например, Принц (2004) рассматривает ревность и несколько альтернативных объяснений наличия у нас такой эмоциональной диспозиции. Эволюционные психологи объясняют распространенность ревности и беспокойства об эмоциональной измене у женщин тем, что их партнеры помогают им воспитывать детей и, значит, распространять их гены. Мужчина же, напротив, не мог быть уверен, что именно он отец ребенка, и потому мужчин больше беспокоит сексуальная неверность. Но все то же можно объяснить, например, экономическими различиями: в большинстве культур мужчины обладают большей экономической властью, и потому эмоциональная неверность дороже обходится женщинам, а значит при прочих равных мужчины предпочтут эмоциональную измену сексуальной.

Эволюционное объяснение каждой из эмоций надо сравнивать с альтернативными, чтобы доказать его превосходство, и это достаточно сложно и затратно. Эволюционный подход, к тому же, вынужден предполагать тесную корреляцию эмоций с классами вызывающих их стимулов: то, что вызывает страх или ревность сейчас должно быть примерно тем же, что и в плейстоцене, например. Но эмоции, утверждает Гриффс (1997), не воплощают такое определенное знание. Эволюционная теория, значит, должна учитывать, что даже если механизмы оценки определенным образом настроены на небольшую группу стимулов, эти стимулы могут меняться в ходе жизни.

Эволюционная психология, конечно, допускает незначительные межкультурные различия эмоциональных реакций, но, похоже, они различаются куда существенней. Как указывает Принц (2004), такая эмоция как ревность сильно варьируется от культуры к культуре: люди Тода --- небольшой этнической группы в горах Нилгири в Индии --- практикуют братскую полиандрию, когда женщина выходит замуж за своих братьев, и почти не проявляют сексуальную ревность внутри групп (Гупка 1991). Также Самоанцы, похоже, живут жизнью, свободной от гнева и такой ревности (Мид 1928). Подобные примеры являются свидетельствами в пользу социального конструктивизма, к рассмотрению которого мы сейчас и перейдем.

\section{Социальный конструктивизм}

Социальный конструктивизм следует отличать от другого вида конструктивизма, а именно --- психологического, теории, которую защищает Джеймс Рассел (2003; см. также Мурс 2017). Согласно Расселу, эмоции состоят из специфических психологических ''строительных блоков'', и базовым таким блоком является центральный аффект --- состояние, характеризуемое, во-первых, гедонистической ценностью и варьирующееся, значит, от крайней приятности до крайней же неприятности, а во-вторых, уровнем возбуждения, варьирующемся от отсутствия активности, что соответствует сну, до неистовства. Другие такие блоки --- оценка, мимика, изменения в вегетативной нервной системе и действия. Такой подход явно отличается от социального конструктивизма, потому что механизмы, реализующие составные блоки эмоций --- общие для всех людей. В рамках разных обществ возможны, конечно, разные концептуализации блоков, но функционирование блоков независимо от таких концептуализаций.

Социальные конструктивисты, чья теория возникла в 1980-х годах выступают против обеих версий биологического детерминизма, утверждая, что эмоции допускают куда большую изменчивость. Эмоции являются социальными конструктами, продуктом социальных практик, так же как деньги и социальные роли. Главные сторонники этого подхода --- психолог Джеймс Аверилл (1980), антрополог Кэтрин Лутц (1988) и философы Ром Харре (1986) и Клэр Армон-Джонс (1985, 1986). Социальный конструктивизм касательно эмоций до сих пор релевантен, на что указывает недавняя работа Миртина Арангурена (2017), утверждающего, что социальному конструктивизму не обязательно предполагать зависимость эмоций от языка, ведь и культура не обязательно зависит от языка, а значит, эмоции животных также могут быть культурно сконструированы.

Есть две основные версии социального конструктивизма. Согласно социально-ролевому конструктивизму, эмоции соответствуют [временным] социальным ролям, определяющим, как человек должен вести себя в определенных сценариях (Аверилл 1980). Программа такого поведения состоит из нормативных убеждений и ожиданий, разделяемых обществом, благодаря которым принятие определенной социальной роли приводит к соответствующему ''эмоциональному синдрому'', включающему оценки, поведение и переживания. Часть реакций, привносимых синдромом, биологически детерминированы, но считаются эмоциями только в случае их интеграции в социальную роль. Ключевое отличие от роли в обычном (не метафорическом) понимании --- в пассивном переживании эмоций, соответствующих социальной роли, как если бы эти реакции были вне контроля. Такой глубокий самообман --- часть общей стратегии, эффективно способствующей выживанию в обществе. Этот подход напоминает предложенный философом-экзистенциалистом Жаном-Полем Сартром (1962), согласно которому эмоции, вообще говоря, не более чем самообман, ''магические'' действия, совершаемые для изменения мира. Так, чтобы удалить объект страха, мы можем потерять сознание.

Теория социальных ролей допускает огромную межкультурную вариативность. Согласно правилам нашего общества, гнев соответствует ситуации, в которой происходит оскорбление, то есть если вас оскорбляют, вы должны злиться в ответ. Содержание таких правил, значит, зависит от того, что в социальной группе считается неким правонарушением. В Главе 1 рассматривалась культура утку, людей, избегающих гнева, усвоив правило, что существует очень мало того, на что следует злиться (Бриггс 1970).

Разыгрывание роли может включать концептуализацию себя как испытывающего соответствующую эмоцию. Так, Клэр Армон-Джонс (1986) и Ром Харре (1986) защищают версию социального конструктивизма, согласно которой переживание гнева, например, требует усвоения концепции гнева и применения ее к своему состоянию. То есть ''нельзя сказать, что некто чувствует эмоцию E, если он не усвоил концепцию E и не может применить эту концепцию к собственному опыту в соответствующих ситуациях'' (1986, 37; см. также Баррет 2017).

Согласно этой версии, то, что считается гневом в одном обществе, может не считаться им в другом, ведь не усвоивший --- в силу социальных причин --- понятие гнева, не может испытывать эту эмоцию. Нечеловеческие животные, не способные на концептуальное мышление, значит, не способны и переживать эмоции. Правила, регулирующие применение эмоциональных концепций, могут --- но не обязательно --- ссылаться и на телесные изменения. Таким образом, обе версии конструктивизма утверждают отсутствие у эмоций какой-либо сущности, потому что ни одно из конкретных применений эмоциональных концепций не неизбежно.

Давайте теперь рассмотрим свидетельства в пользу этих теорий.

\section{Свидетельства в пользу социального конструктивизма}

Стэнли Шахтер и Джером Сингер (1962; см. также Линдквист и Баррет 2008) провели исследование для проверки концептуальной версии социального конструктивизма (согласно которой, эмоции включают не-специфическое физиологическое возбуждение и некий когнитивный ярлык, приписываемый себе переживающим эмоции). В исследовании испытуемым вводили адреналин --- вещество, вызывающее сильную вегетативную реакцию: учащенное дыхание, повышение артериального давления и частоты сердечных сокращений. Всем испытуемым было сказано, что вводимое вещество --- витаминная добавка для проверки зрения, но некоторым из них (''адреналин-информированной'' группе) сообщили также о возможных побочных эффектах. После инъекции участников помещали в комнату ожидания, и там с ними взаимодействовал, демонстрируя либо гнев, либо эйфорию, подставной человек-провокатор.

Влияние марионетки оказалось разным в зависимости от информации, предоставленной участникам. Так, испытуемые из адреналин-неинформированной группы реагировали отрицательно на гнев и положительно --- весельем, например --- на эйфорию, тогда как испытуемые из адреналин-информированной группы почти не реагировали ни на гнев, ни на эйфорию. То есть при примерно одинаковых раздражителях участники по-разному интерпретировали свои состояния и, следовательно, испытывали разные эмоции.

Но дело в том, что ничто не исключает ошибочного приписывания себе соответствующих эмоций испытуемыми из адреналин-неинформированной группы, в результате которого они хотя и вели себя так, будто их испытывают. Хорошо известно, отмечает Гриффтс (1997), что мы склонны к таким конфабуляциям (выдумкам), и эти испытуемые, возможно, так приписывали себе эмоции, не испытывая их. Еще одна проблема, подчеркиваемая Принцем (2004) --- сообщение всеми участниками о себе как о счастливых в анкете, заполняемой в конце эксперимента --- из-за чего неясно, было ли поведение, соответствующее гневу, действительным его проявлением. Так что данное исследование нельзя считать определяющим.

Обе версии социального конструктивизма объясняют существенные межкультурные различия в эмоциональных реакциях. Антропологические и кросс-культурные лингвистические исследования показывают также трудность перевода терминов, обозначающих некоторые эмоции, на другие языки --- что, кажется, говорит о специфичности этих эмоций для определенных культур. Примером такой предположительно культурно-специфичной эмоции является японская амаэ --- приятное чувство зависимости от доброй воли и привязанности другого человека, испытываемое, например, детьми рядом с родителями (Дои 1973). В английском языке нет эквивалентного термина, так что, похоже, амаэ нельзя концептуализировать, не зная японский, а значит, его нельзя и испытать. Еще один пример --- то, что жители Ифалуке, микронезийского кораллового острова в Тихом океане, называют \textit{сонг}: своего рода гнев, имеющий моралистический подтекст и не подразумевающий жажды мести (Лутц 1988) --- английского эквивалента этому термину нет, да и гнев в нашей культуре обыкновенно предполагает желание отомстить. Наконец, русское слово ''тоска'' также не имеет строгого английского эквивалента, потому что включает множество аспектов, отчего может переводиться как ''anguish'', ''pangs'', ''depression'', ''ennui'', ''boredom'', ''longing'', ''yearning'', ''nostalgia'' и ''home sickness'' (Вежбицка 1992, 1999).

Эти примеры показывают, что есть эмоции, незнакомые нашей культуре. Разнообразие эмоциональных реакций огромно. Так, народ Ифалук, который изучал Лутц (1988), использует кроме ''сонг'' множество других эмоциональных терминов, не имеющих английского эквивалента: баю (романтическая любовь/счастье), кер (счастье/волнение), фаго (любовь/сострадание/печаль), метагу (страх/тревога), рус (паника/испуг/удивление), тан (горе/разочарование), люгумент (дискомфорт/вина) и босу (ревность/возбуждение) --- эмоциональный ландшафт ифалуков существенно отличается от знакомого нам, так что социальный конструктивизм кажется вполне правдоподобной позицией.

Но поддерживают ли эти свидетельства социальный конструктивизм касательно всякой эмоции, или только касательно некоторого их подмножества? Ведь у народа ифалук есть и вполне прямо переводящиеся на английский эмоциональные термины: ниябут (отвращение), бобо (разочарование) и малевелу (спокойствие), например (Лутц 1988). То есть по крайней мере некоторые эмоции все же могут быть межкультурными, а значит свидетельства эти поддерживают лишь ограниченную версию социального конструктивизма.

Даже тогда, когда однозначный перевод эмоционального термина отсутствует, это может означать лишь разницу в словарном запасе, так что различия могут оказаться не в переживаемых эмоциях, но в лишь возможностях их описания на определенных языках (Принц 2004). Похоже, что амаэ, например, можно испытывать и не зная соответствующего термина, который однозначно бы на него указал, испытывая, значит, просто приятные чувства по отношению к людям, от которых зависим.

Эта разница между западной и японской культурами, возможно, состоит в значимости, придаваемой этой эмоции, поскольку амаэ считается одним из центральных элементов японской культуры. Такой акцент на амаэ резко контрастирует с западными акцентами на независимости и автономии. Здесь антрополог Роберт Леви (1984) предлагает различать \textit{гиперосознанные} и \textit{гипоосознанные} эмоции: те, что \textit{гиперосознаны} в рамках культуры --- подчеркиваются и выделяются, тогда как \textit{гипоосознанные} --- почти игнорируются, обесцениваются или даже подавляются. То есть эпизоды \textit{амаэ} \textit{гиперосознанны} в японской культуре и \textit{гипоосознанны} --- в западной. Впрочем, игнорирование, обесценивание и подавление определенного вида эмоций, конечно, не исключает их переживание.

Итак, терминологические различия не обязательно влекут различия на уровне эмоций, поэтому данных антропологии и лингвистики недостаточно для подтверждения социально-конструктивной природы эмоций. Однако конструктивисты, особенно те, что говорят о социальных ролях, в ответ могут указать, что эпизоды эмоций связаны с множеством общепринятых правил.

Так, романтическая любовь, какой мы ее знаем, заключается в поведении, соответствующем набору культурно-специфичных правил, имеющих свою историю возникновения и развития (Левис 1936, Аверилл 1985, Билл и Штернберг 1995). Романтическая любовь подразумевает участие в ней строго двух людей --- мужчины и женщины, а также исключительность, из-за чего предполагает ревность. Но ни одно из этих правил не является панкультурным. В викторианской Англии, например, влюбленность подразумевала восхищение, а сексуальное влечение считалось неуместным и даже постыдным, тогда как в современности любовь рассматривается как в том числе интимная привязанность и по определению включает сексуальное влечение. Также и далеко не общепринято современное убеждение, что романтическая любовь --- подходящая основа для брака: множество культур рассматривают ее как даже угрозу долгосрочным отношениям. Сценарии романтической любви, таким образом, варьируются от одного культурного контекста к другому, и то же можно сказать и о других эмоциях, считаемых базовыми. Так, сценарий, допускающий гнев для утку, например, существенно отличается от знакомого нам.

Но значат ли межкультурные различия в таких сценариях --- социальную сконструированность эмоций. Ведь наличие норм в отношении чего-либо еще не означает, что оно социально сконструированно. Есть, например, социальные и моральные нормы касательно причинения боли, но сам акт причинения боли --- не социальный конструкт. Точно так же и существование норм касательно любви, не значит, что это социальный конструкт. Эти нормы могут влиять на переживание и выражение любви, но необязательно, ведь нормы могут нарушаться. Так, один из утку говорил Бриггс: ''гнев утку --- нечто памятное --- но каплуна (белый американец) может злиться с утра до полудня'' (1970, 261).

Итак, следует ли заключить, что есть некоторое ядро эмоций, полностью изолированное от культурного влияния?

\section{Пластичность эмоций}

Как мы увидели, ни аргументы в пользу биологического детерминизма, ни аргументы в пользу социального конструктивизма --- не являются вполне убедительными. Возможно, истина где-то посередине. Эмоциональные диспозиции и эпизоды, похоже, формируются как биологическими, так и культурными факторами. Чтобы понять, как они взаимодействуют, рассмотрим участие культуры в формировании наших эмоций.

Примечательно, что последователи каждой из рассмотренных теорий признают пластичность эмоциональных диспозиций. Через обучение правильным причинам страха или гнева культура частично определяет то, на что будут возникать соответствующие эмоциональные реакции --- и этот тезис принимают даже сторонники теории базовых эмоций. Как именно это происходит --- уже другой вопрос, который мы рассмотрим в Главе 12. Споров не вызывает и то, что культурно-специфичные правила проявления эмоций влияют на их выражение с помощью мимики и телодвижений. В плейстоцене, например, нельзя было взять мобильный телефон чтобы позвать на помощь. Культура также может поощрять определенные типы мотивации через распространение примеров конкретных эмоциональных действий. Культура, значит, формирует причины и следствия наших эмоций.

Были выявлены два механизма, играющие важную роль в формировании наших эмоций. Первый, предложенный Робертом Плутчиком (1980; см. также Принц 2004) --- смешение, то есть процесс объединения двух или более эмоциональных реакций в одну, разделяющую аспекты каждой из исходных. Можно, конечно, согласиться с Гриффтсом (1997) в том, что \textit{не все} небазовые типы эмоций возникают в результате смешения базовых, но некоторые, возможно являются именно таковыми. Фаго, например, сочетает в себе любовь, сострадание и печаль, а тоска представляет смесь печали, страха и скуки. Культура, значит, может если не формировать такие эмоции-смеси --- амаэ, фаго или тоску можно, пожалуй, переживать и не имея соответствующих терминов --- то по крайней мере делать их легко доступными.

Второй же механизм формирования эмоций --- когнитивный. Эмоция может объединяться с такими когнитивными компонентами, как убеждения, приводя к появлению нового типа эмоции. Именно с помощью этого механизма небазовые эмоции формируются из базовых (Джонсон-Лэрд и Оатли, 2000), то есть каждая небазовая эмоция состоит из базовой и некоторого убеждения. Амаэ, например, подразумевает радость \textit{от мысли о зависимости}, негодование --- гнев \textit{из убеждения в несправедливости}, вина --- стресс \textit{из убеждения, что сделал что-то неправильно}. Согласно Джонсон-Лэрду и Оатли (2000), некоторые когнитивные процессы еще более сложны: романтическая любовь, например, включает сексуальное желание и счастье в сочетании с идеализацией любимого, альтруизмом и тоской в разлуке.

Но всегда ли когнитивный компонент состоит в том числе в наличии убеждения? Джонсон-Лэрд и Оатли (2000) считают, что небазовые эмоции всегда включают когнитивный компонент. Принц (2004) же указывает, что когнитивный компонент нужен лишь для приобретения предрасположенности испытывать такие эмоции, а как только процесс обучения завершится, эту эмоцию можно испытывать независимо от убеждения, которое было частью когнитивного компонента. И действительно, такие эмоции, как негодование и вина могут даже сопротивляться убеждениям: можно чувствовать вину и в то же время быть убежденным, что не сделал ничего плохого. Представьте, что поднимаетесь на Эверест и на вашу группу обрушивается лавина, и вы оказываетесь единственным выжившим. В такой ситуации вполне можно почувствовать так называемую вину выжившего даже понимая, что предотвратить смерть товарищей было невозможно. К таким случаям мы еще вернемся в Главе 6.

Итак, культура играет центральную роль в когнитивном механизме формирования эмоций. Это верно еще и потому, что многие убеждения, используемые этим механизмом, в значительной степени именно культурой формируются. Так, понятие несправедливости мы усваиваем в рамках социализации, и то, что ей считается, варьируется от культуры к культуре. Культура также способствует образованию определенных комбинаций эмоций и убеждений. Так, культура, в которой концепции несправедливости нет или она не играет значимой роли, вряд ли будет содержать в своем эмоциональном репертуаре негодование.

Биологические детерминисты, впрочем, могут согласиться со всеми этими пунктами и все же настаивать, что базовые эмоции не получены в результате смешения или когнитивной разработки, а --- напротив --- являются исходным материалом для этих механизмов. Культурные влияния, значит, лишь поверхностны, поскольку касаются только причин и следствий базовых эмоций. Но дело в том, что такие влияния могут быть отнюдь не поверхностны. Эмоции не обязательно подобны авокадо --- с мягкой культурной массой, организованной вокруг твердого природного ядра (Голди 2000; см. также Фаушер и Тапполет 2008).

Независимо от того, влияет ли культура только на причины и следствия базовых эмоций или формирует все эмоции, ее влияние может быть значительно. Формируя объекты страха и отвращения, культура вносит крайне существенный вклад в наши переживания. Эмоциональные реакции также в значительной степени формируются культурой, ведь она определяет, что мы мотивированы делать, испытывая определенные эмоции. Культура определяет наши эмоциональные диспозиции и соответствующие эмоции. И эта пластичность может быть полезна, ведь допускает воспитание эмоций --- тема, к которой мы еще вернемся в Главе 12.

\section{Подведение итогов}

Мы рассмотрели аргументы в пользу биологического детерминизма и социального конструктивизма. Мы пришли к выводу, что эмоции --- продукт как природы (биологии), так и культуры. Например, склонность чувствовать страх развилась, потому что оказалась полезной для выживания, но \textit{то, чего мы боимся}, \textit{как мы выражаем} страх и что при чувстве страха \textit{мотивированы делать} во многом зависит от культуры. Биологическим детерминистам кажется корректным утверждать, что культурные влияния лишь поверхностны, потому что затрагивают только причины и следствия эмоций, но не сами эмоции. Но, похоже, эмоции представляют собой цельный продукт взаимодействия нашей биологии и культуры. Вопрос, который следует рассмотреть --- что считать сущностью эмоций, если не их причины и следствия.

\section{Краткое содержание}

\begin{itemize}
  \item\ Биологический детерминизм в отношении эмоций предполагает, что по крайней мере некоторые их виды зависят от панкультурных психологических черт, закрепленных в нашей генетической структуре и являющимися эволюционными адаптациями, позволившими нашим предкам наиболее успешно распространять эти свои гены.
  \item\ Согласно базовой теории эмоций, только некоторые их виды детерминированы биологически. Такие базовые эмоции состоят из рефлексивных каскадов скоординированных действий, строго определенных отдельными нейронными связями (hardwired neural circuits).
  \item\ Эволюционные психологи --- напротив --- считают все эмоции следствием таких адаптаций.
  \item\ Свидетельства, приводимые приверженцами теории базовых эмоций, включают исследования мимики, нейронных структур и телесных изменений, а также сходства эмоций взрослых, младенцев и нечеловеческих животных, но показывают --- в лучшем случае --- что биология значительно влияет на наши эмоции.
  \item\ Свидетельства эволюционных психологов также неубедительны: такие эмоции, как любовь и вина не могут, кажется, быть такими адаптациями и в то же время различаться от культуры к культуре.
  \item\ Социальные конструктивисты утверждают, что эмоции --- социальные конструкты, потому что являются либо способом исполнения социальных ролей, либо результатом применения к себе определенной эмоциональной концепции.
  \item\ Свидетельства в пользу концептуальной версии конструктивизма, полученные из исследования Шахтера и Сингера, подвергаются критике в том числе за то, что не исключают надуманность и недействительность предполагаемых эмоций испытуемых.
  \item\ Свидетельства от межкультурной изменчивости впечатляют, но, возможно, указывают лишь на различный словарный запас в разных культурах.
  \item\ Свидетельства в пользу конструктивизма, основанные на определенности эмоций распространенными в обществе правилами так же не вполне убедительны, ведь правила могут игнорироваться.
  \item\ Наши эмоциональные диспозиции формируются и биологическими, и культурными факторами.
\end{itemize}

\begin{tcolorbox}
\section{Учебные вопросы}
  \begin{enumerate}
    \item\ Является ли любовь социальным конструктом, смесью или когнитивной разработкой?
    \item\ Могут ли животные переживать эмоции? Если да, то какое значение это имеет в дебатах между биологическим детерминизмом и социальным конструктивизмом?
    \item\ Возможно ли испытывать эмоцию, не концептуализируя ее? Чувствовали ли вы когда-нибудь амаэ или сонг?
  \end{enumerate}
\end{tcolorbox}

\part{Теории эмоций}

\chapter{Чувственные теории}

\section{Введение}

Что лежит в основе эмоций? В результате обсуждения в Главе 2 выяснилось, что эпизоды таких эмоций, как страх или отвращение, обычно довольно кратковременны, и включают чувства, мотивации, интенциональные объекты, оценки и формальные объекты. Эта характеристика помогает не только подчеркнуть статус эмоций как аффективных состояний, но и отличить их от других таких состояний --- чувств и настроений, например. Всякая теория эмоций, пожалуй, должна объяснять эти свойства, а также ответить на вопрос сущности эмоций --- вопрос, на который приведенная характеристика не отвечает.

Следует установить, что из этой предварительной характеристики (и, возможно, не только из нее) существенно для эмоций, что мы и попытаемся сделать в трех следующих главах, рассмотрев \textit{чувственные}, \textit{мотивационные} и \textit{оценочные} теории эмоций.

Должны ли мы к этим трем теориям добавить также теории базовых эмоций и ролевую конструктивистскую теорию, например? Их, конечно, можно считать теориями эмоций, но их тезисы находятся на другом уровне, потому что совместимы с несколькими уже конкретными теориями эмоций. Например, можно придерживаться теории базовых эмоций и в то же время считать, что эмоции, по сути --- мотивации, возникающие в результате развития соответствующих заложенных в нашей генетической структуре механизмов в ходе естественного отбора. Так же можно понимать эмоции как конституированные чувствами, составляющими, значит, основу каскада программ аффектов, заложенных в ходе эволюции.

С теориями же социального конструктивизма все не так однозначно, но, в принципе, он совместим с каждой из трех теорий. Можно, например, утверждать, что оценки возможны благодаря определяемым культурой оценочным концепциям, и если сущность эмоций в оценках, то эмоции социально сконструированны. Не менее успешно можно указывать на сконструированность мотиваций или чувств, и, значит, сконструированность и эмоций, если они главным образом в них состоят. Впрочем, как мы видели в предыдущей главе, социальные конструктивисты склонны отрицать наличие у эмоций хоть какой-то [неизменной] сущности, тем самым отрицая, похоже, возможность разработки позитивной теории эмоций.

Мы начнем с чувственных теорий --- известного и достаточно интуитивного подхода. Чувственные теории, однако, подвергаются серьезной критике, и сегодня мало тех, кто бы их придерживался. Рассмотрение этих возражений позволит лучше понять, что такое эмоции. Прежде чем перейдем к изложению чувственных теорий, давайте взглянем на основные проблемы, с которыми сталкиваются все теории эмоций.

\section{Задачи теорий эмоций}

Всякая теория эмоций должна объяснять их сущность, в связи с чем сталкивается с как минимум двумя демаркационными проблемами. Во-первых, нужно показать, чем эмоции отличаются от таких состояний, как зуд или мигрень, например. Если же отличий, согласно теории, нет, то она должна, по крайней мере, объяснить интуитивность представления об их наличии. Во-вторых, следует выяснить, благодаря чему эмоции различны. Если же различий между ними нет, то такой неинтуитивный вывод объяснить.

Итак, теория эмоций должна учесть или объяснить наши интуитивные представления о них. В дополнение к интуициям о существовании и различности эмоций теория должна, значит, учесть интуиции об их особенностях, то есть учесть, что эмоцию обычно кратковременны, включают чувства и мотивации, а так же связаны с интенциональными и формальными объектами и их оценкой.

Создание такой теории, конечно, непростая задача. Тема эмоции, кажется, слишком запутанная, чтобы надеется на сколь нибудь стройную ее теоретизацию. Эмоции бывают самые разные, и если базовые и прочие эмоции в целом достаточно схожи чтобы говорить и о тех и о других как об именно эмоциях, то страх и, скажем, восхищение отличаются разительно. Даже одна и та же эмоция --- страх, например --- может существенно разниться в зависимости от объекта: бояться нападающего медведя совсем не то же, что бояться за Анну Каренину (см. Главу 12). Во всяком случае, будет продуктивнее начать с рассмотрения ограниченного числа эмоций, как мы сделали это в Главе 2, и уже затем проверить, распространяется ли полученная теория на прочие, более сложные и редкие, эмоции.

Эта стратегия, кстати, не определяет, должны ли какие-то эмоции быть прототипическими и должна ли концепция эмоций быть \textit{концепцией-прототипом} --- то есть должна ли принадлежность к ней зависеть от сходства с парадигматическими ее примерами (Фер и Рассел 1984; Бен-Зеев 2001). Поэтому то, что считается эмоцией, не имеет четких границ, а такие состояния как трепет или спокойствие слишком отличны от парадигматических примеров, чтобы считаться эмоциями.

Трудность разработки теории эмоций также в их сложности как явлений. Рассмотрим, например, эпизод печали. В аэропорту вы наблюдаете за самолетом, на котором возлюбленная улетает на год за границу, и когда он исчезает в облаках, вас охватывает сильная печаль. Даже в таком относительно простом примере задействовано несколько элементов. Во-первых --- вид самолета, уносящего возлюбленную и осознание, что в ближайшее время вы ее не увидите --- то есть вы оцениваете ситуацию как существенную потерю. Во-вторых --- ряд физиологических изменений: кровяное давление повышается, начинают течь слезы, лицо принимает типичное выражение: брови образуют перевернутую букву ''V'', нижняя губа подталкивает верхнюю губу, а уголки верхней опускаются вниз. Поза тела также изменяется: голова наклоняется, а плечи выгибаются вперед. Вы также чувствуете будто что-то разрывает вас на части. Проносятся мысли, что не стоило соглашаться на расставание и что так долго вы не выдержите, возникают воспоминания о похожих случаях из детства. Внимание в целом фокусируется на печальном и мир становится серым и пустым.

Итак, типичный эпизод печали, как и любой другой типичный эмоциональный эпизод, включает следующие компоненты:

a) информационный (осознание, что любимый человек долго не будет с вами)

б) оценочный (расставание понимается как потеря)

в) физиологический (повышение артериального давления и др.)

г) экспрессивный (мимика и поза тела)

д) феноменологический (ощущение, будто что-то разрывает изнутри)

е) когнитивный (навязчивые мысли и смещение фокуса внимания)

ж) мотивационный (стремление абстрагироваться)

Некоторые из них уже указывались в Главе 2: эмоции обычно включают чувства, мотивации и оценки. Теория эмоций должна определять, какой из компонентов является необходимым и достаточным для эмоций, какой составляет их сущность --- то, что Принц (2004) называет ''проблемой частей''.

Ответы предлагаются самые разные. Социальные конструктивисты, например, скептичны к наличию какой-либо сущности у эмоций. Так, Аверилл (1980) не считает необходимым ни один из компонентов так называемого им ''эмоционального синдрома'' --- и так же можно отрицать необходимость каждого из вышеприведенных, и не обязательно делать это в рамках конструктивизма. Например, можно говорить о концепции эмоций как держащейся на семейном сходстве. Семейное сходство --- понятие, введенное Людвигом Витгенштейном (1953) и указывающее на подобие предметов в категории как среди членов семьи, где встречается сходный цвет глаз, черты лица и темпераменты, но ни один из этих наборов черт не относится ко всем членам семьи сразу. То есть можно сказать, что нет характеризующего все эмоции, но есть лишь сходство между некоторыми из них, перечислить которые теоретик эмоций только и может. Джон Элстер (1999) так и делает: его описание эмоций состоит из списка характеристик, аналогичного приведенному выше списку компонентов, но ни один из них не считается необходимым. Но прежде, конечно, стоит рассмотреть более амбициозные подходы, и только если все они потерпят неудачу, будет оправдано взяться за разработку такой скромной версии.

Эти амбициозные подходы можно разделить на чистые и гибридные теории. Чистые подразумевают единственный компонент в качестве необходимого, остальные же считаются необязательными --- контингентными. Так, чувственная теория эмоций считает чувства таким необходимым компонентом, а мотивации, например --- контингентным. Гибридные же теории предполагают необходимыми несколько компонентов. Так, учение Аристотеля об эмоциях, рассмотренное в Главе 1, принято считать гибридной теорией, потому что эмоции, согласно нему, включают удовольствие/неудовольствие, оценку, мотивацию и телесные изменения. Любая гибридная теория, в отличие от чистой, должна объяснять, что удерживает эти необходимые компоненты вместе. Так, приверженцы теории базовых эмоций и гибридной теории, обращаются для ответа на этот вопрос к эволюционным подходам к эмоциям и их функциям.

В этой главе мы рассмотрим теории, понимающие эмоции как чувства. Эмоции, конечно, обычно обладают феноменальными свойствами, но защитники этих теорий утверждают, что только в чувствах эмоции и состоят. Чувства тоже можно понимать по-разному: разделяют обычно теории, имеющие в виду телесные чувства (Джеймс 1884) и теории, имеющие в виду что-то другое (Уайтинг 2009). Мы будем рассматривать чувственную теорию Джеймса, поскольку она наиболее влиятельна, но затронем также те, что отрицают телесность чувств.

\section{Чувственно-телесная теория эмоций Уильяма Джеймса}

Согласно теориям чувств, эмоции --- это особые сознательные переживания, то есть ощущения. Такие теории защищали видные философы. Так, принято считать, что Декарт (1649/1989) придерживался именно чувственной теории эмоций, утверждая, что они лишь восприятия, ощущения или волнения души, вызванные движениями животных духов --- мельчайших частиц, сообщающих душе о телесных изменениях. Юм (1739-1740/2000) схоже понимал эмоции --- как ощущения, вызванные переживанием или мыслью об удовольствии или боли.

Относительно современные обсуждения в этой области сосредоточены на чувственной теории, предложенной Уильямом Джеймсом (1884, 1890). Несмотря на многочисленные и довольно серьезные возражения, эта теория оказалась чрезвычайно влиятельной. Нейробиолог Антонио Дамасио (1994), которого считают нео-Джеймсианцем, использовал идеи Джеймса в построении своей теории эмоций. Интересно, что среди нео-Джеймсианцев есть и защитники мотивационных теорий --- Жюльен Деонна и Фабрис Терони (2012), и оценочных теорий --- Джесси Принц (2004) и Урия Кригель (2014).

Джеймс, подобно Декарту, подчеркивает роль телесных изменений, выделяя грубые эмоции --- такие как горе, страх, ярость, любовь --- очевидно связанные с телесными изменениями, и тонкие --- разные моральные, интеллектуальные и эстетические чувства --- ключевым ингредиентом которых, так же как и всех других эмоций, являются чувства. Эмоции, таким образом --- телесные чувства или ощущения особого рода.

Чувства, о которых говорит Джеймс, связаны с телесными изменениями, вызванными восприятием, воспоминанием или мыслью о чем-то ''захватывающем''. Эти изменения не требуют особой обработки мозгом и касаются внутренних органов, сердечно-сосудистой системы, внешних проявлений, мотиваций и даже прямых импульсов к действиям. Джеймс пишет: ''\textit{организм --- это резонатор}, в котором каждое изменение сознания, сколь бы незначительным не казалось, может оказать значительное на него влияние'' (1890, 450, курсив оригинала). Разнообразие связанных с эмоциями телесных изменений и их комбинаций позволяет, согласно Джеймсу, описать любой ''эмоциональный оттенок''.

Джеймс, значит, разворачивает привычную причинно-следственную связь, согласно которой первично восприятие какого-либо факта, и уже это восприятие ответственно за возникновение эмоции и, затем, телесных изменений: вы сначала видите приближающегося медведя, от этого у вас возникает страх, и от страха повышается пульс, например. Джеймс же утверждает, что восприятие напрямую вызывает телесные изменения: ''Моя теория [...] в том, что \textit{телесные изменения следуют непосредственно за восприятием захватывающего факта, и что наше ощущение этих изменений --- И ЕСТЬ эмоция}'' (1890, 450, курсив и заглавные буквы оригинала, см. также 1884). Джеймс иллюстрирует свою теорию, указывая, что мы испытываем сожаление потому что плачем, злимся потому что наносим удары, боимся потому что дрожим --- а не наоборот.

Здесь стоит пояснить, что теория Джеймса --- не о \textit{связи} чувств с эмоциями, эмоции здесь --- и грубые, и тонкие --- состоят из ощущений телесных изменений. То есть причиной является восприятие, воспоминание или мысль, за ней следуют телесные изменения, которые и ощущаются особым образом, и эти ощущения мы называем эмоциями.

Утверждение, что все телесные изменения так или иначе ощущаются, может создать впечатление, что теория Джеймса --- гибридная, и говорит лишь о важности для эмоций как телесных ощущений, так и изменений. Однако Джеймс допускает переживание эмоции от галлюцинации телесных ощущений. Согласно Дамасио (1994), эмоции возникают не напрямую из-за изменений, то есть могут переживаться в их отсутствие, потому что центры мозга, реагирующие на них, могут активироваться и по другим причинам, создавая впечатление телесных изменений. Именно поэтому достаточно представить, например, грозного медведя, чтобы вызвать эмоции через активацию таких вот центров, в обычной ситуации реагирующих на повышение пульса и т.п.. Учитывая возможность спутать телесные ощущения и телесные изменения, следует считать теорию Джеймса теорией именно чувств, а не гибридной.

Поскольку Джеймс допускает ощущения, не коррелирующие с изменениями, его теория защищена от критики Уолтера Кэннона (1927, см. также Кобос и др. 2002), указывающей, что пациенты с травмами спинного мозга, способны испытывать эмоции, не имея обратной связи со своими телесными изменениями. Теория телесных ощущений допускает своего рода зацикленность причинной цепочки при таких травмах. Конечно, речь не о том, что все эмоции так возникают, что согласуется с ослабленностью некоторых (что интересно --- не печали) эмоциональных переживаний пациентов с поврежденным спинным мозгом пропорционально ослабленности обратной связи (Хоманн 1966).

Телесные изменения, согласно Джеймсу, обычно автоматичны и принадлежат к тому же классу импульсов, что и инстинктивные реакции: такие эмоции, как страх высоты --- это реакции, возникающие раньше и часто противоречащие рациональным суждениям. Джеймс допускает значительные индивидуальные различия в таких реакциях --- они зависят от нашего опыта и личных интересов, а также от социального окружения: например, от нашего воспитания зависит, что мы считаем оскорблением.

Джеймс считал, что даже не имея прямого контроля над эмоциями, мы можем контролировать их косвенно --- через соответствующие контролируемые телесные изменения. Так, мы можем вызвать панику, убегая от опасности, усилить печаль, рыдая, и вызвать меланхолию, сидя в сгорбленной, закрытой позе. Подобным образом эмоции можно и ослаблять, например, считая до десяти, чтобы уменьшить интенсивность гнева. Согласно Джеймсу, ''если мы хотим преодолеть нежелательные эмоциональные тенденции, нам следует научиться усердно и хладнокровно проходить через внешние их проявления, не усиливая их, как мы привыкли'' (1884, 198, см. также 1890).

\section{Аргументы в пользу чувственно-телесной теории}

Привлекательность теории Джеймса состоит главным образом в фокусе на наиболее очевидном, и, казалось бы, важном, компоненте эмоций --- чувствах. Это, кстати, эмпирически подтверждено: Панксепп (2000) проверял обыденные представления разных людей о важности различных компонентов эмоций: мимики, голосовых выражений, чувственных состояний, когнитивных изменений и вегетативных изменений, и выяснил, что чувства имеют самый высокий рейтинг --- близкий к 4 по 5-балльной шкале, немного обгоняя когнитивные и вегетативные изменения. Но являются ли чувства вообще и телесные ощущения в частности --- не только типичными, но и существенными для эмоций?

По мнению Джеймса, одним из ключевых соображений в пользу его теории являются корректные предсказания с помощью нее. Так, из нее следует, что телесные изменения могут влиять на эмоции, и это, похоже, так и есть. Более того, как уже было указано в Главе 3, исследования лицевой обратной связи показывают возникновение телесных изменений и эмоций от имитаций определенных выражений лица. Например, поднятие уголков губ заставляет нас чувствовать себя более счастливыми, и в целом мимическая обратная связь может как модулировать, так и инициировать эмоциональные переживания (Зайонс 1984, см. также недавний метаанализ Коулз, Ларсен И Ленч, 2019). Впрочем, стоит отметить, что исследования лицевой обратной связи касаются лишь ограниченного диапазона эмоций. Николас Коулз и его коллеги (2019), например, говорят о гневе, отвращении, страхе, счастье, печали и удивлении, что недостаточно для подтверждения чувственно-телесной теории эмоций. Более того, эти результаты согласуются с альтернативными теориями эмоций, например с той, согласно которой эмоции и есть телесные изменения и, значит, намеренные телесные движения могут вызывать эмоции, тождественные телесным изменениям, которые уже затем вызывают телесные ощущения. Результаты исследований обратной связи вполне совместимы и с оценочными теориями, ведь лицевая обратная связь может вызывать чувства через активацию оценки через ассоциации (Смит и Кирби, 2004): улыбка, например, вызывает положительную оценку вещей, что и переживается как счастье. Таким образом, свидетельств исследований намеренных движений вообще и обратной лицевой связи в частности, недостаточно для обоснования чувственно-телесной теории.

Джеймс, впрочем, не считает концептуальной истиной свой тезис о телесных ощущениях как сущности эмоций, допуская, что и бестелесные духи обладают эмоциями. Он позиционирует свое объяснение как эмпирическую гипотезу, хотя основной аргумент в его пользу --- не эмпирический. Так называемый аргумент от вычитания состоит в мысленном эксперименте, в ходе которого следует представить сильную эмоцию и мысленно абстрагировать от нее всякие телесные ощущения. У нас, считает Джеймс, останется уже не эмоция, а нейтральное состояние бесчувственного познания:

\smallskip

\quad Что осталось бы от эмоции страха, если бы мы исключили учащенное сердцебиение и поверхностное дыхание, дрожание губ, ослабления конечностей, внутреннее возбуждение, а так же прочие, более тонкие, ощущения? [...] Что останется от горя, убери мы слезы, ощущение удушья и боли в грудине, и другие, не столь заметные чувства? Останется только осознание плачевности определенных обстоятельств, и ничего более.

\smallskip

Здесь, прежде чем переходить к обсуждению мысленного эксперимента, следует отметить, что, согласно Джеймсу, после устранения телесных чувств остается то, что можно было бы назвать оценочным суждением. Это интуитивное предположение предвосхищает оценочные теории эмоций, но Джеймс, похоже, не замечает его проблемности для своей теории, и ничто в его описании эмоций и их причин не содержит и намека на оценочность. Так что совершенно загадочно, как оценочное суждение может остаться от эмоций после исключения из них всех чувств.

Итак, насколько убедителен этот эксперимент? Можно возразить, что утверждая ощущаемость всех телесных изменений, Джеймс не может провести важное разграничение между тем, что останется, если исключить из эмоций телесные ощущения, и тем, что останется, если исключить из них телесные изменения. А ведь, кажется, сильный гнев, например, труднее представить вне телесных движений, чем вне телесных изменений. Без этого разделения аргумент может быть использован как для защиты чувственно-телесной, так и для защиты двигательно-телесной теории.

Еще одна проблема рассматриваемого эксперимента касается сильных эмоций. Как отмечает философ Дэвид Айронс (1894), теория Джеймса может быть обоснована только для таких эмоций. И действительно, мысленный эксперимент выглядит куда менее убедительно при рассмотрении в нем состояния легкого гнева или страха. Защитник теории чувств может ответить, конечно, что интуитивность необязательности телесных ощущений для эмоций опосредованна корреляцией интенсивности эмоций с важностью роли телесных изменений в них, и корреляция эта хорошо объясняется именно чувственно-телесной теорией. Но вовсе не очевидно, что даже проявления сильного сожаления или надежды должны сопровождаться телесными ощущениями --- и рассматриваемый эксперимент здесь не убедителен. Следует, впрочем, указать, что Джеймс рассматривает также моральные, интеллектуальные и эстетические удовольствия и утверждает, что отсутствие их связи с телесными чувствами превращает их в только лишь суждения. Однако, недавняя серия исследований интуиции в отношении аргумента вычитания показала, что большинство людей склонны считать такие эмоции как страх, печаль и гнев, сохраняющимися и даже после ''вычитания'' телесных чувств (Диас 2022). То есть телесные ощущения не столь уж важны для эмоций.

Давайте теперь рассмотрим основные доводы против теории Джеймса.

\section{Проблемы демаркации}

Первое возражение против чувственно-телесной теории эмоций состоит в указании на ее неспособность отличить эмоции от прочих состояний, в частности --- от телесных ощущений, никак с эмоциями не связанными (см., среди прочего, Айронс 1894). Мы можем чувствовать повышение пульса во время пробежки, но не испытывать при этом никаких эмоций, а значит, такие чувства ''эмоциональны'' лишь в той мере, в какой они включены в эмоции, понятые как в общем случае независимые от чувств. Попытка отождествить эмоции с телесными ощущениями, похоже, обречена на провал.

Отвечая на это возражение, можно привести исследования, посвященные корреляциям между телесными и нервными изменениями и эмоциями: в предыдущей главе рассматривалась классификация с помощью алгоритмов машинного обучения, показывающая существование паттернов телесных реакций, соответствующих ряду видов эмоций, таких как удивление, страх, гнев и печаль (Крагель и ЛаБар 2013, см. также Сигел и др. 2018). Эти результаты, конечно, не охватывают все виды эмоций, но можно утверждать, что фигурирующие в них эмоции составляют и все прочие. Если это верно, то можно попытаться отделить эмоции от других состояний сочетаниями конкретных телесных паттернов, что, впрочем, уже эмпирический вопрос.

Однако, этот ответ содержит концептуальные проблемы. Во-первых, установив корреляцию эмоций с телесными паттернами, мы не установили, что эти паттерны всякий раз вызывают эмоции. Во-вторых, такая ориентированная на телесных паттернах стратегия не может сама по себе утверждать демаркацию: что рассматриваемое множество телесных паттернов соответствует особым чувствам --- отдельный эмпирический вопрос. Наконец, поскольку чувственно-телесная теория допускает несоответствие чувств изменениям, то телесные паттерны не получится использовать в качестве демаркационного критерия для признаваемых теорией эмоций.

Но, может, приведенные выше исследования отвечают на возражение, что чувственно-телесные теории не могут отличать разные виды эмоций (см. Кэннон 1927; Бедфорд 1956-1957)? Речь о возможность связи разных видов эмоций с одними и теми же телесными изменениями. Так, согласно Уолтеру Кэннону (1927), телесные ощущения, характерные для страха, не отличаются от характерных для ярости, и эти телесные изменения столь распространены, что страх и ярость должны бы были ощущаться точно так же, как озноб, гипогликемия и лихорадка. Стоит также вспомнить эксперимент Шахтера и Зингера (1962), рассмотренный в предыдущей главе, показывающий, что состояния не-специфического возбуждения должны быть дополнительно эмоционально классифицированы для понимания их как определенных эмоций, и критикуемый за возможность ошибочного приписывания себе эмоций.

Итак, тезис Кэннона, согласно результатам исследований, использующих классификацию паттернов, неверен, потому что есть паттерны, соответствующие (каждый) нескольким видам эмоций: страх и гнев, например, очевидно разные эмоции, но не различаются в телесных изменениях. И даже если некоторые виды эмоций можно различить на основе таких паттернов, неясно, получится ли распространить это на хотя бы большую часть эмоций. Но, возможно, не стоит смешивать телесные паттерны и телесные ощущения. Как подчеркивает Эррол Бедфорд (1956-1957), с точки зрения феноменальных качеств есть лишь небольшая разница между ощущением негодования и ощущением раздражения. Наконец, разграничение видов эмоций на основе телесных паттернов дает неверный ответ в случае возникновения чувств независимо от телесных изменений.

\section{Контраргумент от бессознательных эмоций}

Одно из важных возражений против теорий, утверждающих необходимость чувств для эмоций --- интуитивная возможность эмоций без таких чувств. Речь именно об эмоциональных эпизодах, ведь даже сопровождаемые чувствами эмоциональные диспозиции не представляют проблемы для чувственных теорий --- эти теории говорят именно об эпизодах.

Эмоции без чувств, пожалуй, должны быть бессознательными. Существование таких эмоций, конечно, спорно. Отец психоанализа Зигмунд Фрейд объяснял многое из нашего поведения бессознательными мотивациями, но сомневался в возможности бессознательных эмоций, указывая, что ''суть эмоции в ее осознании для переживания'' (1915/1984, 179). Множество психологов, нейробиологов и философов последовали за Фрейдом в этом сомнении (см., например, Клор 1994, Леду 1994, Деонна и Терони 2012). Другие считают, что хотя эмоции обычно и включают чувства, могут все же быть и бессознательными (см., например, Дамасио 1994, Принц 2004, Шерер 2004).

Стоит отметить, что речь о бессознательных эмоциях в смысле феноменального сознания, а не сознания доступа (Блок 1995). Как это объясняет Нед Блок, феноменальный аспект сознания сводится к тому, каково находиться в данном состоянии, тогда как сознание доступа подразумевает доступ к информации. В известном случае ''слепого зрения'' --- состояния, вызванного повреждением первичной зрительной коры --- пациенты имеют доступ к информации о своем зрительном поле, но лишены феноменального его сознания, из-за чего у них создается впечатление, будто они ничего не видят, но способны достоверно отвечать на вопросы о визуальных раздражителях.

Вообще, обе стороны дискуссии согласны, что эмоции возможны без сознательного доступа к ним в том смысле, что может не хватать информации о собственных эмоциях. Так, можно злиться на друга и не замечать этого, или неосознанно бояться определенных вещей. Но это не значит отсутствие чувств, так что не угрожает чувственным теориям эмоций. Дискуссионный же вопрос --- обязательно ли феноменальное сознание для переживания эмоций.

Как уже обсуждалось в Главе 2, эмоции, казалось бы, не обязательно включают чувства, и в этом заключается их отличие от, например, зуда. Одна из стратегий защиты необходимости чувств для эмоций --- предположить это по умолчанию, и затем прийти к выводу, что это так и есть, объясняя все очевидные случаи бессознательных эмоций. Так, Деонна и Терони (2012) утверждают, что эмоции могут быть по-разному бессознательными даже не будучи бессознательными в феноменальном смысле: можно, например, испытывать гнев, когда внимание направлено на что-то с ним не связанное, или переживать чувство вины, даже не осознавая, что это оно, и эти сценарии, конечно, совместимы с чувствами. Такая стратегия, однако, не исключает случаев эмоций, не связанных с чувствами.

Более амбициозный аргумент был предложен Джеральдом Клором (1994). По мнению Клора, эмоции должны быть сознательными, потому как этого требует их функция --- информировать человека о природе и важности происходящего и соответствующим образом мотивировать его. Чтобы информировать, нужно сначала хотя бы привлечь внимание информируемого, и сделать это можно лишь через сознательный опыт. Этот аргумент, впрочем, терпит неудачу, даже если тезис о такой функции эмоций верен. Как утверждает Принц (2004), эмоции вполне могут информировать и мотивировать и без привлечения к себе внимания и затрагивания чувств, ведь передача информации и формирование мотивации может производиться и за порогом сознательного опыта. Так, эмоции могут играть важную роль в неосознанном направлении внимания.

В ответ может возразить, что испытываемые в ходе переживания эмоций чувства имеют специфическую функцию, отличную от постулируемой Клором. Возможно, ''эмоциональные'' чувства играют особую роль, потому что позволяют сознательно, рефлексивно, контролировать эмоции (Шерер 2004), и, кажется, только через чувства эмоции могут оказать влияние на сознательное мышление. Таким образом, тезис о выполнении чувствами функции контроля --- вполне правдоподобен. Но дело в том, что некоторые эмоциональные эпизоды все еще могут обходиться без чувств и не выполнять такую функцию.

Здесь стоит остановиться, чтобы задаться вопросом: являются ли чувства, связанные с эмоциями --- телесными? Кэролайн Прайс (2015; см. также Соломон 2007, Кригель 2014) указывает, что подход Джеймса к феноменологии эмоций слишком узок --- помимо телесных изменений, переживание страха характеризуется концентрацией внимания на объекте страха и скоростью и спутанностью мыслей. Джеймс Рассел (2003) так же утверждает, что эмоциональные чувства включают не только таковые связанные с телесной активацией, но и удовольствия и неудовольствия --- состояния, очевидно отличные от телесных чувств.

Но, возможно, такое разнообразие эмоциональных чувств --- не проблема для приверженца чувственной теории? Ведь он может утверждать, что эмоции состоят из не только лишь телесных ощущений, но и из всякого чувственного опыта. Но такое расширение концепции несовместимо с первоначальным описанием эмоций, ведь оно не учитывает их интенциональность. Так, когда вы боитесь, ваше внимание направлено на то, чего вы боитесь --- а не на телесные ощущения, повышенный пульс и т.п..

Приведенные рассуждения подводят нас к наиболее широко распространенному набору аргументов против чувственно-телесных теорий.

\section{Контраргументы от интенциональности и мотивации}

Основной аргумент против чувственно-телесных теорий --- их неспособность учесть интенциональность эмоций (см., среди прочего, Питчер 1965, Соломон 1976/1993). Аргумент этот в равной степени эффективен как против теорий, говорящих о телесных изменениях, так и тех, что прослеживаются до Юма (1739-1740/2000, см. Уайтинг 2009) и предполагают эмоциями чисто феноменальные состояния, свободные от интенциональных объектов. Чувственные теории, значит, либо отрицают интенциональность эмоций, либо неверно ее истолковывают. Отрицать ее они вынуждены, чтобы говорить о чувствах как неинтенциональных состояниях. Так, согласно одной из интерпретации Юма, эмоции не имеют репрезентативного содержания, и потому ваш страх при встрече с медведем, например --- не страх медведя, а просто страх. Согласно теории Джеймса, впрочем, можно утверждать, что телесные ощущения направлены --- но не на то, что обычно понимается объектом страха, а на соответствующие телесные изменения, поскольку именно эти изменения они и репрезентируют. Но это, конечно некорректно, ведь объект страха из примера с медведем --- медведь, а не наши телесные изменения при его виде.

Таким образом, как отрицание, так и некорректное понимание интенциональности, сами по себе --- проблематичны. Но они создают также еще одну проблему: из-за них эмоции нельзя оценить с точки зрения рациональности. Например, страх при виде безобидного плюшевого мишки --- явно иррационален. Иррациональность здесь обычно понимается как некорректная характеризация объекта как опасного. При этом ни одна из двух версий чувственных теорий не допускает какой-либо оценки: если у эмоций нет интенционального объекта, то характеризовать нечего, если же эмоции направлены на телесные изменения, то верно или неверно характеризовать можно только сами эти изменения.

Можно попытаться ответить, что эмоции как таковые являются неинтенциональными состояниями, хотя и могут выступать частью сложных состояний, включающих направленные мысли (Уайтинг 2009). Боязнь медведя тогда --- такое сложное состояние, состоящее из убеждения, что поблизости находится медведь и неинтенционального эмоционального состояния. Интересно, что страх перед медведем тогда --- не эмоция, а лишь включает ее. Мы, значит, принимаем сложное состояние за эмоцию потому, что смешиваем его с его эмоциональным компонентном. Однако такое понимание термина ''эмоция'' не соответствует его обыденному использованию. Более того, оно не позволяет говорить об их рациональности/иррациональности.

Более уступчивый ответ на возражение от интенциональности предполагает рассмотрение причин эмоций. Согласно чувственно-телесной теории, то, что вызывает телесные изменения, само вызвано восприятием, воспоминаниями или мыслями, а значит и страх в примере все же связан с медведем, потому что является результатом его восприятия вами. Таким образом, эмоция ''наследует'' свой объект от своей когнитивной основы. Этот ответ, впрочем, не слишком убедителен, ведь, как указывает Прайс (2015), мы не склонны использовать одну и ту же такую стратегию в схожих ситуациях: тот факт, что, например, жаркая погода вызывает у вас мигрень, не значит, что ваша мигрень --- о погодных условиях. Мигрень очевидно не интенциональна.

Думаю, этих контраргументов достаточно, но давайте рассмотрим еще один, поскольку он не менее силен: чувственно-телесные теории искажают связь эмоций с действиями. Так, попытку убежать от медведя можно объяснить эмоциями: они \textit{мотивируют} действия, и, значит каузально объясняют их. Айронс (1894) утверждает, что такое объяснение возможно лишь благодаря наличию у эмоций интенциональных объектов, потому что то, что мы делаем --- так же интенционально, направлено на вещи, и совершенно непонятно, как чувственная теория может объяснить связь таких направленных действий с ненаправленными (или направленными, но на телесные изменения) эмоциями. Джеймс, впрочем, осознает, что чувственно-телесная теория эмоций переворачивает их обыденное понимание: восприятие медведя вызывает телесные изменения, а эти изменения вызывают чувства, которые и составляют страх, и страх, значит, сам по себе никак не объясняет то, что вы затем делаете. В следующей главе мы рассмотрим мотивационные теории, основная привлекательность которых заключается в естественном объяснении действий с точки зрения эмоций.

\section{Подведение итогов}

Похоже, что у нас мало оснований принимать чувственно-телесную теорию. А раз чувства понимаются либо как связанные с телесными изменениями, либо как неинтенциональные состояния, то и чувственные теории вообще (а не только говорящие о телесных чувствах) терпят неудачу из-за неспособности учесть интенциональный и мотивационный компонент эмоций.

\section{Краткое содержание}

\begin{itemize}
\item\ Теории эмоций, такие как чувстввенные, мотивационные и оценочные, пытающиеся ответить на вопрос о сущности эмоций, находятся на другом уровне, нежели, например, теории базовых эмоций и социальный конструктивизм, потому что совместимы с любыми из них.
\item\ Теория эмоций должна решать несколько задач. Во-первых, отличать эмоции от прочих состояний, и одни виды эмоций --- от других. Во-вторых, она должна учитывать наши интуиции касательно эмоций.
\item\ Понимание эмоций затруднено из-за их разнообразия и сложности. Эмоции обычно включают несколько компонентов: информационный, оценочный, физиололгический, экспрессивный, феноменологический, ментальный и мотивационный.
\item\ Теории, утвеждающие необходимым и достаточным компонентом только какой-то один, называются чистыми, а говорящие о нескольких таких компонентах --- гибридными.
\item\ Согласно самой известной чувственной теории --- чувственно-телесной теории Джеймса --- эмоции являются телесными чувствами, то есть ощущениями внутренних и внешних телесных изменений, в том числе мотивационных, и изменения эти вызываются восприятием, воспоминанием или мыслью о чем-либо значимом.
\item\ Основной аргумент в пользу этой теории --- мысленный эксперимент, предлагающий ''вычесть'' чувства из эмоционального эпизода. Проблема этого аргумента в отсутствии четкого различия между телесными ощущениями и телесными изменениями, из-за чего он правдоподобен разве что при рассмотрении сильных эмоций.
\item\ Даже если признать, что телесных изменений достаточно для отделения эмоций от прочих видов состояний и их разделения между собой, чувственно-телесная теория, похоже, не способна решить проблему демаркации.
\item\ Возможность бессознательных эмоций означает необязательность чувств --- телесных и любых других --- для эмоций.
\item\ Наиболее сильные аргументы против чувственных теорий вообще и чувственно-телесных теорий в частности --- указывают на неспособность этих теорий учесть интенциональность эмоций: они вынуждены либо отрицать ее, либо закреплять на телесных изменениях.
\item\ Чувственные теории неспособны также объяснить связь эмоций и действий.
\end{itemize}

\begin{tcolorbox}
\section{Учебные вопросы}
\begin{enumerate}
\item\ Что останется, если вычесть из эмоционального эпизода телесные ощущения, но не телесные движения?
\item\ Могут ли роботы испытывать эмоции, например, страх? Почему?
\item\ Может ли гибридная теория включать чувственный компонент эмоций в качестве одного из ключевых?
\end{enumerate}
\end{tcolorbox}

\chapter{Мотивационные теории}

\section{Введение}

Давайте посмотрим, как выглядит эпизод страха у сурка. Когда он видит кружащего в высоте орла, то тут же замирает и свистит, чтобы предупредить остальных сурков в колонии. Когда орел уже приближается, сурок срывается с места и бежит в нору. Если орел догонит его, он станет сопротивляться, а если схватит --- впадет в своего рода паралич, внешне схожий со смертью. На этом примере хорошо видно, что страх включает в себя мотивацию, и мотивация эта достаточно гибка и может адаптироваться к меняющемуся контексту.

В предыдущей главе указывалось, что теории чувств не могут объяснить действия эмоциями --- например, почему мы убегаем при сильном страхе. Мотивационные же теории, относительно недавно снова обретшие популярность в философских дебатах, такое объяснить, конечно же, способны. Но действительно ли мотивация столь важна для эмоций? Способны ли мотивационные теории отделить эмоции от прочих состояний и разделять их между собой? Наконец, могут ли такие теории учесть такие интуитивно присущие эмоциям свойства, как интенциональность?

Множество знаменитых философов считали мотивацию одним из ключевых компонентов эмоций. Так, Аристотель в своей ''Риторике'' (2007) утверждал, что гнев предполагает желание мести. Фома Аквинский (1268-1273/1981) и Кант (1785/1996) предполагали, что эмоции состоят из мотивационных состояний. Джон Дьюи вслед за ними предположил, что вместо Джеймсианского подхода к эмоциям как к чувствам, их следует рассматривать как ''целенаправленный способ поведения'' (1985, 15), предполагающий готовность действовать определенным образом. Современные психологи также подчеркивают роль мотивации в эмоциях. Роберт Плутчик (1980), например, утверждает, что каждой базовой эмоции соответствует образец адаптивного поведения --- сдача или бегство в случае страха, например. И как указывалось в Главе 3, приверженцы теорий базовых эмоций считают, что составляющие эмоции аффективные программы включают тенденции к действию (см., например, Экман 1999). Мы начнем с рассмотрения так называемой простой мотивационной теории (в отличие от сложных --- включающих среди ключевых компонентов эмоций также и оценки), изложенной психологом Нико Фридой (1986, 2007, 2010), поскольку она представляет наиболее разработанную версию мотивационного подхода. Эта теория оказала сильное влияние на современные теории эмоций, такие как мотивационная теория Андреа Скарантино (2014) и теория установок Деонны и Терони (2012). Заключением этой главы будет вывод, что эмоции не обязательно включают мотивацию, и даже когда включают --- включают так же и оценки.

Фрида в своей теории исходит из того, что эмоциональный опыт вообще и эмоциональное поведение в частности, говорят о наличии побуждений к действию в эмоциях. Поэтому лучший способ понять такие эмоции, как страх или гнев --- рассматривать их как такие ''тенденции к действиям'' (термин, заимствованный Фридой у пионера в оценочных теориях эмоций --- психолога Магды Арнольд (1960)). Согласно Арнольд, эмоция --- это ''ощущаемое влечение к чему-либо, интуитивно оцениваемому как хорошее (полезное) или отторжение к оцениваемому как плохое (вредное)'' (Арнольд 1960, 171). Согласно Фриде, тенденции к действию --- это ''состояния готовности к выполнению этого виду действий'' (1986, 70), где вид действий определяется ''конечным результатом, намеченному или достигнутому''.

Каждая такая тенденция предполагает готовность к множеству действий, направленных на одну и ту же цель. Так, гнев заключается в ''готовности к нападению, оскорблению, отказу/отворачиванию или клевете --- смотря что из этого возможно и уместно в данном контексте'' (1986, 71). Целью гнева при этом является ''восстановление контроля или свободы действий --- как правило, устранение препятствий'' (1986, 89). Результат конкретных действий зависит от контекста, так что объяснение с точки зрения тенденций к ним допускает значительную гибкость, что делает теорию правдоподобней тех, что говорят об эмоциях как жестко определенных, рефлексивных, реакций. Эмоциям ведь иногда можно даже сопротивляться, так что им не всегда удается полностью проявиться в действиях, хотя обычно они именно к полноценным таким действиям приводят.

Фрида предлагает идентифицировать эмоции по тенденциям к действию: гнев направлен на устранение препятствий, страх --- на избегание для защиты, отвращение --- на отклонение для защиты. То есть тенденции к действию соответствуют конкретным целям. И поскольку эмоции направлены на такое установление, поддержание или разрушение отношений с окружающей средой, Фрида подчеркивает реляционный их аспект: ''сфокусированные состояния готовности к действию --- тенденции к действию --- подразумевают готовность устанавливать, поддерживать или изменять отношения с миром'' (Фрида 2007, 33).

Дело, впрочем, в том, что некоторые эмоции --- в отличие от гнева, страха или отвращения --- не преследуют конкретных целей. Радость по поводу выздоровления друга, например, не обязательно должна выражаться в каком-то конкретном действии, направленном на конкретную цель. Согласно Фриде, эта эмоция подразумевает готовность к действию без конкретной цели. В более поздней работе (2010, см. также Скарантино 2014) она говорит об эмоции радости как тенденции к установлению связи с миром в целом. Печаль и горе тоже трудно объяснить определенными тенденциями --- они, похоже, имеют тенденцию уменьшать всякие тенденции к действиям. Фрида называет такие эмоции ''нулевыми состояниями'' и утверждает, что они соответствуют противоположной цели --- отстранению от мира. То есть эти эмоции все еще можно понимать в терминах готовности к действию, хотя и направленных на такую нетипичную цель.

Итак, согласно Фриде, ''эмоции [...] можно определить как способы готовности к реляционным действиям --- либо в форме тенденций устанавливать, поддерживать или разрушать отношения с окружающей средой, либо в форме готовности к реляционным действиям как таковой'' (Фрида 1986, 71). Как указывает Скарантино (2014), поскольку соответствующие таким эмоциям, как радость или печаль, модусы готовности включают общую, неконкретизированную цель, эмоции можно понимать как либо конкретные, либо такие обобщенные, тенденции к действию или бездействию.

Если эмоции --- состояния готовности к действию, то что отличает их от, например, склонности спать после обеда? В ответе Фрида предлагает сосредоточиться на побуждающей природе эмоций (см. также Скарантино 2014): состояния готовности, присущие эмоциям, характеризуются так называемым ''приоритетом контроля'' --- они ''требуют внимания и исполнения'' (1986, 78), превосходя в приоритете прочие, неэмоциональные, тенденции. Приоритет контроля предполагает телесную подготовку к осуществлению хотя бы части соответствующих действий. Так, испугавшись, сурок прерывает всякие занятия и внимательно смотрит за орлом, а его сердце начинает биться достаточно часто чтобы он в любой момент мог убежать или дать отпор (в зависимости от расстояния до хищника).

Итак, согласно простой мотивационной теории, эмоции состоят из относительно гибких тенденций к действию или бездействию, понимаемых в терминах готовности, направленных на конкретные или общие цели и обладающих приоритетом контроля. Насколько убедителен такой подход? Давайте сначала рассмотрим аргументы в его пользу.

\section{Аргументы в пользу простой мотивационной теории}

Мотивационная теория способна объяснить кроме мотивационности и ряд других типичных особенностей эмоций. Связь эмоций с чувствами, например --- из-за связи с ними тенденций к действию как состояний готовности. Чувства, впрочем, не обязательны, ведь эмоции могут мотивировать ментальные действия, происходящие без каких-либо телесных признаков. То есть в рамках простой мотивационной теории эмоции могут обыкновенно, но не обязательно включать чувства. Кратковременность эмоций также вполне объяснима --- приоритетные тенденции действий ведь тоже кратковременны.

Первый из аргументов, которые мы рассмотрим, апеллирует к функции эмоций. Как уже отмечалось при рассмотрении бессознательных эмоций, функция эмоций, похоже, заключается в информировании переживающего их о природе и важности происходящего, а также соответствующем мотивировании. Именно так Фрида и считает: ''функция эмоций состоит в том, чтобы сигнализировать о событиях, имеющим отношение к проблемам их переживающего, а также чтобы мотивировать поведение, направленное на корректное реагирование на эти события'' (Фрида 1994b, 121). И, как мы видели, эмоции не обязательно должны быть сознательными, чтобы выполнять такую функцию, особенно учитывая их мотивационный компонент.

Но действительно ли функция эмоций такова? И даже если так, работает ли этот аргумент только в пользу мотивационных теорий? Ведь эмоции могут играть значительную роль в мотивации, сами мотивациями не являясь. Например, они могут состоять из оценок --- состояний, способных мотивировать, и выполнять, значит, мотивационную функцию. Примечательно также, что и Клор, и Фрида постулируют двойную функцию эмоций --- включающую как мотивационный компонент, так и оценочный: эмоции, помимо мотивирования, информируют о важности событий и сигнализируют о тех из них, что имеют отношение к интересам испытывающего эмоции.

Второй аргумент разработан Скарантоно (2014) и допускает, что мотивационную функцию эмоций могут объяснить и другие к ним подходы. Цель этого аргумента --- показать, что, однако, только мотивационные теории, говорящие о гибких тенденциях действий с приоритетом контроля, могут объяснить наши действия при переживании эмоций. Вслед за Фридой он объясняет импульсивность чувством побуждения, предпочтением действовать как можно скорее, а также ''близорукостью'' таких действий --- их основанностью на ограниченной информации и предвзятости в ее обработке. Для гибкости эмоциональных действий характерно присутствие ограничений на то, что можно считать действием из-за страха или гнева, например. Эмоциональные действия также обычно сопровождаются телесными изменениями --- в мимике и интонациях, а также вегетативными изменениями.

Скарантино, аргументируя свою позицию, предлагает рассмотреть двух людей: Обычного Мэтта, чувствующего эмоции так же, как и мы, и Близнеца Мэтта --- морского пехотинца, обученного не бояться в опасных обстоятельствах. Во время посещения ими цирка, из вольера сбегает тигр. Действия Обычного Мэтта импульсивны, не слишком гибки и по большей части телесны: он дрожит и потеет, забывает, где находится ближайший выход, принимает тень за другого тигра, а оказавшись перед тигром и, казалось бы, осознавая, что должен стоять и не двигаться --- не может удержаться и забирается на дерево, на котором тигр может его достать. Близнец Мэтт же, понимая, конечно, что тигр опасен и что нужно поскорее выйти из помещения, спокойно оценивает ситуацию и идет к ближайшему выходу, а столкнувшись с тигром, остается неподвижным даже при его приближении.

С точки зрения тенденций к действию разницу в действиях Обычного Мэтта и Близнеца Мэтта объяснить нетрудно, но показывает ли это существенность тенденций к действиям для эмоций? Дело в том, что этот аргумент, похоже, плохо обобщается на прочие, не столь сильные, эмоции --- будь даже это такие простые как страх, гнев или отвращение. Так, действия вследствие легких форм этих эмоций не обязательно импульсивны, ограниченно гибки и существенно связаны с телесными движениями. Если бы в рассматриваемом примере сбежал тигренок, а не взрослый тигр, можно было бы разве что немного опасаться, что он может прыгнуть и поцарапать вас, и потому даже Обычный Мэтт действовал бы спокойно. Эти спокойные действия, конечно, не эмоциональны \textit{в смысле импульсивности}, но, строго говоря, они вызваны эмоцией. Защитник мотивационной теории мог бы ответить, что всякое эмоциональное действие должно характеризоваться пусть не крайней, но, все же, импульсивностью. Но некоторые виды эмоций, похоже, провоцируют вовсе не импульсивные действия. Восхищение красивым пейзажем может побуждать сфотографировать его, но соответствующие действия вовсе не обязательно должны быть импульсивными и негибкими. Точно так же, если вы мотивированы кому-то помочь, ваши действия могут выглядеть вовсе не ''эмоционально''.

Рассмотренным аргументам в пользу простой мотивационной теории, таким образом, достаточно легко противостоять. Давайте теперь перейдем к аргументам против этого подхода.

\section{Контраргументы от демаркации}

Как простая мотивационная теория решает проблемы демаркации? Может показаться, что она хорошо объясняет разницу между эмоциями и другими психическими состояниями, ведь тенденции к действию, помимо прочего, характеризуются приоритетом контроля, отличающим их от склонности засыпать после обеда, например, и ограниченной гибкостью, отличающей их от жажды, голода и рефлексов. Так, жажда мотивирует только одно действие --- питье, голод провоцирует есть все, что попадется, а рефлекс --- закрытие глаз при приближении к ним предмета, например --- совершенно негибок.

Но тогда легкий страх, например, не может считаться эмоцией, ведь не предполагает достаточного приоритета контроля --- мотивации, проистекающие из таких слабых эмоций могут и не преобладать над другими склонностями к действию, не требовать, значит, срочной соответствующей когнитивной и телесной подготовки. Скарантино (2014) настаивает, что это лишь нечеткая граница между эмоциями и тем, что ими не является. Но ведь даже такие легкие переживания обыкновенно считаются эмоциями: вы можете немного бояться, что ваш начальник отреагирует на ваш уничижительный комментарий, и все же --- желая его унизить --- не сдерживаться, но ведь даже такой, не достигающей своей цели, страх, можно считать эмоцией.

Более того, некоторые эмоции --- такие как восхищение и сожаление, могут быть и не связаны с мотивацией, а значит, не считаться эмоциями по рассматриваемому критерию. Так, можно восхищаться пейзажем, но не мотивироваться от этого к чему-либо. Можно, впрочем, решить пойти посмотреть на него с других ракурсов, то есть восхищение, возможно, мотивирует продолжать созерцать то, что восхищает. Деонна и Терони (2012) утверждают, что восхищение побуждает исследовать свой объект. Но, похоже, восхищение возможно и без этих сопутствующих мотиваций --- можно просто любоваться пейзажем (Тапполет 2016).

Аналогичные замечания применимы и к облегчению и сожалению. Чуть не потеряв друга из-за тяжелой болезни, вы будете думать, как хорошо, что он выжил, но какую тенденцию к действию может подразумевать такое переживание? Защитник мотивационной теории может ответить, что достаточно уже тенденции к такой мысли, ведь и мысль можно считать действием. Но такие связанные с оценками действия предполагали бы связь эмоций с оценками, а значит, мотивационная теория не полностью описывает их. Когда вы сожалеете, что, например, пять лет назад решили изучать право, тогда как философия оказалась интересней, вы и вовсе не можете что-либо с этим сделать --- то есть для сожаления нет тенденции к действию (или бездействию) и соответствующей цели. Оно, конечно, может побудить вас задуматься о принятом решении, но и то не обязательно.

При разграничении между видами эмоций также проблематично апеллировать только к тенденциям к действию: мы склонны нападать и когда боимся, и когда злимся, и склонны избегать и того, чего боимся, и того, что нам отвратительно. Избегание может быть даже способом выражения гнева. Защитник мотивационной теории может в ответ утверждать, что говорить нужно о целях, на достижение которых направлены различные тенденции действий: страх подразумевает защиту себя от вреда, тогда как цель гнева --- устранение препятствий (Фрида 1986, см. также Скарантино 2014), а отвращения --- устранение отвратительного объекта. То есть даже если действия, связанные со страхом, гневом и отвращением, могут пересекаться, их цели различны.

Такие не связанные с мотивацией эмоции простая мотивационная теория не сможет разграничивать, ведь не сможет говорить о них в терминах тенденций к действиям и целей. Еще одна проблема в существовании, согласно этой теории, эмоций, преследующий одну и ту же цель: устранение препятствий как цель может быть понята и как защита, отчего цели гнева и отвращения становятся неразличимы, точнее, их различение опирается на только лишь произвольное описание цели действий. В случае страха мы ищем защиты от опасности, в случае отвращения --- защиты от осквернения, а в случае гнева --- от препятствий. И раз ''препятствующий'', ''опасный'' и ''загрязненный'' --- оценочные выражения, то для разграничения эмоций, похоже, следует апеллировать не к мотивациям, а к оценкам. Возможна ли вообще тенденция к действию, не отражающая такую оценочную особенность, то есть, можно ли быть мотивированным защитить себя от опасности, например, не представляя то, от чего защищаешься --- опасным? Кажется, нет. Но тогда эмоции должны включать, помимо мотиваций, и оценки.

Роберт Робертс (2003), рассматривая пример гнева на коллегу, указывает, что гнев этот может спровоцировать множество самых разных действий, начиная от презрительного взгляда или улыбки и заканчивая сталкиванием в сточную канаву, и ни одно из этих действий само по себе не привязано к гневу: презрение может быть всего лишь презрением, а столкнуть в канаву можно чтобы спасти от падения чего-нибудь тяжелого. Робертс утверждает, что связанные с гневом действия следует понимать как способы наказания правонарушителя, то есть именно наказание, а не любое устранение препятствий, как считает Фрида, объединяет действия, совершаемые в гневе. Но, согласно Робертсу, для понимания действий в таких терминах как наказание и правонарушение, необходимо уже привлекать оценки, ведь гнев вызывается тогда оскорбительной ситуацией, а ''определение эмоций через готовность к действию переворачивает концептуальные приоритеты [...]'' (2003, 169).

Вывод о необходимости оценок для эмоций легко сделать также и из их интенциональности, что мы сейчас и увидим.

\section{Контраргументы от интенциональности}

Простая мотивационная теория, как демонстрировалось, объясняет ряд типичных особенностей эмоций, таких как их связь с чувствами, например. Но учесть их интенциональность куда сложнее. Когда вы боитесь тигра, ваш страх связан с тигром --- то есть тигр является его интенциональным объектом. Но тенденции к действию, являясь лишь поведенческими, ни на что такое не направлены: ваша тенденция бежать не направлена на бег, так же как и тенденция камня падать --- на падение. В ответ можно указать, что тенденции к действию должны пониматься как психические состояния вроде желаний, имеющие, значит, интенциональные объекты: когда вы хотите пить, ваше желание направлено на питье, которое как процесс и является интенциональным объектом желания. Аналогично, тенденцию убегать от тигра можно понимать как желание убежать от него, а значит бегство будет здесь интенциональным объектом.

Приверженцы мотивационной теории, похоже, понимают тенденции к действию как нечто сугубо телесное, а потому вряд ли одобрят такой ход. Но даже согласись они с ним, полученный интенциональный объект был бы некорректен: когда вы боитесь тигра и хотите убежать, ваши эмоции о тигре, а не о бегстве, хотя желание убежать, конечно же, направлено на бегство. В ответ можно сослаться на цель бегства --- защиту от опасности, то есть страх тигра подразумевает желание такой защиты. Но даже эта защита не выглядит действительным интенциональным объектом страха, такой объект --- это, очевидно, тигр.

Второй контраргумент указывает на связь эмоций с такими вещами, как опасности, загрязнения, препятствия или оскорбления --- тем, что мотивационная теория учесть не может. Страх предполагает не просто мотивацию к борьбе с опасностью, а, прежде всего, оценку нечта как опасного. Как только при описании страха речь заходит об опасного, приходится говорить и о представлении нечта как такого опасного.

Указать на это можно и через так называемые \textit{направления соответствия} (Серль 1983). Желания направлены от мира к разуму, то есть мир должен измениться, чтобы соответствовать желанию: если в середине лета вы захотите, чтобы пошел снег, ваше желание будет удовлетворено только если мир изменится так, что пойдет снег, то есть оно нацелено на реализацию, требующую изменений от мира. Убеждения же направлены противоположно, и должны соответствовать миру: вера в то, что идет снег, соответствует миру, если снег действительно идет. Направления соответствия тесно связаны с условиями истинности.

Итак, согласно простой мотивационной теории, эмоции в смысле соответствия направлены от мира к разуму, как и желания. Но, похоже, это не так. Даже не говоря об истинности или ложности эмоций (но см. Де Соуза 2011), можно говорить о них с точки зрения корректности, то есть отражения корректных оценочных фактов --- фактов об опасности, например. Маленького безобидного паука не стоит бояться, ведь он не опасен, а страх был бы связан с ложным представлением о нем как об опасном. Эмоции, значит, подобно убеждениям, направлены в своем соответствии от разума к миру, и, включая мотивацию, должны включать также оценки.

Отрицая связь эмоций с оценками, нужно объяснить, почему мы часто говорим о них исходя из их соответствия оценочным фактам. Можно, впрочем, признать необходимость оценок, но отрицать их существенность для эмоций, утверждая, как Фрида (1986), что оценки --- лишь причины эмоций, а не сами эмоции. Но ведь исходя из оценок мы говорим о рациональности/иррациональности именно эмоций, а значит, они должны быть их частью.

Еще более уступчивым ответом был бы переход к гибридной теории, сочетающей в качестве существенного для эмоций как мотивации, так и оценки: страх, например, состоит из веры в существование опасности и желания оставаться в безопасности. Однако, согласно Скарантино (2014), убеждения и желания обычно не связаны с чувствами, а значит такой подход не учитывает обыкновенную связь эмоций с ними. В ответ можно утверждать, что мотивационные элементы --- тенденции к действию --- обычно связаны с чувствами. Согласно усложненной версии этого подхода (обратите внимание, впрочем, что неясно, насколько можно считать его гибридным) эмоции состоят из одновременно оценочных и мотивационных состояний, которые представляют вещи согласно оценкам и соответствующе мотивируют (Скарантино, 2014). То есть эмоции --- так называемые ''тяни--толкай''--состояния (Рут Милликен, 2004), имеющие оба направления соответствия, представляя факты и включая цели. Эта теория удачно отвечает на второй контраргумент от интенциональности, а также на контраргументы от демаркации. Впрочем, как далее будет показано, она разделяет некоторые проблемы простой мотивационной теории.

Третий же ответ состоит в отрицании необходимости говорить об оценочном компоненте в терминах оценочного представления. Так, теория установок, разработанная Деонной и Терони (2012), утверждает, что если эмоции и связаны с оценками, то только из-за их связи с мотивацией. К рассмотрению этой теории мы сейчас и перейдем.

\section{Теория установок}

Деонна и Терони (2012, 2014, 2015) вслед за Фридой говорят о состояниях готовности к действию как о ключевом для эмоций, но принимают во внимание и акцент Джеймса на чувствах. Они исходят из того, что ''эмоции тесно связаны с типами готовности к действию, или, точнее, с ощущаемой готовностью к действию''. Так, в гневе мы чувствуем, что наше тело готово к битве с врагом, а от стыдясь, чувствуем, что готовы спрятаться от взглядов других. В отличие от простой мотивационной теории, теория Деонны и Терони предлагает понимать эмоции как состоящие из телесных установок по отношению к интенциональным объектам. То есть чувства готовности к действию --- это телесные ощущения, но направлены они не на тело, а на вещи в мире. Иначе говоря, ''связанные с эмоциями телесные изменения ощущаются как особые установки, обычно направленные на внешние объекты'' (Деонна и Терони 2014, 27). Таким образом, теория установок учитывает интенциональность эмоций.

Деонна и Терони также связывают эмоции с оценками: чувство готовности к действию не только направлено на интенциональный объект, но и (потому) предполагает его видение как опасного, оскорбительного и т.д.. Так, ''страх перед собакой --- это видение собаки как опасной именно потому, что есть ощущение готовности организма действовать в направлении уменьшения вероятного ее воздействия (драка, упреждающее нападение и т.д.) [...]'', и ''чувственное отношение корректно тогда и только тогда, когда собака действительно опасна'' (Деонна и Терони, 2012). То есть теория установок учитывает возможность говорить о рациональности/иррациональности (или корректности) эмоций, исходя из соответствия формального объекта интенциональному.

Стоит еще раз подчеркнуть самую важную часть подхода Деонны и Терони --- оценочный компонент заключается не в репрезентации. То есть боязнь собаки, например, не предполагает представления о ней как об опасной --- страх заключается в оценочной установке, а не оценочном суждении. Установка --- не то же, что содержание, а психическое состояние, на содержание направленное: вера, желание или предположение, например. Касательно одного и того же содержания можно иметь разные установки: можно верить, что идет снег, желать, чтобы он пошел, или предполагать, что он идет. Каждая установка характеризуется своей конститутивной целью. Так, установка убеждения нацелена на истину, отчего убеждение может быть неправильным, если не соответствует ей. Аналогично, мы можем испытывать разные эмоции по поводу того, что идет снег: бояться, что идет снег, злиться, что идет снег.

Деонна и Терони, значит, предлагают говорить об эмоциях как об особом семействе таких установок --- оценочных установках. Страх, например --- оценочная установка по отношению к таким содержаниям, как собаки и тигры. Установки эти оценочны поскольку имеют оценочную цель в качестве конститутивной. Конститутивная цель страха --- опасность, гнева --- обида, отвращения --- осквернение или отвратительность. Такое понимание контрастирует со стандартными оценочными теориями эмоций, где эмоции являются имеющими оценочное содержание установками или включают их.

Таким образом, теория установок утверждает, что эмоции --- чувства готовности к действию, которые следует понимать как оценочные установки, направленные на интенциональные объекты. Это позволяет ей разделять достоинства простой мотивационной теории, учитывая, например, возможную импульсивность мотивации, но теория оценок даже лучше, благодаря способности корректно говорить об интенциональности и оценках. Однако она разделяет с простой мотивационной теорией и некоторые слабости, к рассмотрению которых мы перейдем сразу после обсуждения аргументы в ее пользу.

\section{Аргумент в пользу теории установок}

Теория Деонны и Терони предлагается как более успешная альтернатива оценочным теории. Она, как заявляется, привлекательна казалось бы беспроблемным объяснением оценочного компонента эмоций. Деонна и Терони утверждают, что в отличие от конкурентов, теория отношений учитывает возможность направленности разных эмоций на одно и то же содержание.

Под конкурирующими теориями подразумеваются оценочные, но проблему можно сформулировать и как касающуюся простой мотивационной теории тоже, ведь два человека могут по-разному относиться к одному и тому же объекту: один бояться собаки, а второй --- злиться на нее. Это, согласно простой мотивационной теории, говорит о наличии разных мотиваций с разными, значит, целями: страх направлен на защиту от опасности, а гнев --- на защиту от обиды или правонарушения. То есть эти мотивации отличаются содержанием, но относятся к одному и тому же объекту, что, похоже, некорректно.

Но могут ли разные эмоции иметь одно и то же содержание (Росси и Тапполет, 2019). Конечно, в приведенном примере страх и гнев относятся к одному и тому же объекту --- собаке, но объединяет ли эти эмоции что-то еще? В эмоциях, похоже, содержится больше, чем в их интенциональных объектах. Согласно оценочным теориям, эмоции включают представление, имеющее оценочное содержание, то есть представление объекта как обладающего некоторым оценочным свойством. Содержание эмоции, значит, состоит из интенционального объекта и оценочного свойства: страх состоит из представления собаки как опасной, а гнев --- из представления ее как оскорбительной или мешающей. Интенциональный объект один и тот же, но полное содержание за счет включения также оценочных характеристик, отличается. Аргумент же в пользу теории установок предполагает отсутствие чего-либо в эмоциях помимо интенционального объекта, предвосхищая таким образом основание.

\section{Аргументы против теории установок}

Первая проблема теории установок касается ее акцента на чувствах. Эмоции обычно включают чувства, но далеко не ясно, всегда ли это так. То есть оценочные теории, похоже, могут понимать оценочные установки как мотивационные состояния готовности к действию, которые обычно, но не обязательно ощущаются. Вопрос об учете интенциональности эмоций при таком пересмотре интересен, но оставим его в стороне.

Вторая проблема серьезней: некоторые эмоции, такие как восхищение, облегчение или сожаление, похоже, лишь отдаленно связаны с тенденциями к действию и мотивацией (см. Митчел 2020). Восхищение, как уже указывалось, возможно и как просто восхищение, вне мотивированности даже созерцать или исследовать объект. Такие эмоции проблемны для всех мотивационных теорий, включая теорию ''тяни-толкай'' Скарантино.

Разграничение между видами эмоций также проблематично, ведь для различения тенденций к действиям нужно обращаться к соответствующим оценочным терминам: при страхе мы ищем защиты от опасного, а при отвращении --- защиты от грязного, и ''опасное'' и ''грязное'' являются оценочными терминами. То есть такие мотивации могут возникнуть только при представлении объекта как обладающего соответствующим оценочным свойством --- как можно бояться того, что не представляется опасным? Теория установок, казалось бы, не сталкивается с этой проблемой, ведь она может апеллировать к конститутивной цели разных оценочных установок.

Но, согласно Деонне и Терони, оценочная природа установок вытекает из лежащих в их основе состояний готовности к действию, то есть, например, боязнь собаки --- оценочная установка потому, что заключается в ощущении готовности организма уменьшить воздействие собаки на него. Можно ли тогда назвать такую готовность мотивацией, которая уже включает оценку (Баллард 2021)? Похоже, что да, ведь прежде чем пытаться избежать воздействия собаки, ее нужно посчитать опасной. Но именно это теория установок отрицает.

Следующая проблема касается направления соответствия эмоций. Теория установок, казалось бы, предполагает направленность эмоций от разума к миру, благодаря чему можно оценивать их корректность. Эмоции --- оценочные установки, имеющие оценочные конститутивные цели, и поэтому оценочные свойства используются в условиях их корректности. Аналогично устроена вера: ее конститутивной целью является истина, и поэтому некорректным считается то убеждение, что не соответствует действительности, то есть убеждение корректно тогда и только тогда, когда его содержание истинно. Истина, таким образом, входит в условия корректности веры потому, что является конститутивной целью этой установки. Страх же имеет в качестве своей конститутивной цели опасность, и именно поэтому она используется для определения его корректности: страх корректен тогда и только тогда, когда его содержание (интенциональный объект) --- опасно, то есть совпадает с формальным объектом.

Так может ли теория оценок учесть направленность соответствия эмоций от разума к миру, то есть их корректность в зависимости от того, каков он в терминах оценочных свойств? Можно утверждать, что из обусловленности корректности эмоций их соответствию миру следует включенность в них оценочных репрезентаций. Аналогично, корректность убеждений значит их соответствие действительности: если небо голубое, то убеждение о том, что небо голубое --- корректно, если страх перед собакой корректен, то только потому, что соответствует ее опасности.

Теория установок основана на понимании эмоций как ощущаемых состояний готовности к действию, как мотивационных, значит, состояний, и боязнь собаки --- видение ее страшной, поскольку страх состоит в ощущаемой готовности справиться с опасностью. Раз это мотивация, то готовность справиться с опасностью имеет направленность от от мира к разуму. Мир должен соответствовать желанию, чтобы оно было удовлетворено --- чтобы опасность миновала. Причем мотивация эта корректна не в смысле соответствия миру, а в смысле полезности --- борьба с опасностью полезна, сопутствует нашему благополучию. Но тогда получается, что корректность эмоции определяется не оценочными свойствами ее объекта, а ее полезностью.

Еще одна проблема теории установок и прочих теорий, говорящих об эмоциях в терминах готовности к действию --- существенная ограниченность мотивационных эффектов эмоций. Так, страх предполагает готовность действовать в направлении уменьшения воздействия объекта --- то есть включает поведенческие склонности, приводящие к бегству, борьбе или замиранию. Этого достаточно, чтобы объяснить поведение сурков, но люди ведут себя при страхе куда более разнообразно.

Как отмечает Джеральд Клор, ''Прямые эффекты эмоций носят скорее мотивационный, чем поведенческий характер: можно достичь большего согласия касательно целей тех, кто испытывает страх, чем касательно их конкретного поведения'' (1994a, 111; см. также Де Соуза 1987). Страх, согласно Клору, сопровождается, конечно, телесной активацией, облегчающей некоторые действия, но это не то же, что тенденция к действию. Страх предполагает желание избежать вреда, но как этого достичь --- зависит от контекста, а также от решения того, кто испытывает страх. В терминах Принца (2004), разделяя диспозиции на ''команды к действию'' и ''мотивы'', эмоции --- именно мотивы, ведь не дают повода для конкретных действий. Эмоции могут подготовить нас к действию, но решение остается за нами. Эмоции не всегда импульсивны, не всегда телесная активация столь уж сильна. Таким образом, теория установок, простая мотивационная теория и теория ''тяни-толкай'' не могут объяснить мотивационный компонент эмоций.

\section{Подведение итогов}

Мотивационные теории, в отличие от чувственных, могут объяснить влияние эмоций на наши действия. Но аргументы против них указывают на необходимость тогда и оценочных представлений. Следующая глава посвящена теориям, согласно которым центральную роль в эмоциях играют оценки

\section{Краткое содержание}

\begin{itemize}
  \item\ Согласно простой мотивационной теории, защищаемой Фридой, эмоции --- имеющие приоритет контроля, сравнительно гибкие тенденции к действию или бездействию, преследующие конкретные или общие цели.
  \item\ Теория Фриды хорошо объясняет некоторые особенности эмоций, но мотивационный компонент эмоций так же убедительно объясняется и оценочными теориями.
  \item\ Импульсивность, на которую можно указывать при защите простой мотивационной теории, далеко не всегда присуща эмоциям.
  \item\ Говоря лишь о импульсивных и мотивирующих эмоциях, простая мотивационная теория не может отличить слабые эмоции и эмоции не связанные с мотивацией (восхищение, например), от прочих психических состояний.
  \item\ По той же причине она не может различить и некоторые виды эмоций, а пытаясь различить их по их целям, должна использовать оценочные термины, и, значит, не может уже быть чистой теорией.
  \item\ Простая мотивационная теория едва ли способна учесть интенциональность и формальные объекты эмоций.
  \item\ Теория ''тяни-толкай'' Скарантино решает эти проблемы, указывая, что эмоции включают и репрезентации, и мотивации. Но объясняя эмоции как тенденции к действию, она сталкивается с некоторыми трудностями простой мотивационной теории.
  \item\ Теория установок стремится учесть интенциональность эмоций и пытается говорить об оценках, не говоря об оценочных репрезентациях, утверждая, что эмоция --- это чувство готовности к действию, направленная оценочная установка.
  \item\ Согласно защитникам теории установок, она может объяснить различие эмоций, имеющих одинаковое содержание, но аргумент неубедителен, поскольку предполагает, что содержание эмоции ограничивается ее интенциональным объектом.
  \item\ Теория установок, как и простая мотивационная теория, уязвима для аргумента от не-мотивирующих эмоций.
  \item\ Также теория установок, похоже, не может различить виды эмоций, не апеллируя к оценочным репрезентациям.
  \item\ Мы можем оценивать эмоции с точки зрения их корректности, но теория установок не способна предоставить соответствующий фундамент для этого.
  \item\ Теория установок, говоря про эмоцию как про готовность к действию, предлагает слишком жесткую модель действий, соответствующих эмоциям.
\end{itemize}

\begin{tcolorbox}
  \section{Учебные вопросы.}
  \begin{enumerate}
    \item\ Может ли эмоциям о прошлом соответствовать какая-либо мотивация? Если да, то верно ли это для всех таких эмоций?
    \item\ К чему побуждает вас восхищение произведением искусства? Всегда ли это какие-то определенные действия?
    \item\ Соответствует ли гневу какая-либо цель? Если да, то какая: устранение препятствий, наказание нарушителя или что-то еще?
  \end{enumerate}
\end{tcolorbox}

\chapter{Оценочные теории}

\section{Введение}

Давайте вернемся к примеру про сурка из предыдущей главы. Рассмотрев мотивационные теории, мы поняли, что следует учитывать оценки, связанные с эмоциями --- сурок оценивает орла как опасного. Первый набор вопросов касается связи между оценкой и эмоцией. Как в этом примере связана оценка со страхом? Является ли эта связь каузальной, где оценка выступает причиной эмоции? Или это конститутивное отношение, когда эмоция полностью или частично состоит в оценке? Может ли оно быть и каузальным, и конститутивным? Второй набор вопросов --- про характер этой оценки. В предыдущей главе указывалось, что оценки состоят в репрезентациях. Но это еще далеко не вся история: оценка может состоять в суждении или в восприятии. Оценочные теории отличаются ответами на перечисленные вопросы.

Принято считать, что оценочные теории восходят к стоикам. Как мы видели в Главе 1, они понимали эмоции как суждения о добре и зле, и считали такие суждения по большей части глубоко ошибочными и иррациональными --- что отличает их от прочих приверженцев оценочных теорий. Платон и Аристотель так же считали, что эмоции зависят от оценок. Платон рассматривал гнев как вскипание крови при мысли, что с кем-то поступили \textit{несправедливо}, Аристотель же определял гнев как болезненное желание мести, вызванное мыслью о \textit{неуважительном} отношении. Действительно, по крайней мере некоторые эмоции имеют оценку одним из необходимых компонентов или по крайней мере своей причиной. С этим соглашались также Фома Аквинский, Декарт и Юм.

Оценочные теории вновь приобрели популярность в конце 1950-х годов. Отвергая чувственные теории, такие философы как Ч.Д. Броуд (1954), Эррол Бедфорд (1956-1957) и Энтони Кенни (1963/2003) защищали понимание эмоций как определенного рода оценок, Уильям Лайонс же утверждал, что оценки являются причинами эмоций (1980). Независимо от них, в рамках смены бихевиоризма когнитивизмом, такие психологи как Магда Арнольд (1960) и Ричард Лазарь(1991) защищали так называемые теории отношений и утверждали, что эмоции вызываются оценками. И конститутивная версия, и каузальная, имеют своих современных сторонников.

Мы начнем рассмотрение оценочных теорий с теорий, защищаемых Арнольдом (1960), Лазарем (1991) и Клаусом Шерером (1984, 2005, 2009), поскольку они оказали наибольшее влияние как на психологию, так и на философию. Тезис о каузальности, лежащий в основе этого подхода, совместим со всеми тремя видами теорий эмоций. После мы обратимся к сужденческим теориям, отстаиваемым Робертом Соломоном (1976/1993, 2003, 2007) и Мартой Нуссбаум (2001), а также рассмотрим квазисужденческие теории, разработанные Патрисией Гринспен (1988) и Робертом Робертсом (2003). Несмотря на многочисленные преимущества, сужденческие теории сталкиваются с серьезными возражениями, поэтому мы сосредоточимся больше на теориях, ориентированных на восприятие, защищаемых Сабиной Деринг (2003) и Джесси Принц (2004). К концу главы станет понятно, что нам нужна теория, допускающая различия между чувственными восприятиями и эмоциями. Таким подходом может стать рецептивная теория, краткое описание которой я приведу в конце главы.

\section{Теории отношений}

Согласно Арнольд (1960), относиться к чему-либо --- значит видеть, что это каким-то важным образом на вас влияет. Процесс такой оценки, значит, определяет значимость ситуации для оценивающего. Арнольд считает эмоции результатом отношения к вещам как к хорошим / полезным или плохим / вредным. Она выделяет три измерения этой оценки: хорошее / плохое, присутствующее / отсутствующее и легко достижимое / избегаемое. Так, отношение, вызывающее страх, представляет собой оценку события как плохого, отсутствующего, но возможного в будущем и с трудом избегаемого. Отношение, вызывающее радость подразумевает понимание нечта как хорошего, уже начавшегося и легко поддерживаемого.

Впрочем, и гнев, и отвращение соответствуют плохому, которое есть и которое трудно прекратить, так что трех измерений, постулированных Арнольд, недостаточно. Поэтому Лазарь (1991) предложил дополнительный параметр --- благополучие. Рассматривая благополучие как зависящее от окружающих условий, а также от убеждений и целей, он утверждает: ''задачей оценок является интеграция двух наборов предшествующих переменных --- личности и среды --- в реляционную систему согласно релевантности происходящего для благополучия личности'' (1991, 39). Рассматриваемые оценки, значит, осознанные или нет --- оценки происходящего с точки зрения благополучия оценивающего.

Лазарь выделяет целых шесть измерений таких оценок, разделенных на две категории. К первичным компонентам относятся (а) релевантность цели, (б) соответствие или несоответствие цели и (в) тип вовлечения эго. Ко вторичным --- (г) вина или заслуга, (д) потенциал преодоления трудностей и (е) ожидания. Итак, процесс оценки начинается с релевантности вашим целям --- в противном случае никаких эмоций не возникает. Следующий вопрос --- мешает или помогает оцениваемое достигать ваши цели --- формирует валентность, то есть положительность или отрицательность эмоции. Третий вопрос --- ваша вовлеченность в ситуацию. Лазарь говорит о разных типах вовлеченности эго, среди которых самоуважение и уважение со стороны других. Вина или заслуга же подразумевают ответственность. Интересно, что Лазарь подчеркивает важность потенциала преодоления трудностей. Вопрос об ожиданиях --- вопрос о том, изменится ли ситуация в лучшую или худшую сторону с точки зрения достижения ваших целей.

Согласно Лазарю, сложные оценки, соответствующие разным видам эмоций, фиксируются ''основными реляционными темами'' --- концепцией, играющей важную роль в современных дискуссиях (см. Принц 2004, Соломон 2007, Скарантино 2014). Основная реляционная тема ''обобщает вред и пользу, присущие каждому взаимодействию человека и среды'' (Лазарь 1991, 39). У каждого вида эмоций есть своя основная реляционная тема. Для гнева, например, она определяется как ''унижение или оскорбления меня и/или моего'', для испуга --- как ''столкновение с непосредственной, конкретной и подавляющей физической опасностью'', для печали --- ''переживание безвозвратной утраты'' (Лазарь 1991, 122).

Одно из преимуществ теории Лазаря --- возможность различать 15 видов эмоций: гнев, тревога, испуг, вина, стыд, печаль, зависть, ревность, отвращение, счастье, гордость, облегчение, надежда, любовь и сострадание. Но теория отношений Шерера еще более амбициозна. В своей статье 2009 года он выделяет дюжину измерений, разбитых на четыре категории: актуальность, последствия, потенциал преодоления и нормативное значение. Так называемые ''оценки посредством стимулов'', выполняемые нами при оценке событий, могут быть сложными, но происходят быстро, ''на нескольких уровнях обработки, от автоматического и неявного до сознательного, концептуального или пропозиционального'' (Шерер 2005, 701). Примечательно, что Шерер не предполагает ограниченного набора видов эмоций --- напротив, множество способов объединения и взаимодействия аспектов оценки с течением времени позволяет ему говорить о существовании ''бесконечного множества различных типов эмоциональных эпизодов'' (2009, 1316). Разумеется, некоторые особенно адаптивные модели оценок встречаются чаще. Так, оценочный профиль гнева включает ''новизну, высокую значимость цели, наличие других агентов и намерений, высокую срочность, высокий уровень контроля, а также несправедливость и аморальность'' (2009, 1316).

Важно различать оценку как процесс и оценку как результат. Оценочные теории объединяет решающая причинная роль процесса оценки. Но чем тогда являются сами эмоции? Ответы могут быть разными. Арнольд (1960), например, придерживается мотивационной теории, считая, что эмоции --- ощущаемые тенденции к действию, вызванные оценками: оценки порождают влечение или отвращение, а эмоции отождествляются с такой ощущаемой тенденцией к хорошему или от плохого. Лазарус (1991) отстаивает гибридную теорию, согласно которой эмоции --- сложный процесс, включающий оценку как причину и комбинацию тенденции к действию, физиологического изменения и субъективного аффекта как следствие. Структурно схожая теория предложена Дженефером Робинсоном (2005), считающим эмоцию процессом, начинающимся с некогнитивной и автоматической оценки, запускающей характерные физиологические и поведенческие изменения, а также когнитивный мониторинг ситуации. Шерер определяет эмоцию как ''эпизод взаимосвязанных, синхронизированных изменений'' (2005, 697) во всех или большинстве из пяти подсистем, благодаря чему эмоциональный эпизод включает до пяти компонентов: оценку, вегетативные физиологические изменения, тенденции к действию, мимику и интонации и субъективные ощущения.

Тезис о важности процесса оценки для эмоций, вообще, совместим со всеми тремя видами теорий эмоций.

Считать же что эмоции вызываются не процессами оценки, а самими оценками --- отдельный подход, согласно которому эмоции --- реакция на то, как мы оцениваем вещи. Этот подход близок к мотивационной теории Арнольд, но его в основном защищают философы (Муллиган 2009, Мюллер 2019). Так называемые ''реактивные теории'' утверждают, что эмоции --- установки, возникающие в ответ на предыдущие оценки, причем эти оценки понимаются как чувства или восприятие ценностей. Можно ли говорить об этих оценках как отдельных, неэмоциональных состояниях (см. Митчелл 2019)? Чем может быть, например, чувство или восприятие опасности, если не страхом? В любом случае, оценки как оценочные представления должны рассматриваться как части эмоций, а не только их причины, ведь только тогда можно будет оценивать сами эмоции: лишь потому что страх состоит в репрезентации опасности, его можно оценить как корректный в случае ее действительного присутствия.

Давайте рассмотрим теории, утверждающие, что эмоции состоят из оценочных репрезентаций. Есть несколько способов делать это, и мы начнем с так называемых сужденческих и квазисужденческих подходов.

\section{Сужденческие и квазисужденческие теории}

По мнению таких философов, как Соломон (1976/1993, 2003) и Нуссбаум (2001), эмоции заключаются в оценочных суждениях. Соломон пишет: ''Эмоция --- это оценочное (или ''нормативное'') суждение [...]'' (1976/1993, 126). Аналогично Нуссбаум утверждает, что эмоции --- это ''оценки или оценочные суждения, придающие вещам и переживающему, вне его, воли, значимость для его собственного развития'' (2001, 4). Так, злиться --- значит выносить суждение, что кто-то унизил или оскорбил тебя, бояться --- значит выносить суждение о наличии опасности, грустить --- выносить суждение, что понес утрату.

Делая оценку центральным элементом эмоций, сужденческие теории получают преимущество перед чувственными и мотивационными теориями. Например, они могут учесть наличие у эмоций интенциональных объектов: ваш страх касается тигра, ведь его касается суждение о его опасности. Также благодаря оценочности суждений, из которых состоят эмоции, можно оценивать корректность эмоций. Таким образом, сужденческие теории, в отличие от мотивационных, полагают направленность эмоций от разума к миру.

Сужденческие теории успешно объясняют связь эмоций с мотивациями. Как отмечают Деонна и Терони (2012), оценочные теории могут указывать на зависимость того, что мы собираемся сделать, от наших оценочных суждений. Так, вы желаете быть в безопасности от тигра, когда оцениваете тигра как опасного.

Но обладают ли сужденческие теории средствами для решения двух демаркационных задач? Они могут апеллировать к оценочным суждениям и концепциям, используемым в них, для различения эмоций, утверждая, например, что страх соответствует опасному, гнев --- оскорбительному, отвращение --- отвратительному и т.д.. Здесь у сужденческих теорий, похоже, нет проблем.

Но отличить эмоции от прочих психических состояний оказывается куда сложнее. Дело в том, что оценочное суждение можно выносить и не испытывая никаких эмоций: не обязательно испытывать страх, чтобы судить об опасности. Как же тогда отличать эмоции от таких безэмоциональных оценочных суждений? Можно указывать, что эмоции также включают чувства. Питер Голди (2000, 40), впрочем, критикует такую теорию ''дополнения'' за недостаточно основательную роль чувств: ведь обычно эмоциональные чувства неотделимы от оценочного содержания эмоции, и именно благодаря чувству страха перед тигром мы осознаем его опасность. Аналогично, глядя на красный мак, мы осознаем его цвет благодаря соответствующей ему феноменологии. Этот момент, как дальше будет ясно, говорит в пользу теорий восприятия.

Стандартное возражение против сужденческих теорий --- указание на возможность противоречия между эмоциями и оценочными суждениями (Рорти 1978, Гринспен 1988, Хелм 2001). Такое противоречие происходит в случае так называемых ''непокорных'' эмоций: когда вы, например, боитесь крошечного паука, в то же время не считая его хоть сколь-нибудь опасным. Конечно, вы можете ошибаться, но конфликт есть в любом случае. И если эмоции обязательно включают оценочные суждения, то вы, похоже, страдаете острой иррациональностью. Мы регулярно испытываем непокорные эмоции, и даже если они связаны с иррациональностью, она не столь остра, как соответствующая противоречию в суждениях.

Для ответа на это возражение и были разработаны квазисужденческие теории, фокусирующиеся не на оценочных суждениях, а на состояниях, не предполагающих приверженности истинности их истинности. По мнению Патрисии Гринспен (1988), эмоции ''состоят из двух элементов: аффективных состояний комфорта или дискомфорта и оценочных утверждений, выражающих их интенциональное содержание. Так, страх можно рассматривать как дискомфорт от факта --- действительного, предполагаемого или воображаемого --- от, значит, \textit{мысли} о надвигающейся опасности'' (1988, 4). Эмоции, значит --- чувство комфорта или дискомфорта от оценочных мыслей, которые могут быть навязчивы и не восприниматься как корректные. Точно так же Роберт Робертс (2003) считает, что эмоции --- оценочные конструкции, основанные на беспокойстве о чем-то и являющиеся способом увидеть его в определенных терминах.

Преимущество таких квазисужденческих теорий в возможности конфликта: вы можете легко прийти к выводу, что тигр не опасен, но продолжать навязчиво думать о его якобы опасности. Это так же непротиворечиво, как воображать, что дождь прекратился и светит солнце, веря и даже зная, что он все еще идет.

Как отмечает Беннетт Хелм (2001), проблема такого решения в неучтенности некоторой иррациональности непокорных эмоций. Она, конечно, не так остра, как соответствующая противоречащим суждениям, но она есть. Что-то явно не так, когда вы боитесь паука, которого не считаете опасным.

Другое известное возражение против сужденческих и квазисужденческих теорий --- слишком большие требования от эмоций к когнитивным способностям (см., например, Дей 1994). Вынесение суждения предполагает владение понятиями, которые в нем содержатся. Так, чтобы судить о том, что озеро лазурное, нужно владеть и понятием озера, и понятием лазурного цвета --- например, знать, что лазурный --- разновидность темно-синего. И это же необходимо для квазисуждений. Младенцы и нечеловеческие животные, например, не обладают лингвистическими способностями, обычно коррелирующими с владением понятиями, а значит, видимо, не могут выносить суждения и квазисуждения, оценивать нечто как опасное, оскорбительное или отвратительное. Но младенцы очевидно боятся громких звуков, а животные --- такие как сурок из примера в начале главы --- хищников.

Успех такого возражения зависит от предпочитаемых оппонентом теорий суждений и концепций, что оставляет пространство для дискуссий. Возможны несколько ходов. Например, можно отстаивать минимальное понятие суждения или квазисуждения, позволяющее приписывать их младенцам и нечеловеческим животным. Так, Соломон предполагает, что телесные ощущения, связанные с эмоциями, представляют собой так называемые ''суждения тела'' (2003, 191-192, см. также Нуссбаум 2001, Робертс 2003). Но, как отмечает Скарантино, такая стратегия грозит нефальсифицируемостью результирующей теории. Другой возможный шаг --- отрицать все, или по крайней мере подобные нашим, эмоции у доязыковых существ (см. Наар 2019). Но если мы допускаем общие у нас и младенцев или животных эмоции, то мы можем обратиться к перцептивным теориям, разделяющим, возможно, достоинства оценочных и квазиоценочных подходов, и не страдающих от присущих им недостатков.

\section{Перцептивные теории}

В основе перцептивных теорий лежит интуиция о сходстве эмоций с чувственным восприятием. Сравним зрительный опыт красного мака и чувство гнева на соседа за то, что он сорвал цветок в вашем саду без разрешения. Вот важные для нас их сходства:

\smallskip

a) Наличие интенционального объекта: в случае визуального опыта это мак, в случае гнева --- сосед.

б) Репрезентативное содержание: согласно стандартным репрезентационалистским взглядам на восприятие, такое содержание имеют и визуальные, и эмоциональные переживания --- они определенным образом представляют свой объект.

в) Направленность соответствия от разума к миру, а значит, и наличие условий корректности: визуальное восприятие мака как красного корректно, если он [действительно] красный; гнев корректен если действие соседа [действительно] оскорбительно.

\smallskip

Эти первые три пункта, в отличие от следующих, согласуются с сужденческими и квазисужденческими теориями.

г) Переживания обоих рассматриваемых видов обладают феноменальными свойствами (это идет вразрез со стандартным пониманием суждений). Есть, значит ''то, каково это'' --- увидеть красный мак или испытать гнев.

д) Ни зрительные, ни эмоциональные переживания, неподвластны воле (по крайней мере, непосредственно): вы не можете увидеть красный мак просто потому что решили его увидеть, не можете и разозлиться просто потому что решили разозлиться. То есть переживания обоих типов вызываются автоматически и переживаются пассивно, хотя некоторый косвенный контроль, конечно, есть: можно прийти на маковое поле чтобы увидеть маки или намеренно сосредотачиваться на неприятном и раздражающем чтобы спровоцировать гнев.

у) И чувственный опыт, и эмоции могут противоречить суждениям: если вы носите тонированные красным очки, вы все еще можете сделать верное суждение, что мак, который видите --- белый, хотя видите его красным, так же как и можете считать, что сосед, сорвав цветок в вашем саду, вам не навредил и навредить не намеревался, и все же злиться на него.

ж) Возможность конфликта между восприятием и суждением --- одна из причин считать содержание сенсорных переживаний частично неконцептуальным (см. Крейн 1992). Так, зрительный опыт мака не требует обладания понятиями ''мак'' и ''красный''. Суждения, как состояния, включающие концептуальное содержание, подобны предложениям, состоящим из слов, а чувственное восприятие похоже на изображения, содержание которых не может целиком быть передано словами. Частичная неконцептуальность сенсорного опыта правдоподобна по нескольким причинам (см. Бермудес и Каэн 2015), в том числе потому что успешно объясняет возможность конфликта между ним и суждениями.

Так, знаменитая иллюзия Мюллера-Лайера выглядит так, будто линии имеют разную длину, и выглядят они так даже если вы уверены --- и даже знаете --- что их длины одинаковы. В этом нет ничего загадочного, если учесть, что два этих состояния --- сенсорное восприятие и суждение --- предполагают разные способы репрезентации. Неконцептуальный характер зрительного восприятия --- это то, что делает его схожим с эмоциями и объясняет предполагаемое наличие эмоций у младенцев и нечеловеческих животных (Тапполет 2020).

Этот аргумент (см. де Соуза 1987, Принц 2006, Тапполет 2012) породил множество подходов, указывающих на сходство эмоций и чувственных восприятий. Эмоции, согласно ним, относятся к оценочным характеристикам так же, как восприятие цвета --- к цветам, или, вообще, как восприятие --- к сенсорным характеристикам. Эту идею, впрочем, можно сформулировать совершенно по-разному.

Разделяют, во-первых, буквальные версии, согласно которым эмоции --- буквально виды восприятия, и небуквальные. Некоторые --- Принц (2004, 2006), например --- считают эмоции даже видом чувственного восприятия. Принц утверждает, что эмоции обладают свойствами, характерными для парадигмальных случаев чувственного восприятия, а потому должны считаться формой чувственного восприятия. Это, впрочем, слишком сильное утверждение, ведь есть существенные различия между чувственными восприятиями и эмоциями --- эмоции, например, не требуют органов чувств (де Соуза 1987). Различия между эмоциями и чувственным восприятием можно и не отрицать, но утверждать при этом, что общих характеристик у эмоций достаточно, чтобы считаться чувственными восприятиями. Так, Сабина Деринг (2003, 2007, см. также Тапполет 2016) считает, эмоции --- не чувственные, а аффективные восприятия. Согласно Деринг, аффективным восприятие делает ''его чувственное измерение, называемое также ''аффект'''' (2003, 223).

Небуквальные же версии утверждают хоть и сходство эмоций с чувственными переживаниями, но в конечном итоге кардинальное от них отличие (де Соуза 1987, Голди 2000, Хелм 2001, Митчелл 2021). Рональд де Соуза, например, пишет: ''эмоции следует рассматривать как своего рода восприятие, объектом которого являются аксиологические свойства'' (1987, 45). Но даже рассматривая эмоции как предоставляющие информацию об оценочных (аксиологических) характеристиках, де Соуза отрицает, что они буквально являются перцептивным опытом. Голди утверждает, что эмоции нужно понимать как ''чувства к'', то есть как интенциональные психические состояния, имеющие особый вид эмоционально насыщенного содержания (Голди 2000, 4, см. также Кригель 2014): ''Чувствовать страх по отношению к чему-либо --- значит думать об этом как об опасности в том особом смысле, который предполагает чувство страха'' (2000, 36). Донатан Митчелл разъясняет эту идею, указывая, что эмоциональные переживания --- это чувства предпочтения или недовольства, испытываемые нами по отношению к ценностям. Наконец, Беннетт Хелм утверждает, что эмоции --- это ощущаемые оценки: ''чувствовать эмоцию --- значит чувствовать, что цель этой эмоции имеет определенное значение, то есть радоваться, что дела идут хорошо, или огорчаться, что дела идут плохо'' (2001, 60). Таким образом, центральным небуквальным версиям теорий восприятия является понимание эмоций как квази-восприятий.

Различают также прямые и косвенные версии перцептивного подхода. Прямые говорят об эмоциях как о прямых перцептивных (или квази-перцептивных) представлениях оценочных свойств: бояться, например, значит воспринимать нечто как опасное. То есть, поскольку эмоции подвержены ошибкам в смысле несовпадения с верными суждениями --- можно бояться безобидных вещей или смеяться над плохими шутками --- то эмоции следует описывать с точки зрения перцептивного опыта. Страх тогда заключается в наличие опыта о тигре как об опасном (Тапполет 2016). Содержание этого опыта при этом неконцептуально, так что не обязательно владеть понятием тигра или опасности, чтобы бояться тигра.

Непрямые же версии опосредуют репрезентацию оценочных признаков каким-то другим типом репрезентации. Так, известная буквальная версия этого подхода, разработанная Принцем (2004, 2006), объединяет описание эмоций Джеймсом с точки зрения телесных ощущений и теорию оценки Лазаря. Согласно этой версии, эмоции --- это ''воплощенные оценки'' (Принц 2004, 51). Принц утверждает, что эмоции являются оценками потому что являются телесными чувствами, то есть восприятием телесных изменений: ''в некотором смысле эмоции позволяют нам воспринимать интересующие нас вопросы, и делают это через регистрацию изменений в теле'' (2006, 158). Центральное место в этой теории занимает концепция репрезентации, согласно которой психическое состояние репрезентирует то, что призвано надежно обнаруживать (Дрецке 1988). На основании этого Принц утверждает, что эмоции репрезентируют основные реляционные темы, то есть оценочные характеристики, потому что телесные изменения, которые мы воспринимаем, испытывая их, призваны надежно обнаруживать такие вещи как опасность, загрязнение или потерю. Опираясь на эволюционные представления об эмоциях, Принц утверждает, что соответствующие телесные ощущения были отобраны эволюцией для корреляции с такими вот вещами: страх репрезентирует опасность потому что телесные ощущения, связанные с ним, имеют биологическую функцию корреляции с опасностью, что и позволяет надежно ее обнаруживать и выживать.

Теория воплощенных ошибок Принца подвергается критике по нескольким причинам. Первая проблема --- основанность на связи эмоций с телесными ощущениями. Вторая, описанная Деонной и Терони (2012) --- связана с феноменологией эмоций. Дело в том, что, чувствуя страх, мы представляем опасность как имманентную ему, а не как нечто отдельное: ''Разве опасность не дается нам через страх именно как опасность, разве феноменология страха не есть ровно то, что дает нам этот оценочный аспект мира?'' (2012, 73). Важность феноменологии эмоций для осознаваемого посредством них принимается большинством защитников перцептивных теорий. Именно чувство страха позволяет нам осознать опасность ситуации, и именно чувство гнева демонстрирует нам обидность ситуации. Но какова бы ни была опасность, она определенно не ощущается как телесные изменения, так что вряд ли можно утверждать, что телесные ощущения позволяют осознавать оценочные свойства вещей.

Так что, похоже, прямые версии перцептивных теорий более корректны. Но насколько убедителен перцептивный подход вообще? Давайте перейдем к его буквальным версиям.

\section{Аргументы за и против буквальных перцептивных теорий}

Рассмотрим утверждение, что эмоции --- это перцептивные переживания вещей, обладающие оценочными свойствами. Ясно, что такая теория эмоций может объяснить краткосрочность эмоций и их связь с чувствами, мотивациями, интенциональными и формальными объектами и оценками. Последние три особенности очевидны. Мотивационность можно объяснить, указав на оценочность содержания эмоций.

Один из недостатков буквальных теорий --- необходимость отрицать возможность бессознательных эмоций. Перцептивные теории, впрочем, совместимы с такими эмоциями, если понимать их аналогично случаям слепого зрения, при которых испытуемые могут несознательно регистрировать цвета и формы и затем верно ''угадывать'' их (см. Вайскранц, 1986).

Похоже, что буквальные перцептивные теории обладают всеми достоинствами сужденческих, и при том не страдают от их недостатков. Необязательность концептуальной составляющей опыта позволяет им говорить об эмоциях у младенцев и нечеловеческих животных, также они легко объясняют конфликты между эмоциями и суждениями.

В том, что касается демаркации, буквальные перцептивные теории имеют столь же мало проблем, как и сужденческие --- они так же апеллируют к различным оценочным характеристикам эмоций. От прочих видов оценочных состояний эмоции отличаются ощущениями.

Все это, конечно, звучит слишком хорошо. Есть очень много возражений против буквальных перцептивных теорий. В частности, существует ряд поразительных различий между эмоциями и чувственными восприятиями (см. Салмела 2011, Деонна и Терони 2012, Брэйди 2013).

Во-первых, как уже упоминалось, эмоциям, в отличие от чувственных восприятий, не соответствуют органы. Эмоции, в отличие от таких восприятий, имеют когнитивный базис: чтобы почувствовать гнев на соседа за сорванный в вашем саду мак, вам нужно осознать, что он сделал, например, увидев, как он срывает мак --- это и будет когнитивной основой вашего на него гнева. Но чтобы увидеть, как он срывает мак, ничего предварительно не нужно. Еще одно отличие --- возможность неточности чувственного восприятия: оно должно быть вызвано воспринимаемыми объектами и их свойствами. Но вы можете сожалеть, например, о том, что динозавры вымерли, не имея с ними каузальной связи, вы можете даже испытывать сострадание к вымышленным персонажам, таким как Анна Каренина или Гарри Поттер.

И даже если и чувственное восприятие, и эмоции обладают феноменологией, эта феноменология существенно различается. Для описания того, каково видеть вазу с маками, мы рассказываем о том, что видим, то есть об этой вазе, а не о свойствах опыта как такового --- это называется прозрачностью восприятия. Пытаясь же описать каково испытывать гнев, мы можем описывать не только его причину, но и свои ощущения.

Еще одно отличие --- зависимость эмоций от таких психологических факторов, как ожидание, проекты или предпочтения. Планируя отправиться в поход, вы можете опасаться, что пойдет дождь, но беспокоясь, что сад пострадает от засухи, вы будете надеяться на противоположное --- что пойдет дождь. Более того, еще в Главе 3 мы видели, что объекты страха и гнева могут варьироваться в зависимости от социокультурного контекста, тогда как формы и цвета невосприимчивы к такого рода влияниям. То, что мы склонны видеть, едва ли меняется в течение жизни, а эмоциональные диспозиции пластичны и их, значит, можно воспитывать.

Эмоции и чувственные восприятия различаются и возможностью их оценки с точки зрения рациональности. Как уже указывалось, возможность непокорных эмоций используют в качестве аргумента против сужденческих теорий: такие теории вынуждены говорить об их явной иррациональности, как если бы переживающие их утверждали противоречивые суждения. Но как утверждает Хелм (2001, 2015), непокорные эмоции проблемны и для перцептивных теорий, ведь эмоции, противоречащие оценочным суждениям, все же в некоторой мере иррациональны. Они не подобны зрительным иллюзиям --- в иллюзиях нет иррациональности. Мы можем спросить у переживающего эмоцию, почему он ее переживает, но странно задавать подобные вопросы о чувственном восприятии.

Такое количество существенных различий, похоже, означает, что следует отвергнуть по крайней мере буквальные версии перцептивных теорий. Но не слишком ли это поспешно? Некоторые из описанных различий могут быть преувеличены, к тому же, приведенный список опирается на концепцию чувственного восприятия, которую можно и не принимать. Так, некоторые утверждают, что чувственное восприятие когнитивно проницаемо, так что зависимость от психологических факторов не уникальна для эмоций (Макферсон 2012, Сигел 2017). Некоторые различия, впрочем, верны даже для такой альтернативной концепции, так что следует еще, например, аргументировать возможность отличия чувственного опыта от чувственных состояний (Тапполет 2016). Восприятия тогда не обязательно должны зависеть от органов и могут иметь когнитивную основу, а также не ограниченны каузально, пластичны и могут объясняться и оцениваться.

Можно, впрочем, признать частичность аналогии между эмоциями и чувственными восприятиями, и говорить об эмоциях как о квази-восприятиях. В следующем разделе мы рассмотрим новую версию такого подхода.

\section{Рецептивная теория}

Итак, прямые и буквальные версии перцептивного подхода утверждают, что эмоции --- аффективные и неконцептуальные чувственные переживания вещей, имеющих оценочные характеристики, однако между чувственным опытом и эмоциями есть множество различий. Но неконцептуальные психические состояния не обязательно должны быть чувственными в строгом смысле и лежать на сенсорной периферии --- они могут располагаться на более высоком уровне когнитивной архитектуры.

Эмоции можно моделировать на основе аналоговых представлений величин, таких как аналоговые представления чисел, продолжительности или температуры. Джейкоб Бек (2012, 2019) указывает на веские основания полагать, что животные вроде голубей, крыс и обезьян, да и люди тоже, способны представлять числа аналогово, не обращаясь к явному счету. Голуби, например, могут довольно точно подсчитать, сколько раз клюнули кнопку, хотя вряд ли используют для этого концепции.

Аналоговые репрезентации можно продемонстрировать на примере чисел с вращающимися стрелками --- это \textit{аналоговые} часы, в отличие от цифровых, всегда отображающих время с помощью конкретных чисел. Аналоговые представления имеют тенденцию быть непрерывными и тем самым отражать непрерывность представляемого: время идет, и стрелки на часах постепенно поворачиваются. Аналоговые представления не чисто концептуальны потому что не допускают систематической рекомбинации своего содержания.

Как утверждает Бек (2012), неконцептуальные представления величин могут возникать и вне восприятия в строгом смысле. Это верно, если, например, вы считаете причинный контакт с объектами или их свойствами признаком восприятия, ведь представления аналоговых величин могут запоминаться и активироваться в отсутствие соответствующих стимулов, к тому же, они амодальны --- не привязаны к определенной сенсорной модальности: зрению, слуху, обонянию или вкусу. Мы можем оценить количество визуально репрезентированных точек, не имея времени на явный их подсчет, что указывает на аналоговость таких амодальных репрезентаций.

Защитники перцептивных теорий могут и должны согласиться с аналоговостью эмоций как представлений вещей через их оценочные характеристики. Вещи могут быть более или менее опасными, оскорбительными или отвратительными, и эмоции, отражая эти характеристики, также варьируются в своей интенсивности (Тапполет 2020).

Разница с буквальными подходами в удаленности эмоций с сенсорной периферии --- их нельзя, значит, считать восприятиями в строгом смысле. Подобно представлениям аналоговых величин, эмоции могут активироваться и в отсутствие причинного контакта со стимулами: мы можем испытать страх лишь вспомнив о страшном тигре или услышав, что где-то рядом есть страшный тигр. Эмоции, к тому же, не привязаны к определенной сенсорной модальности: ваш страх может быть основан на том, что видите или на том, что слышите, обоняете или трогаете, он также может быть основан на убеждении или воспоминании.

Таким образом, эмоции --- это не лежащие на сенсорной периферии аналоговые представления вещей, имеющих оценочные характеристики. Это и есть рецептивная теория.

Рецептивная теория наследует основные достоинства буквальный перцептивных теорий: эмоции должны быть репрезентациями, оцениваемыми с точки зрения корректности. А раз они касаются оценочных характеристик, то понятно, как они могут мотивировать. Рецептивная теория, значит, так же легко различает виды эмоций. Если содержание репрезентаций соответствует их формату, то эмоции имеют неконцептуальное содержание, так что и не владеющие концепциями существа могут испытывать эмоции.

Непокорные эмоции также не составляют проблемы для рецептивной теории. Даже для буквальных перцептивных теорий они создавали неудобства, ведь не может быть так, чтобы и эмоция, и суждение представляли вещи такими, какие они есть. Аналогичное неудобство составляют и сенсорные иллюзии: глядя на погруженное в воду весло, вы видите его так, как если бы оно было надломленно, хотя и знаете, что это не так. Эмоции, в отличие от чувственного восприятия, пластичны, поэтому мы можем научиться не бояться пауков, но не можем научиться видеть погруженные в воду весла прямыми (Тапполет 2016). Конечно, может иногда оказываться, что эмоции были корректней суждений, с которыми конфликтовали, но конфликт то все равно был. Рецептивная теория, значит, позволяет говорить о непокорных эмоциях, не беспокоясь о создаваемом различии между эмоциями и восприятиями.

Рецептивная теория разработана с учетом нескольких важных различий между эмоциями и чувственным восприятием. Будучи удаленными от сенсорной периферии, эмоции могут и не иметь соответствующих органов и не обязательно должны причинно соотноситься со своими объектами. Культурное и социальное влияние на эмоции и, вообще, их пластичность легко объясняются благодаря акценту на их когнитивных основах.

Рецептивная теория привлекательна и тем, что позволяет вразумительно вопрошать о причинах эмоций, ведь когнитивную основу эмоций можно рассматривать как такое их объяснение. Так, увидев тигра, вы чувствуете страх потому что увидели тигра --- этот зрительный опыт объясняет вашу эмоцию. Этот опыт может и \textit{оправдать} эмоцию (см. Главу 7).

Все это кажется весьма многообещающим, но как объяснить, что эмоции обыкновенно \textit{ощущаются}? Могут ли аналоговые представления, удаленные от сенсорной периферии, обладать феноменологией? Почему бы и нет? Но как тогда охарактеризовать эти чувства? Аффективные чувства строятся из ряда коррелирующих феноменальных свойств: вовлеченных в когнитивные основы эмоций, или возникающих в результате телесных изменений, а также возникающих в результате определенного направления внимания. Такие переживания по крайней мере частично прозрачны, ведь мы склонны описывать свои эмоции через оценочные характеристики. Как указывает Митчелл (2021), вполне естественно определять гнев как осознание чего-то раздражающего. То есть да, чувства позволяют нам осознавать оценочные свойства вещей, и рецептивная теория может утверждать, что эмоции обычно представляют собой такие ощущаемые оценки.

\section{Подведение итогов}

Сущность эмоций составляют оценочные репрезентации. Есть несколько способов изложить эту идею, наиболее привлекательный --- говорить об эмоциях как об аналоговых представлениях, имеющих неконцептуальное оценочное содержание. Эти представления во многом схожи с чувственными восприятиями, но не являются ими, потому что зависят от когнитивных основ и, значит, расположены на более высоком уровне когнитивной архитектуры. В следующих главах станет ясно, что перцептивные теории вообще и этот подход в частности, хорошо сочетаются с рядом утверждений об оценке эмоций.

\section{Краткое содержание}

\begin{itemize}
  \item\ Среди оценочных теорий выделяют несколько видов. Так, теории отношений фокусируются на этапах процесса оценки, таких как соотнесение с целью и понимание потенциала к преодолению трудностей. Защитники этих теорий расходятся в том, что считать эмоциями, но согласны, что они вызываются процессами оценки, и это утверждение совместимо и с чувственными, и с мотивационными, и с оценочными теориями.
  \item\ Реактивные теории утверждают, что эмоции --- это реакции на оценки, а чувства и восприятие ценностей отдельны от эмоций. Такие теории неспособны объяснить возможность оценки самих эмоций с точки зрения их репрезентационной корректности. Они не могут отличить эмоции от не-эмоций и учесть возможность непокорных эмоций, а также вынуждены отрицать наличие эмоций у младенцев и нечеловеческих животных.
  \item\ Перцептивные теории основаны на сходствах эмоций с чувственным восприятием, и говорят об эмоциях как об оценочных восприятиях (буквальные версии) или квази-восприятиях (небуквальные версии). Выделяют прямые и непрямые версии этих подходов, где непрямые утверждают наличие промежуточного шага.
  \item\ Теория воплощенных оценок --- буквальная непрямая версия перцептивного подхода, согласно которой эмоции, будучи телесными ощущениями, репрезентируют оценочные характеристики.
  \item\ Буквальные и прямые перцептивные теории утверждают, что эмоции --- буквально аффективный перцептивный опыт. Такие теории легко объясняют связь эмоций с чувствами, мотивациями, оценками и интенциональными и формальными объектами. Они также успешно различают эмоции. Но дело в том, что между эмоциями и чувственными восприятиями есть много существенных различий.
  \item\ Можно говорить об эмоциях как о подобных восприятиям своей аналоговостью. Эмоции тогда --- аналоговые представления, имеющие неконцептуальное содержание, но в строгом смысле восприятиями не являющиеся, потому что обладают когнитивными основаниями и тем самым удалены от сенсорной периферии. Это рецептивная теория, и она разделяет все достоинства буквальный перцептивных, но не имеет их недостатков, поскольку не обязана приравнивать эмоции и чувственные восприятия.
\end{itemize}

\begin{tcolorbox}
\section{Учебные вопросы}
  \begin{enumerate}
    \item\ Могут ли нечеловеческие животные испытывать все те виды эмоции, которые может испытывать человек? Если нет, то почему?
    \item\ Как бы вы описали чувство отвращения к испортившейся еде в холодильнике?
    \item\ Вы когда-нибудь испытывали эмоции, противоречащие вашим оценочным суждениям? Если да, то как бы вы это описали?
  \end{enumerate}
\end{tcolorbox}

\part{Нормативность и эмоции}

\chapter{Эмоции и теоретическая рациональность}

\section{Введение}

Один из самых драматичных древних мифов --- история Медеи, рассказанная Еврипидом. Медея делала все, чтобы помочь своему мужу Ясону в поисках золотого руна, но несмотря на это он предал ее и женился на другой женщине. Ярость Медеи побуждает ее совершить кровавую месть: она отравляет свою соперницу и ее отца и убивает двух своих с Ясоном сыновей. Особая трагичность заключается в ее полном осознании делаемого: ''Ужас моего зла охватываает меня, ужас того что я сделаю. Мною овладели гневные страсти --- эмоции беззакония, губящие людей'' (Еврипид 2011, стихи 1053-1056).

В таких историях эмоции --- негативные силы, которые необходимо контролировать разумом или даже полностью подавлять. Философы традиционно обвиняли эмоции в том, что они мешают разуму и тем самым способствуют разного рода иррациональности и безнравственности. Как мы видели еще в Главе 1, и Платон, и Аристотель считали, что эмоции могут порождать акратические --- совершаемые вопреки здравому смыслу --- действия. Стоики и вовсе настаивали на освобождении от большинства эмоций, происходящих, согласно ним, из ошибочных суждений.

Эмоции все еще часто понимаются как искажающие правильное мышление. В Главе 2 рассказывалось про два вида когнитивных процессов: рефлекторные, которые медленны и обдуманны, требуют контролируемого внимания, и интуитивные, которые такого внимания не требуют, возникают автоматически и происходят быстро. Интуитивные эмоциональные процессы имеют тенденцию давать только быстрые и плохо подходящие результаты, нарушая при этом процессы рефлексивные. Функция рефлексивных процессов, таким образом --- проверка и, при необходимости, подавление, интуитивных процессов.

В 1950-х гг., когда оценочные теории эмоций стали популярны, философы стали говорить о необходимости замены такой негативной концепции эмоций более позитивной. После работы философы Рональда де Соузы (1987) и нейробиолога Антонио Дамасио (1994) большинство теоретиков эмоций считают, что они не препятствуют рациональности и морали. Эмоции нужны для правильной работы мотивации и как теоретической, так и практической рациональности, они позволяют нам лучше думать и действовать.

Но ярость Медеи в приведенном примере --- сильный и неконтролируемый гнев --- плохо вписывается в предлагаемую радужную картину, страх Мэтта в примере из Главы 5 мешает ему справиться с опасностью. Речь, значит, о нормативном вопросе: плохи или хороши эмоции вообще и конкретные их виды для нас? Этот вопрос можно разбить на несколько подвопросов.

Во-первых: как эмоции влияют на наше мышление? Вроде бы очевидно, что мешают здравому рассуждению, мешая, значит, теоретической рациональности. Интересно, впрочем, что ужас, испытываемый Медией при мысли об убийстве своих детей, похоже, играет ключевую роль в ее понимании зла, которое она собралась совершить. Очень многие теоретики эмоций указывают на их центральную роль в приобретении эпистемического обоснования, знания и понимания ценностей (см. Эльгин 1996, Доринг 2003, Деонна и Терони 2012a, Сильва 2021 и др.). Аффективные состояния вообще и эмоции в частности необходимы для познания ценностей (см., например, Джонстон 2001) --- одного только разума недостаточно, чтобы полностью осознать опасность ситуации или восхитительную природу горного пейзажа.

Второй вопрос --- как эмоции влияют на нашу практическую рациональность, то есть на рациональность делаемого нами. Случай Медеи здесь хороший пример, ведь представляет конфликт между мотивацией и суждением, характерный для акратического действия. Философы, впрочем, указывают на возможность эмоций помогать быть более рациональными (см., например, Арпали 2000, Джонс 2003, Тапполет 2003).

Третий вопрос --- как эмоции вообще влияют на то, что мы делаем, будь оно связано с благоразумием или моралью. Ярость Медеи приводит к поступкам, вредящим как ей самой так и другим, то есть они плохи и сточки зрения благоразумия, и с точки зрения нравственности. Впрочем, многие философы утверждают о необходимости эмоций для принятия правильных решений и нравственности.

В этой главе мы рассмотрим, как эмоции влияют на рациональность нашего мышления, а в следующей --- как они влияют на рациональность наших действий. Изучение дебатов по этим вопросам позволит нам глубже понять, что такое эмоции.

\section{Обоснование оценочных убеждений}

Даже если эмоции не обязательно связаны с оценочными суждениями, они часто их вызывают, приводя этим к оценочным убеждениям --- диспозиционным состояниям, склоняющим к вынесению определенного рода оценочных суждений. Так, если вы чем-то восхищаетесь, то это восхищение, скорее всего, приведет к суждению или даже убеждению, что его объект его достоин. В последнее время стало популярно говорить, что эмоции не только вызывают оценочные убеждения, но и оправдывают их. Речь, конечно, об эпистемическом обосновании, связанном со знанием и, значит, истиной, а не об обосновании согласно с благоразумием или моралью.

Возможность такого обоснования зависит от предполагаемой общей эпистемологической структуры. Эпистемолог-когерентист Кэтрин Элгин, пишет: ''при отсутствии особых причин не доверять эмоциям, выводы на их основе по умолчанию приемлемы'' (1996, 150, см. также 2008). Согласно Элгин, обоснованность зависит от интеграции и вклада эмоции в целостную когнитивную систему, то есть систему, находящуюся в рефлективном равновесии, так что различные ее компоненты поддерживают друг друга. Эмоции по умолчанию разумны, но должны способствовать рефлексивному равновесию, чтобы быть полностью обоснованными (Элгин 1996, 155). Деонна и Терони (2012a) также считают, что эмоции обладают условной обосновывающей силой, но эта сила зависит не от вклада в целостную когнитивную систему, а от некоторых состояний, которые оправдывают эмоции. Так, обоснованных эмоций достаточно, чтобы обосновать связанные с ними оценочные суждения (2012a, 120). Эмоции, значит, приводят к обоснованным убеждениям в той мере, в какой эти убеждения обосновываются когнитивным базисом эмоций.

Наиболее обсуждаемое предположение состоит в том, то эпистемическая сила эмоций того же рода, что и таковая чувственного восприятия. Деринг, говоря об эмоциях как об аффективных восприятиях (см. Главу 6), утверждает, что эмоция, как и чувственное восприятие, может служить обоснованием, делая это посредством своего репрезентативного содержания (Деринг 2003, 215, см. также Милона 2016, Тапполет 2016).

Рассмотрим пример чувственного восприятия. Принято считать, что чувственные восприятия способны обосновывать убеждения. Обоснование это основного, непосредственного рода, потому что не зависит от какого-либо другого состояния, которое обосновывало бы само чувственное восприятие. Так, чувственный опыт красного мака делает веру в красноту видимого мака по умолчанию обоснованной (см. Ауди 2011). Если помимо этого опыта у вас есть убеждение, что цветок освещен красным светом, то вера эта обоснована по умолчанию, но не в конечном счете. Убеждение об освещении здесь --- дефитер от (слова defeat --- побеждать, подрывать). То есть окончательно обоснованное совпадает с обоснованным по умолчанию при отсутствии дефитеров. Но и окончательно обоснованное может не быть знанием, ведь можно ошибаться даже учитывая все имеющиеся сведения.

Можно, впрочем, утверждать сходство эмоций с чувственным восприятием, отказывая им в той же обосновывающей силе. Быть может, эмоции чем-то существенно отличаются от чувственного восприятия, что лишает их такой силы. Хотя, конечно, трудно говорить об эмоциях как о восприятиях, не наделяя их соответствующей обосновывающей силой. Возражения против такой силы эмоций обычно используются как возражения против перцептивных теорий (см. Брэди 2013).

Но обо всем по порядку: почему эмоции обеспечивают непосредственное предварительное обоснование оценочных убеждений? Если ваше убеждение в, например, отвратительности куска сыра, основано на испытываемом к нему отвращении, то оно сразу предварительно обоснованно. А в отсутствие дефитеров --- таких как вера в то, что вы нездоровы и что любая еда вызовет у вас такое отвращение --- эта вера обоснованна и в конечном счете, и для вас эпистемически допустимо верить, что сыр отвратительный. Заметьте, что дефитер здесь --- убеждение в отношении определенных условий, а не сами эти условия. Если же убеждения не соответствуют действительности, придерживаться их не стоит.

Таким образом, ваша вера в отвратительность x предварительно оправдана, если

а) Вы испытываете отвращение к x

б) Ваша вера основана на вашем отвращении

И если дефитеры отсутствуют, то эта вера окончательно обоснована (см. Тапполет 2016). Это не означает, впрочем, что убеждение истинно, ведь эмоции могут вводить вас в заблуждение.

Заметьте, что кондициональность (а не бикондициональность) здесь по той причине, что могут быть другие способы прийти к обоснованным оценочным убеждениям. Так, вы можете иметь полное право предварительно полагать, что нечто отвратительно поскольку оно на днях вызвало у вас отвращение или потому что так сказал тот, кому стоит доверять.

Как отмечает Роберт Коуэн (2016), эмоции могут раскрыть нам оценочные характеристики вещей, например чувство вины --- способ осознать моральную значимость сделанного. Тезис об обосновании, как мы видели, поддерживается и перцептивными подходами, в том числе рецептивными их версиями. То есть даже настаивая на различиях между эмоциями и чувственными восприятиями, можно утверждать, что эмоции --- ощущаемые представления оценочных характеристик. Рецептивная теория утверждает, что испытываемые нами чувства, заставляют нас осознавать оценочные свойства вещей. Так что вполне оправдано считать, что эмоции играют центральную роль в обосновании оценочных убеждений --- подобную той, которую играют сенсорные восприятия в обосновании перцептивных.

Тезис об обосновании обычно комбинируют с утверждением, что эмоции эпистемически необходимы для познания ценностей --- если бы не они, мы бы были слепы к ценностям. Марк Джонстон пишет: ''[если] на человека никогда не влияли способы вещей быть красивыми, банальными, возвышенными, ужасными или привлекательными, то и соответствующие ценности ему неведомы'' (2001, 183, см. также Деринг 2003). То есть оценочно различать можно только посредством эмоций.

В известном мысленном эксперименте Фрэнк Джексон (1982) предлагает нам представить ученую Мэри, вынужденную жить в черно-белой комнате. Мэри обладает всеми мыслимыми научными знаниями о цветах: знает нейрофизиологию зрения и физику длин волн, но поскольку она не видела никаких цветов, ей что-то все же неизвестно, ведь когда ее выпустят из черно-белой комнаты, она наконец узнает, как \textit{выглядят} цвета. Она не знала, как выглядит синий цвет, пока не увидела небо, и, значит, не могла приписывать цвета на основе своего зрительного опыта. Аналогично и никогда не испытывавший удовольствия не знает чего-то существенного.

Тезис эпистемической незаменимости предполагает, что эмоции обеспечивают непосредственное предварительное обоснование: не прочувствовав, что шутка смешная, вы не можете предварительно обосновать убеждение, что она смешная. Вам, ведь, придется полагаться тогда на информацию, предоставленную другими, или же использовать какие-то сложные принципы, соотносящие неоценочные свойства с оценочным. Эмоции же поставили бы вас в, казалось бы, лучшее положение --- благодаря им вы бы имели предварительно обоснованное оценочное суждение. Но, возможно, эмоции не единственны в такой роли. У ученой Мэри могут быть обоснованные представления о цвете неба, основанные на научных знаниях. Возможно, эмоции даже недостаточно надежны для такой роли.

Недавно было выдвинуто огромное количество возражений против обеспечения эмоциями предварительного непосредственного обоснования, не говоря уже об окончательном обосновании. В следующих разделах мы эти возражения и рассмотрим.

\section{Возражение от ненадежности}

Легко найти примеры эмоций, делающих нас предвзятыми в оценочных убеждениях. Так, согласно Питеру Голди, наши эмоции могут искажать (и действительно глубоко искажают) наш взгляд на вещи (2004, 249), заставляя нас сосредотачиваться на подтверждающем наши чувства. Например, когда мы боимся, мы склонны сосредотачиваться на том, что оправдало бы наш страх, оставляя в стороне то, что могло бы доказать нашу неправоту в предположении опасности. В широко известном исследовании Талия Уитли и Джонатан Хайдт (2005) испытуемые, которым гипнотически внушили испытывать отвращение всякий раз, когда они читают слово ''часто'', были более суровы в моральной оценке тех, чье описание содержало это слово --- то есть отвращение мешало им выносить корректные суждения.

Эмоции, конечно, могут иногда обманывать нас, но являются ли они тотально ненадежными, то есть приводят ли чаще всего к неверным оценкам? Но как проверить, что наши чувства чаще всего верны? Для ответа на этот вопрос, похоже, придется обращаться к чувственному восприятию, видя что-то как лучшее чем что-то другое. Непонятно, какие условия позволили бы избежать предвзятости. В случае чувственного восприятия мы понимаем, что способствует точности наблюдения, но с эмоциями все куда сложней: есть ли что-нибудь достаточно объективное для обоснования хоть какого-нибудь ответа на вопрос надежности эмоций? К этому последнему вопросу мы вернемся в следующей главе, а пока отложим его.

Можно предусмотреть две основные стратегии для преодоления возражения о ненадежности. Первая состоит в утверждении, что эмоции или по крайней мере какие-то их виды --- достаточно надежны, чтобы иметь обосновывающую силу. На надежность эмоций могут указать эволюционные соображения (см. Элгин 2008): если они позволяли нашим предкам противостоять опасностям, то, наверное, предоставляли достаточно точную информацию об оценочных характеристиках. Некоторые эмоции, значит, можно защитить от обвинения в ненадежности. Но надежны ли эмоции, которые мы испытываем в современности, сейчас? Наш мир существенно отличается от мира наших плейстоценовых предков, так что придется полагать, что формирование наших эмоций культурой делает их точными, но почему?

Можно снизить ожидания и указать, что нет необходимости постоянно утверждать надежность всех видов эмоций вне зависимости от контекста, ведь можно утверждать ограниченную надежность некоторых видов эмоций к некоторых случаях. Деонна и Терони (2012b) предполагают, что наша склонность к гневу, например, может быть в целом ненадежной, но вполне надежно отслеживать определенные виды правонарушений.

Более того, следует отметить, что требуемая надежность может быть вовсе минимальной. Чувственное восприятие тоже далеко на стопроцентно надежно, и мы совершаем всевозможные ошибки из-за его неточности. Но даже эта низкая надежность не значит отсутствие верной информации вообще. И так же достаточно минимальной надежности эмоций. Если так снизить планку надежности, то возражение от ненадежности станет не столь уж и весомым.

Но этих соображений, конечно, недостаточно для доказательства возможности эмоций обеспечить предварительное обоснование, так что рассмотрим вторую стратегию. Пусть эмоции как правило ненадежны. Заметим, что это не противоречит тезису об обосновании, ведь можно подвергнуть сомнению релайабилистское предположение, что только надежные источники информации могут обеспечить предварительное обоснование. Это сложные и интенсивно обсуждаемые эпистемологические проблемы, но следует указать, что многие эпистемологи говорят о необязательности зависимости обоснования от чего-то внешнего. Так, согласно эвиденциализму, наличие перцептивного опыта нечта достаточно для предварительной обоснованности существования этого нечта (см. Лайонс 2016). Эвиденциализм является интерналистским подходом, в котором обоснование обеспечивается опытом, то есть вашими состояниями, о которых вы осведомлены.

Но, возможно эвиденциалистская эпистемология не оставляет средств для объяснения надежности эмоций. Соображения надежности, впрочем, будут занимать видное место в условиях обоснования, учитывая, что для возникновения дефитера достаточно наличия оснований для ненадежности эмоций. Если у вас есть основания полагать, что ваш страх при виде паука ненадежен (вы знаете, например, что у вас арахнофобия), то ваше убеждение в опасности паука не оправдано в конечном счете. Да, вам кажется что он опасен, но этого не достаточно.

Итак, если, как вам известно, все ваши эмоции ненадежны, то дальше предварительно обоснования вы с их помощью не продвинетесь, ведь для каждого такого обоснования будет дефитер. Это, конечно, было бы печально, но не противоречило бы тезису обоснования. Но все даже и не столь плохо: в зависимости от случая у нас могут быть причины считать, что эмоция вводит нас в заблуждение, или считать ее вполне достоверной. И даже в трудных случаях мы можем определить, что противоречит нашим эмоциям: например, знание о своей арахнофобии позволяет понять, что своему страху паука вам доверять не следует.

Мы, заметьте, заранее, априорно, не знаем, какие условия представляют собой дефитеры --- скорее, именно благодаря опыту и сравнению чувствуемого с предположительно чувствуемым другими, мы можем понять, какие из условий благоприятны, а какие нет (см. Тапполет 2016). Таким образом, если вы испытываете отвращение к куску синего, заплесневелого сыра, но ваши друзья-французы говорят, что он так выглядеть и должен, то вы можете заключить, что отвращение вводит вас в заблуждение, и что этот сыр не столь уж отвратителен. То же и в случаях чувственного восприятия, когда приходится полагаться на утверждения других для проверки надежности своего опыта: войдя в теплое помещение из холода, вы почувствуете, будто в помещении очень жарко, но друзья скажут вам, что это вовсе не так и что так кажется лишь потому что вы пришли с улицы, а на улице гораздо холоднее.

Если мы позволяем эмоциям быть нашей пусть и сомнительной, но отправной точкой в размышлениях о таких вещах, как отвращение, опасность или обидность, то непонятно, что могло бы быть альтернативной такой точке. Давайте рассмотрим еще одно известное возражение против тезиса об обосновании.

\section{''Почему''-возражение}

Одно из важных различий между эмоциями и чувственными восприятиями --- вторые не допускают вопросов ''Почему?''. Можно спросить, почему вы гневаетесь на друга: что он такого сделал, чтобы вызывать этот гнев? Вы, возможно, злитесь, потому что он взял ваш велосипед без спроса. Нас, впрочем, может интересовать и нормативная сторона этого вопроса: действительно ли то, что сделал ваш друг, оправдывает гнев. Мы можем также спросить, обоснован ли ваш гнев, учитывая ваше знание его причин и возможность этого знания ошибочности. Но все эти нормативные ''Почему?'' не применимы к чувственному восприятию. Мы обычно не задаемся вопросом, что обосновывает наше визуальное переживание мака как красного, когда мы на него смотрим --- у нас просто есть опыт, и он не нуждается в каких-либо обоснованиях.

Учитывая эту асимметрию, Майкл Брэди (2013) утверждает, что эмоции не могут служить даже предварительным обоснованием, ведь, в отличие от чувственного восприятия, они не могут блокировать регресс --- не могут остановить поиск обоснования.

Брэди (2013) признает, что в некоторых случаях вопросов ''Почему?'' не возникает: чувство отвращения, вызванное куском оленины, дает достаточно оснований полагать ее неаппетитность, но обычно все куда менее благоприятно: когда вы чувствуете страх, услышав шум в квартире посреди ночи и прислушиваетесь, чтобы узнать, насколько он оправдан (грабитель это или просто муж вернулся домой из паба), то вопросы ''Почему?'' вполне актуальны и полагаться на свой страх в утверждении опасности было бы опрометчиво.

Как отмечает Джонатан Митчелл (2017), эмоции отличаются от чувственного восприятия сложностью. Так, чтобы бояться за себя, нужно иметь некоторое представление о том, что вредно именно для вас. Чувствуя эмоцию и формируя оценочное убеждение, мы часто не совсем уверены, что именно это эмоцию вызвало. Скажем, вы идете домой в темноте и вдруг начинаете чувствовать страх: непонятно, что именно вас пугает, и есть ли то, чего стоило бы бояться.

Бреди считает, что мы обычно стремимся выяснить, оправданы ли наши эмоции, потому что в нормальных обстоятельствах на непозволительно принимать их за чистую монету (Бреди 2013, 88). Клор и Гаспер (2000) предполагают, что чувства иногда воспринимаются как свидетельство убеждений, с ними согласующихся: ''свидетельства чувственных ощущений можно рассматривать как сенсорные свидетельства из внешней среды, так что то, во что верят как пропозиционально, так и эмоционально, может показаться особенно достоверным'' (2000, 25). Адам Пелсер (2014), впрочем, утверждает, что мы часто принимаем наши эмоции за чистую монету --- хоть даже нам и не следовало бы так делать.

Как уже отмечалось, эмоции, в отличие от чувственных переживаний, имеют когнитивную основу. Из-за этого может показаться, что они имеют эпистемические условия обоснования: если вы основываете свой страх на необоснованном убеждении, что где-то рядом притаился тигр, то ваш страх эпистемически необоснован. Но если о том, что где-то рядом разгуливает тигр, вам сообщил заслуживающий доверия человек, то страх уже вполне оправдан даже если тигра на самом деле рядом нет.

Учитывая это, можно утверждать, что эмоции не могут выступать блокаторами регресса: если они нуждаются в обосновании, то не могут сами быть непосредственным обоснованием --- только передавать его. Эмоции, таким образом, эпистемически излишни (см. Голди 2004, Брэди 2013): оценочные убеждения с тем же успехом могли обосновываться тем, что обосновывает соответствующие эмоции.

Уступчивый ответ состоит в согласии, что обоснование эмоций значит обоснование и оценочных убеждений, но что эмоции все же играют особую роль в обосновании вторых. Деонна и Терони (2012a) утверждают, что эмоции представляют собой привилегированный эпистемический путь к познанию ценностей --- обоснование эмоций зависит от релевантности того, на что вы реагируете эмоцией, для ее формального объекта. Если эмоция сама по себе обоснована, то она сможет обосновать и соответствующее оценочное убеждение. Поскольку эмоции позволяют нам осознавать значимость фактов, которые сложно сформулировать сознательно и которые могут быть недоступны вне ощущений, то эмоции предоставляют привилегированный путь к обоснованным убеждениям.

Можно было бы, впрочем, добавить в качестве условия и обоснованность когнитивных оснований: если вы считаете, что где-то рядом ходит тигр потому что прочитали это на Facebook, то ваш страх необоснован. Пельсер (2014) утверждает, что мы можем критиковать эмоции исходя из их основанности на таких вот необоснованных утверждениях. То есть когнитивная основа эмоций не лишена необходимости обоснования. Можно также задаться вопросом, не обязательно ли обоснование этой эмоции будет включать общее убеждение, что тигры на свободе опасны (см. Смот 2014, Эчеверри 2019). Во всяком случае, это может показаться необходимым для обоснования убеждения, что находящийся рядом тигр опасен, и если все, что обосновывает эмоции, может обосновать и соответствующее оценочное убеждение, то условия обоснования, изложенные Деонной и Терони, неполны.

Но возможно эмоции и убеждения не одинаковы в смысле обоснования. Допустимо верить только в то, во что есть веские эпистемические основания верить --- то есть убеждения должны подкрепляться чувственным восприятием или обоснованными убеждениями. Но есть ли аналогичное требование к эмоциям, действительно ли допустимо переживать только те из них, что подкреплены другими психическими состояниями? Допустим, вы видите приближающегося тигра, и это заставляет вас его бояться. Далеко не очевидно, что ваш страх нуждается в эпистемической поддержке со стороны зрительного опыта приближающегося тигра --- вы можете и не осознавать, что видите его. Ваш страх также не нуждается в поддержке со стороны каких-либо убеждений --- не требуется, например, верить, что приближающиеся тигры опасны, поскольку в большинстве случаев эмоции не вызывают таких убеждений. У вас просто есть предрасположенность бояться приближаться к тиграм и подобным им животным, и даже если допустить, что предрасположенности включают [имплицитные] репрезентации, трудно понять, как неосознаваемое состояние может иметь обосновывающую силу.

Может, эмоции и вовсе не имеют условий обоснованности? Так бы не было асимметрии между ними и чувственными восприятиями. Но как же тогда некоторым эмоциям не хватает обоснования, например, когда вы основываете их на необоснованном убеждении?

Вторая стратегия ответа на обвинение в эпистемической избыточности заключается в указании на возможность эмоций блокировать регресс даже при некоторой их открытости вопросу ''Почему?''. Так, можно отрицать, что чувства могут быть оправданы или необоснованны в том же смысле, что и убеждения. Майкл Милона (2016, см. также Кован 2016) утверждает, что в каком бы смысле эмоции ни были необоснованны, чувственное восприятие может быть необоснованно ровно так же. Так, глядя на иллюзию, вы вполне можете знать, что это лишь иллюзия, и тогда вашему визуальному опыту не хватает ''обоснования'', ведь у вас есть основания считать, что он неверен. Важно, что это совместимо со способностью перцептивного опыта блокировать регресс. Эмоции, значит, аналогично могут блокировать его, ведь так же обеспечивают предварительное обоснование. Но не слишком ли такой ход расширяет понятие обоснования?

Третий и наиболее правдоподобный ответ --- согласиться с отличием эмоций от чувственных восприятий условиями обоснования, но утверждать, что эмоции по умолчанию оправданы. Тогда наличие оснований считать свой опыт неточным не повлияет на условия предварительного обоснования. Так, осознание основанности ваших эмоций на необоснованных утверждениях будет отказывать им лишь в окончательной обоснованности, но не в предварительной. Но не будут ли они тогда эпистемически излишни? Ведь, кажется, все, что их предварительно обосновывает, должно обосновывать также и соответствующие оценочные убеждения. Такое рассуждение, однако, ошибочно: достаточное для предварительного обоснования эмоций, похоже, сильно отличается от необходимого для обоснования соответствующих оценочных убеждений.

Но не следует ли нам предпочесть другой, неэмоциональный путь к обоснованию оценочных убеждений? Конечно, если у вас есть обоснованное убеждение, что тигр на свободе и что тигры на свободе опасны, то корректно верить, что этот тигр опасен. Но как формируется убеждение, что тигры на свободе опасны? Возможно, оно индуктивно основано на конкретных представлениях об опасности конкретных тигров или подобных им животных. Можно предположить, что именно благодаря конкретным случаям эмоций, испытываемых вами или другими людьми, вы можете сделать вывод, что тигры на свободе опасны. Аналогично убеждение, что спелые помидоры красные, основано на зрительных впечатлениях, испытываемых вами при взгляде на спелые помидоры. Возможно, есть еще какой-то путь, но если эмоции хотя бы частично подобны чувственным восприятиям, то, кажется, такой путь тоже приведет обратно к эмоциям.

В любом случае, тезису об обосновании, похоже, не угрожает допущение эмоциями вопросов ''Почему?''. И поскольку эмоции по умолчанию обоснованны, они могут как обосновывать убеждения, так и быть объектом таких вопросов. Давайте рассмотрим третью линию аргументации против тезиса об обосновании.

\section{Возражение от фокуса внимания}

Психологи (см. Уэллс и Мэтьюз 1994), нейробиологи (см. Дамасио 1994) и философы (см. де Соуза 1987, Фохер и Тапполет 2002) в целом признают, что эмоции тесно связаны с вниманием --- когда вы боитесь нападающего тигра, ваше внимание сосредоточено на этом тигре. Внимание часто сравнивают с прожектором, который освещает некоторую часть сцены, оставляя остальную часть в темноте, прожектор этот можно удерживать на месте или перемещать, а также менять его сфокусированность. Уильям Джеймс (1890) указывал: ''фокализация, концентрация сознания составляют его суть. Они позволяют не учитывать одно, чтобы эффективнее взаимодействовать с другим'' (1890, 403-404). Следуя Джеймсу, внимание обычно понимают как избирательность в обработке информации: когда вы визуально фокусируетесь на маке, вы замечаете больше его деталей, чем если бы были сосредоточены на всем поле цветов перед вами.

Ожидаемо, что эмоции направляют наше внимание на объекты, представленные как имеющие определенные оценочные характеристики. Брэди (2013) утверждает, что, поскольку эмоции не только направляют внимание, но и имеют тенденцию захватывать и поглощать его, то их функция заключается не в обосновании оценочных убеждений, а в мотивации поиска доводов за и против них. В этом он видит важную эпистемическую роль эмоций: ''Одна из важных вещей, на которую способно упорное внимание --- улучшение нашего представления потенциально значимых объектов и событий через предоставление возможность обнаружить причины, влияющие на точность наших первоначальных эмоциональных оценок'' (2013, 93). Эмоции, значит, способствуют переоценке своего объекта, что позволяет перейти от грубых первоначальных оценок к более точным.

Чем же помощь эмоций в сборе информации о своем объекте противоречит тезису об обосновании? Согласно Брэди, эмоции сами по себе исключают принятие нашей первоначальной эмоциональной оценки за чистую монету: ''эмоции сами по себе [...] поднимают, а не замалчивают, вопрос обоснования, удерживая наше внимание сосредоточенным на своем объекте и побуждая нас искать соображения, обосновывающие или опровергающие эмоциональную кажимость'' (2013, 97). Страх от шума посреди ночи, таким образом, сигнализирует о \textit{возможной} опасности, вызывая состояние бдительности и позволяя вам полностью сосредоточиться на определении, есть ли в доме злоумышленник.

Но действительно ли эмоции и связанные с ними явления внимания поднимают а не замалчивают вопрос обоснования? Как отмечает Брэди (2013), смысл устойчивости внимания именно в отслеживании считающегося опасным: когда вы фокусируетесь на приближающемся тигре, вас наверняка не интересует, может ли видимость острой опасности быть ложной. Скорее, вы пытаетесь определить, что он задумал и как, значит, вам можно спастись. То есть вы уже исходите из его опасности. Брэди предвосхищает это возражение и утверждает, что рассматриваемая гипотеза совместима с мотивирующей на переоценку объекта функцией эмоций: утверждение о переоценке следует из тезиса об облегчении наблюдения вниманием. Можно, конечно, согласиться, что сосредоточенность внимания на тигре может способствовать переоценке ситуации, например, когда он уйдет, но неясно, как наблюдение за предположительно опасным может поднимать вопрос об этой уже предполагаемой опасности.

Следует также подчеркнуть, что разные эмоции сопровождаются разными моделями внимания. Когда вы чувствуете интерес к чему-то, ваше внимание сосредотачивается на этом предмете, но когда вам скучно на концерте, ваше внимание рассеивается и вы стараетесь отвлечься (Иствуд и др. 2012). Отвращение же коррелирует с концентрацией внимания (Чараш и Маккей 2002). Но, казалось бы, испытывая отвращение к чему-либо, мы переживаем конфликт между желанием наблюдать его и не наблюдать. Страх, похоже, может иметь тот же эффект. Кеннет Хугдаль и Кьелл Мортен Стормарк (2003) считают, что мы склонны активно избегать дальнейшей обработки стимулов, воспринимающихся как отталкивающие, перемещая внимание на что-нибудь другое. Наконец, даже если некоторые положительные эмоции действительно привлекают наше внимание, они делают это особым образом, поскольку имеют тенденцию расширять фокус внимания (Фредриксон и Браниган 2005, Гупта 2019). Так, если вы рады, что рабочий день закончился, вы не сосредотачиваете внимание на нем, и в голову приходят всевозможные слабо связанные между собой идеи и образы. Во всех этих случаях трудно обоснованно утверждать, что эмоции помогают собирать информацию о своих объектах, поддерживая фокус внимания на них. Скука на концерте не ставит вас в хорошую эпистемическую позицию по отношению к его деталям, равно как и отведение взгляда от чего-либо отвратительного или страшного или размышление о самых разных вещах от радости что закончился рабочий день. Таким образом, влияние эмоций на концентрацию внимания не представляет угрозы для их роли в качестве предварительного обоснования.

\section{Подведение итогов}

В этой главе мы увидели, что тезис об обосновании, следующий из перцептивных подходов, не только сам по себе правдоподобен, но и хорошо противостоит основным возражениям. Следующий вопрос, к которому мы обратимся --- хорошо ли обоснование наших оценочных убеждений на эмоциях с практической точки зрения, помогает ли это нам действовать рационально или мешает.

\section{Краткое содержание}

\begin{itemize}
  \item\ Нормативный вопрос об эмоциях делиться на несколько подвопросов: как эмоции влияют на рациональность того, что мы думаем, как они влияют на рациональность того, что мы делаем, и как они влияют на благоразумность и моральность наших действий.
  \item\ В попытке ответить на первый из этих вопросов следует разобраться, могут ли эмоции обеспечить эпистемическое обоснование оценочных убеждений.
  \item\ Из перцептивных представлений об эмоциях следует предварительная обоснованность убеждений относительно таких свойств, как опасность, отвратительность или оскорбительность. Окончательно они оправданы только если нет дефитеров.
  \item\ С теориями восприятия также хорошо сочетается эпистемическая необходимость эмоций для нашего знания оценочных характеристик. Поскольку эмоции могут быть недостаточно достоверны, оно не означает, что эмоции предварительно обоснованны.
  \item\ Являются ли эмоции достаточно надежными чтобы служить источниками обоснованности? Можно либо утверждать, что эмоции достаточно надежны, либо отвергнуть релайабилистское предположение, лежащее в основании этого вопроса, и принять эвиденциализм.
  \item\ Эмоции, в отличие от чувственного восприятия, похоже, допускают нормативные вопросы ''Почему?'', так что можно подумать, что они имеют эпистемические условия обоснования, находящиеся на одном уровне с условиями оценочных убеждений. Некоторые философы считают, что эмоции могут иметь условия обоснованности, предоставляя привилегированный путь к знанию оценочных свойств. Другая возможность состоит в непротиворечии чувства оправданности эмоций их функционированию в качестве блокаторов регресса.
  \item\ Имеют ли эмоции эпистемические условия обоснования, и если да, то какие? Согласно некоторым философам, эмоции обоснованны своей когнитивной основой, но на самом деле неясно, нуждаются ли эмоции в эпистемической поддержке со стороны прочих состояний.
  \item\ Третья возможная проблема --- о тесной связи между эмоциями и вниманием: может показаться, что, фокусируя внимание на их объектах, эмоции поднимают, а не замалчивают вопросы обоснования. В ответ можно указывать, что функция этого фокуса в помощи справиться с формальным объектом, и что совпадение интенционального и формального объекта здесь предполагается, а не ставится под вопрос.
\end{itemize}

\begin{tcolorbox}
  \section{Учебные вопросы.}
  \begin{enumerate}
    \item\ Можете ли вы поверить, что фильм хорош, на основании чьих-то отзывов, или вам нужно самим его посмотреть?
    \item\ Как кофеин и алкоголь влияют на ваши эмоции? Как вы могли бы определить, мешают ухудшают они точность ваших эмоций или наоборот?
    \item\ Почему у нас есть термометры, но нет аналогичных устройств для оценки оскорбительности или отвратительности?
  \end{enumerate}
\end{tcolorbox}

\chapter{Эмоции и практическая рациональность}

\section{Введение}

История Медеи ясно показывает, что эмоции могут заставить нас делать вещи, которые плохи как для нас самих, так и для других. Эмоции могут вызывать действия, противоречащие нашему здравому смыслу --- она в ужасе от того, что планирует сделать, но все же делает это. В более общем плане эмоции часто рассматриваются как источник акратических и даже иррациональных действий.

Практическая иррациональность касается того, что мы делаем. Точно так же, как мы должны не иметь противоречащих друг другу убеждений, мы не должны совершать и иррациональных действий. Часто указывается, что рациональность требует делать то, что считаем лучшим действием с учетом всех обстоятельств. Акратические действия, пренебрегая этим требованием, значит, практически иррациональны. Еще одним примером может быть прокрастинация. Она заключается в свободном откладывании того, что, казалось бы, имеет достаточные основания быть сделанным сейчас. Другое фундаментальное требование практической рациональности заключается в необходимости использовать те средства, которые считаются необходимыми для достижения целей.

Такие требования можно считать структурными, потому что они зависят от когерентности действий агентов: непоследовательно действовать вопреки здравому смыслу или не делать сейчас то, что, казалось бы, следует сделать сейчас. Другая группа требований связанных с рациональностью касается только внутренней согласованности. Например, требование верить только в то, что истинно. Удовлетворение этого требования зависит от того, как устроен мир, а не только от нашего о нем представления --- это, значит, яркий пример \textit{существенного} требования. Еще одно существенное требование --- действовать на основании фактов а не на основе убеждений о них.

Ни одно из этих требований не идентично требованию благоразумия: действия в соответствии со здравым смыслом могут вполне способствовать вашим интересам, и, значит, быть полезными для вас, но можно действовать исходя из полезности, но делая при этом не лучшие, с учетом всех обстоятельств, вещи. Это может быть по, например, моральным причинам. Можно также использовать средства, необходимые для достижения цели, не делая при этом то, что полезно для вас лично.

Итак, эмоции иногда вызывают акратические действия: гнев Медеи заставил ее убить своих сыновей, страх перед прыжком в озером может помешать прыгнуть чтобы добраться до другой его стороны, страх не справиться с задачей может стать причиной ее откладывания, в результате чего вы с ней действительно не справитесь. Но мешают ли эмоции вообще нашей практической рациональности или же они могут иногда помочь в ней? Некоторые философы, например Номи Арпали (2000) и Карен Джонс (2003), утверждали, что эмоции иногда могут помочь нам преодолеть практическую иррациональность. Некоторые также предполагают, что без эмоций принятие решений сильно затруднено, если вообще возможно. Так, согласно де Соузе (1987) и Дамасио (1994), эмоции играют ключевую роль в принятии решений.

Поскольку акрасия занимала центральное место в философских дискуссиях с античных времен, то сосредоточимся на роли эмоции касательно акрасии. Сначала, впрочем, мы рассмотрим так называемые арациональные действия --- экспрессивные преднамеренные, но совершаемые без какой-либо причины, действия. После попробуем выяснить, как эмоции могут вызывать такие действия, а также могут ли эмоции хотя бы иногда приводить к моральным, пусть и иррациональным, действиям, и могут ли быть акратические действия иногда не только моральными, но и более рациональными, чем соответствующие лучшему с точки зрения агента решению.

\section{Арациональные действия}

Как уже обсуждалось в Главе 5, некоторые эмоции тесно связаны с мотивацией, в том числе потому, что включают желание избежать вреда. Так, попытку убежать от опасного тигра можно считать рациональным поступком, ведь он способствует достижению цели --- избежать вреда. Но некоторые вызванные эмоциями действия, похоже, не служат какой-либо цели. Как указывает Розалинда Херстхаус (1991), примеры таких экспрессивных действий включают бег, прыжки от радости, выдергивание волос, закрытие лица руками от стыда или закрытие глаз от страха --- во всех этих случаях может показаться, что у агента нет веской причины так делать, и, значит, эмоции приводят здесь к практической иррациональности.

Как пример часто приводится Джейн, которая из ненависти к Джоан выкалывает глаза на ее фотографии. Согласно Херстхаус (1991), действие Джейн преднамеренно и добровольно, но бессмысленно. Но поскольку оно совершается не вопреки мотивам агента, то не иррационально. Намеренные действия, значит, не всегда осмысленны. Осмысленное действие подразумевает действие исходя из желания и убеждения об эффективности средств к его удовлетворению, и арациональные действия соответствуют наличию первого, но не второго.

Подобные действия, конечно, граничат с практической иррациональностью. Будучи преднамеренны, они, вероятно, совершаются в свете какого-то воспринимаемого блага --- трудно отрицать, что намеренное действие совершается исходя из достижения посредством него блага. Но что хорошего для Джейн в выкалывании глаз на фотографии ненавистного человека? Она, конечно, могла считать такое действие инструментально полезным, но ведь не обязательно. И она точно не верит, что вред, нанесенный фотографии, распространится на изображенного на ней человека. И даже если она делает это, полагая, что это наилучший способ ''выплеснуть'' гнев, причиной может быть не это. Практическая рациональность требует делать только то, для чего есть веские причины, так что подобные арациональные действия, конечно, не рациональны. И все же, действия Джейн вполне понятны, так же как и прыжки от радости, например --- это просто то, что мы делаем, когда испытываем соответствующие эмоции.

Майкл Смит (1998) указывает, что из-за ненависти Джейн хочет выколоть глаза фотографии, и средство для этого попросту совпадает с целью. То есть действие желается как таковое, а не как средство для достижения какой-то дальнейшей цели. Непонятно впрочем, почему гнев побудил Джейн сделать такую странную вещь. То есть интуитивно ее действие все же не кажется рациональным.

Согласно Питеру Голди (2000), который, как и Смит, защищает модель объяснения действий исходя из желания и убеждения, действие Джейн объясняется не обычным желанием (desire), а таким, которое, как она сама знает, не может быть удовлетворено (wish). Из ненависти Джейн хотела бы выцарапать глаза Джоан, но сделать это настоящей Джоан она боится и вообще считает это аморальным, так что делает то же с ее фотографией. Вне возможности совершить реальное действие, она совершает символическое --- "Экспрессивно действуя, Джейн символическим образом разыгрывает то, что никогда бы не сделала в действительности" (Голди 2000, 131).

Такое объяснение оставляет место для ощущения странности действий Джейн при их все же понятности. Но не слишком ли оно рационализирует эти действия? Ведь верить, что действие с фотографией повлечет то же, что и действие с изображенным на ней человеком --- явно иррационально.

Деринг (2003, см. также Бетцлер 2009) предлагает понимать такие случаи вне модели желания и убеждения, с точки зрения оценочного содержания эмоций. Деринг, как уже указывалось в Главе 6, отдает предпочтение перцептивной теории эмоций. Деринг утверждает, что эмоции ''могут рационализировать действия, ведь, подобно восприятиям, имеют репрезентативное интенциональное содержание''. Предположим, Джейн боится Джоан. Действия Джейн по ее избеганию тогда вполне понятны --- Джоан представляется Джейн опасной. Аналогично Джейн может набросится на Джоан из ненависти.

Но как перенести эти рассуждения на арациональные действия вроде выцарапывания глаз на фотографии? Согласно Дерингу (2003), действия Джейн заключаются в символическом выражении ее ненависти. Но нужно ли Джейн для этого рассматривать свои действия как хороший способ символически выразить свои эмоции? Ведь рассматривая их так, она уже стремится к некоторому благу. Если же она их так не рассматривает, то, возможно, ее действие и не преднамеренно?

Итак, нам нужен ответ, учитывающий следующую интуицию: преднамеренность действия значит что оно в некотором смысле ориентированно на благо, однако вообще Джейн его так ориентированным не считает. Можно рассматривать этот случай как конфликт эмоциональной и рациональной точек зрения. Джейн кажется, что выцарапывание глаз на фотографии --- хорошее, полезное действие, хотя она с этим не согласна. То есть она не согласна с доводом к своим действиям, но этот довод навязчиво присутствует, из-за чего действие кажется понятным, но бессмысленным.

Важное замечание Деринга --- возможность экспрессивных действий вне символической замены. Херстхаус (1991) указывает, что экспрессивные действия весьма разнообразны: потливость и дрожь, например, непроизвольны, а сжимание ладоней в кулаки может быть и непроизвольным, и произвольным. Некоторые действия делаются намеренно, но именно потому что во власти эмоций мы хотим их сделать. Некоторые действия можно объяснить с точки зрения оценочных характеристик, а в случае конфликта могут возникнуть арациональные действия. Далее мы рассмотрим, как эмоциональны действия могут противоречить суждению о правильном действии с учетом всех обстоятельств.

\section{Эмоции и акрасия}

Акратические (от греческого kratein, что означает контроль и привативное ''a'') или безвольные действия обычно подразумевают конфликт между делаемым нами и тем, что, как считаем, должны делать: например, когда вы выпиваете еще один бокал вина, хотя считаете, что, учитывая все обстоятельства, следует от него воздержаться --- что этот бокал приведет к сильной мигрени.

Различают свободные и несвободные акратические действия. Первые значат намеренное действие вопреки здравому смыслу, вторые не столь намерены. Убийство Медеи --- пример несвободной акрасии, ведь, судя по ее словам, гнев побудил ее совершить убийство против ее воли. Медея в некотором смысле не была собой, когда действовала согласно своему гневу. Это касается вопроса соотношения эмоций и автономии, к которому мы еще вернемся.

Возможность свободно и намеренно действовать вопреки здравому смыслу создает ряд загадок о природе намеренного действия и связи между практическим суждением и действием (см. Штрауд 2008). Как мы можем свободно и намеренно выбирать образ действий, считаемый нами же неполноценным? Лишь сравнительно недавно, начиная с конца 1960-х годов, такая возможность стала допускаться (см. Дэвидсон 1969). Сейчас большинство философов согласятся с ее наличием. Этот тип акратического действия можно определить так:

\smallskip

(Строгая акрасия) Действие является строгим акратическим тогда и только тогда, когда (а) оно является свободным и преднамеренным и (б) при его выполнении агент считает, что существуют значимые причины его, учитывая все обстоятельства, не выполнять

\smallskip

Следует отметить, что лучшее (учитывая все обстоятельства и цели агента) решение агент принимает именно в момент действия, и это отличает строгую акрасию от другой формы практической иррациональности --- слабости воли, как когда, например, просыпаясь, вы решаете все же не идти на пробежку (см. Холтон 2009).

Как эмоции могут вызвать строгие акратические действия? Возможно, они включают в себя мотивацию, вызывающую это действие. В случае с Медеей, похоже, причиной ее злого поступка является желание мести, сопровождающее ее гнев и заставляющее ее действовать против своей воли. Но эмоции не обязательно должны включать такую непреодолимую мотивацию. Сильный гнев может мотивировать действовать и не против своей воли. Кажется, может быть желание мести, не лишающее свободы, но достаточно сильное, чтобы мотивировать к соответствующим действиям.

Возможно, эмоции способствуют акрасии через управление вниманием. Согласно де Соузе, привлекая внимание к определенным соображениям, эмоции мешают нашим практическим рассуждениям, так что мы приходим к иррациональному выводу, что именно акратическое действие следует совершить, полагая при том, решение это все корректно учитывает и, значит, рационально. Эмоции, указывает он, ''идеально приспособлены для роли арбитров среди доводов'', ведь ''их основная роль заключается в установлении паттернов значимого в доводах''. То что эмоции так влияют на наши решения --- вполне правдоподобно, это также совместимо с ролью эмоциональных мотивов в мотивировании акратических действий.

Следуя Дональду Дэвидсону (1969), де Соуза предполагает, что акрасия обязательно включает в себя ошибку рассуждения, происходящую еще до действия. Это, впрочем, неочевидно. Акрасия конечно может, но не то что бы должна, содержать такую ошибку. Думая, стоит ли выпить еще один бокал вина, вы тщательно обдумываете все свои доводы за и против и приходите к выводу, что следует воздержаться, однако все же выпиваете еще один бокал. Между рассуждением и мотивацией (и именно между ними), значит, лежит акратический разрыв. Это, впрочем, не отменяет роли эмоций в возникновении такого разрыва. Привлекая ваше внимание к некоторым конкретным аспектам ситуации, эмоция вполне может повлиять на мотивацию, вне зависимости от вывода из рассуждений.

Роль эмоций в акрасии, возможно, не только каузальная. Как мы видели при рассмотрении рациональных действий, составленность эмоций из оценочных представлений позволяет объяснить действия с точки зрения представленных оценочных свойств. Оценочные свойства, будучи связаны с практическими доводами, могут объяснять ваши действия. Так, страх при мысли о прыжке в озеро представляет этот прыжок опасным. Опасность здесь --- причина, пусть и не единственная, от таких прыжков воздерживаться. Страх, значит, сообщает вам причину, по которой вам не следует прыгать. Если вы воздерживаетесь от прыжка, несмотря на рациональный вывод прыгнуть, то это объясняется его сообщенной страхом опасностью. Это, конечно, не значит, что не прыгать --- рационально, ведь учитывая все обстоятельства, вы пришли к выводу, что прыгать надо.

Акрасию можно сравнить с непокорными эмоциями. Непокорность значит конфликт между эмоциями и конкретным оценочным суждением, например когда вы боитесь паука, считая что он вовсе не опасен. В случае же акрасии мы имеем дело с конфликтом между чувствуемым и результатом оценки с учетом всех обстоятельств. Вы, скажем, боитесь опасности и не совершаете действие, которое, считаете, рационально было бы совершить. При страхе вы в любом случае будете склонны испытывать общее негативное чувство по отношению к действию и будете неохотно его совершать.

Итак, эмоции являются значимыми триггерами акратических действий. Но ограничивается ли их влияние на наше поведение созданием мотивации для таких действий?

\section{Обратная акрасия}

Акратические действия всегда представляли особый интерес для философов, поскольку считалось, что они связаны не только с практической иррациональностью, но и с неблагоразумными и, прежде всего, аморальными, действиями. Часто рассматривались случаи акрасии, когда агенты действуют вопреки своим моральным суждениям. Акрасия тогда рассматривалась как одна из главных источников безнравственных поступков. Можно ли свободно и намеренно действовать вопреки суждению о моральной неприемлемости такого образа действий? Р.М.Хэйр (1952), например, отрицает это. Он утверждает, что за моральными суждениями обязательно следуют соответствующие действия, и что это концептуальная истина. То есть если вы не смогли действовать в соответствии со своими моральными суждениями, то это потому, что вы, как и Медея, не были свободны так действовать, ну или потому что на самом деле не считали что следует так поступать.

Относительно недавно акцент сместился на случаи акрасии, мотивирующие действовать вопреки ошибочному суждению. Если вы ошибочно решите делать что-нибудь опасное, но из-за страха остановитесь, то сделаете разумный выбор. Так же и в случае действий, совершаемых вопреки ошибочному морально суждению, акратическое действие вполне может быть моральным, будучи иррациональным. Такие случаи соответствуют ''обратной акрасии''.

В пример обычно приводят случай с персонажем Марка Твена Гекльберри Финном (см. Беннет 1974, Арпали и Шредер 1999). Гек помог рабу Джиму сбежать от своего хозяина на плоту по реке Миссисипи. Когда Джим уже вот-вот будет свободен, у Гека возникают сомнения: сдать Джима властям, думает он, будет правильным поступком. Но когда появляется такая возможность, Гек обнаруживает, что не может сдать Джима и лжет охотникам за рабами, чтобы защитить его. Джонатан Беннетт (1974), интерпретирует этот случай как конфликт между ошибочным моральным суждениям и сочувствием. Моральные принципы Гека, усвоенные им без всякой критики в сельской местности штата Миссури и согласно которым рабы являются всего лишь еще одним видом собственности, глубоко ошибочны, и чувство симпатии к Джиму заставляет его действовать морально.

Гек здесь описывает себя как слишком слабого человека --- что он не может следовать совести и действует несвободно. Можно сказать, что симпатия к Джиму заставила Гека действовать против своей воли. Но в отличие от случая с гневом Медеи, странно говорить о сочувствии как такой одолевающей и лишающей свободы эмоции. Можно сказать, что Гек был волен поступить иначе, но решил не следовать своей совести --- то есть решил действовать вопреки своему здравому смыслу --- из симпатии к Джиму.

Обычной реакцией на эту историю является моральная похвала не только действия Гека, но и его самого за им сделанное. Арпали и Шредер (1999) утверждают, что обратная акрасия несовместима с рационалистическими представлениями о самости. Они отмечают, что ''читатель романа Твена склонен рассматривать поступок Гекльберри не как случайное доброе дело, совершенное аморальным человеком, а как показатель этого человека моральности'' (1999, 163). Если мы готовы похвалить Гека за то, что он сделал, то мы не можем в то же время считать, что сделанное им было ему навязано и совершенно не соответствовало его воле, ведь мы считаем его источником его действий.

Интересно, что в отличие от рационалистического подхода к самости, согласно которому ответственные действия проистекают из здравого суждения агента, удобный для принятия здесь подход допускает свободное и, значит, ответственное, действие согласно эмоциям. Согласно Арпали и Шредер (1999), любое описание личности, отождествляющее ее с ее рациональными способностями, не может объяснить случай Гека. Действия, являющиеся результатом эмоций, а не здравого смысла агента, можно считать автономными оттого что их источник находится в самом агенте, а не вовне (см. Тапполет 2016). Автономия здесь понимается как самоуправление: вы автономны если определяете свои мысли и действия, а не находитесь во власти внешних по отношению к вас сил, таких как угрозы и принуждения (см. Басс и Вестлунд 2018).

Итак, акрасия значит только иррациональность, но не обязательно ошибочность действия: когда рациональное суждение ошибочно, эмоции могут привести к благоразумным действиям. Но обязательно ли действия вопреки ошибочному решению соответствуют практической иррациональности?

\section{Рациональная акрасия}

До недавнего времени считалось самим собой разумеющимся, что действовать вопреки здравому смыслу --- иррационально. Однако по мнению ряда авторов, таких как Элисон Макинтайр (1990, см. также Арпали 2000 и Тапполет 2003), акратические действия не обязательно должны быть иррациональными, во всяком случае они могут быть рациональнее, чем некоторые действия, основанные на здравом смысле.

Несмотря на то, что действия Гека противоречит его здравому смыслу, оно основано на его практических соображениях. Его суждение не только ложно, но и пренебрегает важными соображениями, которые Гек мог и должен был учитывать. То есть, подумай он получше, он бы понял, что Джим был его другом, доверял ему и что его стремление к свободе вполне законно. Приняв это во внимание, Гек бы увидел, что должен помочь Джиму. Гек, к сожалению, не учел все это, и потому его здравый смысл не отображает всех должных соображений. Но сочувствие позволило ему таки не выдать Джима. Эмоции, значит, помогли Геку принять во внимание доводы, которыми он пренебрег, сделав его более рациональным, чем если бы он действовал в соответствии со своим здравым смыслом.

Арпали (2000) предлагает рассмотреть аналогичный случай: чувства Эмили, студентки химического факультета, заставляют ее отказаться от программы химии вопреки здравому смыслу. Она считает, что чувства эти беспочвенны, но на самом деле они вызваны факторами, являющимися вескими причинами задуматься о смене направления: например, химия не позволит ей в полной мере реализовать свои таланты. Как пишет Арпали, ''Эмили [...] действует гораздо более рационально, решаясь сменить направление обучения, ведь у нее не просто есть веские причины это сделать --- она по этим причинам и действует'' (2000, 504-505). Арпали сравнивает случай Эмили со случаем Алисы, схожие чувства которой вызваны заниженной самооценкой --- Алиса действует дважды иррационально, решаясь сменить направление, ведь ее убеждение о наличие у нее на это веских причин ложно. Единственная же иррациональность Эмили --- ее неспособность отказаться от убеждения, что ей следует продолжить образование химика. Карен Джонс пишет: ''эмоции агента могут быть привязаны к соответствующим причинам таким образом, что позволят ему отслеживать их даже когда здравый смысл на это неспособен'' (2003, 186).

Согласно Макинтайру (1990), такие возможности возникают потому, что агенты не обязательно ясно видят все свои причины для тех или иных действий, и потому в своих размышлениях не принимают во внимание все такие практические доводы. Если понимать практические причины как внешние факты, то конечно в суждениях о них можно ошибаться: если моральные факты представляют собой практические причины и выдавать Джима морально недопустимо (это моральный факт), то конечно насчет этого можно ошибаться --- Гек, значит, просто не смог учесть этот моральный факт в своих рассуждениях. Однако, как утверждает Макинтайр (1990, см. также Арпали 2000), можно также следовать Бернарду Уильямсу (1981) и предположить, что практические причины внутренни в том смысле что зависят от мотивации агента, но что мы все же можем в отношении них ошибаться. То есть мы можем ошибаться в отношении наших мотиваций и доводов к действиям, или же иметь ложные убеждения о средствах, которые для достижения наших целей следует использовать, и эмоции могут здесь помочь. Согласно Макинтайру, ''чувствительность агента к доводам может в некоторых случаях превосходить его более интеллектуальную способность приходить к ним в рассуждениях'' (1990, 390).

Конечно, возможность действий вопреки здравому смыслу быть более рациональными, чем здравому смыслу соответствующие --- спорный вопрос. С точки зрения агента правильно было бы следовать здравому смыслу. Заманчиво проследовать на Дерингом и заявить: ''практическую рациональность [...] следует понимать с точки зрения субъективных доводов: агент рационален в той степени, в которой руководствуется своими такими доводами'' (2010, 287). Субъективные доводы здесь --- те, которые агенты уверены что имеют, то есть руководство такими причинами и есть руководство своим здравым смыслом.

Требует ли практическая рациональность от агента следовать своему лучшему суждению или требование ее в том, чтобы действовать исходя из практических доводов? Должны ли мы делать то, на что, согласно нашим рассуждениям, есть веские причины, или же следует учитывать и мотивации вне рассуждений? Нормативная сила первого требования можно считать следующей из второго (Колодный 2005). Чтобы быть рациональным, а не только таковым себе казаться, лучше не следовать иррациональному лучшему суждению. Следование своим эмоциям, таки образом, в некоторых случаях может оказаться более надежным путем к практической рациональности, чем следование здравому смыслу.

Итак, эмоции могут быть причиной далеко не только хаоса иррациональности --- в некоторых случаях они также могут позволить нам действовать более рационально. В следующем разделе мы обсудим роль эмоций в эффективном принятии решений.

\section{Эмоции и принятие решений}

Как мы уже неоднократно замечали, эмоции тесно связаны с вниманием: когда вы чего-то боитесь, ваше внимание сосредотачивается на объекте страха и возможных путях защиты от него. Согласно де Соузе (1987), это позволяет эмоциям играть решающую практическую роль. Разум как таковой не способен определить, какая информация важна для принятия решения: нам нужно знать, как из всей имеющейся информации извлечь только необходимое. Де Соуза называет это ''проблемой рамок'' (1987, 193). Эмоции, утверждает он, являются естественным решением такой проблемы. Робот, продолжающий анализировать многочисленные последствия своих действий в ситуации, когда следует срочно обезвредить бомбу, будет уничтожен именно из-за недостатка страха.

Эмоции способны решить проблему рамок именно благодаря своей связи с вниманием --- они ''ограничивают диапазон информации, которую организм будет принимать во внимание'' (1987, 195). Эмоции так позволяют ''контролировать решающий фактор значимости среди того, что в противном случае было бы неуправляемым множеством объектов, интерпретаций и стратегий мышления и поведения'' (1987, 249). Такова, считает де Соуза, биологическая функция эмоций.

Рами Маджид (2019) указывает, что роль эмоций в определении значимости информации обеспечивается их оценочным компонентом: ''Эмоции ощущаются хорошими или плохими, и в силу этого их интенциональные объекты также представляются как хорошие или плохие, как положительные или отрицательные'' (2019, 58). Наше внимание привлекается к частям информации, представленным как положительные или отрицательные, представленное же ''нейтрально'' игнорируется.

Схожую гипотезу предлагает Дамасио (1994): эмоциональные чувства функционируют как ''соматические маркеры'', которые ''путем обучения связаны с предсказанием будущих результатов определенных сценариев'' (1994, 174). Соматические маркеры помогают нам в размышлениях, выделяя варианты как хорошие или плохие и фокусируя внимание на отрицательном или положительном их результате, так что число рассматриваемых вариантов ограничивается. Дамасио пишет: ''соматическое состояние, негативное или позитивное, вызванное появлением данной репрезентации, действует не только как маркер ценности репрезентированного, но также и как усилитель непрерывной рабочей памяти и внимания'' (1994, 197-198). Соматические маркеры также облегчают анализ затрат на варианты и выгод от них: положительные чувства побуждают нас отдавать предпочтение соответствующим вариантам, а негативные побуждают соответствующих избегать. Вполне вероятно, что эмоции играют центральную роль в оценке наших возможностей. Согласно Брейди --- ожидание стыда или удовольствия, комфорта или дискомфорта --- ''ожидание таких вещей является центральным фактором в [нашем] присвоении полезности вариантам'' (2018, 77).

Дамасио приводит исследования пациентов, перенесших поражения головного мозга в префронтальной доле, имеющей решающее значение для обработки эмоций. Особенно показательным является случай Эллиота. Эллиот был умным, квалифицированным и трудоспособным человеком, имел завидным личный и профессиональный статус, но после операции по удалению опухоли головного мозга не смог больше работать и поддерживать отношения со своей семьей. Он не мог больше правильно распределять свое время и принимать разумные решения --- ''его способность принимать решения была нарушена, как и его способность составлять эффективный план на предстоящие часы, не говоря уже о планировании на месяцы и годы своего будущего'' (Дамасио 1994, 37). Удивительно, но в его прочих способностях не произошло никаких изменений: он все еще был способен очень хорошо решать задачи на рассуждение, в том числе морального и практического характера. Но даже рассмотрев несколько подходящих вариантов, он все равно не мог решить, что делать. Эта трудность, по мнению Дамасио, возникла из-0за изменения его чувств: то, что когда-то вызывало сильные положительные или отрицательные эмоции, больше не вызывало никаких эмоций вообще. Отсутствие эмоциональных реакций не позволяло Эллиоту присвоить ценность рассматриваемым им вариантам.

Биологическая эволюция и обучение определяют, какие эмоциональные предубеждения будут нами руководить. Кто-то реагирует страхом даже на безобидных собак, а кто-то не боится даже крупных. Некоторые эмоциональные предубеждения могут приводить к действиям, не имеющим веских причин, и, значит, иррациональным. Но, если указанное де Соузой и Дамасио верно, то иррациональность, порождаемая эмоциями --- цена любого решения. Без эмоций мы не смогли бы выбирать, так что ошибочно считать, что эмоции должны быть исключены из наших практических рассуждений.

Джон Элстер (1999) выступает против такой картины, обвиняя де Соуза и Дамасио в карикатурности их понимания теории рационального выбора. Эта теория, утверждает он, не предполагает рассмотрение рациональным агентом всех возможных результатов всех возможных вариантов. Агенты могут полагаться на рациональные методы выбора соответствующей информации, а не на эмоции. Впрочем, похоже, что мы обычно полагаемся именно на эмоции. Следует также отметить, что вопрос интерпретации гипотезы о зависимости принятия решений от соматических маркеров --- крайне спорный (см. Бартол и Линквист 2015). В любом случае, будучи связаны с оценками, эмоции должны играть центральную роль в определении вариантов действий агента.

Итак, эмоции, как правило, выполняют функцию определителя значимости, облегчая принятие решений за счет сокращения числа рассматриваемых вариантов. Но эмоции по-разному влияют на внимание. Как уже отмечалось в Разделе 5 Главы 7, скука отвлекает наше внимание, а ощущение счастья приводит к расширению его фокуса. Из-за разнообразия воздействия на внимание, влияние эмоций на принятие решений довольно сложно --- оно не может состоять в просто-таки сосредоточении нашего внимания на объекте эмоции.

\section{Подведение итогов}

Мы начали с вопроса оправданности традиционной концепции эмоций как источников практической иррациональности. Рассматривая экспрессивные действия --- преднамеренные, но беспричинные, то есть арациональные --- мы увидели, что эмоции и соответствующие им оценочные репрезентации могут объяснять наше поведение. Арациональные действия, значит, можно понимать как конфликт между репрезентацией, создаваемой эмоциями, и здравым смыслом. Такова, по-видимому, структура акратических действий.

Случаи акратических действий показывают, что эмоции могут вызывать практическую иррациональность за счет своей способности к мотивации и влиянию на внимание. Но эмоции могут иметь и положительный эффект --- они могут привести к акратическим действиям, которые лучше, чем соответствующие здравому смыслу. Эмоции могут заставить действовать более соответствующим практическим соображениям агента образом, делая результирующее поведение рациональней соответствующего здравому смыслу.

Роль эмоций в производстве практической рациональности, таким образом, неоднозначна --- иногда они приводят к иррациональности, но иногда помогают ее достичь. В любом случае, эмоции играют ключевую практическую роль --- сосредотачивая наше внимание на важном и отмечая варианты действий как хорошие или плохие, они компенсируют недостаток разума как такового.

\section{Краткое содержание}

\begin{itemize}
  \item\ Арациональные действия, такие как выкалывание глаз на фотографии ненавистного человека --- экспрессивные действия, преднамеренные но беспричинные, поскольку их нельзя объяснить с точки зрения желаний или убеждений агента. Поскольку эмоции имеют репрезентативное содержание, можно предположить, что они помогают объяснить такие действия.
  \item\ Эмоции также могут вызывать акратические действия --- выполняемые вопреки здравому смыслу агента. Такие действия не обязательно несвободны. Имея репрезентативное содержание, эмоции объясняют акратические действия с точки зрения оценочных свойств.
  \item\ Эмоции могут предоставить более точные суждения, не основанные на доводах агентов, и позволить им действовать в соответствии с действительно имеющимися у них доводами.
  \item\ Эмоции, возможно, играют ключевую роль в практической рациональности, позволяя нам делать выбор --- они фокусируют наше внимание на важном для принятия решения, сводя бесконечность возможных вариантов до конечного числа. Эмоции отмечают варианты как хорошие и плохие, что неспособен сделать разум сам по себе.
\end{itemize}

\begin{tcolorbox}
  \section{Учебные вопросы.}
  \begin{enumerate}
    \item\ Эмоции часто считают причиной акратических действий. Но могут ли эмоции помочь придерживаться здравого смысла?
    \item\ Что бы вы сделали на месте Гека, и какие действия, на ваш взгляд, были бы правильными?
    \item\ Как эмоции влияют на другие формы практической иррациональности, такие как прокрастинация?
  \end{enumerate}
\end{tcolorbox}

\chapter{Сентиментализм}

\section{Введение}

Эмоции, очевидно, связаны с оценкой --- страх, отвращение и стыд связаны со страшным/опасным, отвратительным и постыдным соответственно. Этими свойствами мы характеризуем вещи как хорошие или плохие: отвратительное это что-то плохое, восхитительное --- хорошее.

Более того, как было показано в Главе 7, эмоции, похоже, могут предварительно обосновывать оценочные убеждения. И если эмоции репрезентируют оценочные свойства, то воздействие на первые можно объяснить с точки зрения вторых, в свете которых и совершается действие.

Являются ли оценочные свойства объективными, наравне с формой или массой? Если да, то каково их отношение к натуральным свойствам, постулируемым естественными и социальными науками? Возможно, оценочные свойства зависят от нашей субъективности так же, как цвета зависят от нашего их восприятия. Как станет ясно, эту аналогию можно понимать по-разному.

Прояснение природы оценочных свойств может вынудить нас пересмотреть отношение между эмоциональным и оценочным. Что, если оценочные --- наиболее, казалось бы, многообещающие --- теории привержены ошибочной концепции оценочных свойств как объективных. Без понимания природы оценочных свойств у нас не будет и полноценного понимания природы эмоций.

Для ответа на вопрос о природе оценочного нам придется обратиться к метаэтике --- части моральной философии, занимающейся наиболее абстрактными ее вопросами, такими как вопрос о моральной объективности и истинностности моральных суждений. Ряд метаэтических теорий подчеркивают эмоции в своем понимании оценочного --- моральные факты, концепции и суждения так или иначе связаны с аффектами вообще и эмоциями в частности. Метаэтические дебаты в значительной степени строятся вокруг противостояния сентиментализма и рационализма. Моральный сентиментализм восходит к Юму (1739-1740/2000) и утверждает чувства источником морали. Моральный же рационализм связан с Кантом (1785/1996) и утверждает основой морали рациональность.

Недавно стало ясно, что противоречие между сентиментализмом и рационализмом не следует переоценивать: подчеркивание роли эмоций в этике вполне совместимо с истинностностью моральных суждений. А некоторые сентименталистские теории, говорящие о моральных утверждениях как о выражающих лишь эмоции, понимают, похоже, эмоции как просто чувства.

Первый раздел этой главы разъясняет понятие сентиментализма, далее мы обсудим различные теории оценочных свойств, начиная с простого субъективизма (Принц 2007) и заканчивая теориями идеального наблюдателя (Льюис 1989, Кауппинен 2014). Затем мы сосредоточимся на неосентименталистских взглядах, восходящих к Францу Брентано (1889/2009) и до сих пор актуальных как в редукционистских, так и в не-редукционистских (известных также как теории чувствительности) своих версиях.

\section{Виды сентиментализма}

Один из вопросов, традиционно интересовавших философов --- вопрос о природе ценности. Термин ''ценность'' имеет несколько значений, так что вопрос можно понимать по-разному. Говоря о ценности, часто говорят о таких благах как удовольствие, свобода или справедливость. Затрагиваются при этом также идеалы, направляющие наши действия. Например, о демократии и автономии можно говорить как о западных ценностях, а надежность и честность считать чьими-то личными ценностями.

Ценности в смысле благ или идеалов довольно интересны сами по себе, но метаэтические дебаты касаются ценностей в несколько другом смысле. Метаэтическая теория пытается охватить фундаментальный вопрос о природе приписываемого нами объекту высказываниями о его хорошести, плохости, достойности восхищения, отвратительности и т.д.. Действительно ли это приписываемое существует и какова его природа?

Нормативные концепции обычно разделяют на два больших семейства: оценочные и деонтические. Для деонтического семейства центральными понятиями являются обязательное, дозволенное и запрещенное. Деонтические концепции используются в суждениях относительно должного --- будь то с моральной или с, например, юридической точки зрения. Оценочные же концепции используются для оценки и не касаются должного напрямую. Тем не менее, хорошесть или плохость зависит от от того, что нам следует делать --- согласно любым моральным теориям, должное хотя бы иногда зависит от последствий. Так что, похоже, деонтические и оценочные концепции тесно связаны.

Моральный сентиментализм касается только морального подмножества оценочной и деонтической области. Но что можно считать моральным? Распространено предположение, что мораль связана с тем, что мы должны другим или с тем, что для них хорошо или плохо, но не с тем, что хорошо или плохо для нас. Вопрос о наличии моральных обязательств перед самим собой довольной спорный. В любом случае, оценочные свойства не обязательно должны быть моральными: можно быть достойным морального восхищения за спасение тонущего ребенка ценой риска собственной жизнью, но можно вызывать восхищение и вовсе не связанными с моралью способами --- будучи великим пианистом, например. То есть сентиментализм в отношении оценочного не то же, что моральный сентиментализм.

Метанормативные вопросы вообще и метаэтические вопросы в частности принято разделять на четыре группы. Например, вопросы о природе моральных фактов. В качестве такого факта можно рассмотреть "боль --- это плохо". Моральная онтология отвечает на вопрос объективности этого факта, его независимости от наших реакций и суждений. Так, согласно моральному реализму, моральные факты могут быть объективными, моральный антиреализм же это отрицает. Моральный конструктивизм значит, что объективность моральных фактов --- результат процесса субъективного или интерсубъективного конструирования. Вопросы онтологии связаны с вопросами семантики, но не тождественны им. Один из вопросов --- отличие моральных утверждений от описывающих положение вещей. Вопросы о функциях морального языка не следует путать с вопросами о природе моральных суждений, понимаемых как психологические состояния, а не суждения, а также о концепциях, подразумевающих такие суждения. Моральная психология занимается вопросами моральных суждений: например, обязательно ли они связаны с мотивацией, то есть значит ли уверенность в должности нечта на это нечто мотивированность. Связь моральных суждений с мотивацией значит невозможность их вынесения без этой мотивации, и, значит, невозможность акратических действий. Еще одна группа вопросов касается возможности знания о морали и этого знания обоснования --- ими занимается моральная эпистемология.

Сентиментализм принимает самые разные формы в зависимости от того, к какой конкретной части метаэтик он относится (см. обзор Кауппинен 2022). Так, касательно моральной онтологии сентименталист считает, что моральные факты так или иначе зависят от аффектов. Например, что моральные факты --- проекции наших моральных установок, иначе говоря, что в отличие от разума ''вкус обладает продуктивной способностью и окрашивает все объекты цветами, заимствованными из внутренних чувств'' (Юм 1751/1948, 269, см. также Блэкберн 1998). Семантический же моральный сентиментализм говорит о значительной роли эмоций в моральной семантике. Экспрессивисты вообще и эмотивисты в частности, такие как А.Дж.Айер (1936), Чарльз Стивенсон (1937), а в последнее время Саймон Блэкберн (1998) и Аллан Гиббард (1990), например, утверждали, что моральные суждения находятся на том же уровне, что и междометия, то есть выражают субъективные реакции, а не истинные или ложные утверждения. Некоторые же утверждают, что моральные суждения вызваны эмоциональными реакциями (Хайдт 2001) или состоят из аффектов вообще и эмоций в частности (Блэкберн 1998, Принц 2007) --- это один из видов психологического морального сентиментализма. Эпистемический же моральный сентиментализм утверждает зависимость морального знания и/или его обоснования от эмоций.

Каждый из этих теоретических вариантов может быть защищен как в отношении нормативности в целом, а не только в рамках оценочного или деонтического подхода. Так, эмоции, как мы видели в Главе 7, можно рассматривать как источник эпистемического обоснования, и, возможно, оценочного знания --- согласие с этим значит принятие эпистемического сентиментализма. Эпистемический сентиментализм, впрочем, не значит сентиментализм в отношении ценностей. Можно, например, считать, что эмоции позволяют нам войти в контакт с независимо существующей оценочной реальностью, не имеющей, однако, ничего общего с нашими реакциями. Таким образом, несмотря на некоторое сходство, различные формы сентиментализма логически независимы.

Давайте теперь сосредоточимся на одном конкретном типе сентиментализма --- в отношении оценочных фактов и свойств, которые они затрагивают. Речь, значит, пойдет об онтологии оценочных фактов --- моральных и не только.

\section{Простой субъективизм}

Как мы видели, оценочные теории эмоций говорят об эмоциях как о включающих репрезентации вещей как наделенных определенными оценочными свойствами или ценностью. Рассмотрим пример восхищения. Из-за лексической связи между словами ''восхитительный'' и ''восхищение'' кажется очевидным, что:

\smallskip

1. x восхитителен тогда и только тогда, когда x вызывает восхищение у какого-либо наблюдателя.

\smallskip

Похоже, что это и имел ввиду Томас Гоббс, говоря ''Но что бы ни было предметом желания человека --- это то, что он, со своей стороны, назовет добром, а объект же ненависти и отвращения он назовет злом'' (1651/2012, Книга 1, Глава 6).

Такое представление о ценностях называют простым субъективизмом, ведь говорить о восхитительном как о том, чем некто восхищается --- самый простой способ связать восхитительное с восхищением. Субъективность же здесь оттого, что восхитительность так совершенно зависит от субъекта. Это делает подход релятивистским --- восхитительность нечта релятивна к эмоциональным реакциям субъектов.

Простой субъективизм обычно считается совершенно безнадежной позицией, ведь, кажется, пейзаж может оставаться восхитительным даже уже не вызывая у нас восхищение. Наши эмоциональные реакции слишком непостоянны, чтобы их можно было рассматривать как определяющий фактор ценности. Можно включить в определение психологические диспозиции и утверждать, что достойное восхищения зависит исключительно от диспозиций к восхищению им, и ценность, значит, тоже состоит в диспозиции к видению ценным. Итак,

\smallskip

2. x восхитительно тогда и только тогда, когда наблюдатель склонен восхищаться x.

\smallskip

Подобную теорию касательно добра и зла защищал Джесси Принц (2007): моральные свойства можно определить с точки зрения диспозиции испытывать ряд связанных с ними эмоций, то есть с точки зрения одобрения или неодобрения. Например, сентимент неодобрения лжи состоит в склонности испытывать стыд или вину, говоря ложь, и разочарование, презрение и отвращение, когда это делают другие. Принц утверждает: ''Действие имеет свойство быть морально правильным/неправильным только в том случае, если есть наблюдатель, который испытывает к нему чувство одобрения/неодобрения'' (2007, 92).

Другое решение проблемы непостоянтсва состоит в определении ценностей с точки зрения того, что \textit{вещи} склонны вызывать. Так, мы можем перестать испытывать восхищение от пейзажа, отвернувшись от него, но он все еще восхитителен, ведь склонен вызывать у нас восхищение. То есть

\smallskip

3. x достоин восхищения тогда и только тогда, когда x способен вызывать его у наблюдателя.

\smallskip

Эдвард Вестермарк, похоже, придерживается такого рода теории, когда пишет ''назвать поступок хорошим или плохим --- значит подразумевать его способность вызвать соответствующую эмоцию у так его называющего'' (1906, 4), хотя и имеет в виду скорее семантику, чем онтологию. Диспозиционный подход к ценностям остается субъективистским, ведь склонность нечта вызывать восхищение у человека зависит от этого человека (это, впрочем, предвосхищение основания --- прим. переводчика). Теория эта, значит, является релятивистской: то, что способно вызвать восхищение у меня --- альпийские пейзажи, например, --- не обязательно должно вызывать восхищение у вас.

Если бы не Дж.Э.Мур (1912/2005), определение ценностей Вестермарка, вероятно, было бы забыто. Иронично, что Мур упоминает Вестермарка только для серьезного к его подходу возражения. Оно направлено не на онтологический тезис о ценностях, а на определение добра и зла через диспозиции вызывать определенные чувства. Мур указывает на тогда невозможность несогласия: моя уверенность в диспозиции действия вызывать у меня неодобрение совместимо с отсутствием такой его диспозиции для вас. Оба суждения, значит, истинны, и разногласий быть не может.

Это возражение применимо и к тезису о природе ценностей: если ценность заключается в склонности вызывать определенные чувства, то моя уверенность в хорошести нечта совместима с отсутствием такой уверенности у вас, и, значит, разногласия и здесь невозможны.

Та же проблема затрагивает и подход Принца (2007): если неправильность определяется неодобрением утверждающего неправильность, то утверждение неправильность одним не противоречит утверждению неправильности другим. Справедливости ради следует сказать, что общая концепция Принца несколько сложнее, поскольку он считает, что суждения о добре и зле на самом деле являются эмоциями. То есть суждение о неправильности чей-то лжи заключается в стыде или вине за свою ложь, например. Но это все не касается возможности разногласий --- моральное суждение определяется чувством по отношению к своему объекту и может не разделятся всеми субъектами.

Другая проблема --- необходимость учета возможности ошибки касательно ценностей. Если ценности определяются нашими диспозициями, то ошибаться мы можем только касательно этих диспозиций. Так, по мнению Принца, пытать младенцев ради и считать это правильным вполне корректно, если это было весело.

Еще одна проблема --- одновременно восхитительность и невосхитительность одной и той же вещи (это, впрочем, предвосхищение основания абсолютизма --- прим. переводчика): если я склонен восхищаться пейзажем, а вы --- нет, то пейзаж, значит, и восхитительный, и не восхитительный. Субъективист может отвечать, что вещи восхитительны не сами по себе, а только для кого-то. Возможно, это и верный путь, но давайте сначала изучим менее ревизионистские подходы.

\section{Теории идеального наблюдателя}

Считается, что так называемые вторичные качества --- цвета, например --- состоят в склонности вызывать определенные реакции, например, зрительные переживания. Эти реакции определяются как те, что \textit{нормальные} испытуемые пережили бы в \textit{нормальных} обстоятельствах --- при достаточном освещении, например. Например, красный цвет определяется диспозицией вызывать соответствующие переживания у нормального субъекта в нормальных обстоятельствах. И то же можно утверждать о ценностях: что это диспозиции вызывать определенную реакцию у нормальных субъектов в нормальных обстоятельствах. Итак,

\smallskip

4. x восхитителен тогда и только тогда, когда x способен вызывать восхищение у нормального наблюдателя в нормальных обстоятельствах.

\smallskip

В обычном прочтении, основанном на общем понимании вторичных качеств, нормальность понимается эмпирически. Например, нормальным считается статистически наиболее распространенное. Как, впрочем, отмечают Д'Армс и Джейкобсон (2000b), ни один такой подход не может быть успешным --- они ''не могут отразить нормативные аспекты суждений'' (2000b, 726). То есть склонность нечта вызывать восхищение у нормального наблюдателя в нормальных обстоятельствах не имеет нормативной силы, способной заставить нас тратить время на созерцание и изучение этого нечта, например --- способность вещи вызвать восхищение у большинства людей ничего не говорит о конкретном человеке.

Одна из заманчивых поправок здесь восходит к Адаму Смиту (1759/2002) и заключается в обращении не к \textit{нормальному}, а к \textit{идеальному} наблюдателю (см. Ферт 1951, Льюис 1989, Кауппинен 2014):

\smallskip

5. x восхитителен тогда и только тогда, когда x способен вызвать восхищение у идеального наблюдателя в нормальных условиях.

\smallskip

Теории идеального наблюдателя допускают некоторую объективность, поскольку нет зависимости от субъективных реакций \textit{случайного} наблюдателя. Возможны, значит, и разногласия: восхитительность определяется не утверждающим ее, а идеальным наблюдателем, вещи, значит, не могут быть одновременно восхитительными и не восхитительными. Но что если разные идеальные наблюдатели утверждают разное? В любом случае, идеальность наблюдателя как будто придает нормативной силы его утверждениям --- конечно же его вердикт имеет некоторый вес в ответе на вопрос о том, что мне следует чувствовать и делать.

Как здесь понимать термин ''идеал''? Идеальность наблюдателя как его правильность понимания ценностей вещей можно сразу исключить --- это сделало бы определение круговым, было бы равносильно признанию восхитительности нечта вне зависимости от наблюдателя.

Но такое эпистемическое прочтение очень естественно. Как отметил Дэвид Енох (2005), трудно оправдать введение идеализирующих условий чем-то кроме лучшей приспособленности идеального наблюдателя к отслеживанию независимо существующих ценностей --- идеализация нужна для избегания ошибок о них. Енох приводит пример: чтобы узнать время, можно посмотреть на часы, но нужно также чтобы, например, батарейка в них была заряжена, то есть нужна определенная доля идеализации, и даже такое утверждение предполагает достаточную точность отслеживания времени этими часами. Енох пишет: ''естественное обоснование идеализации применимо только к случаям, когда соответствующая процедура или реакция считаются отслеживающими независимую от нее истину'' (2005, 764). Теории идеальных наблюдателей не могут апеллировать к такому обоснованию, и, значит, опираются на необоснованный ход.

В ответ Антти Кауппинен (2014) утверждал, что идеализация может быть практична и оправдываться необходимостью решения практических проблем, а не установлении фактов, и что проблема неидеализированных ответов в практических проблемах ими вызываемых --- для их решения нам приходится игнорировать некоторые из наших ответов, тем самым рассматривая определенные ответы как авторитетные и тем самым утверждая их как идеальные в смысли оптимальности для решения проблем. Какие практические проблемы видит Кауппинен? Если мы будет определять ценности с точки зрения наших реакций, не пытаясь эти реакции скорректировать, то получим нежелательные результаты: действия исходя лишь из желаний легко приводят к ухудшению здоровья, ярким обоям, распаду брака и недоверчивому сообществу. Впрочем, все эти вещи кажутся плохими безотносительно неодобрения их идеальными наблюдателями.

Есть еще два пути к пониманию, что такое идеальный наблюдатель. Первый состоит в разъяснении того, что его делает идеальным: например, доброта и беспристрастность. Второй состоит в его определении в неоценочных терминах: например, как информированного обо всех неоценочных фактах. Проблема, как отметил Расс Шафер-Ландау (2003) заключается в следующей дилемме: для избежания бесконечного регресса нам придется предположить не могут быть определены в терминах реакций идеального наблюдателя, но неясно, будут ли реакции идеального наблюдателя иметь право определять оценочные свойства. Почему, например, мы должны доверять наблюдателю, который не является добрым и беспристрастным? Можно также усомниться в возможность сохранения нормативности ценностей при выражении характеристик наблюдателя в терминах неоценочных --- почему нас должно волновать, что такой наблюдатель чем-то восхищается? Он, казалось бы, может и ошибаться.

Очевидным шагом будет явное использование нормативного понятия в определении ценностей, что знаменует переход к неосентименталистским взглядам.

\section{Неосентиментализм}

Неосентиментализм обычно формулируют как тезис, касающийся оценочных концепций и суждений, но он может также касаться и природы оценочных свойств --- эта онтологическая форма и будет интересовать нас. Согласно нему, ценности следует понимать в терминах нормативно подтверждающихся субъективных реакций --- то есть подтверждающихся как уместные, правильные, рациональные, обоснованные и т.п..

Неосентиментализм иногда подается как объяснение общих представлений о добре и зле. Брентано писал: ''В самом широком смысле, добро --- то, что достойно любви, то, что можно любить правильной любовью'' (1889/2009, 11). Брентано понимал любовь не как определенную эмоцию или чувство, а как включающую в себя все положительные эмоции, характеризуемые им как принадлежащие к сфере склонностей или удовольствий. Современные неосентименталисты, однако, чаще всего сосредотачиваются на определенном диапазоне ценностей, иногда называемых ''сентиментальными'' или ''аффективными'', и утверждают возможность их понимания аналогично свойству вызывать восхищение. Таким образом,

\smallskip

6. x восхитителен тогда и только тогда, когда в восхищение им уместно.

\smallskip

Поскольку такие ценности как восхитительность объясняются с точки зрения наших реакций, и поскольку реакции эти играют центральную роль в определении восхитительного, то такая теория утверждает субъективность (зависимость от наших реакций) ценностей. Но поскольку такое объяснение допускает разногласия и ошибки, то оно допускает и объективность ценностей. Джон Макдауэлл (1985) указывает, что сентиментальные ценности следует объяснять с точки зрения соответствующих своим объектам аффективных реакций, и что так понимаемые ценности могут быть как в некотором смысле субъективными, так и в некотором смысле объективными: они объективны потому что не ''лишь фрагменты наших субъективных состояний'' (Макдауэлл 1985, 114). То есть нечто вполне может вызвать у вас восхищение, достойным восхищения не являясь --- восхищение ваше не будет, значит, уместным.

В целом существует два основных способа понимать неосентиментализм. Первый --- ''неприоритетный'' --- характерен для теорий чувствительности (Макдауэлл 1985, Виггинс 1987) и не пытается редуцировать ценности. Таким образом, поскольку утверждение не состоит в редуцировании восхитительности к уместности восхищения, не следует отдавать приоритет правой части бикондиционала --- вместо этого предлагается понимать бикондиционал как разъяснение, а не как анализ.

Второй же способ предполагает редуктивный анализ ценностей. Одобрение, значит, не может быть выражено в оценочных терминах --- уместная реакция не может считаться хорошей, иначе возникает цикличность. Одобрение, значит, должно пониматься деонтически, и уместными реакциями обычно считаются те, что в том или ином смысле мы \textit{должны} иметь. Либо же можно свести ценности к ответам, доводы в пользу которых нас есть.

Поскольку неосентиментализм понимает ценности с точки зрения аффективных реакций вообще и эмоций в частности, он не тождественен теориям установки, говорящим о ценностях с точки зрения нормативно подтвержденных установок независимо от этих установок аффективности. Более того, теории уместного эмоционального отношения вообще и неосентиментализм в частности --- близкие родственники так называемой теории перекладывания ответственности, согласно которой ценности объясняются через нормативно одобренные \textit{действия}, а не \textit{эмоции}. Таким образом, неосентиментализм --- всего лишь одна из версий более общего подхода, объясняющего ценности деонтически.

Неосентиментализм привлекателен демистификацией ценностей и объяснением их связи с тем, что мы должны делать. Влодек Рабинович и Тони Реннов-Расмуссен указывают, что такой подход устраняет атмосферу таинственности, ведь ''нет ничего странного в предписывающих импликациях приписывания ценностей, если ценность эксплицируется в деонтических терминах''. Сводя ценности к таким установкам как эмоции, мы получаем объяснение способности ценностей направлять действия через порождение мотивации --- это происходит из-за тесной связи между эмоциями и мотивацией. Таким образом, даже если неосентиментализм не привержен моральному интернализму, он предполагает по крайней мере частую связь между ценностями и мотивацией.

Основное возражение против неосентименталистского анализа --- возражение о ''неправильных причинах''. Быть восхитительным --- значит предоставлять достаточные основания для восхищения, но у нас может быть множество причин за и против, и только некоторые из них связаны непосредственно с ценностью объекта. Представьте, что злой демон угрожает вам страшными страданиями если вы не будете им восхищаться --- у вас ведь тогда есть веская, благоразумная причина восхищаться им, но не потому что он восхитителен. Этот пример можно модифицировать для включения моральных соображений: если демон обещает положить конец болезням если вы будете восхищаться им, у вас есть моральное основание им восхищаться, не связанное с действительной его восхитительностью --- согласно неосентименталистскому подходу, у вас нет достаточных оснований восхищаться им. Как подчеркивают Д'Армс и Джейкобсон (2000a), было бы ''моралистическим заблуждением'' заключить из моральной правильности/неправильности испытывать эмоцию, уместность этой эмоции согласно ценности своего объекта.

Казалось бы, легко отличить благоразумные и моральные причины от относящихся к уместности восхищения. Но неосентименталисты должны объяснить, как отличать неправильные причины от правильных, если только правильные касаются вещей как таковых. И поскольку неосентименталист стремится свести к минимуму оценочные свойства, он не может просто апеллировать к ценности вещей для указания правильных причин --- не может, например, сказать, что правильными причинами для восхищения являются основанные на восхитительности.

Один из ответов заключается в разграничении причин на объектные --- являющиеся правильными --- и связанные с состоянием (Парфит 2001/2018). Причина, связанная с состоянием, касается ценности реакции на объект --- например, ценности восхищения демоном. Но в некоторых случаях неправильные причины зависят от самого объекта: причиной для восхищения демоном могут быть его тщеславие и мстительность, и, следовательно, склонность к угрозам и принуждению к восхищению им, но эти свойства не делают его восхитительным, и, значит, являются неправильными причинами.

Эндрю Рейснер (2009) отмечает основанность такой методологии на нашей интуиции касательно ценностей: правильные причины --- те, что касаются ценности вещей как таковых, а демон, несмотря на наличие причин им восхищаться, восхищения недостоин. Таким образом, на первом месте стоят ценности и интуиция в отношении них, и если и есть причина восхищаться человеком, то только потому, что он восхитителен --- а не наоборот. Но тогда не получается свести ценности к нормативно одобренным установкам, ведь, согласно редуктивному анализу, у восхитительности нечта есть причина, не заключающаяся в его восхитительности.

В знаменитом ''Евтифроне'' Платона (2002) Сократ задается вопросом: любят ли боги благочестивых потому что они благочестивы, или же они благочестивые благочестивы потому что их любят боги? Аналогично: восхищаемся ли мы нечтом потому что оно достойно восхищения, или же оно достойно восхищения потому что мы им восхищаемся? В рамках неосентиментализма: восхищение нечтом уместно потому что оно восхитительно или же нечто восхитительно потому что восхищение им уместно? Неосентименталисты придерживаются второго, но далеко не очевидно, что это верный ответ. Скорее, уместно восхищение тем, что восхитительно --- именно это сказали бы о ценностях реалисты. Таким образом, ответ неосентименталиста на дилемму Евтифрона противоречит предположениям о ценностях, что приводит к возражению от неправильных причин.

Должны ли мы просто отказаться от неосентименталистского подхода и заявить, что ценности --- естественные или, может быть, даже не-естественные объективные свойства вещей, не имеющие ничего общего с нашими чувствами? Но восхитительность и постыдность явно связаны с эмоциями. В следующем разделе мы рассмотрим другой способ понять связь между эмоциями и ценностями.

\section{Теории чувствительности}

До сих пор мы фокусировались на редуктивных подходах к ценности. Дэвид Виггинс (1987) и Джон Макдауэлл (1985) предлагают нередуктивную версию неосентиментализма: такие свойства, как восхитительность, постыдность или отвратительность неразрывно связаны с проявлением человеческой чувствительности, но не сводятся к какому-либо реляционному свойству, такую чувствительность включающему. Нередуктивный и редуктивный неосентиментализмы, значит, согласны в восхитительности нечта тогда и только тогда когда восхищение им уместно. Нередуктивный неосентиментализм, однако, отрицает определение восхитительности через уместность восхищения и предлагает неприоритетное прочтение бикондиционала, лишь указывающее на связь, но не редуцирующее. Макдауэлл (1987) указывает, что реакции и оценочные свойства --- ''братья и сестры, но не родители и дети'' --- ни те, ни другие не имеют метафизического приоритета.

Как может оценочное свойство, не будучи определено в терминах наших реакций, быть с этими реакциями все же связанным? Если восхитительность --- не склонность вызывать восхищение, то как оно может быть тесно связано с восхищением? Возможно, свойства вроде восхищения, хотя и не сводятся к возникновению реакций, все же очерчиваются концепциями, от реакций зависящими, то есть они концептуально связаны с соответствующими эмоциональными реакциями. Важно, что зависимость концепции от реакции не означает этой концепции к этой реакции сведения --- в неприоритетном прочтении не концепция ценности, ни концепция адекватного ответа не являются более первичными или более фундаментальными. То есть восхитительность, отвратительность, постыдность и т.д. --- \textit{нереляционные} или \textit{монадические} свойства, но соответствующие понятия --- \textit{восхитительность}, \textit{отвратительность}, \textit{постыдность} и т.д. --- зависят от реакции.

Виггинс (1976) рассматривает пример красноты и указывает, что и она может являться монадическим, но в некотором смысле ''релятивным'' свойством. Относительно оно потому, что категория цвета антропоцентрична --- ''соответствует интересу, укорененному только в существах, чей сенсорный аппарат схож с нашим'' (1976, 107). Это, впрочем, совместимо с монадичностью красноты: мы видим вещи красными потому что они красные. Таким образом, оба рога дилеммы Евтифрона верны: мы видим вещи красными потому что они красные, и вещи считаются красными потому что мы их такими видим.

То же и касательно ценностей: мы восхищаемся нечтом потому что оно восхитительно, и восхитительно оно потому что мы им восхищаемся. Восхитительность --- монадическое свойство, на которое мы стремимся реагируем восхищением, но и восхитительность зависит от реакции, ведь соответствующая концепция антропоцентрична --- она есть у нас только потому, что мы восприимчивы к цвету. Перефразируя Виггинса (1976), можно сказать, что чувство восхищения необходимо для существования перспективы, с которой можно было бы воспринять восхитительность.

Так интерпретированный, неосентиментализм без труда справляется с возражением от неправильных причин. Правильная причина --- та, что зависит от объективной ценности вещей. Реакция восхищения уместна тогда и только тогда, когда объект восхитителен, так что злой демон, заставляющий им восхищаться, сам вовсе не восхитителен. Итак,

\smallskip

7. x восхитителен тогда и только тогда, когда восхищение им уместно, и восхищение x уместно потому что x восхитителен.

\smallskip

Это, конечно, круговой ответ. Но защитники теорий чувствительности и не пытаются предложить редуктивный анализ --- только прояснить концепцию восхитительного, указав на ее связь с прочими концепциями, и неприоритетность конечно приводит к такому кругу. И все же, чем подобное утверждение лучше, чем ''нечто является кошкой тогда и только тогда, когда утверждение ''нечто является кошкой'' --- истинно''? Это ведь тоже верное утверждение, но оно не привносит никакой новой информации и не представляет интерес.

Виггинс отвечает: ''пытаясь понять, является ли x хорошим, правильным или красивым, мы не можем обратиться к чему-либо более фундаментальному, чем наши чувства'' (Виггинс 1987, 189). То есть, размышляя о восхитительности нечта, мы можем апеллировать лишь к нашему чувству восхищения им --- более фундаментального обоснования мы найти не сможем. Это, впрочем, не значит, что чувства всегда правильны --- скорее, мы должны начать с чувств чтобы хотя бы задать вопрос о восхитительности.

Чтобы развеять возможные сомнения, нужно будет сравнить нынешнее это ценностей с теориями, согласно которым ценности никоим образом с эмоциями не связаны и прояснить связь между оценочными и естественным свойствами: идентично ли вызывающее восхищение определенному естественному свойству или их дизъюнкции? Или же рассматриваемая связь --- супервентность, то есть невозмножность различия на оценочном уровне без различий на уровне естественном? В любом случае, нередуктивная версия неосентиментализма --- привлекательная концепция.

\section{Подведение итогов}

Мы начали с вопроса необходимости пересмотра выводов касательно природы эмоций после выяснения природы ценностей. Если нередуктивный неосентиментализм верен, то такая необходимость отпадает. Нередуктивный неосентиментализм и перцептивные (включая рецептивные) теории эмоций хорошо согласуются (см. Тапполет 2016). Поскольку эмоции --- оценочные репрезентации, верные при совпадении формального и интенционального своих объектов, уместность можно понимать как такую верность.

Не придерживаются ли перцептивные теории неверного --- объективистского --- подхода к ценностям? Наиболее правдоподобная теория ценностей утверждает, что оценочные свойства объективны в строгом смысле монадических свойств, так что, похоже, защитникам перцептивных теорий не о чем беспокоится. Какие еще теории ценностей совместимы с перцептивными теориями? Хорошим кандидатом является теория ошибок, согласно которой мы говорим и чувствуем так, будто объективные ценности существуют, хотя на самом деле их нет (см. Маки 1977).

Редуктивные неосентименталисты пытаются избежать круга в объяснении, и в этом лучше всего могут помочь мотивационные теории. Значит, принимая редуктивный подход, следует пересмотреть аргументы против мотивационных теорий.

Итак, подход сентименталистов кажется вполне верным. Однако следует подчеркнуть, что рассмотренная его версия допускает несоответствие эмоций действительности и, значит, ошибки в отношении ценностей. Возможность репрезентативной ошибки обычно считается отличительной чертой рационализма, так что разрыв между этими двумя подходами --- сентиментализмом и рационализмом --- похоже, не столь существенен.

\section{Краткое содержание}

\begin{itemize}
  \item\ Термин ''ценность'' может относиться к хорошим вещам или идеалам, а также к оценочным свойствам вещей или людей, таким как их восхитительность или доброта. Оценочные концепции отличаются от деонтических, таких как концепция должного.
  \item\ Сентиментализм в отношении оценочного или деонтического вообще и моральный сентиментализм в частности, принимают разные формы в зависимости от затрагиваемых ими частей метаэтики: онтологический сентиментализм, семантический (эмотивизм, например), психологический или эпистемический.
  \item\ Простой субъективизм утверждает определенность ценности нечта соответствующей в связи с ним эмоцией или к этой эмоции склонности у наблюдателя. Основная проблема этого подхода --- он не учитывает возможность разногласий и ошибок.
  \item\ Теории идеального наблюдателя утверждают определенность ценности нечта его способностью вызывать соответствующую эмоцию у идеального наблюдателя. Таким образом, возможны разногласия и ошибки. Но такие теории с трудом обосновывают использование идеализации. Более того, перед ними стоит дилемма: либо они допустить характеризацию идеальных наблюдателей в оценочном ключе и, значит, прийти к кругу, либо отрицать это и растерять весь возможный авторитет таких наблюдателей.
  \item\ Неосентиментализм в редуктивной версии утверждает сводимость ценностей к эмоциям. Этот подход привлекателен объяснением нормативности ценностей. Неосентиментализм сталкивается с возражением о неправильных причинах, указывающем на необходимость связывать эмоций с независимыми ценностью их объектов как таковых.
  \item\ Теории чувствительности не пытаются свести ценности к эмоциям. Ценности, согласно ним --- монадические свойства, даже если очерчивающие их концепции антропоцентричны и зависят от в том числе от реакций.
  \item\ Теории чувствительности хорошо сочетаются с оценочными теориями эмоций и, в частности, с перцептивными теориями.
\end{itemize}

\begin{tcolorbox}
  \section{Учебные вопросы.}
  \begin{enumerate}
    \item\ Какой подход к эмоциям совместим с пониманием моральных утверждений как лишь выражений эмоций, не являющихся истинными или ложными?
    \item\ Согласятся ли разные идеальные наблюдатели в своих ответах? Если нет, то что это влечет в отношении совместимости теорий идеального наблюдателя и релятивизма?
    \item\ Жестокость шутки --- это своего рода причина не испытывать от нее удовольствие. Но делает ли она ее менее смешной?
  \end{enumerate}
\end{tcolorbox}

\chapter{Этика и эмоции}

\section{Введение}

Эмоции наносят большой вред: из-за них можно убивать, лгать и отказывать в помощи. Даже такие эмоции как сострадание и любовь, могут приводить к аморальным действиям. Например, из-за неуместного сострадания можно лишить кого-то части автономии, печаль может заставить не замечать нужд других, и даже удивление может спровоцировать аморальную реакцию --- действие из предубеждения или обидное замечание. Некоторые эмоции кажутся аморальными даже вне зависимости от последствий --- например, зависть, ревность и ненависть.

Получается, что морально совершенные агенты должны не испытывать никаких эмоций вовсе? Но тогда чувства не влияли бы на их мысли и действия. Даже стоики и Иммануил Кант сочли такое преувеличением: морально совершенные агенты должны были бы совершенно рассуждать, переходя от моральных принципов к практическим соображениям.

Такой идеал кажется как минимум недосягаемым. Большинство современных моральных философов согласны, что для правильных действий эмоции необходимы --- и так считают сторонники всех трех основных этических подходов --- деонтологии, консеквенциализма и этики добродетели. Деонтологи, такие как Марсия Барон (1995), Нэнси Шерман (1997) и Карла Баньоли (2003), стремятся подчеркнуть роль таких эмоций, как уважение к моральному закону, а также любовь и сострадание.

То же верно и касательно консеквенциализма, историческим предшественником которого является классический утилитаризм, защищенный Джереми Бентамом (1789/1961) и Джоном Стюартом Миллем (1863/1998). Некоторые современные консеквенциалисты определяют правильные действия через мотивы с наилучшими последствиями (Адамс 1976), но даже консеквенциалисты действий подчеркивают важность характера для моральности агента. Они, значит, по большей части согласны с центральным значением эмоционально индуцированной мотивации для моральной агентности.

Эти выводы можно рассматривать как результат растущей важности третьего основного подхода --- этики добродетели. Если не следовать Энском и не отказываться от понятия правильного действия вообще, то правильным действием называется то, что совершил бы добродетельный агент. Вообще, добродетели необходимы для моральной агентности, и раз добродетели связаны с эмоциями, то и эмоции для нее необходимы.

Несогласны все эти теории в том, какую именно роль играют эмоции. Мы рассмотрим вопрос влияния эмоций на моральную агентность, сосредоточившись на их роли в добродетелях, начиная с аристотелевской их концепции, разработанной Розалиндой Херстхаус (1999) и Джулией Аннас (2011). Как станет ясно, стандартная интерпретация Аристотеля ставит знание в основу всяких добродетелей. Мы обсудим также два возражения, которые можно выдвинуть против такой сократической интерпретации, а затем рассмотрим подход Гопал Шринивасан (2020), согласно которому эмоции являются ядром добродетелей, а добродетели необходимы для счастья и процветания.

\section{Аристотелевская теория добродетели}

Аристотель определяет добродетели как укоренившиеся черты характеры, заслуживающие похвалы. Добродетели необходимы для счастья или процветания (эвдемонии) людей, что, согласно Аристотелю, является высшим благом. Он различает этические и интеллектуальные добродетели, где последние касаются теоретического и практического рассуждения. Отрицая, что добродетели --- черты, с которыми мы рождаемся, Аристотель утверждает возможность их воспитания через приобретение привычек. Он рассматривает такие добродетели, как честность, мужество, щедрость и справедливость.

Этические, да и прочие добродетели --- черты характера, заставляющие агентов действовать и реагировать определенным образом. Добродетели не являются эмоциями, но связаны с эмоциональными диспозициями, точнее, подразумевают диспозиции к правильным эмоциям в правильных ситуациях. Правильное --- то, что не является крайностью:

\smallskip

И страх, и уверенность, и гнев, и жалость, и вообще удовольствие и боль могут ощущаться и слишком сильно, и слишком слабо, и в обоих случаях неправильно; но ощущать их в нужное время, по отношению к нужным предметам и людям и с правильным мотивом и правильным способом --- это то среднее и то лучшее, что характерно для добродетели. (Никомахова этика 1106b20-25)

\smallskip

Так, смелый не чувствует ни слишком сильного, ни слишком слабого страха, в отличие от труса или самоуверенного и опрометчивого.

Со связью добродетелей с правильными эмоциями согласны современные аристотелики. Так, Херстхаус говорит, что добродители подразумевают предрасположенность испытывать эмоции, при этом ''эмоции эти ощущаются в нужных случаях, по отношению к нужным людям и объектам и по нужной причине [...]'' (1999, 108). Добродетели состоят не только в диспозициях к правильным эмоциям, но и, разумеется, в диспозициях к правильным --- вызываемым правильными эмоциями --- действиям. Херстхаус пишет:

\smallskip

Полная аристотелевская концепция добродетели --- концепция сложной черты характера, то есть диспозиции, прочно укорененной в своей обладателе и, так сказать, целиком его пронизывающей. Эта диспозиция многонаправленная и включает также множество других действий [sic]: желания, эмоции и эмоциональные реакции, восприятия, отношения, интересы и ожидания.

\smallskip

Так, честный человек будет стараться избегать нечестных поступков, таких как мошенничество --- будет всегда поступать честно, например, говорить правду. Делать он это будет характерным образом, то есть с готовностью и без колебаний. Он также будет огорчен, если близкие и дорогие ему люди будут нечестны, не будет обижаться на честную критику, будет возмущен нечестным поступком и радоваться честному.

Как различные части этой диспозиции соотносятся друг с другом? Являются ли какая-то из них центральной или они равны? Этот вопрос относит к вопросу о роли знания для добродетелей. Аристотель считал, что этические добродетели требуют особых познавательных способностей, которые он рассматривал как добродетель интеллекта, связанного с действиями --- практическую мудрость (фронезис). Она предполагает правильное понимание благ, делающих жизнь счастливой и полноценной и их друг с другом сочетания. Этические добродетели, в отличие от врожденных склонностей, нуждаются в практической мудрости. Как она взаимодействует с прочими компонентами добродетелей? Давайте рассмотрим так называемую сократическую интерпретацию аристотелевской теории.

\section{Сократическая интерпретация}

Практическая мудрость помогает правильно рассуждать о практических вопросах и составляет основу этических добродетелей (Херстхаус 1999; Аннас 2011) --- согласно этой концепции, диспозиции, составляющие добродетель, следуют из практического рассуждения. Если вы правильно рассуждаете о практических вопросах, то можете поступать правильно и испытывать правильные эмоции. Херстхаус пишет: ''Каждая из добродетелей предполагает умение делать все правильно, поскольку каждая включает в себя фронезис, то есть способность правильно рассуждать о практических вопросах'' (1999, 12). То есть даже если добродетели включают в себя как эмоциональный, так и интеллектуальный компоненты, последний является фундаментальным.

Для проявления каждой конкретной добродетели требуется особое суждение касательно конкретного блага. Как указывает Лизл Ван Зил: ''Добродетельный человек не просто заботится об определенных вещах --- он заботится о правильных вещах и имеет правильные желания. Это требует суждения об этих вещах --- что они действительно хороши или важны, и, значит, заслуживают защиты или стремления'' (2018, 23). Точно так же и Майкл Брейди считает, что добродетельная мотивация необходимо включает в себя оценочный элемент (2005, 87), имеющих приоритет над мотивационным --- желаниями и эмоциями. Так, доброжелательный человек --- тот, кто ценит благополучие других и потому мотивирован действовать соответствующим образом. Добродетельные агенты действуют на основе не просто случайных оценочных суждений --- эти их суждения составляют знание. Как утверждает Брейди: ''доброжелательный человек --- тот, кто знает, что содействие благополучию других ценно само по себе, правдивый человек --- тот, кто знает, что истина хороша сама по себе, а ложь --- плоха'' (2005, 99).

Добродетель, значит --- знание, а порок --- неведение, добродетели же --- достоинства характера, заключающиеся, главным образом, в знаниях. Они включают диспозиции к эмоциям, но диспозиции эти не фундаментальны --- оценочные знания важнее, потому что эмоции возникают через них. Диспозиции к эмоциям --- диспозиции реагировать на определенное оценочное знание. Так, столкнувшись с кем-то, кого истина совершенно не волнует и кто готов других вводить в заблуждение, честные агенты формируют о нем негативное оценочное суждение, согласно которому испытывают возмущение или презрение.

Основной аргумент этого подхода основан на необходимости добродетельных действий совершаться по правильным причинам. Добродетельный агент, утверждает Аристотель, ''должен действовать ради самого этого выбранного образа действия [...]'' (Никомахова этика, 1105a28-34). Так, мужественный спасет тонущего ребенка не потому, что это удачный способ впечатлить других, а потому что это удачный способ спасти ребенка. Добродетельные агенты, значит, не только выбирают правильное, но и знают причину, по которой это правильное правильно. И без оценочных знаний никак не понять, какие причины правильны.

Насколько эта интерпретация близка к утверждаемому Аристотелем? Действительно ли он считал, что добродетели --- достоинства характера, заключающиеся в знаниях? Это, как мы увидим, не единственный способ понять Аристотеля.

\section{Возражение от скромности}

Возможно, не все добродетели требуют знания. Джулия Драйвер утверждает, что некоторые из них --- ''добродетели невежества'' --- требуют игнорирования ряда вещей. Она говорит о слепой благотворительности, импульсивной смелости, прощении и доверии, а также скромности. Скромность, согласно Драйвер --- добродетель, требующая некоторой недосамооценки: скромному необходимо считать себя не столь хорошим, каков он на самом деле есть. Так, лучший пианист в мире, будучи скромным, не должен до конца осознавать, что он лучший пианист в мире --- то есть должен иметь ложное представление о себе. Это также значит, что нельзя знать, что поступаешь скромно, не хвастаясь своим талантом, например. То есть невозможно действовать скромно по правильным причинам. Драйвер делает вывод, что знание факторов, создающих добродетель, не является необходимым ее условием (2001, xv).

Защитники сократической интерпретации могут ответить, что скромность, конечно, предполагает знание --- например, точное понимание значимости своих достижений. Лучшему в мире пианисту, чтобы быть скромным, нужно знать, что он лучший в мире пианист, но также и понимать ограниченность этого факта: например, что удача и возможности сыграли важную роль в его достижениях и что это не единственный имеющий значение аспект их оценки. Возможно, он лучший в классической музыке, но есть те, кто лучше играет джаз или блюз. И, в любом случае, он хорош не во всем --- не хорош как повар или как философ. Он, значит, не считает себя лучшим в мире вообще, его самооценка не зависит от достижений в какой-то области. Скромность, таким образом, совместима со знанием о своем превосходстве в какой-то области.

Но нельзя ли быть скромным, не имея убеждений о границах своих достижений? Возможно, для скромности важно не невежество и не знание пределов своих достижений, а забота об определенных вещах. Так, Фред Шуелер утверждает, что по-настоящему скромный человек --- тот, кого не волнует, впечатляются ли люди его достижениями (1997, 478). Драйвер, впрочем, отмечает, что не заботиться о впечатлениях других можно и потому, что испытываешь к ним презрение --- что уже, конечно, не похоже на скромность. Арпали и Шредер (2013) утверждают, что скромный может знать, какими качествами обладает, но его не волнует, насколько высоко эти качества оцениваются --- он не ценит свои высокие достижения. То есть даже если скромность часто приводит к незнанию собственных достоинств в качестве побочного эффекта (если вас не волнует ваш ''рейтинг'', вы, конечно, с меньшей вероятностью будете искать информацию, подтверждающую или опровергающую вашу его оценку), это незнание не обязательно. Скромность тогда зависит от объектов заботы а не от мыслей о пределах своих достижений.

Что в этом контексте значит забота? Арпали и Шредер (2013, см. также Франкфурт 1982) считают, что забота --- это желание определенных результатов. Возможно, забота --- это многонаправленная диспозиция испытывать ряд эмоций. Так, заботящийся о своем месте в рейтинге пианистов, склонен радоваться высокому такому своему рейтингу и бояться, что он станет ниже, а также завидовать тем, у кого он выше и т.д.. Дэвид Шумейкер отмечает, что ''заботясь о нечте, я становлюсь уязвимым для приобретений и потерь --- эмоциональных взлетов и падений в соответствии с изменениями нечта'' (2003, 91).

Таким образом, ошибка нескромного носит аффективный характер --- его заботит нечто не столь уж важное. Скромный же такой ошибки не допускает --- его может волновать ценность музыки, то, что он опытный музыкант и т.д., но не волнует победа над другими пианистами.

\section{Возражение от чрезмерной интеллектуализации}

Итак, главная проблема сократической интерпретации --- чрезмерная интеллектуализация добродители. В большинстве случаев, действуя из смелости или великодушия, мы не размышляем, следует так делать или нет. Как отмечает Аннас, добродетельные агенты, похоже, ''реагируют прямо и немедленно на ситуации, что часто контрастирует с более сдержанной и вдумчивой реакцией не столь добродетельных'' (2011, 28). Так, щедрый человек сразу откликается на нужды других, не задумываясь, и когда его спрашивают о мотивации, он указывает на соответствие ситуации. То есть добродетельные агенты обычно не руководствуются явными мыслями о ценности того, что делают, а рассматривают ситуации как просто требующие действий, независимо от соответствующих размышлений.

Аннас (2011) опирается на предположение Аристотеля о подобности добродетелей навыкам --- и те и другие ''практичны и обоим можно научиться только на практике, делая соответствующее''. Роль знаний в добродетелях аналогична таковой в навыках вроде игры на фортепиано или футболе. Для развития добродетелей необходимы сознательное практическое рассуждение, точно так же, как когда учимся играть на фортепиано: мы учимся у образцов для подражания и учителей, повторяя то, что они делают, следуя инструкциям, а также понимая, почему они делают именно это и именно так. Все это, конечно, требует сознательных усилий и рассуждений. Овладев навыком принятия добродетельных решений, человек может применять его более не задумываясь, но всегда сможет дать отчет о том, что и почему он делает. Таким образом, добродетель --- склонность не только действовать правильно, но и делать это по правильным причинам.

Но можно ли при этом говорить об оценочном знании как ядре добродетели? Возможно, не следует предполагать пропозициональное знание добродетельного агента внутренне ценным. Учась играть на фортепиано, вы не усваиваете общие пропозициональные истины о том, почему то или иное должно быть сделано именно таким образом. Вы, конечно, тренируетесь определенным образом двигать пальцами, но тренируетесь также переживать определенные состояния во время исполнения соответствующих частей музыкального произведения. Как отмечает Аннас (2011), есть много навыков, применение которых не связано с чувствами. Но ведь наличие этих навыков подразумевает привычку заботиться о соответствующем.

Аристотель в Никомаховой этике утверждает, что моральное воспитание --- в значительной степени воспитание эмоций. Как пишет Бернард Уильямс:

\smallskip

если моральное воспитание не касается вопросов о том, чего следует бояться, на что злиться, что презирать --- то непонятно, чего оно вообще может касаться. Аристотель верно понял, что нас интересует не столько привитие принципов, сколько воспитание эмоций. (1973, 225)

\smallskip

Потому основной альтернативой сократической интерпретации является та, согласно которой эмоции составляют ядро добродетелей.

\section{Теория добродетелей Шринивасана}

Следуя утверждению Аристотеля, что правильные эмоции характеризуют добродетель, многие говорили об основанности добродетелей на эмоциях. Так, Кристьян Кристьянссон считает, что ''эмоциональные реакции являются важными ингредиентами добродетелей'' (2018, 20), и приписывает это утверждение Аристотелю. Поскольку добродетели --- диспозиции, то эмоциональные диспозиции составляют их суть. Это утверждение можно развивать по-разному. Мы сосредоточимся на недавно разработанном подходе, согласно которому в основе ряда добродетелей лежат эмоциональные диспозиции. Шринивасан рассматривает доброту, смелость и сострадание, сомневаясь, что справедливость можно описать в его подходе.

Шринивасан стремится ответить на вопрос, каким должен быть психологический склад добродетельных агентов. Ответ заключается в наличие соответствующих эмоций, которые необходимы для таких добродетелей, как доброта, смелость и сострадание --- склонность к переживанию соответствующих эмоций ''необходима для того, чтобы образцы добродетели прошли самый основной тест центрального ее измерения'' (2020, 33). Этот центральный тест включает последовательное вынесение правильных практических суждений о ситуациях, требующих проявления добродетели --- нечто, что заставляет агента стабильно действовать правильно. Непример, обладать состраданием --- значит постоянно принимать правильные суждения о том, что делать в соответствующих (требующих проявления сострадания) ситуациях. То есть добродетель должна содержать определенный эмоциональный компонент, включать в себя морально скорректированную эмоциональную диспозицию.

Такая морально скорректированная черта --- то, что позволяет добродетельному агенту выносить правильные практические суждения. В этом заключается важное отличие от сократической интерпретации. Но если черты морально скорректированных эмоций не соответствуют вердикту морального суждения, то что они такое?

Даже утверждая, что не имеет каких-либо особых возражений против перцептивных теорий эмоций, Шринивасан официально занимает нейтральную позицию по этому вопросу и говорит главным образом об отношении эмоций к вниманию. Де Соуза пишет: ''эмоция ограничивает диапазон информации, которую организм будет принимать во внимание'' (1987, 195). Опираясь на идеи де Соузы, Шринивасан предполагает две пары относительно бесспорных утверждений о функциях эмоций: одна касается значимости входных данных, то есть того, что эмоции вызывает, другая же --- значимости выходных данных, то есть реакций на вызвавшее эмоции.

Вот первая пара:

а) Эмоции контролируют значимость входных данных.

б) Эмоции пластичны в контроле этой значимости.

Эмоции, значит, имеют функцию избирательного фокусирования внимания и подчеркивания значимости определенных вещей, таких как опасность, вред, потеря и т.д.. В дополнение к этому может измениться и круг вещей, на которых эмоция фокусирует внимание. Шринивасан пишет: ''триггеры эмоции не фиксированы, но допускают изменение в определенной степени'' (2020, 53). Так, можно научиться бояться одного и перестать бояться другого. С этим согласны даже биологические детерминисты --- даже они утверждают, что эмоциональные диспозиции частично формируются культурой.

Вторая пара утверждений касается роли эмоций в реакциях:

в) Эмоции контролируют значимость выходных данных.

г) Эмоции пластичны в контроле этой значимости.

Контролируя значимость выходных данных, эмоции контролируют ''заметность некоторых конкретных реакций в диапазоне возможных ответов'' (2020, 60). Например, в случае страха такими реакциями обычно является бегство, борьба или замирание. Еще в Главе 3 подчеркивалось, что окружающая среда и культура формируют наши эмоциональные диспозиции, влияя на связанные с ними тенденции действий.

Эти четыре утверждения значат, что эмоциональные черты идеально подходят в качестве основного компонента добродетелей. Контролируя значимость входных и выходных данных, они существенно влияют на то, как мы реагируем на конкретные ситуации, и, будучи пластичными, допускают воспитание лучшей на эти ситуации реакции.

\section{Аргументы Шринивасана}

Добродетельное действие требует осознания агентом морально значимых свойств ситуации. Добродетельные агенты постоянно ''настроены на моральный мир'' (2020, 133). Переживание эмоций, значит, помогает агенту стать добродетельнее. Как отмечает Шринивасан, эмоции, контролируя входную и выходную значимость, отлично подходят для объяснения значимости определенных свойств мира для данного агента (и следования определенных реакций из этих свойств значимости) (2020, 134).

Эмоцию, лежащую в основе сострадания, Шринивасан называет сочувствием. Опираясь на размышления Питера Голди (2000), Шона Николса (2004) и Дэниела Бэтсона (2011), он противопоставляет его эмпатии и личным страданиям. Эмпатия представляет собой процесс, приводящий к опосредованному разделению гомологичного аффекта --- разделению чужой печали или радости, например, что предполагает образное представление себя на месте другого. Личное страдание же не всегда гомологично и эгоцентрично --- можно быть расстроенным горем или бедой других. В отличие от эмпатии, сочувствие не требует гомологии: то, что вы испытываете, может, но не должно, быть того же рода, что и испытываемое объектом сочувствия: можно сочувствовать тому, кто попал в беду, даже когда он находится без сознания. Сочувствие сопровождается альтруистической мотивацией, конечной целью которой является улучшение не собственного благосостояния, а благосостояния другого. Сочувствие, значит --- эмоция, вызванная бедственным положением другого, она мотивирует агента помочь ему. Испытывая сочувствие, вы сосредоточены на попавшем в беду и его страданиях, и понимаете, что ему нужна помощь. Таким образом, склонный испытывать сочувствие склонен лучше замечать, что другие нуждаются в помощи, что делает сочувствие центральным компонентом сострадания.

Второй аргумент в пользу центральной роли эмоций для таких добродетелей как доброта, смелость или сострадание, заключается в необходимости переживания эмоций для правильного совершения добродетельных действий. Так, в случае добродетели сострадания акт сострадания ''должен искренне выражать сочувствие'' (2020, 15).

Таким образом, сострадание имеет основным компонентом морально скорректированную черту сочувствия --- благодаря ей получается лучше замечать, когда другие находятся в беде и помогать им.

Но что делает черту сочувствия морально корректной? Утверждается, что ''обычное сочувствие должно быть скорректировано таким образом, чтобы реже приводить к моральным ошибкам'' (Шринивасан 2020, 148). Связывая этот момент с утверждениями о пластичности эмоций, Шринивасан утверждает, что на входной стороне склонность обычного сочувствия к ошибкам ''может быть морально скорректирована путем добавления подходящих и устранения неподходящих триггеров'' (2020, 155). Так, морально скорректированное сострадание --- склонность испытывать сочувствие ко всем, кто находится в бедственном положении.

Морально скорректированная эмоциональная черта, значит --- эмоциональная диспозиция к соответствующему ситуации. Шринивасан не хочет связывать себя с какой-либо конкретной теорией эмоций, но есть веские основания утверждать, что лучше всего его теории соответствует оценочная: ''подчинение чувствам образца добродетели иногда является лучшим способом выяснить, что является морально правильным поступком в данной ситуации'' (2020, 168). Это, утверждает он, объяснятся проистеканием авторитета добродетельных агентов из протоавторитета их эмоциональных реакций. Например, вам следует прислушиваться к советам образца сострадания, ведь они основаны на опыте сочувствия, которое он испытывает, воображая описанный вами случай. Шринивасан настаивает на важности безошибочности протоавторитетных эмоций: сочувствие не должно возникать в отсутствие кого-то находящегося в бедственном положении --- такие ошибки являются эмоциональным эквивалентом иллюзии (2020, 187). Очевидно, предполагается оценочное репрезентативное содержание эмоций.

Таким образом, добродетели связаны с диспозицией к переживанию подходящих эмоций: добродетель сострадания, например, заключается в диспозиции испытывать сочувствие к тому, кто его действительно достоин. Как понимать оценочное содержание, соответствующее сочувствию? Возможно, сочувствие возникает только тогда, когда его интенциональный объект находится в бедственном положении. Сочувствие также должно представлять это положение как нечто плохое, и может, значит, моделироваться на основе боли (см. Бэйн).

Не игнорируется ли в теории Шринивасана связь добродетелей с многонаправленными эмоциональными диспозициями?

\section{Важность заботы}

В описании сострадания Шринивасаном утверждается, что эта добродетель включает главным образом морально скорректированную черту сочувствия, то есть ядром сострадания является \textit{однонаправленная} эмоциональная диспозиция. Некоторые считают, что добродетели связаны с \textit{многонаправленными} диспозициями. Так, честность предполагает склонность чувствовать огорчение при столкновении с нечестностью, спокойное отношение к честной критике, гнев по поводу нечестных действий, радость по поводу честных и т.д..

Теории добродетелей, значит, делятся на

1. Однонаправленные: где добродетель включает в первую очередь однонаправленную эмоциональную диспозицию.

2. Многонаправленные: где добродетель включает в первую очередь многонаправленную эмоциональную диспозицию.

Я попытаюсь показать совместимость интуиций, лежащих в основе этих моделей.

Согласно многонаправленной модели, наличие добродетели означает наличие чувства, состоящего в уязвимости к приобретениям и потерям какого-либо объекта, будь то истина или благосостояние. Наличие добродетели, таким образом, означает заботу об этом объекте. Например, честный заботится об истинности, нечестного же она не заботит. Нечестный человек, значит, совершает аффективную ошибку, не заботясь о стоящем заботы.

Забота фактически представляет форму оценки (Сейдман 2009): честный заботится об истине потому что ценит ее. Доброжелательный же ценит благополучие других, что и заставляет его о нем заботиться и чувствовать радость, когда у них все хорошо и грусть, когда плохо.

Не следует ли понимать все добродетели многонаправленно? Например, понимать сострадание как включающее многонаправленную эмоциональную диспозицию. Можно даже подумать, что ''доброжелательность'' и ''сострадание'' --- два названия одной и той же добродетели. И согласно Шринивасан, добродетель сострадания ''характерно ориентированна [...] на благополучие другого'' (2020, 141). То есть ценность, на которую призвано реагировать сострадание --- та же, на которую реагирует доброжелательность --- благополучие других. Разве мы не ожидаем, что сострадательные агенты --- те, кто заботится о благополучии других, поскольку не только испытывает сочувствие к ним, когда их дела идут плохо, но и радуется, когда их дела идут хорошо? Похоже, добродетель сострадания должна включать многонаправленную эмоциональную диспозицию.

Очевидно, что сочувствие играет особую роль в добродетели сострадания. Не все компоненты многонаправленной эмоциональной диспозиции должны иметь равное положение. Возможно, эти диспозиции состоят из структурированных связок однонаправленных диспозиций, и особую роль играет какая-то одна из них. В случае сострадания черта сочувствия играет такую особую роль и лежит в основе этой добродетели. Иначе говоря, сочувствие характерно для добродетели сострадания.

Какова может быть эта особая роль? Возможно, черта сочувствия фундаментальна онтологически. В любом случае, сочувствие и его проявления имеют эпистемический приоритет: мы говорим, что некто обладает добродетелью сострадания на основании его проявлений склонности к сочувствию к находящимся в беде --- это, похоже, потому, что мы считаем сочувствие фундаментальным для этой добродетели онтологически, т.е. что остальные однонаправленные диспозиции необязательны для сострадания.

В любом случае есть основание утверждать, что добродетели в основном связаны с многонаправленными эмоциональными диспозициями, в которым диспозиция испытывать определенный вид эмоций играет особую роль. Добродетели можно рассматривать как чувства, состоящие в заботе о какой-то ценности, например, об истине или благополучии, причем эта забота, по крайней мере иногда, включает в себя эмоциональную предрасположенность, например сочувствие в случае сострадания, являющееся его отличительной чертой.

Все ли добродетели включают какую-то особую черту как часть многонаправленной эмоциональной диспозиции? Все ли добродетели вообще связаны с эмоциональными диспозициями? Рассмотрение этих вопросов потребует от нас понимания того, какие черты характера являются добродетелями и что вообще делает черту добродетелью. Согласно известному предположению, восходящему к Аристотелю, добродетели --- черты характера, необходимые для счастья или развития.

\section{Добродетели и счастье}

Многие теоретики добродетели считают, что она обязательно способствует счастью носителя (например, Херстхаус 1999, Аннас 2011). Некоторые считают, что быть добродетельным одновременно необходимо и достаточно для счастья. Но даже добродетельным агентам приходится преодолевать различные трудности, так что необходимость, но недостаточность добродетели для счастья, похоже, проще защищать. Но даже это утверждение может показаться сомнительным, если вспомнить о злонамеренных агентах --- неочевидно, что недостаток честности, сострадания, смелости и т.д. способствует их несчастью. Может даже показаться, что проявление пороков способствует их счастью. Дэниел Хейброн (2008) рассматривает случай императора Монгольской империи Чингисхана (ок. 1158-1227), который, судя по всему, вел счастливую и процветающую жизнь, хотя и был ответственен за убийство десятков миллионов людей. В любом случае действия, проистекающие из добродетелей, могут привести к самопожертвованию, а не к счастью или развитию. Добродетель и счастье, похоже, лишь условно связаны: в некоторых обстоятельствах проявление добродетели приводит к счастью, но обратной связи между ними нет.

Если говорить о добродетелях как о состоящих в основном из эмоциональных диспозиций, то можно аргументировать, что добродетель и счастье обязательно связаны: добродетели в силу своей природы имеют тенденцию приносить пользу людям. Термин ''счастье'' используется в двух разных смыслах (см. Хейброн 2008). В одном смысле счастье --- положительное психологическое состояние. Есть разные варианты пояснения, в чем оно заключается: в положительном балансе удовольствия и неудовольствия или в удовлетворенности своей жизнью, например. В другом смысле --- ''счастье'' синоним ''благополучия'' и ''процветания'', что соответствует эвдемонии у Аристотеля. Таким образом вопрос о соотнесении добродетели со счастьем распадается на два вопроса: как добродетели соотносятся с психологическим счастьем, и как соотносятся с процветанием.

Как добродетели могут быть связаны с психологическим счастьем? Дэниел Хейброн защищает теорию эмоционального состояния, согласно которой психологическое счастье состоит в положительном балансе аффективных состояний вообще и эмоциональных состояний в частности. Люди счастливы, когда положительные такие состояния --- будь то радость, надежда, восхищение, восторг или спокойствие --- уравновешивают негативные.

Поскольку ядром добродетелей являются эмоциональные диспозиции, то всякий раз, когда агент испытывает эмоцию в результате своей диспозиции, это влияет на его состояние психологического счастья / несчастья.

Раз добродетель приводит к возникновению эмоций, то эмоции, составляющие психологическое счастье или несчастье агента, являются устоявшимися --- когда психологическое счастье является результатом добродетелей, оно само по себе может рассматриваться как общее состояние в смысле соответствия оценочной реальности. Поэтому можно утверждать, что добродетели необходимо относятся к воспитанию, и уместное счастье является необходимым компонентом процветания, причем это совместимо с разными пониманиями процветания. Процветание можно понимать как, например, доступ к определенному множеству благ, таких как удовольствия, дружба и знания. Уместное счастье, значит, можно рассматривать как одно из благ, необходимых для процветания. Но это совместимо также и с теориями самореализации, ведь уместное счастье играет важную роль и в них (Хейборн 2008).

Таким образом:

1. Добродетели состоят в склонностях испытывать соответствующие эмоции.

2. Уместные эмоции составляют уместное счастье.

3. Уместное счастье является составное частью воспитания.

Значит добродетели, в основе которых лежат эмоциональные диспозиции, приводящие к положительным эмоциям, как правило делают жизнь человека лучше. Чтобы связать добродетели с воспитанием, необходимо требовать чтобы эти диспозиции приводили к уместным положительным эмоциям. Следует тогда уточнить, что только при благоприятных внешних условиях обладание добродетелями может привести к процветанию. Например, мир, в котором добродетельные агенты будут сталкиваться с нечестностью или страданиями, будет иметь тенденцию делать их психологически несчастными и негативно влиять на их развитие. Но мир, в котором они окружены честностью и заботой, сделает их счастливыми и процветающими.

Мы обсудили первую посылку аргумента. Вторая и третья содержат существенные и спорные утверждения. Тем не менее, аргумент показывает плодотворность подхода к добродетелям, ориентированного на эмоциональные диспозиции.

\section{Подведение итогов}

Эмоции могут быть как разрушительны, так и полезны. В Главе 7 мы рассмотрели возможность положительного влияния эмоций на рациональность нашего мышления и пришли к выводу, что наши убеждения насчет оценочных характеристик предварительно оправданы, когда основаны на соответствующей эмоции. Такие убеждения могут быть даже окончательно оправданы, если нет свидетельств, что эмоции вводят нас в заблуждение. Эмоции могут привести к действиям, противоречащим здравому смыслу, но даже эти действия могут быть благоразумными и даже морально правильными. В некоторых случаях акратические действия, вытекающие из эмоций, могут согласовываться с практическими соображениями агента. Есть также основания полагать, что эмоции полезны влиянием на внимание и провоцирование выбора.

В этой главе мы рассмотрели моральность и аморальность эмоций. Добродетели можно понимать как черты характера, в основе которых лежат эмоциональные диспозиции, причем эти диспозиции представляют собой заботу об определенных ценностях, таких как истина или благополучие.

Вывод таков: эмоции двулики, и только будучи уместными, могут быть частью добродетельной моральной агентности. В свете этого интересно, как начать испытывать уместные эмоции. Это вопрос, к которому мы обратимся в двух заключительных главах.

\section{Краткое содержание}

\begin{itemize}
  \item\ Приверженцы аристотелевской концепции добродетели считают, что добродетели --- сложные черты характера, заключающиеся в склонности испытывать правильные эмоции, а также иметь множество других соответствующих психических состояний, включая желания.
  \item\ Согласно сократической интерпретации, ядро добродетели составляет знание того, что хорошо и что плохо.
  \item\ Одна из проблем сократической интерпретации --- некоторые добродетели, например, скромность, требуют как раз-таки незнания. Похоже, дело в том, что добродетель вроде скромности зависит скорее от того, что нас волнует, нежели от того, что мы знаем.
  \item\ Другое возражение --- обвинение в чрезмерной интеллектуализации. Но добродетели можно понимать как виды навыков.
  \item\ Многие философы считают, что эмоциональные диспозиции являются ядром добродетелей. Шринивасан, например, утверждает, что эмоция сочувствия и соответствующая нравственно скорректированная черта занимает центральное место в добродетели сострадания.
  \item\ Многонаправленность добродетелей как соответствующих эмоциональных диспозиций и особая роль некоторых однонаправленных диспозиций не противоречат друг другу.
  \item\ Основанность добродетелей на эмоциональных диспозициях позволяет говорить об их (добродетелей) способствовании процветанию и психологическому счастью.
\end{itemize}

\begin{tcolorbox}
  \section{Учебные вопросы.}
  \begin{enumerate}
    \item\ Хотели бы вы жить с тем, кто соответствует радикальному рациональному идеалу моральной агентности и потому не испытывающему эмоций? Если нет, почему?
    \item\ Есть ли что-то неправильное в таких неприятных эмоциях как зависть, ревность или злорадство, если это никому не вредит? Если да, то почему?
    \item\ Вспомните морально достойный поступок Гека, когда он из сочувствия помог бежать Джиму из рабства. Считаете ли вы его морально добродетельным, или чего-то не хватает?
  \end{enumerate}
\end{tcolorbox}

\part{Искусство и формирование эмоций}

\chapter{Эмоциональная регуляция и музыка}

\section{Введение}

Вспомните последний фильм, который вы смотрели. Возможно, вы были заинтригованы, или же вас тронули невзгоды, которые пришлось пережить персонажам. Художественные произведения вообще и фильмы в частности могут заставить нас испытывать самые разные эмоции: от страха, печали и гнева до надежды и восхищения, и то же касается музыки. Однако, даже если некоторые эмоции, индуцируемые созерцанием искусства, негативны, мы продолжаем получать от этого созерцания удовольствие.

Как произведения искусства вообще и художественные произведения в частности могут вызывать эмоции? Как нас может волновать судьба персонажа, если мы знаем, что он вымышлен, и, значит, никаких страданий в действительности не переживает? Это так называемый парадокс вымысла. Произведения искусства не только вызывают эмоции, но и выражают их. Яркая картина или веселая мелодия могут выразить радость, тогда как темные цвета или мрачная музыка часто используются для выражения страха и печали. Но как и почему можно называть музыкальное произведение грустным или радостным? Что за печаль или радость оно выражает, чьи это эмоции? Считая, что для выражения состояния необходимо в нем находиться, мы все же не считаем, что произведения искусства буквально переживают эмоции. Еще одна странность --- привлекательность искусства, выражающего негативные эмоции. Грустная музыка, казалось бы, должна вызывать грусть, и оттого отталкивать, но пользуется большой популярностью. Это так называемый парадокс грустной музыки, он, похоже, аналогичен парадоксу трагедии, впервые сформулированному Аристотелем, и вопрошающему о причинах возможности привлекательности трагедий, вызывающих такие отталкивающие эмоции как ужас.

Рассмотрение связи искусства и эмоций важно для ответа на вопрос об эмоциональной регуляции. Искусство --- важный ее инструмент.

Как мы уже видели ранее, вопрос регуляции, наряду с вопросом о сущности и нормативным вопросом, является одним из трех фундаментальных вопросов философии эмоций. Цель этой главы --- изучить проблему эмоциональной регуляции, сосредоточив внимание на музыке, а также упомянув литературу и сентиментальное воспитание. Мы начнем с обзора науки об эмоциональной регуляции, затем обратимся к тому, как прослушивание музыки в ней помогает, рассмотрев различные способы индуцирования эмоций музыкой. Мы рассмотрим парадокс грустной музыки, а так же ее способствование возникновению ''состояния потока''.

\section{Наука об эмоциональной регуляции}

Философы прошлого прекрасно осознавали, что некоторые из наших эмоций вредны и что их необходимо регулировать, формировать. Аристотель считал, что приобретение добродетели требует воспитания эмоций, что предполагает обретение привычек: нам нужно привыкнуть испытывать правильные эмоции в правильных случаях. Стоики, гораздо более критично относясь к эмоциям и считая, что большая их часть укоренена в глубоко ошибочных суждениях, предлагали множество методов их искоренения: когнитивную терапию, выработку привычек и даже музыкальную терапию, которая, как утверждали, имеет свойство делать душу более рациональной. Фома Аквинский и Декарт указывали, что мы можем контролировать свои эмоции посредством размышления: например, представляя их причину в ином свете.

Современные психологи дополнили этот список. Джеймс Гросс (1998, 2014), ведущий деятель науки об эмоциональной регуляции, подчеркивает разнообразие ее типов. Он предлагает понимать эмоции как последовательность ситуация -- внимание -- оценка -- реакция. Любая эмоция, значит, начинается с ситуации, которая привлекает внимание и получает нашу оценку, что вызывает мультисистемную реакцию, включающую эмпирические, поведенческие и нейробиологические изменения. Так, обнаружение змеи в палатке приведет к эмоциональной реакции при условии, что вы обратите на нее внимание и расцените ее как угрозу. Во многих случаях этот эмоциональный отклик вызовет изменение ситуации: вы прогоните змею, например. Таким образом, эмоциональная последовательность характеризуется петлей обратной связи.

Эта последовательная модель совместима с рядом теорий эмоций. Она, например, может быть сформулирована в терминах рецептивной теории, согласно которой эмоции состоят в ощущаемых оценках, имеющих когнитивную основу (см. Главу 6). Согласно этой теории, последовательность начинается с реальной или воображаемой ситуации, которая привлекает ваши внимание. Понимание ситуации --- будь то визуальное восприятие, убеждение или какое-то другое подобное состояние, составляет когнитивную основу вашей ощущаемой оценки. В зависимости от типа эмоции и обстоятельств, эта оценка вызовет ряд реакций, таких как изменение выражения лица или мотивация.

Гросс различает несколько типов регуляции эмоций на основе этой последовательности. Начнем с разделения на внутриличностную --- направленную на свои эмоции, и межличностную --- направленную на эмоции других, регуляцию. Внутриличностная такая регуляция особенно важна для развития детей, для их нравственного воспитания, но важна и в нашем повседневном социальном взаимодействии.

Интересно, что эмоции, или, по крайней мере, их выражения, могут служить инструментами, с помощью которых мы пытаемся регулировать эмоции других. Например, вы можете намеренно проявить гнев на друга за то, что он разбил вазу, заставив его чувствовать вину, сожаление или даже страх. Подобные случаи не ограничиваются человеческим взаимодействием --- собака может агрессивно лаять на другую собаку с намерением вызвать страх. Внутриличностная регуляция также может начинаться с эмоций: если вам удастся посмеяться над собой из-за страха сделать что-то, вы сможете значительно этот страх уменьшить и даже преодолеть.

Другой набор различий, предложенный Гроссом (2014), касается влияния характеристик эмоций на процесс регуляции. Вот несколько аспектов, на которые можно повлиять: интенсивность, продолжительность и смещение мультисистемной реакции. Регулирование эмоций может, например, увеличить или уменьшить интенсивность гнева, но может также повлиять на выражение вашего лица или совершаемые действия.

По мнению Гросса (2014), эти различия проявляются во всех пяти семействах процессов регуляции эмоций: выбор ситуации, модификация ситуации, развертывание внимания, когнитивные изменения и модуляция реакции. Выбор и модификация ситуации относятся к первому элементу последовательности. Стратегии развертывания внимания --- ко второму. Примером такой стратегии является отвлечение, например, когда вы пытаетесь забыть о событии, которое вас раздражает, просматривая фильм. Ярким примером межличностной версии этой стратегии являются успешные попытки Яго вызвать ревность Отелло. Рональд де Соуза указывает, что хитрость Яго в направлении внимания Отелло на казалось бы доказательства неверности Дездемоны. Третье семейство процессов включает когнитивные изменения, целью которых является воздействия на элемент оценки. Такие изменения могут заключаться в изменении вашего мышления о ситуации, например, когда вам удается убедить себя, что экзамен --- не угроза, а интересное испытание, в результате чего происходит переоценка ситуации. Наконец, пятое семейство стратегий --- модуляция реакций --- направлено на прямое изменение мультисистемной эмоциональной реакции. Так, прогулка, глубокое дыхание или употребление алкоголя, могут быть использованы для модуляции физиологических изменений и ощущений. Как мы видели ранее в Главе 3, выражения лица зависят в том числе от правил поведения в обществе, что объясняет межкультурные между ними различия. Управление выражением лица, значит, является одной из возможных целей стратегий эмоциональной регуляции: подавление экспрессивности через подавление мимики и прочего соответствующего поведения --- распространенная стратегия.

Учитывая большое разнообразие стратегий, следует рассмотреть их преимущества и недостатки. Как отмечает Гросс (1998, 2014), некоторые стратегии имеют тенденцию быть более трудозатратными и проблемными. Например, подавление эмоционального поведения имеет более проблематичные последствия, чем переоценка --- стратегия, направленная на более ранние стадии развертывания эмоции. Подавление может ослаблять положительные, но не отрицательные эмоции, тогда как переоценка --- наоборот.

Цели усилий по эмоциональной регуляции могут быть очень разными. Есть, конечно, гедонистические мотивы: мы часто стремимся снизить интенсивность и продолжительность отрицательных эмоций и увеличить --- положительных. Но иногда нам может захотеться усилить свой гнев, чтобы легче противостоять сопернику, или отвлечься от смешной мысли во время похорон.

Вопрос, который не находится в центре внимания науки об эмоциональной регуляции --- как она может быть направлена на ощущение эмоций вообще. Однако, он может быть важен, если принять аристотелевское понимание добродетели, согласно которому эмоциональная регуляция равносильная своего рода обучению. Его мы рассмотрим в Главе 12.

\section{Эмоциональные эффекты музыки}

Когда мы пытаемся избавиться от негативной эмоции --- например, от грусти --- мы часто смотрим фильм, читаем роман или слушаем музыку. Как подчеркивают Патрик Юслин и Джон Слобода (2010), наиболее распространенным мотивом для прослушивания музыки является ее влияние на эмоции, то есть возможность с помощью нее регулировать и, значит, улучшать, свое психическое состояние.

Каким образом музыка может вызывать у нас эмоции? Давайте рассмотрим, как и почему мы испытываем эмоции от созерцания произведений искусства. Кажется странным, что вымышленное может вызывать у нас эмоции, ведь в обычных обстоятельствах, например, обнаружив, что бояться нечего, мы перестаем бояться. Но в случае с вымышленными персонажами мы испытываем самые разные эмоции --- сострадание и страх, например --- по отношению к, как знаем, вымышленному. Насколько искренними могут быть такие эмоции? Во всяком случае, совершенно ясно, что мы испытываем эмоции по поводу вымышленного.

С музыкой, однако, другая трудность. Музыка --- не изобразительное искусство, потому и непонятно, как ее прослушивание может вызывать эмоции. Ведь, в отличие от фильма или романа, она не рассказывает нам истории о вещах и людях. Музыкальное произведение, конечно, может сопровождаться текстом или названием, но как таковое оно ничего не описывает и не изображает --- даже произведение без названия и прочего контекста вполне способно вызвать эмоции, несмотря на отсутствие всякого повествовательного элемента.

Как музыка может вызывать эмоции? Как отмечает Стивен Дэвис (2010), некоторые способы, которыми она вызывает эмоции, вовсе не загадочны: музыкальное произведение может вызывать случайные личные ассоциации, например, когда его играли на вашей свадьбе. Ассоциации могут быть и общими для группы или культуры --- у хоккейных болельщиков есть наверняка стойкие ассоциации между победой и национальным гимном. Более того, оценочные свойства музыкального произведения, такие как его репитативность или неожиданная смена ритма, могут вызывать скуку или удивление. Эмоции может вызвать и удачное исполнение трудного отрывка. Такие эмоции, заключаясь в реакциях на вещи и их оценочные свойства, не столь уж загадочны.

Стоит упомянуть еще две причинно-следственные связи музыки и эмоций. Во-первых --- эффект вовлечения: музыка подразумевает ритм, так что ее прослушивание напрямую влияет на периферическую нервную систему. Этот эффект может привести к аффективным состояниям: прослушивание музыки с энергичным ритмом не только подталкивает двигаться соответствующим образом, но и вызывает чувство приподнятого настроения.

Более того, как отмечает Клаус Шерер (2013), прослушивание музыки может вызывать эмоции из-за процесса сопереживания, включающего воображение себя на месте другого. Вызываемые так эмоции особенно важны в случае живой музыки, где вид музыканта, выражающего эмоции, индуцирует эти же эмоции у публики. Чувства может вызвать и воображение себя на месте композитора.

Невролог Изабель Перец (2010) утверждает, что музыка может кооптировать эмоциональные цепи в мозгу, которые развились для эмоциональной реакции на голосовые выражения. Среди людей есть поразительная инвариантность в суждениях о том, является ли музыка радостной или пугающей: даже если индивидуальный опыт и культура играют роль в наших реакциях на музыка, склонность эмоционально реагировать кажется врожденной. Так, эмоциональная реакция младенцев на музыку может быть обусловлена врожденной склонностью определенным образом реагировать на определенный тон голоса.

Есть еще один способ воздействия музыки на нас, касающийся музыкальной выразительности. Так, похоронный марш из Третьей Симфонии Бетховена считается образцовым примером музыкального произведения, выражающего печаль. Симфония, конечно, сама не переживает никаких эмоций, так что не совсем ясно, как она может их выражать. Для ответа на этот вопрос был предложен ряд теорий: теории возбуждения, теории личности и теории сходства.

Согласно теориям возбуждения, выражение эмоций в музыке следует понимать с точки зрения эмоций, вызываемых у аудитории. Грустная музыка, значит --- та, что вызывает печаль, веселая --- та, что радует, страшная --- та, что вызывает страх и т.д.. Основная трудность этого подхода касается случаев, когда аудитория, казалось бы, ничего такого не чувствует, хотя произведение все же явно выражает чувства. Защитники теорий возбуждения отвечают, что реакция возбуждения может быть чрезвычайно слабой и состоять просто из ощущений, связанных с очень тонкими телесными изменениями или даже этих изменений нейронной симуляцией. Тем не менее, трудно поверить, что музыка может выражать печаль только в той мере, в какой мы чувствуем хотя бы минимальную этой печали степень. Параллельно с описанием оценочных свойств, вызывающих диспозиции (см. Главу 9), можно утверждать, что грустная музыка склонна вызывать печаль в определенных обстоятельствах, и тогда следует описать эти обстоятельства. Что тогда вернее: что грустная музыка вызывает грусть в стандартных обстоятельствах у стандартных слушателей, или что нам грустно, когда мы ее слушаем? Но множество людей слушают грустную музыку потому что она улучшает их психологическое состояние.

Теории личности отрицают, что нам нужно что-либо чувствовать, слушая выразительную музыку. Для того, чтобы музыкальное произведение считалось выражающим эмоции, достаточно чтобы мы слышали, что оно выражает эмоции воображаемого человека. Так, слушая грустную музыку, мы можем представить такого ''музыкального персонажа'', который чувствует грусть. Но такого рода воображение не кажется необходимым для понимания выразительности музыкального произведения.

Теории сходства утверждают, что музыка \textit{похожа} на то, как люди выражают свои эмоции: музыкальное произведение считается выражением эмоций в силу сходства его динамической структуры с таковой у поведения людей. Есть не только ощутимое сходство между мелодиями и вокальным выражением, но сходство контуров музыкальных произведений с движениями тела. По словам Дэвиса, ''музыка напоминает походку, манеру поведения или речи, которые обычно выражают возникающие у человека эмоции [...]'' (2010, 33). Согласно Дэвису, это общее явление --- мы часто рассматриваем вещи как выражающие эмоции на основе их внешнего вида: например, свисающие ветви деревьев напоминают грустных людей, так что ивы выглядят грустными.

Таким образом, теория возбуждения совместима с несколькими теориями о том, как выразительность музыки приводит к переживанию эмоций. Так, Том Кокрейн защищает теорию возбуждения, используя элементы теории личности и теории сходства. Он пишет: ''музыка захватывает механизм симуляции, который распознает действия и эмоции других'' (2010, 203), что очень близко к утверждению о кооптировании эмоциональных систем для реакции на человеческие голоса (см. Перец 2010). По мнению Кокрейна, возникновение эмоций от прослушивания музыки подобно эмоциональному заражению, затрагивающему зеркальные нейроны, возбуждающиеся как при совершении определенных действий нами, так и при наблюдении совершения этих действий другими: чувствуя радость или страх от радости или страха другого, мы на самом деле бессознательно имитируем выразительные его реакции. Подобный подход доступен и приверженцу теории личности: испытываемые нами эмоции при прослушивании музыки возникают посредством эмоционального заражения от воображаемого ''музыкального персонажа''. Дэвис (2010, 2013) утверждает, что мы ''подхватываем'' посредством процесса эмоционального заражения эмоции, выраженные в музыкальном произведении, и потому музыка, выражающая печаль, вызывает печаль и у нас --- но, в отличие от ''полномасштабной'' эмоции печали, эта не будет направлена на какой-то конкретный объект. Дэвис пишет: ''если слушатель отражает выразительность музыки, эта его реакция вызвана музыкой и отслеживает ее, но не принимает ее или любую другую вещь в качестве своего эмоционального объекта'' (2010, 37). То есть, как еще утверждает Дженефер Робинсон (2005), аффективные состояния, вызванные музыкой, на самом деле не эмоции, а настроения.

Мы обнаружили несколько способов, которыми музыка может вызывать эмоции и увидели, в частности, что она может вызывать настроения, отражающее выраженное в музыкальном произведении. Следует подчеркнуть, что чувствуемое слушателем, скорее всего, является результатом взаимодействия нескольких аффективных состояний, вызываемых разными причинными механизмами: например, музыкальное произведение может вызвать счастливые воспоминания, но также иметь энергичный ритм. Вообще, мы часто заранее знаем, как прослушивание определенной музыки повлияет на наше психическое состояние, и сознательно используем это для получения желаемого эффекта.

\section{Парадокс грустной музыки}

Зачем нам слушать грустную музыку, если она нас только расстраивает? Мы же обыкновенно избегаем грусти. Парадокс грустной музыки может быть приведен так:

\begin{enumerate}
  \item\ Грустная музыка вызывает грусть.
  \item\ Грусть --- неприятное аффективное состояние.
  \item\ Мы намеренно слушаем грустную музыку.
\end{enumerate}

Есть ряд решений этого парадокса, некоторые из которых соответствуют решениям, предложенным для парадокса отрицательных эмоций в искусстве.

Самый распространенный ход состоит в принятии всех трех допущений и попытке показать, что они совместимы. Основное предположение здесь состоит в компенсации негативного аффективного состояния некоторой от него пользой. Такие решения можно найти еще у Аристотеля, который, как известно, утверждал, что трагедия способствует катарсису. Ужас и жалость, испытываемые при наблюдении за трагической судьбой персонажей, позволяют нам очиститься от негативных эмоций. То же можно и предположить в отношении музыки, конечно. Как указывает Джерролд Левинсон (1982), реакция на музыку может позволить человеку контролируемым образом выпустить определенное количество вредных эмоций, которые он переживает (1982, 338). Левинсон считает, что боль от негативных аффективных состояний, вызванных музыкой, компенсируется множеством других ''вознаграждений''. Фактически, согласно Левинсону, по крайней мере некоторые преимущества, предлагаемых прослушиванием музыки --- переживание определенных эмоций, ''общение'' с композитором и т.д. --- нельзя получить без сопутствующих негативных переживаний. Но, казалось бы, даже если музыка может предложить подобное ''вознаграждение'', это происходит не всегда.

Другое схожее решение начинается с указания на наличие негативных аспектов у множества ценимых нами вещей. Те, кто любит участвовать в марафонах, не отрицают, что подготовка к ним и само в них участие имеет как положительные, так и отрицательные аспекты, но испытываемая боль, например, в значительной степени компенсируется общей ценностью усилий. Нечто подобное, можно предположить, происходит и в случае с музыкой. Дэвис, например, считает, что наша цель --- понять и оценить музыку, и это то, что мы ценим, даже если это требует переживания негативных эмоций. Это решение, похоже, указывает в правильном направлении, поскольку фокусируется на ценности прослушивания музыки. Однако сомнительно, что многие из нас ставят перед собой такую цель --- понять музыку.

Вообще говоря не обязательно принимать все положения, образующие парадокс. Может показаться, что преимуществом взгляда Дэвиса и прочих компенсаторных решений является возможность объяснить взаимодействие не только с грустной музыкой, но и со всякой другой, вызывающей негативные эмоции. Но в случае с грустной музыкой есть нечто особенное --- она особенно привлекательна, и, похоже, есть что-то неправильное в мысли, что ее прослушивание приводит к негативному опыту, который необходимо, значит, как-то компенсировать.

Давайте тогда рассмотрим решения, отрицающие некоторые из его пунктов. Лаура Сайзер (2019) утверждает, что грусть, вызванная грустной музыкой --- не так уж неприятна. Учитывая замечание Левинсона (1997) об отделенности печали от грустной музыки от повседневных забот и ее таким образом ''чистоте'', Лаура считает, что эта печаль не только усиливает сосредоточение на музыке, помогая лучше в нее ''погрузиться'', но и может быть позитивна и утешительна. Вопрос здесь в том, является ли сама грусть составляющей позитивный опыт, или же ее следует рассматривать как часть или результат в целом приятного опыта. Если грусть, вызванная музыкой, сама по себе неприятна, то предлагаемое решение на самом деле относится к компенсаторным. Согласно некоторым исследованиям (Вуоскоски и др., 2012), прослушивание грустной музыки вызывает реакции печали, которые не воспринимаются как негативные или неприятные, в отличие от эмоциональных реакций на пугающую музыку, оценивающихся как явно неприятные. Может показаться, что это усиливает аргументы против второго пункта парадокса, но это может и означать необходимость сомнения в первом его пункте. Может ли музыка вызывать печаль, если пережитое от ее прослушивания не неприятно?

Питер Киви (1990, 2006) считает, что музыка не может вызывать ни того, что он называет ''обычными эмоциями'', ни настроения. Он признает, что красота и другие положительные эстетические свойства могут вызывать чувство трепета, но эти чувства не связаны с аффектом, выраженным в музыке. Поскольку его аргумент предполагает спорную теорию эмоций, согласно которой эмоции являются пропозициональными установками, он не столь уж убедителен. Вообще, трудно поверить, что музыка не может вызывать ни эмоций, ни настроений.

Отрицание способности грустной музыки обыкновенно вызывать грусть противоречит эмпирическим исследованиям, согласно которым оценки переживаемой печали были значительно выше, чем оценки других эмоций (но см. другие результаты, согласно которым грустная музыка воспринимается как выражение грусти, но не как ее переживание вызывающая в Зентнер и др. 2008). Гораздо надежней считать, что грустная музыка вызывает реакцию, которая в некотором роде схожа с печалью, но заключается во все же положительном аффективном состоянии. Недавние эмпирические исследования показывают, что типичной реакцией на грустную музыку является чувство волнения, тронутости. Согласно Вуоскоски и Ээроле (2017), ощущение тронутости играет центральную роль в наслаждении, получаемом нами от прослушивания грустной музыки. В чем состоит это ощущение? Оно обычно вызывается критическими жизненными событиями, такими как рождение, смерть и брак, а также созерцанием произведений искусства и природы.

Более пристальный взгляд на природу чувства тронутости показывает, что его можно рассматривать не как часть механизма, приводящего нас от печали к удовольствию, а как единую реакцию на грустную музыку. То есть грустная музыка обыкновенно вызывает вовсе не грусть, а чувство волнения. Согласно Деонне (2020; см. также 2018), волнение --- отдельная эмоция, состоящая в ''переживании пораженности какой-то конкретной представленной положительной ценностью, ее важностью'' (2020, 190).

Особый интерес представляет то, что люди, сообщающие о своем опыте такого волнения, склонны называть его в целом приятным, даже если он включал как положительные, так и отрицательные эмоции --- как если бы тронутость могла сопровождать некоторую печаль наряду с радостью, испытываемой при созерцании чего-то крайне ценного и великого. Возможно дело в том, что ценное выглядит еще более ценным на в целом негативном фоне. Поскольку тронутость часто связана со смешанными чувствами, есть основания полагать, что это типичная реакция, испытываемая при прослушивании красивой, но грустной музыки. Согласно Деонне, слушая грустную музыку, мы просто таки поражаемся красотой и возвышенностью воплощенного в ней.

Это легко объясняет, почему волнение, включая чувство грусти, заставляет многих думать, что прослушивание грустной музыки заставляет их грустить, хотя грусть эта, казалось бы, не переживается как нечто негативное. Это также объясняет, почему грустная музыка трогает нас сильнее, чем веселая или пугающая: вполне вероятно, что осознание печали, выраженной в ней, может вызвать эмоцию, включающую в себя в том числе и печаль. Действительно, грусть, пронизывающая такую музыку, служит темным фоном, подчеркивающим ее красоту.

Следует отметить, что это решение распространяется и на другие формы искусства, выражающие печаль. Даже если задействованные механизмы, вероятно, будут разными, грустные повествования часто трогают нас так же, как и музыка. Может ли подобное решение разрешить парадокс не только грусти, но и прочих негативных эмоций, переживаемых при созерцании искусства?

\section{Музыка, состояние потока и счастье}

Многим из нас нравится слушать музыку, будь то грустная, радостная или выражающая какие-то другие аффективные состояния. Как указывает Михай Чиксентмихайи (1975, 1990), прослушивание музыки может вызвать ощущение потока, являющееся, похоже, особым видом удовольствия --- получаемым от определенных видов деятельности (Тапполет 2022). Его могут вызвать самые разные занятия --- прослушивание музыки, скалолазание, парусный спорт, шахматы, видеоигры, танцы, рисование, лепка, сочинение и письмо. Прослушивание музыки, значит, может участвовать в эмоциональной регуляции через такое состояние потока.

Музыку связывали с состоянием потока с самого начала его исследований (см. Чиксентмихайи, 1975). Хотя значительная часть работы Чиксентмихайи посвящена сочинению музыки и игре на музыкальном инструменте, он также учитывает и влияние ее прослушивания. Так, он иллюстрирует необходимость принятия некоторого вызова для возникновения состояния потока: ''музыкальное произведение, если слишком простое для слушателя, будет скучным, а слишком сложное будет вызывать фрустрацию, удовольствие же возможно на границе между скукой и тревогой, когда проблемы легко уравновешиваются способностями к их решению'' (1990, 52).

В одном из своих ранних описаний состояния потока Чиксентмихайи называет его ''оптимальным опытом'', ''целостным ощущением, которое испытываешь, действуя с совершенной вовлеченностью'' (1975, 36). Агент в состоянии потока переживает свое действие ''как единое течение от одного момента к другому, в котором он контролирует свои действия и в котором мало различий между собой и окружающей средой, между стимулом и ответом, между прошлым, настоящим, и будущим''.

По мнению Чиксентмихайи и других исследователей (Чиксентмихайи 1975, Накамура и Чиксентмихайи 2012, Энгезер и Шипе-Тиска 2012), переживание потока зависит от следующих двух условий:

\begin{itemize}
  \item\ Воспринимаемые проблемы соответствуют имеющимся навыкам.
  \item\ Есть четкие ближайшие цели и немедленная обратная связь о прогрессе в их достижении.
\end{itemize}

Кроме того, утверждается, что для состояния потока свойственны:

\begin{itemize}
  \item\ Интенсивная и целенаправленная концентрация на том, что делаешь до ощущения отрезанности от прочего мира.
  \item\ Слияние действия и самоосознания, агент не воспринимает себя отдельно от производимого действия.
  \item\ Утрата рефлексивного самосознания, размышления о себе становятся вовсе неактуальны.
  \item\ Ощущение контроля над своими действиями и требованиями окружающей среды, ощущение превосходства, способности справляться с задачей.
  \item\ Искажение временного опыта: время, кажется проходит быстрее.
  \item\ Восприятие деятельности как внутренне ценной, так что нет необходимости в прочих целях и вознаграждениях.
\end{itemize}

Несколько удивительно, но интуитивно важный аспект --- удовольствие --- не фигурирует в текущих списках свойств состояния потока. Чиксентмихайи изначально намеревался изучить природу и условия наслаждения, и в первой своей книге писал, что состояние потока ''является важнейшим компонентом удовольствия'' (1975, 11), и что деятельность, это состояние провоцирующая, воспринимается как очень приятная. Чиксентмихайи допускает, конечно, и другие, более пассивные виды наслаждения, но состояние потока ассоциируется с высшей его формой.

Как уже говорилось выше, решающее значение для состояния потока имеет воспринимаемый баланс между сложностью деятельности и навыками, что приводит к ощущению контроля и мастерства. Когда задача воспринимается как слишком сложная, возникает тревога, когда как слишком простая --- скука. Состояние потока же противоположно тревоге и скуке, так что, похоже, тоже является эмоцией.

То же следует из более внимательного рассмотрения свойств потока --- становится ясно, что опыт потока имеет все основные черты эмоций (см. Главу 2). Он обычно характеризуется феноменальными свойствами, ведь включает не только чувство контроля, но и чувство самоценности деятельности. Он обычно включает и мотивацию --- желание продолжить деятельность или вернуться к ней позже. Интенциональный объект --- деятельность, порождающая состояние потока, она оценивается положительно. Формальным объектом состояния потока, возможно, является не что иное, как приносящее удовольствие от соответствия навыкам. Состояние потока, как положительная эмоция, делает непосредственный вклад в психологическое счастье.

Музыка, значит, когда подходит для вхождения в состояние потока, является приятным времяпрепровождением. Когда вы поглощены музыкой, ваше внимание полностью ей захватывается и остальной мир, кажется перестает существовать. Это слияние приводит к потере интереса к себе как к объекту рефлексивного самосознания, время проходит быстрее.

Как прослушивание музыки вызывает состояние потока? Ведь, кажется, прослушивание недостаточно деятельно, не требует сколь-нибудь серьезного напряжения. В конце концов, прослушивание музыки --- не математическая задача. Возможно, прослушивание музыки включает то, что можно назвать ''внутренним танцем''. Как отмечает Сайзер, ''до изобретения традиционного западного концертного зала, предполагающего, что слушатели будут сидеть неподвижно, музыка и движение были тесно связаны'' (2019, 258). Увлеченность как механизм воздействия музыки на эмоции играет важную роль и часто приводит к телесным движениям. Второй важный механизм --- музыкальное ожидание (см. Мейер 1956, Гурон 2006): слушатели имеют ряд ожиданий от музыки, которые зависят от их знакомства с произведением и его стилем. Эти ожидания либо оправдываются, либо нет, и обыкновенно предпочитают музыку не слишком предсказуемую. При прослушивании музыки работают оба эффекта --- и увлечение, и ожидание, и именно их комбинация позволяет танцевать, а не просто топать ногой или делать какие-то другие простые движения. Поскольку танец --- занятие, требующее усилий и навыков, то может объяснить способствование музыки вхождению в состояние потока.

Мы, конечно, можем не танцевать и даже не притопывать ногой, слушая музыку, но все еще возможен так называемый внутренний танец. Джекендофф считает, что корреляция музыки и танца вызывает ''музыкальный аффект --- эмоциональную реакцию на музыку --- потому что она провоцирует репрезентации телесных движений независимо от их перевода в действительные движения'' (1992, 16). Задача, связанная с прослушиванием музыки, таким образом, состоит в умении удачно исполнить этот по крайней мере внутренний танец, синхронизируясь с музыкальным произведением.

Веселая, грустная, и даже пугающая музыка может вызвать состояние потока, и наличие этого эффекта зависит от знакомства слушателя с музыкой и ее стилем, а также он навыков прослушивания, так что и не удивительно, что музыкальные вкусы могут столь различаться. В некоторых случаях, конечно, эффекты могут комбинироваться: например, когда грустная музыка не только трогает, но и вызывает состояние потока. Так что неудивительно, что музыка является одним из популярнейших инструментов эмоциональной регуляции.

\section{Подведение итогов}

Мы рассмотрели, как искусство вообще и музыка в частности, могут использоваться для эмоциональной регуляции. Музыка способна заставить нас испытывать сильные аффективные состояния, некоторые из которых чрезвычайно приятны. Грустная музыка, например, далеко не всегда вызывает негативное переживание печали --- обычно она вызывает волнение, положительную эмоцию от восхищения ее красотой. Как правило, прослушивание музыки провоцирует вхождение в состояние потока, которые, похоже, является своего рода удовольствием от деятельности. Музыка, значит, инструмент, оказывающий мощное влияние на наше психическое состояние. В следующей главе мы выясним, как она соотносится с другими формами искусства --- кино и литературой, с точки зрения эмоциональной регуляции.

\section{Краткое содержание}

\begin{itemize}
  \item\ Эмоциональная регуляция заключается в формировании наших эмоций, например, в отношении того, когда и как долго мы их испытываем и как мы их выражаем.
  \item\ Она может быть внутриличностной --- направленной на свои эмоции, или межличностной --- направленной на эмоции других.
  \item\ Существует пять семейств стратегий эмоциональной регуляции, касающихся различных элементов эмоциональной последовательности: выбор ситуации, модификация ситуации, развертывание внимания, когнитивные изменения и модуляция реакции.
  \item\ Цели эмоциональной регуляции могут включать гедонистическое наслаждение, но могут его и не затрагивать.
  \item\ Музыка может вызывать эмоции несколькими способами: вызывая личные ассоциации или демонстрируя такие свойства, как скучность или виртуозность, а также посредством вовлечения, сопереживания и заражения.
  \item\ Грустная музыка может вызвать чувство волнения --- в целом приятную эмоцию, испытываемую от красоты музыки, подчеркнутой мрачным фоном, что объясняет, почему так много людей любят ее слушать.
  \item\ В целом, прослушивание музыки часто провоцирует вхождение в ''состояние потока'', заключающееся в своего рода удовольствии от не слишком простой, но и не слишком сложной, деятельности.
\end{itemize}

\begin{tcolorbox}
  \section{Учебные вопросы.}
  \begin{enumerate}
    \item\ Вспомните, как вы пытались избавиться от каких-то негативных эмоций, например от гнева или отвращения. Что это были за эмоции, и какие стратегии вы использовали и насколько успешно?
    \item\ Создает ли страшная музыка парадокс, аналогичный парадоксу грустной музыки, и если да, то как мы можем его разрешить?
    \item\ Какую музыку вы бы порекомендовали другу, чтобы он почувствовал себя более счастливым, и как бы обосновали свой выбор?
  \end{enumerate}
\end{tcolorbox}

\chapter{Сентиментальное воспитание и художественная литература}

\section{Введение}

Музыка --- мощный инструмент для регуляции эмоций, но другие формы искусства --- особенно повествовательные --- имеют сравнимый потенциал. Просмотр фильма или чтение романа могут, так же как и музыка, вызывать волнение и погружать в состояние потока. Таким образом, повествовательное искусство является важным источником счастья, чем не отличается существенно от музыки.

Но повествовательные формы искусства влияют на нас иначе, нежели музыка. Мы склонны реагировать на истории: любопытством по поводу того, что произойдет далее, страхом, надеждой или печалью, в зависимости от того, как разворачивается сюжет. Загадочна здесь возможность испытывать эмоции по поводу вымышленного при совершенном осознании его вымышленности. В этой главе, значит, мы рассмотрим так называемый парадокс вымысла. Мы также рассмотрим, как и почему нам может нравится повествовательное искусство, вызывающее крайне негативные эмоции, такие как страх и отвращение --- парадокс трагедии, а также более общую его версию --- парадокс негативных эмоций, а также сходства и различия их возможных решений с таковыми парадокса грустной музыки.

Предыдущая глава касалась возможности и способов эмоциональной регуляции, в этой же рассмотрим, как искусство может способствовать развитию наших эмоциональных диспозиций и сентиментов. Распространено убеждение, что повествовательное искусство вообще и литература в частности, вносят важный вклад в наше моральной развитие, но давайте начнем со скептическим подходом: на первый взгляд неочевидно, что проводящие много времени за чтением романов или просмотром фильмов, обладают каким-то особым моральным превосходством.

Мы начнем с отличий воспитания эмоций от их регуляции, затем обсудим, в чем заключаются изменения эмоциональных диспозиций, а так же рассмотрим основные стратегии, позволяющие таких изменений добиться. Понимание способности вымысла вызывать эмоции потребует разрешения парадокса вымысла, после чего мы рассмотрим парадокс трагедии.

\section{Что такое сентиментальное воспитание?}

Идея о необходимости воспитания эмоций восходит к Аристотелю и занимает центральное место не только в этике добродетели, но и в теории добродетели (Херстхаус 1999, Кристьянсон 2018, Шринивасан 2020). Она играет важную роль в любом этическом подходе, признающем зависимость наших действий от испытываемых эмоций: для сторонника консеквенциализма, например, вполне естественно утверждать, что часто следует полагаться на чувство сострадания, потому что именно оно способствует преумножению блага. Консеквенциалист, значит, может лишь надеяться, что возможен процесс, заставляющий чаще испытывать сострадание. Что именно включает в себя воспитание эмоций и чем оно отличается от эмоциональной регуляции?

Регуляция эмоций заключается в их формировании: изменении их длительности, интенсивности и соответствующих поведенческих и нейробиологических изменений (Гросс 2014). Такая регуляция может касаться любого элемента эмоциональной последовательности, переводящей посредством концентрации внимания и оценки, от ситуации к мультисистемной реакции, включающей эмпирические, поведенческие и нейробиологические изменения. То есть регуляция имеет дело с эмоциональными эпизодами. Эмоциональные эпизоды, как указывалось в Главе 2, отличаются от эмоциональных диспозиций, которые являются не самими эмоциями, а предрасположенностями к этим эмоциям возникновению. И эти диспозиции также можно изменять.

Способ изменения эмоциональных диспозиций мы и будем называть воспитанием эмоций (см. Робинсон 2005). Обычно сентиментальное воспитание связывают с чувством любви, как это показано во влиятельном романе Густава Флобера XIX века ''Сентиментальное воспитание'', но в этой главе мы будем использовать этот термин для обозначения, конечно, воспитания всех видов эмоциональных диспозиций, как однонаправленных, так и многонаправленных.

Основное различие между эмоциональной регуляцией и воспитанием эмоций касается их целей. Регуляция позволяет достигать самых разных целей --- как кратковременных гедонистических, так и более сложных стратегических. Воспитание же подразумевает стремление испытывать правильные --- соответствующие некоторым нормативным стандартам --- эмоции.

Учитывая эту разницу, можно различить четыре способа формирования того, что мы чувствуем:

\begin{enumerate}
  \item\ \textit{Эмоциональная регуляция} --- регуляция эмоциональных эпизодов, то есть эмоций, с произвольной целью.
  \item\ \textit{Воспитание эмоций} --- регуляция эмоций с целью согласования их с нормативными стандартами.
  \item\ \textit{Сентиментальная регуляция} --- регуляция эмоциональных диспозиций с произвольной целью.
  \item\ \textit{Сентиментальное воспитание} --- регуляция эмоциональных диспозиций с целью согласования их с нормативными стандартами.
\end{enumerate}

Поскольку цель переживания правильных эмоций --- не единственная, которую можно преследовать при регуляции эмоций, изменения, направленные на переживание правильных эмоций --- подкатегория изменений эмоций, к которым мы можем стремиться, так что воспитание эмоций --- своего рода регуляция эмоций. Аналогично, изменения, направленные на возникновение правильных эмоций путем вмешательства на уровне диспозиций --- подкатегория изменений диспозиций, к которым мы можем стремиться., и, значит, сентиментальное воспитание --- одна из форм сентиментальной регуляции. В оставшейся части главы мы сосредоточимся на сентиментальном воспитании.

Что считать правильными эмоциями? Возможно, это те, что рекомендованы социальными нормами группы, к которой принадлежит агент. Такие нормы могут иметь целью социальную кооперацию, например, но не обязательно, и, вообще, подобным образом направленные нормы могут приводить к нежелательным последствиям: возможны социальные нормы, согласно которым определенная группа должна испытывать покорный страх при взаимодействии с доминирующими членами общества.

Возможно, правильные эмоции --- те, что хороши с моральной точки зрения (о какого рода нормах тогда говорилось в абзаце выше? --- прим. переводчика). Сентиментальное воспитание, значит, направлено на формирование морально корректных эмоций (Кристьянссон, 2018). Впрочем, без основательной моральной теории это определение едва ли практически применимо. Впрочем, есть смысл сосредоточиться на эмоциях, правильных в смысле уместности --- то есть репрезентационно корректных, соответствующих действительным оценочным характеристикам. Так, страх уместен в ответ на опасный или страшный объект, а восхищение уместно на нечто восхитительное. Это понятие уместности --- центральное в аристотелевском понимании добродетели (см. Главу 10).

\section{Сентиментальное воспитание как калибровка}

Вопрос сентиментального воспитания часто связывается с воспитанием детей. Оно, конечно, имеет первостепенное значение, ведь их диспозиции более податливы, чем таковые взрослых. Но и эмоциональные диспозиции взрослых людей допускают изменения. Исследования расовых предубеждений с участием взрослых участников показывают, что реальных опыт общения с членами расово чуждой группы может уменьшить расовые предубеждения (Цзо и Хан 2013). Гопал Шринивасан (2020) утверждает, что страдания членов чуждой группы могут, посредством сочувствия, способствовать уменьшению таких предубеждений. Таким образом, сентиментальное воспитание возможно на протяжении всей жизни человека.

Как мы уже видели, как биологические детерминисты (Экман 1999, Гриффс 1997, Туби и Космидес 2000), так и социальные конструктивисты (Аверилл 1980, Харре 1986, Армон-Джонс 1986) --- признают податливость эмоциональных диспозиций (см. Главу 3). Так, биологические детерминисты считают, что эмоции зависят от врожденных свойств, возникающих в результате эволюционного отбора, но могут далее формироваться средой и культурой. Социальные конструктивисты утверждают, что эмоции --- продукт социальных практик, так что допускают еще большую гибкость их диспозиций.

Вопрос сентиментального воспитания отличается от вопроса изначального возникновения эмоциональных диспозиций. Биологические детерминисты указывают, что рождаемся с некоторыми уже имеющимися диспозициями к таким эмоциям, как страх от громких звуков. Но по мере развития мы учимся ассоциировать со страхом и вовсе другое --- пауков и высоту, например.

Социальные конструктивисты, с другой стороны, говорят о социальных факторах. Похоже, что мы рождаемся с предрасположенностью реагировать не полноценными эмоциями, и эти реакции становятся эмоциями только когда связываются с социальными сценариями. Согласно Рональду де Соузе (1987), такие реакции, чтобы стать эмоциями, должны быть связаны с типами ситуаций, пережитыми в детстве --- такие ситуации определяют характеристики переживаемых в будущем полноценных эмоций. Эта связь зависит в том числе от усвоения терминов, обозначающих эмоции: например, дети могут выучить термин ''страх'' и связывать его с грозой и прочими соответствующими сценариями. Соуза утверждает, что сценарии эти, даже если отчасти детерминированы, во многом зависят от контекста и социально информированы: они изначально черпаются из повседневной жизни, и уже затем подкрепляются историями, искусством и культурой, с которыми нам приходится сталкиваться, постепенно обогащаясь (1987, 182).

Все эти вопросы развития, конечно, интересны, но почти не касаются сентиментального воспитания. Ни детерминисты, ни конструктивисты, не говорят, как происходит переход от имеющийся диспозиции касательно класса объектов к той же диспозиции касательно более узкого или широкого класса. Они, конечно, допускают такие изменения, но еще предстоит выяснить, как они возможны.

Культурное формирование эмоциональных диспозиций может происходить через смешение или через когнитивную разработку (см. Главу 3, Раздел 7). Смешение заключается в объединении двух или более эмоциональных реакций для получения новой (Плутчик 1980, Принц 2004): например, эмоцию, которую жители Ифалука называют фаго, можно понимать как смесь любви, сострадания и печали (Лутц 1988). Когнитивная разработка сочетает некий когнитивный компонент --- убеждение, например --- и эмоции, для получения новой эмоции (Джонсон-Лэрд и Оатли, 2000): возмущение, например, подразумевает убеждение в несправедливости и гнев, а вина --- страдание от убеждения в неправильном своем поступке. Эмоции, полученные в качестве когнитивной разработке и достаточно устойчиво закрепившиеся, можно ощущать и уже без соответствующих убеждений --- когда процесс ''обучения'' завершен, они возникают в определенных ситуациях вне зависимости от наличия соответствующих убеждений.

Смешение и когнитивная разработка, похоже, играют важную роль в культурном формировании эмоций, но есть более фундаментальный процесс --- добавление новых или удаление имеющихся триггеров эмоций --- ''калибровка''. Смешивание и когнитивная разработка, в отличие от калибровки, используют имеющиеся эмоциональные диспозиции для получения новой. Так, смешение для получения фаго приводит к слиянию элементов различных эмоциональных диспозиций --- диспозиций чувствовать любовь, сострадание и печаль --- для образования новой, с собственным феноменологическим и мотивационным профилем, а также с собственным набором триггеров, при этом первоначальные диспозиции остаются нетронутыми.

Когнитивная разработка аналогично подразумевает изменения существующих эмоциональных диспозиций: например, когнитивно выработанная склонность испытывать гнев в некоторых ситуациях приводит к диспозиции испытывать его и в некоторых других, не затрагивая исходную диспозицию. То есть агент сохраняет диспозицию испытывать эмоцию при столкновении с определенными стимулами, но приобретает еще одну, связанную с несправедливостью, например. Калибровка же не приводит к выработке новой эмоциональной диспозиции, а изменяет имеющуюся.

Калибровку можно рассматривать как обучение --- приобретение неявного знания об оценочных свойствах вещей. Разумеется, если вы начинаете, например, боятся безобидных вещей, такое изменение нельзя связать со знанием --- это просто ошибка.

Но как именно происходит калибровка?

\section{От обуславливания к социальному обучению}

Андреас Олссон и Элизабет Фелпс в своей влиятельной обзорной статье ''Социальное обучение страху'' (2007) указывают три способа изменения диспозиции к страху: обуславливание, обучение через наблюдение и обучение под руководством. Обуславливание подразумевает переживание аверсивного стимула. Обуславливание страха, изучение которого восходит к работам Павлова об условных рефлексах (1927) состоит в сочетании естественного аверсивного раздражителя с нейтральным, благодаря чему нейтральный раздражитель начинает ассоциироваться со страхом: например, когда вас кусает собака, вам больно и неприятно, и эти переживания, ассоциируясь с собакой, формируют боязнь всяких собак вообще. Возможно, конечно, и так называемое положительное подкрепление, когда аверсивный стимул сочетается с достаточно сильным приятным, что постепенно снижает страх, связанный с первым.

Олссон и Фелпс, впрочем, подчеркивают, что выяснение, чего стоит боятся, на основе только лишь аверсивных стимулов слишком травмоопасно. Более того, из-за того, что вы так полагаетесь на ограниченный объем опыта, часть которого может быть и вовсе неверно истолкована, процесс обучения может оказаться ненадежным. Приобретение же страхов косвенно, через рассказы об опасных и неприятных вещах, например, не предполагает травмоопасности, а также часто основана на интерсубъективном, что делает его более надежным. Социальное обучение, основанное на наблюдении за другими, позволяет избежать недостатков обучения через непосредственное участие.

Так, молодые особи макак-резусов учатся бояться змей, наблюдая за реакциями сородичей. Олссон и Фелпс (2007) отмечают, что способность обнаруживать признаки страха и адекватно реагировать на них обнаружена у многих нечеловеческих животных, включая приматов, птиц, мышей, кошек, и дает существенное эволюционное преимущество, ведь позволяет быстро научаться правильным реакциям. Например, вороны на основе наблюдения за тем, как определенным образом выглядящие люди поймали и связали их сородичей, научились бояться так выглядящих людей, даже если не видели самой сцены поимки (Корнелл, Марзлафф и Пекораро, 2012).

Страх очень заразителен: просмотр фильма, в котором кто-то выражает страх, может вызвать страх у зрителя. В случае людей достаточно даже простого знания о страхах другого, без наблюдения его соответствующей реакции, хотя, конечно, процесс перенимания этих страхов будет обусловлен и прочими факторами --- знакомством или родством с этим человеком, а также его социальным статусом, что будет обосновывать доверие. То есть мы перенимаем эмоциональные диспозиции только у определенной части окружающих людей, что объясняет важность ролевых моделей в эмоциональном обучении.

Третий способ обучения специфичен для существ, говорящих на одном языке: например, когда родители говорят детям, что трогать электрические розетки опасно. Олссон и Фелпс (2007) пишут: ''экспериментальные исследования детей, связанные с рассказыванием историй, показывают, что вербальные инструкции могут быть сильным стимулом для обучения страху'' (Олссон и Фелпс 2007, 1099, см. также Филд, Аргирис и Ноулз 2001). Детям предлагалось представить, что две куклы --- монстры. Им показывали соответствующее видео или рассказывали соответствующие истории, и результаты показали, что обучение посредством вербального общения может быть столь же эффективным, как и обучение через наблюдение.

Конечно же, обучение через вербальное общение так же обусловлено социальной ролью информатора, как и любое другое обучение, так что дети, например, будут больше доверять родителям, чем незнакомым людям. Перечисленные способы обучения работают не только для страха, но и для любой другой эмоции.

\section{Парадоксальная сила художественной литературы}

Просмотр фильмов или чтение романов вызывают у нас самые разные эмоции. Дженефер Робинсон пишет: ''Хорошая история вызывают у нас любопытство и напряжение по поводу дальнейшего развития сюжета, может заставить нас смеяться или плакать, чувствовать страх, гнев и даже ужас, отвращение, любовь и сострадание --- в общем, весь спектр эмоций'' (2005, 105).

Художественные произведения вызывают эмоции отчасти так же, как и музыкальные: фильм или роман делают это благодаря связи с эмоционально заряженными событиями --- ассоциациям с ними и их оценочным свойствам. Но, конечно, в отличие от музыки, здесь наши эмоции провоцируют также события и персонажи произведения. Так, в трагедии Шекспира ''Макбет'' есть ужасающий момент, когда Макдуф узнает, что Макбет убил леди Макдуф и всех их детей --- нас ужасает факт убийств и сам Максет. Во время просмотра фильма мы даже можем почувствовать нечто вроде страха: представьте Чарльза, который смотрит фильм, где монстр внезапно прыгает в его сторону. Дерек Мэтраверс (1998) отмечает, что этот страх довольно необычен: монстр обитает в вымышленном мире, так что не может в действительности угрожать Чарльзу, так что в его реакции есть нечто аномальное, как, впрочем, и в поведении монстра, когда он будто бы бросается на Чарльза.

Наиболее эффективный способ, с помощью которого вымысел вызывает эмоции --- наше сопереживание персонажам. При чтении романа Толстого ''Анна Каренина'', мы, наверняка, будем испытывать массу эмоций по поводу судьбы главной героини: сочувствовать, когда ей не позволяют видеться с сыном, которого она очень любит, бояться за нее, когда она начинает думать о самоубийстве. В целом, отличительной особенностью многих эмоций, испытываемых при просмотре фильмов или чтении романов, является их отношение к персонажам, нам небезразличным --- до определенной степени мы идентифицируем себя с ними, сопереживаем им. Как указывает Кейт Оутли (1999), именно наша творческая вовлеченность в вымысел является источником соответствующих эмоций: мы отчасти разделяем опыт персонажей, отчего склонны испытывать эмоции, сходные с теми, что испытывают они.

Робинсон (2005) подчеркивает, что способ повествования также играет важную роль: используя яркие метафоры, писатель может заставить нас сосредоточиться на определенной особенности, что повлияет на наши эмоции по отношению к персонажам. Вообще от того, как персонаж описывается, зависит наша склонность идентифицировать себя с ним.

Но, все же, как вообще могут вымышленные персонажи вызывать эмоции? Они ведь не существуют. Этот так называемый парадокс вымысла включает три положения:

\begin{enumerate}
  \item\ Мы испытываем эмоции, направленные на вымышленных персонажей и вымышленные истории.
  \item\ Чтобы испытать эмоцию по отношению к объекту, необходимо верить, что он существует.
  \item\ Мы не верим в существование вымышленных персонажей и историй.
\end{enumerate}

Большинство решений парадокса предполагают отказ от одного из трех тезисов. Можно утверждать, например, что эмоции, испытываемые в отношении вымышленного --- не подлинные. Наиболее известный вариант такого решения предложен Уолтоном (1978): просмотр фильма или чтение романа сопровождается притворством, сравнимом с играми детей, когда они, играя в песочнице, притворяются, что пекут печенье. Соответствующие эмоции, значит --- тоже притворство, и Чарльз на самом деле не боится монстра, а лишь участвует в игре, в которой притворяется, что испытывает страх. Это, конечно, не значит, что Чарльз на самом деле ничего не испытывает --- даже квазиэмоция схожа с действительным страхом.

Однако, согласно Уолтону, состояние Чарльза кардинально отличается от такового при настоящем страхе, ведь Чарльз не считает, что есть опасность и не имеет, значит, соответствующей мотивации --- Чарльз не выбегает из кинотеатра в попытках спасти свою жизнь. Но, как мы видели в Главе 6, суждения не требуются для эмоций --- чтобы бояться, не обязательно верить, что грозит опасность. Более того, как указывалось в Главе 5, связь между эмоциями и мотивацией не столь уж сильна, и эмоции, направленные на вымышленное, схожи с испытываемыми по поводу действительного: когда мы смотрим документальный фильм или читаем новости, например --- в большинстве подобных случаев мы не находимся в прямой конфронтации с сообщаемой реальностью, так что неудивительно, что переживаемые эмоции не приводят к действиям.

Основной довод в пользу утверждения Уолтона о квазиэмоциях --- Чарльз не верит в существование монстра: ''Чарльз полностью осознает, что монстр не настоящий, и это [...] веская причина отрицать соответствующую природу его страха'' (1978, 6). То есть Уолтон придерживается второго пункта парадокса и отрицает искренность эмоций, направленных на вымышленное. Но так ли нужно верить в существование объекта, чтобы возникли направленные на него эмоции?

Наиболее правдоподобное решение парадокса состоит в отказе от второго его пункта. Как указывает Робинсон: ''эмоции не требую веры в свои ''объекты'' для возникновения'' (2005, 151). Если не обязательно верить в существование объекта чтобы возникла эмоция о нем, то что заменяет эту веру? Возможно, достаточно лишь представлять, что нечто существует, чтобы по поводу него смогли возникать эмоции (Ламарк 1981, Кэролл 1990). Таким образом, эмоции, направленные на вымышленное, включают воображение этого вымышленного как существующего: Чарльзу достаточно представить, что монстр собирается его уничтожить, чтобы испытать страх по поводу него. В целом возможность воображения приводить к эмоциям общепризнанна (Дамасио 1994, Оутли 1999, Галлезе 2001): одни и те же нейрологические сети активируются независимо от действительности того, на что направлено внимание.

Следующий вопрос касается рациональности эмоций по поводу вымышленного. Как они могут быть рациональны, если их объекты не существуют? Возможно, они основаны на убеждениях о том, что происходит в произведении (Терони 2019). Так, опасаясь за жизнь Анны Карениной, мы не верим, что она действительно существует и что ей действительно угрожает опасность --- мы верим, что она существует в рамках вымысла и что там ей и угрожает опасность. То есть страх вполне рационален, ведь и основан на убеждениях касательно происходящего в рамках вымысла, и направлен на происходящее в нем.

Итак, полноценные эмоции по отношению к вымышленному не только возможны, но и рациональны. Но почему нам так нравится смотреть фильмы и читать художественную литературу, если это вызывает массу в том числе и негативных эмоций?

\section{Парадокс негативных эмоций}

Как мы видели в Разделе 4 Главы 11 парадокс грустной музыки в склонности слушать грустную музыку несмотря на склонность избегать грусти. Парадокс негативных эмоций можно сформулировать в схожих терминах (Левинсон 1997, Смэтс 2009, Наар 2013a):

\begin{enumerate}
  \item\ Такие художественные произведения, как трагедии и триллеры, вызывают негативные эмоции.
  \item\ Негативные эмоции --- аверсивные аффективные состояния, которых мы склонны избегать.
  \item\ Мы намеренно смотрим / читаем трагедии и триллеры.
\end{enumerate}

Многие философы утверждают совместимость этих пунктов и объясняют эту совместимость некоторой выгодой от просмотра или чтения того, что вызывает негативные эмоции. Так, согласно Аристотелю, трагедии способствуют катарсису, то есть очищению от негативных эмоций.

Но, похоже, во многих случаях мы смотрим трагедии не с целью достижения какой-то дальнейшей выгоды или пользы, и находим сам этот опыт просмотра привлекательным. Это же возражение применимо к решению через целесообразность просмотра трагедий из-за знаний о природе и поведении людей, которые они дают (Пакер 1989, Нуссбаум 2001).

Может быть, эмоции вроде страха и отвращения не обязательно аверсивны? Некоторые считают, что следует отвергнуть второй пункт парадокса (Уолтон 1990, Нил 1992, Левинсон 1982). Возможно, эмоции ни приятны, ни неприятны, приятность --- свойство интенционального их объекта, но не их самих. Но обычно мы понимаем страх, например, как неприятное и аверсивное переживание. Быть может, негативные эмоции, направленные на вымышленное --- не аверсивны, поскольку переживаются в отсутствие какого-либо реального взаимодействия (Левинсон 1982). Но хотя страх при столкновении с настоящим медведем сильнее, чем таковой при только лишь этой встречи воображении, негативная составляющая его феноменологии не выглядит иной. Некоторые люди даже избегают произведений искусства, вызывающих негативные эмоции.

Аарон Смэтс (2009) утверждает, что в просмотре трагедий и триллеров нас интересует то богатство опыта, которое они позволяют получить, то множество в том числе и негативных эмоций, что они вызывают. Переживание этих эмоций по поводу вымышленного помогает затем в решении проблем в действительности: ''искусство безопасно дает нам возможность получить богатые эмоциональные переживания, которые либо невозможно, либо слишком рискованно получать в повседневной жизни'' (2009, 74). Но почему мы находим такой богатый опыт привлекательным? Если бы у нас была возможность выбора, разве не предпочли бы мы испытать богатый положительный опыт, свободный от аверсивных эмоций?

Итак остался последний вариант --- отрицать сказанное в первом пункте парадокса. Действительно ли такие произведения, как трагедии и триллеры, вызывают негативные эмоции? Дэвид Юм (1957/1985b) предполагает, что неприятные эмоции от просмотра трагедий могут быть преобразованы в положительный опыт, если учесть получаемое эстетическое удовольствие. Это удовольствие, будучи достаточно сильным, смешиваясь с эмоциями, приводит к в целом положительному опыту. То есть наша реакция на вымышленное частично обусловлена эстетическими свойствами произведения искусства. Вопрос того, как различные реакции объединяются для порождения целостного переживания, сложен. Даже если в некоторых случаях происходит что-то вроде описанного Юмом, трудно видеть в этом решение парадокса.

Похоже, нам придется принять плюралистическое решение (Наар 2013a, Штрол 2019): в зависимости от случая будут уместны разные объяснения. В некоторых случаях это получаемая польза, например, знания, в других --- положительный опыт, которых перевешивает негативные эмоции. Это позволяет объяснить и грустные произведения, которые, как и грустная музыка, могут вызывать в целом положительную эмоцию волнения, и именно эту эмоцию мы стремимся испытать, когда читаем грустные романы или смотрим мелодраму. Поклонники триллеров же, очевидно, хотят испытать острые ощущения, которые в целом приятны, и, хотя коррелируют со страхом, все же отличаются от него (см. Бакли 2016). Возможно, Ноэль Кэрролл (1990) прав, отмечая, что поклонников этого жанра испытывают не только и не столько страх, сколько любопытство и острые ощущения.

\section{Литература и сентиментальное воспитание}

Распространено мнение, что чтение художественной литературы помогает нам стать морально лучше. Так, Марта Нуссбаум (1990, 2001) утверждает, что чтение реалистических романов XIX века, таких как ''Золотая чаша'' Генри Джеймса, дает знания о морали --- вообще, такие романы лучше, чем трактаты в деле получения этих знаний, потому что трактаты о морали дают лишь интеллектуальное понимание некоторых положений, не учитывая сложности реальной жизни. Робинсон (2005) подчеркивает, что усвоение морального знания зависит от переживаемых эмоций, и что поэтому взаимодействие с вымыслом способствует нашему нравственному воспитанию: хороший роман побуждает читателя размышлять над воспитанием эмоций персонажей, а также над их реакциями на происходящее.

Впрочем, сама Робинсон указывает, что речь скорее о ''воспитании эмоциями'' (2005, 156), чем о воспитании эмоций. Даже если описываемый в художественном произведении мир вымышлен писателем, наши реакции на этот мир соответствуют реакциям на мир реальный. Так что калибровка эмоциональных диспозиций, происходящая при взаимодействии с вымышленными персонажами, ничем не отличается от таковой в реальной жизни. Давайте рассмотрим, как процессы калибровки --- обуславливание, обучение через наблюдение и обучение по инструкции --- работают в контексте чтения художественной литературы.

Обуславливание заключается в сочетании аверсивного стимула с нейтральным. Поскольку мы не в состоянии ощутить реальные стимулы при чтении романа, может показаться, что обуславливание не будет работать. Но эмоции по поводу происходящего в произведении очевидно возможны. Если в романе описывается, как собака бросается на ребенка и кусает его, мы сопереживаем ему и испытываем в целом неприятные эмоции, что может привести к ассоциации собак с аверсивным опытом и, следовательно, к приобретению склонности испытывать страх при встрече с ними.

Социальное обучение посредством наблюдения работает при чтении художественной литературы куда проще: читая романы, мы наблюдаем за тем, как персонажи эмоционально реагируют  на ситуации, и, в зависимости от способствующих доверию факторов --- знакомству с персонажем и его социального статуса --- мы скорее всего приобретем диспозицию к аналогичным эмоциональным реакциям. Это, наверное, основной способ воздействия кино и литературы на эмоциональные диспозиции.

Воздействие через обучение по инструкции подразумевает получение информации об оценочных характеристиках: истории, включающие положительную или отрицательную оценку нечта влияют на оценку их читателем / слушателем этого нечта. Такую информацию может давать рассказчик или персонажи: персонаж может высказать оценочное мнение, и если мы ему доверяем (это зависит от того, как он в истории представлен рассказчиком), то наши эмоциональные диспозиции соответствующе меняются.

Но ведь тогда чтение художественной литературы может сильно исказить восприятие, а точнее, оценку, реальности: может возникнуть как адекватная, так и вовсе неадекватная конфигурация диспозиций. Обучение посредством вымышленного, впрочем, сопряжено с не большим риском, чем обучение в реальности. Корректность усваиваемых эмоциональных предрасположенностей не зависит от реальности информатора.

Мы обыкновенно склонны сопереживать персонажам произведения, тем самым перенимая их точку зрения и испытывая эмоции схожие с испытываемыми ими. Чтение романов, значит, особенно тех, что эмоционально вовлекательны, позволяет развить способность сопереживать. Дело в том, что эмпатия связана с альтруистической мотивацией, и чем больше мы склонны к сопереживанию, тем больше склонны помогать тем, кто нуждается в помощи.

Есть исследования, показывающие, что чтение романов может быть мощной стратегией улучшения отношения к чуждой группе. Одно из них включало чтение отрывков из книг Джоан Роулинг о Гарии Поттере и показало, что сильная идентификация с главным героем повысила эмпатию к стигматизированным группам. В целом читатели художественной литературы, как правило, лучше справляются с задачами на эмпатию, чем те, кто читает в основном литературу техническую (Мар, Оатли и Петерсон 2009). Интересно, что основное повышение эмпатии происходит не сразу после чтения, а неделю спустя.

Термин ''эмпатия'' допускает разные значения, но мы будем говорить об эмпатии как о процессе, включающем воображение и приводящем к опосредованному разделению гомологичных аффектов (см. Раздел 7 Главы 10). Например, когда вы сопереживаете чувствующему радость, вы испытываете эмоцию, схожую с его радостью. Эмпатия отличается от сочувствия, потому что сочувствие подразумевает страдание того, на кого направлено и желание это страдание прекратить. Эмпатия также отличается от только склонности испытывать эмоции, направленные на других.

Эмпатия, впрочем, отличаясь от сочувствия, все же способствует ему: как отмечает Нуссбаум (2001), она ''является очень важным инструментом, позволяющим понять, что происходит с другим человеком и установить с ним контакт'' (2001, 330-331). Эмпатия также помогает сосредоточиться на тяжелом положении другого человека, так что разумеется, эмпатия и сочувствие обыкновенно коррелируют. Поскольку сочувствие репрезентирует чужое страдание как нечто плохое, то подразумевает, что страдание, в целом --- это плохо. Эмпатия, значит, играет роль детектора страдания.

Эмпатию как таковую часто критикуют за недостатки (см. Принц 2001, Блюм 2017). Чтение художественной литературы, однако, помогает преодолеть один из них, расширяя спектр известных случаев страдания. Но, конечно, не всякая литература будет иметь такой эффект --- некоторая может, например, усиливать предубеждения об определенных группах людей, усиливая отстраненность от них и оттого уменьшая возможную к ним эмпатию и сочувствие, что означает наличие существенного компонента удачи в формировании эмоциональных диспозиций посредством литературы.

\section{Подведение итогов}

Мы начали эту главу с вопроса: могут ли художественные произведения способствовать не только регулированию эмоций, но и воспитанию эмоциональных диспозиций. Мы увидели, что воспитание этих диспозиций заключается в их калибровке с помощью обуславливания, обучения посредством наблюдения и обучения посредством инструкций. Изучая, как художественные произведения вызывают эмоции, мы пришли к выводу, что они делают это во многом так же, как и реальные ситуации, то есть воображения чего-то, что в реальности вызвало бы эмоции достаточно для возникновения этих эмоций, так что калибровка посредством вымышленного не отличается принципиально от таковой посредством действительного. Калибровка посредством вымысла, значит --- мощный, хотя и рискованный способ изменения эмоциональных диспозиций.

\section{Краткое содержание}

\begin{itemize}
  \item\ Сентиментальное воспитание является важным этапом в нравственном воспитании человека и отличается от регулирования эмоций направленностью на развитие эмоциональных диспозиций для переживания уместных эмоций.
  \item\ И биологические детерминисты, и социальные конструктивисты признают податливость эмоциональных диспозиций.
  \item\ Помимо смешения и когнитивной разработки, есть еще один способ влияния на эмоции --- калибровка --- процесс, посредством которого изменяется набор триггеров эмоции.
  \item\ Есть три способа калибровки: обуславливание, обучение через наблюдение и обучение посредством прямых инструкций.
  \item\ Вымысел вызывает эмоции почти так же, как реальность, и это настоящие, подлинные эмоции, которые могут быть рациональными.
  \item\ Мы иногда смотрим фильмы и читаем книги, вызывающие в том числе негативные эмоции. Это потому, что они помогают правильно откалибровать эмоциональные диспозиции, а также вызывают и положительные эмоции --- любопытство, волнение или восторг.
  \item\ Чтение художественной литературы оказывает существенное влияние на наши эмоциональные диспозиции, поскольку все три способа калибровки работают как в отношении реального, так и в отношении вымышленного.
  \item\ Приведет ли калибровка посредством чтения литературы к сентиментальному воспитанию --- зависит, конечно, от литературы.
  \item\ Чтение эмоционально вовлекающей литературы развивает способность к эмпатии и альтруистическую мотивацию.
  \item\ Эмпатия отличается как от сочувствия, так и от диспозиции к сочувствию, но, поскольку служит своего рода детектором страдания, играет важную роль в возникновении обоих.
\end{itemize}

\begin{tcolorbox}
  \section{Учебные вопросы}
  \begin{enumerate}
    \item\ Приведите пример изменения эмоциональных диспозиций, не являющийся сентиментальным воспитанием.
    \item\ Стоики настаивали на искоренении некоторых эмоций. Что потребуется для искоренения гнева, какие изменения каких сентиментов? Насколько такое искоренение реалистично?
    \item\ Как у нас формируются образцы для подражания в лице вымышленных персонажей и есть ли какие-то особо надежные методы выбора таких образцов?
  \end{enumerate}
\end{tcolorbox}

\end{document}
